\documentclass[12pt,a4paper,openany]{book}
\usepackage[utf8]{inputenc}
\usepackage[T1]{fontenc}
\usepackage[french]{babel}
\usepackage[margin=2.5cm]{geometry}
\usepackage{amsmath,amssymb,amsthm,mathrsfs,mathtools}
\usepackage{tikz}
\usetikzlibrary{arrows.meta,calc,positioning,patterns,decorations.markings,decorations.pathmorphing,cd}
\usepackage{pgfplots}
\pgfplotsset{compat=1.18}
\usepackage[most]{tcolorbox}
\tcbuselibrary{theorems,skins,breakable}
\usepackage[colorlinks=true,linkcolor=blue!60!black,citecolor=green!50!black]{hyperref}
\usepackage{makeidx}
\makeindex
\usepackage{enumitem}
\usepackage{fancyhdr}
\usepackage{caption}
\usepackage{booktabs}
\usepackage{microtype}
\usepackage{xcolor}
\usepackage{minitoc}
\dominitoc

\pagestyle{fancy}
\fancyhf{}
\fancyhead[LE]{\leftmark}
\fancyhead[RO]{\rightmark}
\fancyfoot[C]{\thepage}

% ── Theorem environments (amsthm) ──
\theoremstyle{definition}
\newtheorem{definition}{Définition}[chapter]
\newtheorem{exemple}[definition]{Exemple}
\newtheorem{exercice}{Exercice}[chapter]

\theoremstyle{plain}
\newtheorem{theoreme}[definition]{Théorème}
\newtheorem{proposition}[definition]{Proposition}
\newtheorem{lemme}[definition]{Lemme}
\newtheorem{corollaire}[definition]{Corollaire}

\theoremstyle{remark}
\newtheorem{remarque}[definition]{Remarque}

% ── tcolorbox wrapping (visual only) ──
\tcolorboxenvironment{definition}{enhanced,breakable,colback=blue!5,colframe=blue!70!black,fonttitle=\bfseries,before skip=10pt,after skip=10pt,left=8pt,right=8pt}
\tcolorboxenvironment{theoreme}{enhanced,breakable,colback=red!5,colframe=red!70!black,fonttitle=\bfseries,before skip=10pt,after skip=10pt,left=8pt,right=8pt}
\tcolorboxenvironment{proposition}{enhanced,breakable,colback=orange!5,colframe=orange!70!black,fonttitle=\bfseries,before skip=10pt,after skip=10pt,left=8pt,right=8pt}
\tcolorboxenvironment{lemme}{enhanced,breakable,colback=yellow!5,colframe=yellow!50!black,fonttitle=\bfseries,before skip=10pt,after skip=10pt,left=8pt,right=8pt}
\tcolorboxenvironment{corollaire}{enhanced,breakable,colback=green!5,colframe=green!50!black,fonttitle=\bfseries,before skip=10pt,after skip=10pt,left=8pt,right=8pt}
\tcolorboxenvironment{remarque}{enhanced,breakable,colback=gray!5,colframe=gray!50!black,before skip=8pt,after skip=8pt,left=6pt,right=6pt}
\tcolorboxenvironment{exemple}{enhanced,breakable,colback=purple!5,colframe=purple!50!black,before skip=8pt,after skip=8pt,left=6pt,right=6pt}

% ── Custom commands ──
\newcommand{\C}{\mathbb{C}}
\newcommand{\R}{\mathbb{R}}
\newcommand{\Z}{\mathbb{Z}}
\newcommand{\N}{\mathbb{N}}
\newcommand{\Q}{\mathbb{Q}}
\newcommand{\D}{\mathbb{D}}
\newcommand{\T}{\mathbb{T}}
\DeclareMathOperator{\re}{Re}
\DeclareMathOperator{\im}{Im}
\DeclareMathOperator{\Arg}{Arg}
\DeclareMathOperator{\Log}{Log}
\DeclareMathOperator{\Res}{Res}
\DeclareMathOperator{\ind}{ind}
\newcommand{\dd}{\mathrm{d}}
\newcommand{\abs}[1]{\left|#1\right|}
\newcommand{\norm}[1]{\left\|#1\right\|}
\newcommand{\conj}[1]{\overline{#1}}

\title{\Huge\textbf{Analyse Complexe}\\[0.5cm]
\Large Notes de Cours\\[0.3cm]
\large Licence L3, Mathématiques}
\author{}
\date{}

\begin{document}

% ═══════════════════════════════════════════════════════════════════════════════
% PAGE DE TITRE
% ═══════════════════════════════════════════════════════════════════════════════
\begin{titlepage}
\begin{tikzpicture}[remember picture, overlay]
  \fill[left color=blue!15!white, right color=blue!5!white]
    (current page.south west) rectangle (current page.north east);
  \fill[color=orange!70!yellow]
    ([yshift=-2cm]current page.north west) rectangle ([yshift=-2.3cm]current page.north east);
  \fill[color=blue!80!black, rounded corners=8pt]
    ([xshift=3cm, yshift=-4cm]current page.north west) rectangle
    ([xshift=-3cm, yshift=-10cm]current page.north east);
  \node[anchor=center, text=white, font=\Huge\bfseries]
    at ([yshift=-6cm]current page.north) {Analyse Complexe};
  \node[anchor=center, text=orange!80!yellow, font=\Large]
    at ([yshift=-7.5cm]current page.north) {Notes de Cours};
  \node[anchor=center, text=white!80, font=\large]
    at ([yshift=-8.8cm]current page.north) {Licence L3 --- 2025--2026};
  \node[anchor=center, text=blue!80!black, font=\Large\itshape]
    at ([yshift=-12cm]current page.north) {Ya\'e Ulrich Gaba};
  \draw[orange!70!yellow, line width=2pt]
    ([xshift=5cm, yshift=-13.5cm]current page.north west) --
    ([xshift=-5cm, yshift=-13.5cm]current page.north east);
  \node[anchor=center, text width=12cm, align=center, text=blue!60!black, font=\itshape]
    at ([yshift=-16cm]current page.north) {<<\,Le plus court chemin entre deux v\'erit\'es dans le domaine r\'eel passe par le domaine complexe.\,>>\\[4pt]--- Jacques Hadamard};
  \node[anchor=south, text=gray, font=\small]
    at ([yshift=1.5cm]current page.south) {\today};
\end{tikzpicture}
\end{titlepage}

\tableofcontents

% ═══════════════════════════════════════════════════════════════════════════════
% PRÉFACE
% ═══════════════════════════════════════════════════════════════════════════════
\chapter*{Préface}
\addcontentsline{toc}{chapter}{Préface}

L'analyse complexe est l'une des théories les plus élégantes et les plus puissantes des mathématiques. Là où l'analyse réelle offre une grande liberté --- fonctions continues non dérivables, fonctions dérivables non deux fois dérivables, etc. --- l'analyse complexe impose une rigidité remarquable : une seule dérivée complexe suffit à garantir que la fonction est infiniment dérivable, analytique, et entièrement déterminée par ses valeurs sur n'importe quel ouvert.

Cette rigidité provient de la structure algébrique et géométrique du corps $\C$. La dérivabilité complexe, formalisée par les \emph{équations de Cauchy-Riemann}, lie les dérivées partielles réelles de manière à créer un couplage profond entre la partie réelle et la partie imaginaire d'une fonction holomorphe. Ce couplage a des conséquences spectaculaires : le théorème intégral de Cauchy, la formule intégrale de Cauchy, le développement en série entière, le principe du maximum, le théorème de Liouville, et bien d'autres résultats fondamentaux.

\medskip

\noindent\textbf{Connexions avec d'autres domaines.}
L'analyse complexe n'est pas une discipline isolée. Ses applications irriguent de nombreux champs :
\begin{itemize}[leftmargin=2cm]
  \item \textbf{Mécanique des fluides.} Les fonctions holomorphes décrivent les écoulements plans irrotationnels de fluides incompressibles. La partie réelle donne le potentiel de vitesse, la partie imaginaire la fonction de courant.
  \item \textbf{Électrostatique.} Les fonctions harmoniques, étroitement liées aux fonctions holomorphes, modélisent les potentiels électrostatiques dans le plan.
  \item \textbf{Théorie des nombres.} La fonction zêta de Riemann $\zeta(s) = \sum_{n=1}^{\infty} n^{-s}$, définie comme fonction d'une variable complexe, encode la distribution des nombres premiers.
  \item \textbf{Géométrie et topologie.} Les surfaces de Riemann, les fibrés en droites, la théorie de Hodge trouvent leurs racines dans l'analyse complexe.
  \item \textbf{Théorie du contrôle.} Les transformées de Laplace et de Fourier reposent sur l'évaluation de fonctions dans le plan complexe.
\end{itemize}

\medskip

\noindent\textbf{Prérequis.}
Ce cours suppose acquis :
\begin{itemize}[leftmargin=2cm]
  \item Analyse réelle (suites, séries, continuité, dérivabilité, intégration).
  \item Topologie élémentaire (ouverts, fermés, compacité, connexité dans $\R^n$).
  \item Algèbre linéaire (espaces vectoriels, applications linéaires, matrices).
  \item Calcul différentiel en plusieurs variables (dérivées partielles, différentielle).
\end{itemize}

\medskip

\noindent\textbf{Plan du cours.}
\begin{description}[leftmargin=1.5cm]
  \item[Chapitre 1.] Nombres complexes : rappels algébriques, géométrie du plan complexe, topologie de $\C$, sphère de Riemann.
  \item[Chapitre 2.] Fonctions holomorphes : dérivée complexe, équations de Cauchy-Riemann, conformité, fonctions harmoniques.
  \item[Chapitre 3.] Séries entières et fonctions analytiques.
  \item[Chapitre 4.] Intégration complexe et théorème de Cauchy.
  \item[Chapitre 5.] Formule intégrale de Cauchy et applications.
  \item[Chapitre 6.] Séries de Laurent et théorie des résidus.
  \item[Chapitre 7.] Applications conformes et théorème de Riemann.
\end{description}

% ═══════════════════════════════════════════════════════════════════════════════
% NOTATIONS
% ═══════════════════════════════════════════════════════════════════════════════
\chapter*{Notations}
\addcontentsline{toc}{chapter}{Notations}

\begin{tabular}{@{}ll@{}}
\toprule
\textbf{Symbole} & \textbf{Signification} \\
\midrule
$\N, \Z, \Q, \R, \C$ & Ensembles de nombres usuels \\
$\re(z),\; \im(z)$ & Partie réelle et partie imaginaire de $z \in \C$ \\
$\abs{z}$ & Module de $z$ \\
$\conj{z}$ & Conjugué de $z$ \\
$\arg(z)$ & Argument de $z$ (multivalué) \\
$\Arg(z)$ & Argument principal, dans $(-\pi,\pi]$ \\
$\Log(z)$ & Logarithme principal de $z$ \\
$\D$ & Disque unité ouvert $\{z \in \C : \abs{z} < 1\}$ \\
$\T$ & Cercle unité $\{z \in \C : \abs{z} = 1\}$ \\
$D(a,r)$ & Disque ouvert de centre $a$ et rayon $r$ \\
$\conj{D}(a,r)$ & Disque fermé de centre $a$ et rayon $r$ \\
$C(a,r)$ & Cercle de centre $a$ et rayon $r$ \\
$\hat{\C} = \C \cup \{\infty\}$ & Plan complexe étendu (sphère de Riemann) \\
$\mathscr{O}(\Omega)$ & Fonctions holomorphes sur $\Omega$ \\
$\mathscr{C}^k(\Omega)$ & Fonctions de classe $\mathscr{C}^k$ sur $\Omega$ \\
$f'(z_0)$ & Dérivée complexe de $f$ en $z_0$ \\
$\frac{\partial f}{\partial z},\;\frac{\partial f}{\partial \conj{z}}$ & Opérateurs de Wirtinger \\
$\Res_{z=a} f$ & Résidu de $f$ en $a$ \\
$\ind_\gamma(a)$ & Indice du point $a$ par rapport au lacet $\gamma$ \\
$\dd z = \dd x + i\,\dd y$ & Différentielle complexe \\
\bottomrule
\end{tabular}

% ═══════════════════════════════════════════════════════════════════════════════
% CHAPITRE 1
% ═══════════════════════════════════════════════════════════════════════════════
\chapter{Nombres complexes --- Rappels et géométrie du plan}\label{ch:nombres-complexes}

\section{Construction algébrique de \texorpdfstring{$\C$}{C}}\label{sec:construction-C}

Le corps des nombres complexes peut être construit de plusieurs manières équivalentes. Nous présentons ici la construction algébrique classique.

\begin{definition}[Corps des nombres complexes]\label{def:corps-C}
On définit $\C = \R^2$ muni des opérations suivantes : pour $(a,b), (c,d) \in \R^2$,
\begin{align}
  (a,b) + (c,d) &= (a+c,\, b+d), \label{eq:addition-C} \\
  (a,b) \cdot (c,d) &= (ac - bd,\, ad + bc). \label{eq:multiplication-C}
\end{align}
On note $i = (0,1)$ l'unité imaginaire et on identifie $a \in \R$ avec $(a,0) \in \C$. Tout élément $z = (a,b)$ s'écrit alors $z = a + ib$.
\end{definition}

\begin{theoreme}[Structure de corps]\label{thm:C-corps}
$(\C, +, \cdot)$ est un corps commutatif. Son élément neutre pour l'addition est $0 = (0,0)$, et son élément neutre pour la multiplication est $1 = (1,0)$. De plus, $i^2 = -1$.
\end{theoreme}

\begin{proof}
La vérification des axiomes de corps est un calcul direct. Vérifions les points essentiels.

\emph{Associativité de la multiplication.} Soient $z_1 = a_1 + ib_1$, $z_2 = a_2 + ib_2$, $z_3 = a_3 + ib_3$. On calcule $(z_1 z_2)z_3$ et $z_1(z_2 z_3)$ et on vérifie que les deux expressions sont égales par développement. Ce calcul, bien que fastidieux, ne présente aucune difficulté.

\emph{Inverse multiplicatif.} Soit $z = a + ib \neq 0$. Alors $a^2 + b^2 > 0$ et
\[
  z^{-1} = \frac{a}{a^2 + b^2} - i\frac{b}{a^2 + b^2} = \frac{a - ib}{a^2 + b^2}.
\]
On vérifie que $z \cdot z^{-1} = 1$.

\emph{Unité imaginaire.} $i^2 = (0,1) \cdot (0,1) = (0\cdot 0 - 1\cdot 1,\; 0\cdot 1 + 1\cdot 0) = (-1,0) = -1$.

Les autres axiomes (commutativité, distributivité) se vérifient de manière analogue.
\end{proof}

\begin{remarque}\label{rem:C-non-ordonne}
Le corps $\C$ ne peut pas être muni d'un ordre total compatible avec ses opérations. En effet, si un tel ordre existait, on aurait soit $i > 0$, d'où $i^2 = -1 > 0$, contradiction ; soit $i < 0$, d'où $-i > 0$ et $(-i)^2 = -1 > 0$, même contradiction. Cette impossibilité est fondamentale : elle explique pourquoi de nombreuses techniques d'analyse réelle (comparaison, encadrement) ne se transposent pas directement en analyse complexe.
\end{remarque}

\begin{definition}[Partie réelle, partie imaginaire, conjugué]\label{def:re-im-conj}
Soit $z = a + ib \in \C$ avec $a, b \in \R$.
\begin{enumerate}[label=(\roman*)]
  \item La \emph{partie réelle} de $z$ est $\re(z) = a$.
  \item La \emph{partie imaginaire} de $z$ est $\im(z) = b$ (c'est un réel, non un imaginaire pur).
  \item Le \emph{conjugué} de $z$ est $\conj{z} = a - ib$.
\end{enumerate}
\end{definition}

\begin{proposition}[Propriétés du conjugué]\label{prop:proprietes-conjugue}
Pour tous $z, w \in \C$ :
\begin{enumerate}[label=(\roman*)]
  \item $\conj{z + w} = \conj{z} + \conj{w}$, \quad $\conj{z \cdot w} = \conj{z} \cdot \conj{w}$.
  \item $\conj{\conj{z}} = z$.
  \item $z + \conj{z} = 2\re(z)$, \quad $z - \conj{z} = 2i\,\im(z)$.
  \item $z \conj{z} = \re(z)^2 + \im(z)^2 \geq 0$, avec égalité si et seulement si $z = 0$.
  \item $z \in \R \iff z = \conj{z}$.
  \item La conjugaison $z \mapsto \conj{z}$ est un automorphisme de corps de $\C$ (le seul automorphisme continu non trivial).
\end{enumerate}
\end{proposition}

\begin{proof}
Posons $z = a + ib$ et $w = c + id$ avec $a,b,c,d \in \R$.

(i) $\conj{z+w} = \conj{(a+c) + i(b+d)} = (a+c) - i(b+d) = (a-ib) + (c-id) = \conj{z} + \conj{w}$.

Pour le produit : $zw = (ac-bd) + i(ad+bc)$, donc $\conj{zw} = (ac-bd) - i(ad+bc)$. D'autre part, $\conj{z}\,\conj{w} = (a-ib)(c-id) = (ac-bd) - i(ad+bc)$. D'où l'égalité.

(ii) $\conj{\conj{z}} = \conj{a-ib} = a + ib = z$.

(iii) $z + \conj{z} = (a+ib) + (a-ib) = 2a = 2\re(z)$. De même, $z - \conj{z} = 2ib = 2i\,\im(z)$.

(iv) $z\conj{z} = (a+ib)(a-ib) = a^2 + b^2$.

(v) et (vi) sont immédiats.
\end{proof}

\section{Module et argument}\label{sec:module-argument}

\begin{definition}[Module]\label{def:module}
Le \emph{module} de $z = a + ib \in \C$ est
\[
  \abs{z} = \sqrt{a^2 + b^2} = \sqrt{z\conj{z}}.
\]
Géométriquement, $\abs{z}$ est la distance de $z$ à l'origine dans le plan $\R^2$.
\end{definition}

\begin{proposition}[Propriétés du module]\label{prop:proprietes-module}
Pour tous $z, w \in \C$ :
\begin{enumerate}[label=(\roman*)]
  \item $\abs{z} \geq 0$, avec $\abs{z} = 0 \iff z = 0$.
  \item $\abs{\conj{z}} = \abs{z}$.
  \item $\abs{zw} = \abs{z}\abs{w}$ \quad (multiplicativité).
  \item $\abs{z^{-1}} = \abs{z}^{-1}$ pour $z \neq 0$.
  \item $\abs{\re(z)} \leq \abs{z}$ et $\abs{\im(z)} \leq \abs{z}$.
  \item $\abs{z} \leq \abs{\re(z)} + \abs{\im(z)}$.
\end{enumerate}
\end{proposition}

\begin{proof}
Les propriétés (i), (ii), (v) et (vi) découlent directement de la définition.

(iii) $\abs{zw}^2 = (zw)\conj{(zw)} = zw\conj{z}\,\conj{w} = (z\conj{z})(w\conj{w}) = \abs{z}^2\abs{w}^2$.

(iv) Conséquence de (iii) : $\abs{z}\abs{z^{-1}} = \abs{zz^{-1}} = \abs{1} = 1$.
\end{proof}

\begin{theoreme}[Inégalité triangulaire]\label{thm:inegalite-triangulaire}
Pour tous $z, w \in \C$ :
\[
  \abs{z + w} \leq \abs{z} + \abs{w}.
\]
De plus, $\bigl|\abs{z} - \abs{w}\bigr| \leq \abs{z - w}$ (inégalité triangulaire inverse).
\end{theoreme}

\begin{proof}
On calcule :
\begin{align*}
  \abs{z + w}^2 &= (z+w)\conj{(z+w)} = (z+w)(\conj{z}+\conj{w}) \\
  &= z\conj{z} + z\conj{w} + w\conj{z} + w\conj{w} \\
  &= \abs{z}^2 + 2\re(z\conj{w}) + \abs{w}^2.
\end{align*}
Or $\re(z\conj{w}) \leq \abs{z\conj{w}} = \abs{z}\abs{w}$, d'où
\[
  \abs{z+w}^2 \leq \abs{z}^2 + 2\abs{z}\abs{w} + \abs{w}^2 = (\abs{z} + \abs{w})^2.
\]
L'inégalité inverse s'obtient en écrivant $z = (z-w) + w$, d'où $\abs{z} \leq \abs{z-w} + \abs{w}$.
\end{proof}

\begin{definition}[Argument]\label{def:argument}
Soit $z \in \C^* = \C \setminus \{0\}$. Un \emph{argument} de $z$ est un réel $\theta \in \R$ tel que
\[
  z = \abs{z}(\cos\theta + i\sin\theta).
\]
L'ensemble des arguments de $z$ est $\arg(z) = \{\theta_0 + 2k\pi : k \in \Z\}$ pour tout argument particulier $\theta_0$.

L'\emph{argument principal} $\Arg(z) \in (-\pi, \pi]$ est l'unique argument dans cet intervalle.
\end{definition}

\begin{remarque}\label{rem:argument-multivaluation}
La multivaluation de l'argument est un thème récurrent en analyse complexe. Elle est à l'origine des difficultés liées à la définition du logarithme complexe et des puissances complexes. La \og coupure \fg{} le long de l'axe réel négatif, qui définit la branche principale de l'argument, est un choix conventionnel ; d'autres coupures sont parfois plus adaptées.
\end{remarque}

\section{Forme exponentielle et formule d'Euler}\label{sec:forme-exponentielle}

\begin{theoreme}[Formule d'Euler]\label{thm:euler}
Pour tout $\theta \in \R$ :
\[
  e^{i\theta} = \cos\theta + i\sin\theta.
\]
\end{theoreme}

\begin{proof}
Considérons les séries entières convergentes :
\[
  e^{i\theta} = \sum_{n=0}^{\infty} \frac{(i\theta)^n}{n!}
  = \sum_{n=0}^{\infty} \frac{i^n \theta^n}{n!}.
\]
Séparons les termes pairs et impairs. Pour $n = 2k$ : $i^{2k} = (-1)^k$. Pour $n = 2k+1$ : $i^{2k+1} = (-1)^k i$. Donc :
\begin{align*}
  e^{i\theta} &= \sum_{k=0}^{\infty} \frac{(-1)^k \theta^{2k}}{(2k)!} + i\sum_{k=0}^{\infty} \frac{(-1)^k \theta^{2k+1}}{(2k+1)!} = \cos\theta + i\sin\theta. \qedhere
\end{align*}
\end{proof}

\begin{corollaire}[Forme exponentielle]\label{cor:forme-expo}
Tout nombre complexe $z \neq 0$ s'écrit de manière unique sous la forme
\[
  z = r e^{i\theta}, \quad r = \abs{z} > 0,\; \theta = \Arg(z) \in (-\pi,\pi].
\]
\end{corollaire}

\begin{proposition}[Propriétés de l'exponentielle complexe]\label{prop:exp-proprietes}
Pour tous $\theta, \varphi \in \R$ :
\begin{enumerate}[label=(\roman*)]
  \item $e^{i\theta} \cdot e^{i\varphi} = e^{i(\theta + \varphi)}$ \quad (les arguments s'additionnent).
  \item $\conj{e^{i\theta}} = e^{-i\theta}$.
  \item $\abs{e^{i\theta}} = 1$ pour tout $\theta$.
  \item $\cos\theta = \frac{e^{i\theta} + e^{-i\theta}}{2}$, \quad $\sin\theta = \frac{e^{i\theta} - e^{-i\theta}}{2i}$.
\end{enumerate}
\end{proposition}

\begin{proof}
(i) C'est l'identité $(\cos\theta + i\sin\theta)(\cos\varphi + i\sin\varphi) = \cos(\theta+\varphi) + i\sin(\theta+\varphi)$, qui se vérifie par les formules d'addition trigonométriques.

(ii) $\conj{e^{i\theta}} = \cos\theta - i\sin\theta = \cos(-\theta) + i\sin(-\theta) = e^{-i\theta}$.

(iii) $\abs{e^{i\theta}}^2 = e^{i\theta}\conj{e^{i\theta}} = e^{i\theta}e^{-i\theta} = e^0 = 1$.

(iv) Se déduit immédiatement de la formule d'Euler et de (ii).
\end{proof}

\begin{theoreme}[Formule de De Moivre]\label{thm:de-moivre}
Pour tout $\theta \in \R$ et tout $n \in \Z$ :
\[
  (\cos\theta + i\sin\theta)^n = \cos(n\theta) + i\sin(n\theta).
\]
\end{theoreme}

\begin{proof}
C'est une conséquence immédiate de la propriété $(e^{i\theta})^n = e^{in\theta}$, elle-même conséquence de la propriété d'additivité de l'exponentielle.
\end{proof}

\begin{exemple}[Linéarisation trigonométrique]\label{ex:linearisation}
Exprimons $\cos^3\theta$ en fonction de $\cos(k\theta)$. On utilise les formules d'Euler :
\begin{align*}
  \cos^3\theta &= \left(\frac{e^{i\theta}+e^{-i\theta}}{2}\right)^3 = \frac{1}{8}\left(e^{3i\theta} + 3e^{i\theta} + 3e^{-i\theta} + e^{-3i\theta}\right) \\
  &= \frac{1}{8}\left(2\cos 3\theta + 6\cos\theta\right) = \frac{1}{4}\cos 3\theta + \frac{3}{4}\cos\theta.
\end{align*}
\end{exemple}

\section{Racines \texorpdfstring{$n$}{n}-ièmes de l'unité}\label{sec:racines-unite}

\begin{theoreme}[Racines $n$-ièmes]\label{thm:racines-n-iemes}
Soit $w \in \C^*$ et $n \in \N^*$. L'équation $z^n = w$ possède exactement $n$ solutions dans $\C$, données par :
\[
  z_k = \abs{w}^{1/n}\, \exp\!\left(i\,\frac{\Arg(w) + 2k\pi}{n}\right), \quad k = 0, 1, \ldots, n-1.
\]
\end{theoreme}

\begin{proof}
Écrivons $w = \abs{w}e^{i\alpha}$ avec $\alpha = \Arg(w)$, et cherchons $z = re^{i\theta}$ tel que $z^n = w$. On obtient $r^n e^{in\theta} = \abs{w}e^{i\alpha}$. L'identification des modules donne $r = \abs{w}^{1/n}$ (unique réel positif). L'identification des arguments donne $n\theta = \alpha + 2k\pi$, soit $\theta = \frac{\alpha + 2k\pi}{n}$. Pour $k = 0, 1, \ldots, n-1$, on obtient $n$ valeurs distinctes de $e^{i\theta}$ ; pour $k = n$, on retrouve la valeur $k=0$.
\end{proof}

\begin{definition}[Racines $n$-ièmes de l'unité]\label{def:racines-unite}
Les \emph{racines $n$-ièmes de l'unité} sont les solutions de $z^n = 1$ :
\[
  \omega_k = e^{2ik\pi/n}, \quad k = 0, 1, \ldots, n-1.
\]
On note $\omega = \omega_1 = e^{2i\pi/n}$ la racine primitive canonique. Alors $\omega_k = \omega^k$.
\end{definition}

\begin{proposition}[Propriétés des racines de l'unité]\label{prop:racines-unite-proprietes}
Soit $\mathbb{U}_n = \{\omega^k : k = 0, \ldots, n-1\}$ l'ensemble des racines $n$-ièmes de l'unité.
\begin{enumerate}[label=(\roman*)]
  \item $\mathbb{U}_n$ est un sous-groupe cyclique de $(\C^*, \cdot)$, engendré par $\omega$.
  \item $\sum_{k=0}^{n-1} \omega^k = 0$.
  \item $\prod_{k=0}^{n-1} \omega^k = (-1)^{n+1}$.
  \item $z^n - 1 = \prod_{k=0}^{n-1}(z - \omega^k)$.
  \item Les racines $n$-ièmes de l'unité forment les sommets d'un polygone régulier à $n$ côtés inscrit dans le cercle unité.
\end{enumerate}
\end{proposition}

\begin{proof}
(i) $\omega^j \cdot \omega^k = \omega^{j+k} = \omega^{(j+k) \bmod n} \in \mathbb{U}_n$, et $(\omega^k)^{-1} = \omega^{-k} = \omega^{n-k} \in \mathbb{U}_n$.

(ii) La somme $S = \sum_{k=0}^{n-1} \omega^k$ est une série géométrique. Comme $\omega \neq 1$ (pour $n \geq 2$), on a $S = \frac{\omega^n - 1}{\omega - 1} = \frac{1-1}{\omega-1} = 0$.

(iii) $\prod_{k=0}^{n-1}\omega^k = \omega^{0+1+\cdots+(n-1)} = \omega^{n(n-1)/2} = e^{i\pi(n-1)} = (-1)^{n-1}= (-1)^{n+1}$.

(iv) Le polynôme $z^n - 1$ est de degré $n$ et admet les $n$ racines $\omega^k$ ; d'où la factorisation.

(v) Les points $\omega^k = e^{2ik\pi/n}$ sont régulièrement espacés de $2\pi/n$ sur le cercle unité.
\end{proof}

\begin{center}
\begin{tikzpicture}[scale=2.2]
  % Axes
  \draw[gray!40,->] (-1.4,0) -- (1.4,0) node[right,black] {$\re$};
  \draw[gray!40,->] (0,-1.4) -- (0,1.4) node[above,black] {$\im$};

  % Cercle unité
  \draw[blue!50!black, thick] (0,0) circle (1);

  % Racines 7-ièmes
  \foreach \k in {0,...,6} {
    \pgfmathsetmacro{\angle}{360*\k/7}
    \pgfmathsetmacro{\nexta}{360*(\k+1)/7}
    \coordinate (w\k) at (\angle:1);
    \fill[red!70!black] (w\k) circle (1.5pt);
    \draw[red!70!black, dashed, thin] (0,0) -- (w\k);
    \node[red!70!black] at (\angle:1.2) {\small $\omega^{\k}$};
  }
  % Polygone
  \draw[orange!70!black, thick] (0:1) \foreach \k in {1,...,6} { -- ({360*\k/7}:1) } -- cycle;

  \node[below left] at (0,0) {\small $0$};
  \node[below] at (0,-1.5) {Racines $7$-ièmes de l'unité};
\end{tikzpicture}
\end{center}

\begin{exemple}[Racines cubiques de $-8$]\label{ex:racines-cubiques}
Résolvons $z^3 = -8$. On a $\abs{w} = 8$ et $\Arg(w) = \pi$, donc
\[
  z_k = 2\exp\!\left(i\,\frac{\pi + 2k\pi}{3}\right), \quad k = 0,1,2.
\]
\begin{itemize}
  \item $z_0 = 2e^{i\pi/3} = 2\left(\frac{1}{2} + i\frac{\sqrt{3}}{2}\right) = 1 + i\sqrt{3}$.
  \item $z_1 = 2e^{i\pi} = -2$.
  \item $z_2 = 2e^{i5\pi/3} = 2\left(\frac{1}{2} - i\frac{\sqrt{3}}{2}\right) = 1 - i\sqrt{3}$.
\end{itemize}
Vérifions : $(1+i\sqrt{3})^3 = (1+i\sqrt{3})(1+i\sqrt{3})^2 = (1+i\sqrt{3})(1+2i\sqrt{3}-3) = (1+i\sqrt{3})(-2+2i\sqrt{3}) = -2+2i\sqrt{3}-2i\sqrt{3}+2i^2 \cdot 3 = -2-6 = -8$.
\end{exemple}

\section{Inversion et transformations de Möbius}\label{sec:mobius}

\begin{definition}[Inversion complexe]\label{def:inversion}
L'\emph{inversion complexe} est l'application $z \mapsto 1/z$, définie sur $\C^*$. En forme polaire, si $z = re^{i\theta}$, alors $1/z = r^{-1}e^{-i\theta}$ : le module est inversé et l'argument est changé en son opposé.
\end{definition}

\begin{proposition}[Image de droites et cercles par l'inversion]\label{prop:inversion-images}
L'inversion $z \mapsto 1/z$ transforme :
\begin{enumerate}[label=(\roman*)]
  \item un cercle ne passant pas par l'origine en un cercle ne passant pas par l'origine ;
  \item un cercle passant par l'origine en une droite ne passant pas par l'origine ;
  \item une droite ne passant pas par l'origine en un cercle passant par l'origine ;
  \item une droite passant pas par l'origine en une droite passant par l'origine.
\end{enumerate}
\end{proposition}

\begin{proof}
L'équation générale d'un cercle ou d'une droite dans $\C$ s'écrit
\[
  \alpha z\conj{z} + \conj{\beta} z + \beta\conj{z} + \gamma = 0,
\]
avec $\alpha, \gamma \in \R$ et $\beta \in \C$. C'est un cercle si $\alpha \neq 0$, une droite si $\alpha = 0$.

Si l'on pose $w = 1/z$, soit $z = 1/w$, l'équation devient
\[
  \alpha \frac{1}{w\conj{w}} + \conj{\beta}\frac{1}{w} + \beta\frac{1}{\conj{w}} + \gamma = 0.
\]
En multipliant par $w\conj{w}$ :
\[
  \alpha + \conj{\beta}\conj{w} + \beta w + \gamma w\conj{w} = 0,
\]
soit $\gamma w\conj{w} + \beta w + \conj{\beta}\conj{w} + \alpha = 0$. On obtient une équation de même type, avec $\alpha$ et $\gamma$ échangés. Les quatre cas s'en déduisent.
\end{proof}

\begin{definition}[Transformation de Möbius]\label{def:mobius}
Une \emph{transformation de Möbius} (ou \emph{homographie}) est une application de la forme
\[
  T(z) = \frac{az + b}{cz + d}, \quad a,b,c,d \in \C,\; ad - bc \neq 0.
\]
La condition $ad - bc \neq 0$ assure que $T$ n'est pas constante.
\end{definition}

\begin{proposition}[Propriétés des transformations de Möbius]\label{prop:mobius-proprietes}
\begin{enumerate}[label=(\roman*)]
  \item L'ensemble des transformations de Möbius forme un groupe pour la composition, isomorphe à $\mathrm{PGL}_2(\C)$.
  \item Toute transformation de Möbius est une composée de translations ($z \mapsto z + b$), rotations-homothéties ($z \mapsto az$, $a \neq 0$), et de l'inversion ($z \mapsto 1/z$).
  \item Les transformations de Möbius envoient cercles et droites sur cercles et droites.
  \item Elles sont conformes (préservent les angles).
\end{enumerate}
\end{proposition}

\begin{proof}
(i) Si $T_1(z) = \frac{a_1 z + b_1}{c_1 z + d_1}$ et $T_2(z) = \frac{a_2 z + b_2}{c_2 z + d_2}$, alors $T_1 \circ T_2$ correspond au produit des matrices $\begin{psmallmatrix}a_1 & b_1 \\ c_1 & d_1\end{psmallmatrix}\begin{psmallmatrix}a_2 & b_2 \\ c_2 & d_2\end{psmallmatrix}$, dont le déterminant est $(a_1d_1 - b_1c_1)(a_2d_2 - b_2c_2) \neq 0$.

(ii) Si $c = 0$, alors $T(z) = \frac{a}{d}z + \frac{b}{d}$. Si $c \neq 0$, on effectue la division euclidienne :
\[
  T(z) = \frac{a}{c} - \frac{ad - bc}{c^2}\cdot\frac{1}{z + d/c} = \frac{a}{c} + \frac{bc-ad}{c^2}\cdot\frac{1}{z+d/c}.
\]
C'est bien une composée de translation, inversion, homothétie et translation.

(iii) Découle de (ii) et de la 
ef{prop:inversion-images}.

(iv) Chaque opération élémentaire (translation, rotation-homothétie, inversion) est conforme, donc leur composée l'est aussi.
\end{proof}

\begin{exemple}[Transformation de Cayley]\label{ex:cayley}
La transformation de Cayley $T(z) = \frac{z - i}{z + i}$ envoie le demi-plan supérieur $\{z \in \C : \im(z) > 0\}$ sur le disque unité $\D$. En effet, si $z = x + iy$ avec $y > 0$ :
\[
  \abs{T(z)}^2 = \frac{\abs{z-i}^2}{\abs{z+i}^2} = \frac{x^2 + (y-1)^2}{x^2 + (y+1)^2} = \frac{x^2+y^2-2y+1}{x^2+y^2+2y+1}.
\]
Le numérateur est strictement inférieur au dénominateur (leur différence est $-4y < 0$), donc $\abs{T(z)} < 1$.

La transformation inverse est $T^{-1}(w) = i\frac{1+w}{1-w}$.
\end{exemple}

\section{Topologie de \texorpdfstring{$\C$}{C}}\label{sec:topologie-C}

On identifie $\C$ à $\R^2$ muni de la distance $d(z,w) = \abs{z - w}$.

\begin{definition}[Disques et cercles]\label{def:disques}
Soient $a \in \C$ et $r > 0$.
\begin{enumerate}[label=(\roman*)]
  \item Le \emph{disque ouvert} de centre $a$ et de rayon $r$ est $D(a,r) = \{z \in \C : \abs{z-a} < r\}$.
  \item Le \emph{disque fermé} est $\conj{D}(a,r) = \{z \in \C : \abs{z-a} \leq r\}$.
  \item Le \emph{cercle} est $C(a,r) = \{z \in \C : \abs{z-a} = r\}$.
\end{enumerate}
\end{definition}

\begin{definition}[Ouverts, fermés, bord]\label{def:ouverts-fermes}
Soit $A \subseteq \C$.
\begin{enumerate}[label=(\roman*)]
  \item $A$ est \emph{ouvert} si pour tout $z \in A$, il existe $r > 0$ tel que $D(z,r) \subseteq A$.
  \item $A$ est \emph{fermé} si $\C \setminus A$ est ouvert.
  \item Le \emph{bord} de $A$ est $\partial A = \conj{A} \setminus \mathring{A}$.
  \item $A$ est \emph{borné} si $A \subseteq D(0,R)$ pour un certain $R > 0$.
  \item $A$ est \emph{compact} s'il est fermé et borné (théorème de Heine-Borel).
\end{enumerate}
\end{definition}

\begin{definition}[Connexité]\label{def:connexite}
Un ouvert $\Omega \subseteq \C$ est \emph{connexe} si l'une des conditions équivalentes suivantes est satisfaite :
\begin{enumerate}[label=(\roman*)]
  \item $\Omega$ ne peut pas s'écrire comme réunion disjointe de deux ouverts non vides.
  \item Pour tous $z, w \in \Omega$, il existe un chemin continu $\gamma : [0,1] \to \Omega$ tel que $\gamma(0) = z$ et $\gamma(1) = w$.
\end{enumerate}
Un ouvert connexe est appelé un \emph{domaine}.
\end{definition}

\begin{remarque}\label{rem:connexite-ouverts}
Pour les ouverts de $\C$, connexité et connexité par arcs sont équivalentes. On peut même imposer que les chemins soient des lignes polygonales. En revanche, pour des parties quelconques de $\C$, la connexité par arcs est strictement plus forte que la connexité.
\end{remarque}

\begin{definition}[Connexité simple]\label{def:simplement-connexe}
Un domaine $\Omega \subseteq \C$ est \emph{simplement connexe} si tout lacet continu $\gamma : [0,1] \to \Omega$ (avec $\gamma(0) = \gamma(1)$) est homotope à un point dans $\Omega$. Intuitivement, $\Omega$ est simplement connexe s'il n'a \og pas de trous\fg{}.
\end{definition}

\begin{exemple}[Domaines importants]\label{ex:domaines}
\begin{enumerate}[label=(\roman*)]
  \item $\C$ est un domaine simplement connexe.
  \item $\C^* = \C \setminus \{0\}$ est un domaine \emph{non} simplement connexe.
  \item $D(a,r)$ est un domaine simplement connexe.
  \item La couronne $\{z \in \C : r_1 < \abs{z-a} < r_2\}$ est un domaine non simplement connexe.
  \item Le demi-plan $\{z \in \C : \re(z) > 0\}$ est simplement connexe.
\end{enumerate}
\end{exemple}

\section{Le plan complexe étendu et la sphère de Riemann}\label{sec:sphere-riemann}

\begin{definition}[Plan complexe étendu]\label{def:plan-etendu}
Le \emph{plan complexe étendu} est $\hat{\C} = \C \cup \{\infty\}$, où $\infty$ est un point formel. On munit $\hat{\C}$ de la topologie suivante :
\begin{itemize}
  \item Les ouverts de $\C$ restent ouverts dans $\hat{\C}$.
  \item Les voisinages de $\infty$ sont les ensembles $\{\infty\} \cup (\C \setminus K)$ où $K \subseteq \C$ est compact.
\end{itemize}
L'espace $\hat{\C}$ est compact et homéomorphe à la sphère $S^2$.
\end{definition}

\begin{definition}[Projection stéréographique]\label{def:projection-stereo}
La \emph{projection stéréographique} est l'homéomorphisme $\Phi : S^2 \to \hat{\C}$ défini comme suit. On identifie $S^2 = \{(X,Y,Z) \in \R^3 : X^2 + Y^2 + Z^2 = 1\}$ et on projette depuis le pôle Nord $N = (0,0,1)$ sur le plan équatorial $\{Z = 0\} \simeq \C$.

Pour $(X,Y,Z) \in S^2 \setminus \{N\}$, la droite passant par $N$ et $(X,Y,Z)$ coupe le plan $Z=0$ en
\[
  \Phi(X,Y,Z) = \frac{X}{1-Z} + i\frac{Y}{1-Z}.
\]
On pose $\Phi(N) = \infty$. L'application inverse est
\[
  \Phi^{-1}(z) = \left(\frac{2\re(z)}{\abs{z}^2+1},\; \frac{2\im(z)}{\abs{z}^2+1},\; \frac{\abs{z}^2-1}{\abs{z}^2+1}\right).
\]
\end{definition}

\begin{proof}[Vérification de la formule inverse]
Soit $z = x + iy$ et posons $r^2 = x^2 + y^2 = \abs{z}^2$. Le point $(X,Y,Z)$ sur $S^2$ est paramétré par $(X,Y,Z) = (1-t)N + t\,(x,y,0) = (tx, ty, 1-t)$ pour un $t > 0$. La condition $X^2+Y^2+Z^2 = 1$ donne $t^2r^2 + (1-t)^2 = 1$, soit $t^2 r^2 + 1 - 2t + t^2 = 1$, d'où $t^2(r^2+1) = 2t$, soit $t = \frac{2}{r^2+1}$. On obtient :
\[
  X = \frac{2x}{r^2+1},\quad Y = \frac{2y}{r^2+1},\quad Z = 1 - \frac{2}{r^2+1} = \frac{r^2-1}{r^2+1}. \qedhere
\]
\end{proof}

\begin{center}
\begin{tikzpicture}[scale=1.6]
  % Sphère (ellipse)
  \draw[thick,blue!60!black] (0,0) ellipse (2 and 0.6);
  \draw[thick,blue!60!black] (0,0) ellipse (2 and 2);

  % Pôle Nord
  \fill[red!70!black] (0,2) circle (2pt) node[above right] {$N = \infty$};

  % Plan
  \draw[gray,thick] (-3,0) -- (3,0) node[right] {$\C$};

  % Point sur la sphère
  \coordinate (P) at ({2*cos(50)},{2*sin(50)});
  \fill[green!50!black] (P) circle (2pt) node[above right] {$(X,Y,Z)$};

  % Projection
  \pgfmathsetmacro{\projx}{2*cos(50)/(1-sin(50))*0.6}
  \fill[orange!70!black] (\projx,0) circle (2pt) node[below] {$z = \Phi(X,Y,Z)$};

  % Ligne de projection
  \draw[dashed,red!60!black] (0,2) -- (\projx,0);
  \draw[dashed,red!60!black,thin] (0,2) -- (P);

  \node at (0,-1.2) {Projection stéréographique};
\end{tikzpicture}
\end{center}

\begin{proposition}[Propriétés de la projection stéréographique]\label{prop:stereo-proprietes}
\begin{enumerate}[label=(\roman*)]
  \item $\Phi$ est un homéomorphisme de $S^2$ sur $\hat{\C}$.
  \item $\Phi$ est conforme : elle préserve les angles.
  \item $\Phi$ envoie les cercles de $S^2$ sur les cercles et droites de $\hat{\C}$ (les droites correspondant aux cercles passant par $N$).
\end{enumerate}
\end{proposition}

\begin{definition}[Distance chordale]\label{def:distance-chordale}
La \emph{distance chordale} sur $\hat{\C}$ est définie par
\[
  \chi(z,w) = \frac{2\abs{z-w}}{\sqrt{1+\abs{z}^2}\sqrt{1+\abs{w}^2}},
  \quad \chi(z,\infty) = \frac{2}{\sqrt{1+\abs{z}^2}}.
\]
Elle correspond à la distance euclidienne entre les points correspondants sur $S^2$.
\end{definition}

\section{Exercices}\label{sec:exo-ch1}

\begin{exercice}[$\star$]\label{exo:module-somme}
Montrer que pour $z_1, z_2 \in \C$ :
\[
  \abs{z_1 + z_2}^2 + \abs{z_1 - z_2}^2 = 2(\abs{z_1}^2 + \abs{z_2}^2).
\]
Interpréter géométriquement ce résultat (identité du parallélogramme).
\end{exercice}

\begin{exercice}[$\star$]\label{exo:inegalite-module}
Soient $z_1, \ldots, z_n \in \C$. Montrer l'inégalité triangulaire généralisée :
\[
  \abs{z_1 + z_2 + \cdots + z_n} \leq \abs{z_1} + \abs{z_2} + \cdots + \abs{z_n}.
\]
À quelle condition a-t-on égalité ?
\end{exercice}

\begin{exercice}[$\star$]\label{exo:lieux-geometriques}
Décrire les ensembles suivants et les représenter dans le plan complexe :
\begin{enumerate}[label=(\alph*)]
  \item $\{z \in \C : \abs{z - 1} = \abs{z + 1}\}$.
  \item $\{z \in \C : \abs{z - 1} + \abs{z + 1} = 4\}$.
  \item $\{z \in \C : \re(z^2) = 1\}$.
  \item $\{z \in \C : \abs{z - 1} = 2\abs{z + 1}\}$.
\end{enumerate}
\end{exercice}

\begin{exercice}[$\star\star$]\label{exo:racines-explicites}
\begin{enumerate}[label=(\alph*)]
  \item Calculer les racines quatrièmes de $-4$.
  \item Calculer les racines cinquièmes de $1+i$.
  \item Résoudre $z^6 + z^3 + 1 = 0$.
\end{enumerate}
\end{exercice}

\begin{exercice}[$\star\star$]\label{exo:somme-racines}
Soit $n \geq 2$ un entier et $\omega = e^{2i\pi/n}$.
\begin{enumerate}[label=(\alph*)]
  \item Calculer $\sum_{k=0}^{n-1} \cos\!\left(\frac{2k\pi}{n}\right)$ et $\sum_{k=0}^{n-1} \sin\!\left(\frac{2k\pi}{n}\right)$.
  \item Calculer $\prod_{k=1}^{n-1} \sin\!\left(\frac{k\pi}{n}\right)$. (\emph{Indication} : utiliser la factorisation de $z^n - 1$ et évaluer en $z = 1$ après simplification.)
  \item Montrer que $\sum_{k=0}^{n-1} \abs{1 - \omega^k}^2 = 2n$.
\end{enumerate}
\end{exercice}

\begin{exercice}[$\star\star$]\label{exo:mobius-disque}
Soit $a \in \D$ (c'est-à-dire $\abs{a} < 1$). Montrer que la transformation de Möbius
\[
  \varphi_a(z) = \frac{z - a}{1 - \conj{a}z}
\]
est un automorphisme du disque unité $\D$, c'est-à-dire une bijection holomorphe de $\D$ sur $\D$. Vérifier que $\varphi_a(\varphi_a(z)) = z$ (involution).
\end{exercice}

\begin{exercice}[$\star\star$]\label{exo:apollonius}
Soient $a, b \in \C$ distincts et $k > 0$ avec $k \neq 1$. Montrer que l'ensemble
\[
  \left\{z \in \C : \frac{\abs{z - a}}{\abs{z - b}} = k\right\}
\]
est un cercle (cercle d'Apollonius). Déterminer son centre et son rayon.
\end{exercice}

\begin{exercice}[$\star\star\star$]\label{exo:ptolemee}
Soient $z_1, z_2, z_3, z_4 \in \C$. Démontrer l'inégalité de Ptolémée :
\[
  \abs{z_1 - z_3}\abs{z_2 - z_4} \leq \abs{z_1 - z_2}\abs{z_3 - z_4} + \abs{z_1 - z_4}\abs{z_2 - z_3}.
\]
(\emph{Indication} : considérer le birapport et utiliser l'inégalité triangulaire.)
\end{exercice}

\begin{exercice}[$\star\star\star$]\label{exo:stereo-cercles}
Montrer que la projection stéréographique envoie les cercles de $S^2$ sur les cercles et les droites de $\C$. Plus précisément, un cercle sur $S^2$ est l'intersection de $S^2$ avec un plan $\alpha X + \beta Y + \gamma Z = \delta$ (avec $\alpha, \beta, \gamma, \delta \in \R$). Exprimer l'image par $\Phi$ en fonction de $\alpha, \beta, \gamma, \delta$ et montrer que c'est un cercle si $\gamma \neq \delta$, une droite si $\gamma = \delta$.
\end{exercice}

\begin{exercice}[$\star\star\star$]\label{exo:mobius-triple-transitivite}
Montrer que le groupe des transformations de Möbius agit triplement transitivement sur $\hat{\C}$ : pour tous triplets $(z_1, z_2, z_3)$ et $(w_1, w_2, w_3)$ de points distincts de $\hat{\C}$, il existe une unique transformation de Möbius $T$ telle que $T(z_k) = w_k$ pour $k = 1,2,3$.

\emph{Indication :} Construire d'abord la transformation $T$ envoyant $(z_1, z_2, z_3)$ sur $(0, 1, \infty)$ à l'aide du birapport :
\[
  T(z) = \frac{(z - z_1)(z_2 - z_3)}{(z - z_3)(z_2 - z_1)}.
\]
\end{exercice}

% ═══════════════════════════════════════════════════════════════════════════════
% CHAPITRE 2
% ═══════════════════════════════════════════════════════════════════════════════
\chapter{Fonctions holomorphes --- Conditions de Cauchy-Riemann}\label{ch:holomorphes}

\section{Dérivée complexe}\label{sec:derivee-complexe}

La notion centrale de l'analyse complexe est celle de dérivabilité complexe. Bien que sa définition soit formellement identique à celle de la dérivée réelle, ses conséquences sont incomparablement plus riches.

\begin{definition}[Dérivabilité complexe]\label{def:derivee-complexe}
Soit $\Omega \subseteq \C$ un ouvert et $f : \Omega \to \C$. On dit que $f$ est \emph{dérivable au sens complexe} (ou $\C$-dérivable) en $z_0 \in \Omega$ si la limite
\[
  f'(z_0) = \lim_{h \to 0} \frac{f(z_0 + h) - f(z_0)}{h}
\]
existe dans $\C$, où $h \in \C^*$ tend vers $0$ de manière quelconque. On appelle $f'(z_0)$ la \emph{dérivée complexe} de $f$ en $z_0$.
\end{definition}

\begin{remarque}\label{rem:derivee-directionnelle}
La condition cruciale est que la limite doit être la même quelle que soit la direction dans laquelle $h \to 0$. En analyse réelle, $h$ ne peut tendre vers $0$ que par valeurs positives ou négatives (deux directions). En analyse complexe, $h$ peut tendre vers $0$ dans une infinité de directions --- le long de l'axe réel, de l'axe imaginaire, le long de spirales, etc. Cette contrainte beaucoup plus forte est la source de la rigidité de l'analyse complexe.
\end{remarque}

\begin{definition}[Fonction holomorphe]\label{def:holomorphe}
Soit $\Omega \subseteq \C$ un ouvert.
\begin{enumerate}[label=(\roman*)]
  \item $f$ est \emph{holomorphe en} $z_0 \in \Omega$ si $f$ est $\C$-dérivable en $z_0$.
  \item $f$ est \emph{holomorphe sur} $\Omega$ si $f$ est holomorphe en tout point de $\Omega$.
  \item On note $\mathscr{O}(\Omega)$ l'ensemble des fonctions holomorphes sur $\Omega$.
  \item $f$ est \emph{entière} si elle est holomorphe sur $\C$ tout entier.
\end{enumerate}
\end{definition}

\begin{proposition}[Algèbre des fonctions holomorphes]\label{prop:algebre-holomorphe}
Soient $f, g \in \mathscr{O}(\Omega)$ et $\lambda \in \C$.
\begin{enumerate}[label=(\roman*)]
  \item $\lambda f + g \in \mathscr{O}(\Omega)$ et $(\lambda f + g)' = \lambda f' + g'$.
  \item $fg \in \mathscr{O}(\Omega)$ et $(fg)' = f'g + fg'$ \quad (règle de Leibniz).
  \item Si $g(z_0) \neq 0$, alors $f/g$ est holomorphe au voisinage de $z_0$ et $(f/g)' = \frac{f'g - fg'}{g^2}$.
  \item Si $g : \Omega \to \Omega'$ et $f : \Omega' \to \C$ sont holomorphes, alors $f \circ g \in \mathscr{O}(\Omega)$ et $(f \circ g)' = (f' \circ g) \cdot g'$ \quad (règle de la chaîne).
\end{enumerate}
En particulier, $\mathscr{O}(\Omega)$ est une $\C$-algèbre.
\end{proposition}

\begin{proof}
Les preuves sont identiques à celles du cas réel. Montrons par exemple la règle de Leibniz. On a
\begin{align*}
  \frac{(fg)(z_0+h)-(fg)(z_0)}{h} &= \frac{f(z_0+h)g(z_0+h) - f(z_0)g(z_0)}{h} \\
  &= f(z_0+h)\frac{g(z_0+h)-g(z_0)}{h} + g(z_0)\frac{f(z_0+h)-f(z_0)}{h}.
\end{align*}
Comme $f$ est dérivable en $z_0$, elle est continue en $z_0$, donc $f(z_0+h) \to f(z_0)$. En passant à la limite :
\[
  (fg)'(z_0) = f(z_0)g'(z_0) + g(z_0)f'(z_0). \qedhere
\]
\end{proof}

\begin{exemple}[Fonctions holomorphes élémentaires]\label{ex:holomorphes-elementaires}
\begin{enumerate}[label=(\roman*)]
  \item \textbf{Polynômes.} Tout polynôme $P(z) = a_n z^n + \cdots + a_1 z + a_0$ est entier, et $P'(z) = na_n z^{n-1} + \cdots + a_1$. En effet, $(z^n)' = nz^{n-1}$ se démontre par récurrence à l'aide de la règle de Leibniz.

  \item \textbf{Fractions rationnelles.} Si $P$ et $Q$ sont des polynômes, $f(z) = P(z)/Q(z)$ est holomorphe sur $\C \setminus \{z : Q(z) = 0\}$.

  \item \textbf{Exponentielle.} $f(z) = e^z = e^{x+iy} = e^x(\cos y + i\sin y)$. On vérifie directement :
  \[
    \frac{e^{z_0+h}-e^{z_0}}{h} = e^{z_0}\frac{e^h - 1}{h} \to e^{z_0} \quad \text{quand } h \to 0,
  \]
  car $\frac{e^h-1}{h} \to 1$ (par le développement $e^h = 1 + h + h^2/2 + \cdots$). Donc $f'(z) = e^z$.

  \item \textbf{Fonctions trigonométriques.} $\cos z = \frac{e^{iz}+e^{-iz}}{2}$ et $\sin z = \frac{e^{iz}-e^{-iz}}{2i}$ sont entières, avec $(\cos z)' = -\sin z$ et $(\sin z)' = \cos z$.
\end{enumerate}
\end{exemple}

\begin{exemple}[Fonctions non holomorphes]\label{ex:non-holomorphes}
\begin{enumerate}[label=(\roman*)]
  \item \textbf{Conjugaison.} $f(z) = \conj{z}$ n'est pas holomorphe. En effet,
  \[
    \frac{f(z_0+h)-f(z_0)}{h} = \frac{\conj{h}}{h}.
  \]
  Pour $h = t \in \R^*$, ce quotient vaut $1$. Pour $h = it$ avec $t \in \R^*$, ce quotient vaut $\frac{-it}{it} = -1$. Les limites sont différentes, donc la dérivée n'existe pas.

  \item \textbf{Module.} $f(z) = \abs{z}$ n'est pas holomorphe (sauf peut-être en $z = 0$, mais un calcul montre que la dérivée n'existe pas non plus en $0$).

  \item \textbf{Partie réelle.} $f(z) = \re(z)$ n'est pas holomorphe. En effet, $\re(z) = \frac{z + \conj{z}}{2}$, et si $\re$ était holomorphe, $\conj{z} = 2\re(z) - z$ le serait aussi, ce qui est faux.
\end{enumerate}
\end{exemple}

\begin{remarque}\label{rem:holomorphe-vs-derivable}
Les exemples ci-dessus illustrent un phénomène fondamental : les fonctions $z \mapsto \conj{z}$, $z \mapsto \re(z)$, $z \mapsto \abs{z}^2$ sont toutes de classe $\mathscr{C}^\infty$ en tant que fonctions de $\R^2$ dans $\R^2$. Pourtant, elles ne sont pas holomorphes. La $\C$-dérivabilité est une condition strictement plus forte que la $\R$-différentiabilité.
\end{remarque}

\section{Conditions de Cauchy-Riemann}\label{sec:cauchy-riemann}

Les conditions de Cauchy-Riemann caractérisent les fonctions holomorphes en termes de leurs parties réelle et imaginaire.

\begin{theoreme}[Conditions de Cauchy-Riemann --- condition nécessaire]\label{thm:cauchy-riemann}
Soit $\Omega \subseteq \C$ un ouvert et $f : \Omega \to \C$ une fonction holomorphe en $z_0 = x_0 + iy_0 \in \Omega$. Écrivons $f(z) = u(x,y) + iv(x,y)$ où $u, v : \R^2 \to \R$. Alors les dérivées partielles de $u$ et $v$ existent en $(x_0,y_0)$ et satisfont les \emph{équations de Cauchy-Riemann} :
\begin{equation}\label{eq:cauchy-riemann}
  \frac{\partial u}{\partial x}(x_0,y_0) = \frac{\partial v}{\partial y}(x_0,y_0), \qquad
  \frac{\partial u}{\partial y}(x_0,y_0) = -\frac{\partial v}{\partial x}(x_0,y_0).
\end{equation}
De plus, la dérivée complexe s'exprime par
\begin{equation}\label{eq:derivee-via-CR}
  f'(z_0) = \frac{\partial u}{\partial x}(x_0,y_0) + i\frac{\partial v}{\partial x}(x_0,y_0) = \frac{\partial v}{\partial y}(x_0,y_0) - i\frac{\partial u}{\partial y}(x_0,y_0).
\end{equation}
\end{theoreme}

\begin{proof}
Puisque la dérivée complexe existe, la limite $\frac{f(z_0+h)-f(z_0)}{h}$ est la même quel que soit le chemin emprunté par $h \to 0$.

\emph{Étape 1 : $h$ réel.} Prenons $h = t \in \R^*$, $t \to 0$. Alors
\begin{align*}
  f'(z_0) &= \lim_{t \to 0} \frac{f(x_0+t, y_0) - f(x_0,y_0)}{t} \\
  &= \lim_{t \to 0} \frac{u(x_0+t,y_0)-u(x_0,y_0)}{t} + i\lim_{t \to 0}\frac{v(x_0+t,y_0)-v(x_0,y_0)}{t} \\
  &= \frac{\partial u}{\partial x}(x_0,y_0) + i\frac{\partial v}{\partial x}(x_0,y_0).
\end{align*}

\emph{Étape 2 : $h$ imaginaire pur.} Prenons $h = it$ avec $t \in \R^*$, $t \to 0$. Alors
\begin{align*}
  f'(z_0) &= \lim_{t \to 0} \frac{f(x_0, y_0+t) - f(x_0,y_0)}{it} \\
  &= \frac{1}{i}\left[\frac{\partial u}{\partial y}(x_0,y_0) + i\frac{\partial v}{\partial y}(x_0,y_0)\right] \\
  &= \frac{\partial v}{\partial y}(x_0,y_0) - i\frac{\partial u}{\partial y}(x_0,y_0),
\end{align*}
où l'on a utilisé $1/i = -i$.

\emph{Étape 3 : Identification.} En identifiant les parties réelles et imaginaires des deux expressions de $f'(z_0)$ :
\[
  \frac{\partial u}{\partial x} = \frac{\partial v}{\partial y} \quad\text{et}\quad \frac{\partial v}{\partial x} = -\frac{\partial u}{\partial y}. \qedhere
\]
\end{proof}

\begin{remarque}\label{rem:CR-notation-compacte}
En notation compacte, les équations de Cauchy-Riemann s'écrivent $u_x = v_y$ et $u_y = -v_x$, ou encore sous forme matricielle :
\[
  \begin{pmatrix} u_x & u_y \\ v_x & v_y \end{pmatrix} = \begin{pmatrix} u_x & -v_x \\ v_x & u_x \end{pmatrix}.
\]
La matrice jacobienne de $f$ vue comme application $\R^2 \to \R^2$ est donc de la forme $\begin{psmallmatrix} a & -b \\ b & a\end{psmallmatrix}$, qui représente la composition d'une rotation et d'une homothétie. C'est la clé de l'interprétation géométrique de l'holomorphie.
\end{remarque}

\begin{theoreme}[Conditions de Cauchy-Riemann --- condition suffisante]\label{thm:CR-suffisant}
Soit $\Omega \subseteq \C$ un ouvert et $f = u + iv : \Omega \to \C$. Si $u$ et $v$ sont de classe $\mathscr{C}^1$ sur $\Omega$ et satisfont les équations de Cauchy-Riemann \eqref{eq:cauchy-riemann} en tout point de $\Omega$, alors $f$ est holomorphe sur $\Omega$.
\end{theoreme}

\begin{proof}
Soit $z_0 = x_0 + iy_0 \in \Omega$. Comme $u$ et $v$ sont $\mathscr{C}^1$, elles sont différentiables au sens réel en $(x_0,y_0)$ : il existe des fonctions $\varepsilon_1, \varepsilon_2$ tendant vers $0$ en $(0,0)$ telles que
\begin{align*}
  u(x_0+s, y_0+t) - u(x_0,y_0) &= u_x s + u_y t + \abs{(s,t)}\varepsilon_1(s,t), \\
  v(x_0+s, y_0+t) - v(x_0,y_0) &= v_x s + v_y t + \abs{(s,t)}\varepsilon_2(s,t),
\end{align*}
où les dérivées partielles sont évaluées en $(x_0,y_0)$. Posons $h = s + it$. Alors :
\begin{align*}
  f(z_0+h) - f(z_0) &= (u_x s + u_y t) + i(v_x s + v_y t) + \abs{h}\varepsilon(h),
\end{align*}
où $\varepsilon(h) \to 0$ quand $h \to 0$. En utilisant les équations de Cauchy-Riemann ($u_x = v_y$, $u_y = -v_x$) :
\begin{align*}
  (u_x s + u_y t) + i(v_x s + v_y t) &= u_x s - v_x t + i(v_x s + u_x t) \\
  &= (u_x + iv_x)(s + it) \\
  &= (u_x + iv_x)\,h.
\end{align*}
Donc
\[
  \frac{f(z_0+h)-f(z_0)}{h} = u_x + iv_x + \frac{\abs{h}}{h}\varepsilon(h).
\]
Or $\abs{\abs{h}/h} = 1$, donc $\frac{\abs{h}}{h}\varepsilon(h) \to 0$ quand $h \to 0$. Ainsi $f'(z_0) = u_x(x_0,y_0) + iv_x(x_0,y_0)$ existe.
\end{proof}

\begin{remarque}\label{rem:CR-C1-necessaire}
L'hypothèse $\mathscr{C}^1$ dans le théorème précédent peut être affaiblie : il suffit que les dérivées partielles existent et soient continues en $(x_0, y_0)$ (et non sur tout $\Omega$). Il existe même un résultat plus fin dû à Looman et Menchov : si $f$ est continue et si les équations de Cauchy-Riemann sont satisfaites partout (avec existence des dérivées partielles), alors $f$ est holomorphe.
\end{remarque}

\begin{definition}[Opérateurs de Wirtinger]\label{def:wirtinger}
On définit les \emph{opérateurs de Wirtinger} par :
\[
  \frac{\partial}{\partial z} = \frac{1}{2}\left(\frac{\partial}{\partial x} - i\frac{\partial}{\partial y}\right), \qquad
  \frac{\partial}{\partial \conj{z}} = \frac{1}{2}\left(\frac{\partial}{\partial x} + i\frac{\partial}{\partial y}\right).
\]
\end{definition}

\begin{proposition}[Holomorphie et opérateur de Wirtinger]\label{prop:wirtinger-holomorphie}
Soit $f = u + iv$ de classe $\mathscr{C}^1$ sur un ouvert $\Omega$. Alors :
\begin{enumerate}[label=(\roman*)]
  \item Les équations de Cauchy-Riemann sont équivalentes à $\frac{\partial f}{\partial \conj{z}} = 0$.
  \item Si $f$ est holomorphe, alors $f'(z) = \frac{\partial f}{\partial z}$.
\end{enumerate}
\end{proposition}

\begin{proof}
(i) On calcule :
\begin{align*}
  \frac{\partial f}{\partial \conj{z}} &= \frac{1}{2}\left(\frac{\partial f}{\partial x} + i\frac{\partial f}{\partial y}\right) = \frac{1}{2}\left[(u_x + iv_x) + i(u_y + iv_y)\right] \\
  &= \frac{1}{2}\left[(u_x - v_y) + i(v_x + u_y)\right].
\end{align*}
Cette expression est nulle si et seulement si $u_x = v_y$ et $u_y = -v_x$.

(ii) De même, $\frac{\partial f}{\partial z} = \frac{1}{2}[(u_x+v_y) + i(v_x - u_y)] = u_x + iv_x = f'(z)$ en utilisant les équations de Cauchy-Riemann.
\end{proof}

\begin{exemple}[Vérification par Cauchy-Riemann]\label{ex:CR-verification}
\begin{enumerate}[label=(\roman*)]
  \item Soit $f(z) = z^2 = (x+iy)^2 = (x^2-y^2) + 2ixy$. Donc $u = x^2 - y^2$ et $v = 2xy$.
  \[
    u_x = 2x = v_y, \qquad u_y = -2y = -v_x.
  \]
  Les conditions de Cauchy-Riemann sont satisfaites, $u$ et $v$ sont $\mathscr{C}^1$, donc $f$ est holomorphe et $f'(z) = u_x + iv_x = 2x + 2iy = 2z$.

  \item Soit $f(z) = e^z = e^x\cos y + ie^x\sin y$. Donc $u = e^x\cos y$, $v = e^x\sin y$.
  \[
    u_x = e^x\cos y = v_y, \qquad u_y = -e^x\sin y = -v_x.
  \]
  Donc $f$ est holomorphe et $f'(z) = e^x\cos y + ie^x\sin y = e^z$.

  \item Soit $f(z) = \abs{z}^2 = x^2 + y^2$. Ici $u = x^2+y^2$ et $v = 0$.
  \[
    u_x = 2x, \quad v_y = 0 \implies u_x = v_y \iff x = 0.
  \]
  \[
    u_y = 2y, \quad v_x = 0 \implies u_y = -v_x \iff y = 0.
  \]
  Les conditions ne sont satisfaites qu'en $z = 0$. Mais les dérivées partielles sont continues, donc $f$ est $\C$-dérivable en $z = 0$ seulement, avec $f'(0) = 0$. La fonction n'est holomorphe sur aucun ouvert.
\end{enumerate}
\end{exemple}

\section{Interprétation géométrique : conformité}\label{sec:conformite}

\begin{definition}[Application conforme]\label{def:conforme}
Soit $\Omega \subseteq \C$ un ouvert et $f : \Omega \to \C$ une fonction de classe $\mathscr{C}^1$. On dit que $f$ est \emph{conforme} en $z_0 \in \Omega$ si $f$ préserve les angles orientés entre courbes passant par $z_0$, et préserve l'orientation. Plus précisément, si $\gamma_1$ et $\gamma_2$ sont deux courbes $\mathscr{C}^1$ passant par $z_0$ avec $\gamma_1'(0), \gamma_2'(0) \neq 0$, alors l'angle orienté entre $f \circ \gamma_1$ et $f \circ \gamma_2$ en $f(z_0)$ est égal à celui entre $\gamma_1$ et $\gamma_2$ en $z_0$.
\end{definition}

\begin{theoreme}[Holomorphie et conformité]\label{thm:holomorphe-conforme}
Soit $f$ holomorphe sur un ouvert $\Omega$. Si $f'(z_0) \neq 0$, alors $f$ est conforme en $z_0$.

Réciproquement, si $f$ est $\mathscr{C}^1$, préserve les angles et l'orientation en $z_0$, et sa différentielle en $z_0$ est non nulle, alors $f$ est holomorphe en $z_0$ avec $f'(z_0) \neq 0$.
\end{theoreme}

\begin{proof}
\emph{Sens direct.} Soit $\gamma : (-\varepsilon, \varepsilon) \to \Omega$ une courbe $\mathscr{C}^1$ avec $\gamma(0) = z_0$ et $\gamma'(0) \neq 0$. La courbe image est $\sigma = f \circ \gamma$, de dérivée
\[
  \sigma'(0) = f'(z_0) \cdot \gamma'(0).
\]
L'argument de $\sigma'(0)$ est
\[
  \arg(\sigma'(0)) = \arg(f'(z_0)) + \arg(\gamma'(0)).
\]
L'angle entre deux courbes $\gamma_1, \gamma_2$ en $z_0$ est $\arg(\gamma_2'(0)) - \arg(\gamma_1'(0))$. L'angle entre leurs images est
\[
  \arg(\sigma_2'(0)) - \arg(\sigma_1'(0)) = \arg(\gamma_2'(0)) - \arg(\gamma_1'(0)).
\]
Le terme $\arg(f'(z_0))$ se simplifie : l'angle est préservé.

De plus, $\abs{\sigma'(0)} = \abs{f'(z_0)}\abs{\gamma'(0)}$, donc les longueurs sont multipliées par le facteur constant $\abs{f'(z_0)}$ (localement). C'est pourquoi $f$ est dite \emph{conforme} : elle ressemble localement à une rotation d'angle $\arg(f'(z_0))$ composée avec une homothétie de rapport $\abs{f'(z_0)}$.

\emph{Sens réciproque.} La matrice jacobienne de $f = (u,v)$ en $(x_0, y_0)$ est $J = \begin{psmallmatrix}u_x & u_y \\ v_x & v_y\end{psmallmatrix}$. La préservation des angles et de l'orientation impose que $J$ est de la forme $\lambda R_\theta = \lambda\begin{psmallmatrix}\cos\theta & -\sin\theta \\ \sin\theta & \cos\theta\end{psmallmatrix}$ avec $\lambda > 0$. Posant $a = \lambda\cos\theta$ et $b = \lambda\sin\theta$, on obtient $u_x = a = v_y$ et $u_y = -b = -v_x$, qui sont les équations de Cauchy-Riemann.
\end{proof}

\begin{exemple}[Conformité de $f(z) = z^2$]\label{ex:conformite-z2}
La fonction $f(z) = z^2$ est conforme en tout point $z_0 \neq 0$ (car $f'(z_0) = 2z_0 \neq 0$). En $z_0 = 0$, $f'(0) = 0$ et la fonction double les angles : l'angle entre deux courbes est multiplié par $2$.

Considérons la grille formée par les droites $x = c$ (constante) et $y = c$ (constante). L'image de $x = c$ par $f(z) = z^2 = (x^2-y^2)+2ixy$ est la courbe paramétrée par :
\[
  u = c^2 - y^2, \quad v = 2cy \quad\implies\quad u = c^2 - \frac{v^2}{4c^2},
\]
qui est une parabole d'axe horizontal. De même, l'image de $y = c$ est $u = x^2 - c^2$, $v = 2cx$, soit $u = \frac{v^2}{4c^2} - c^2$, une autre famille de paraboles. Ces deux familles de paraboles se coupent à angles droits, confirmant la conformité.
\end{exemple}

\begin{center}
\begin{tikzpicture}[scale=0.85]
  % Grille originale
  \begin{scope}[shift={(-5.5,0)}]
    \node[above] at (1.5,3.3) {\textbf{Plan $z$}};
    \draw[->] (-0.5,0) -- (3.5,0) node[right] {$x$};
    \draw[->] (0,-3.3) -- (0,3.3) node[above] {$y$};
    \foreach \c in {0.3,0.6,...,3.0} {
      \draw[blue!60,thin] (\c,-3) -- (\c,3);
    }
    \foreach \c in {-3.0,-2.7,...,3.0} {
      \draw[red!60,thin] (0,\c) -- (3,\c);
    }
  \end{scope}

  % Grille image
  \begin{scope}[shift={(4,0)}]
    \node[above] at (1.5,3.3) {\textbf{Plan $w = z^2$}};
    \draw[->] (-4.5,0) -- (4.5,0) node[right] {$u$};
    \draw[->] (0,-4) -- (0,4) node[above] {$v$};

    % Images des droites x = c (paraboles u = c^2 - v^2/(4c^2))
    \foreach \c in {0.3,0.6,...,2.4} {
      \draw[blue!60,thin,domain=-3.5:3.5,samples=80,smooth]
        plot ({\c*\c - (\x)^2/(4*\c*\c)}, {\x});
    }

    % Images des droites y = c (paraboles u = v^2/(4c^2) - c^2)
    \foreach \c in {0.3,0.6,...,2.4} {
      \draw[red!60,thin,domain=-3.5:3.5,samples=80,smooth]
        plot ({(\x)^2/(4*\c*\c) - \c*\c}, {\x});
    }
    % y = -c donne les mêmes courbes
  \end{scope}
\end{tikzpicture}

\medskip
\textit{Transformation conforme $z \mapsto z^2$ : la grille orthogonale est envoyée sur deux familles de paraboles qui se coupent à angles droits.}
\end{center}

\begin{exemple}[Conformité de l'exponentielle]\label{ex:conformite-exp}
La fonction $f(z) = e^z$ est entière et $f'(z) = e^z \neq 0$ pour tout $z$, donc $f$ est conforme en tout point. L'image de la droite horizontale $y = c$ est le rayon $\{re^{ic} : r > 0\}$. L'image de la droite verticale $x = c$ est le cercle $\{e^c e^{iy} : y \in \R\} = C(0, e^c)$. Les rayons et les cercles sont orthogonaux, confirmant la conformité.
\end{exemple}

\section{Fonctions harmoniques et lien avec l'holomorphie}\label{sec:harmoniques}

\begin{definition}[Fonction harmonique]\label{def:harmonique}
Soit $\Omega \subseteq \C$ un ouvert. Une fonction $u : \Omega \to \R$ de classe $\mathscr{C}^2$ est \emph{harmonique} si elle satisfait l'équation de Laplace :
\[
  \Delta u = \frac{\partial^2 u}{\partial x^2} + \frac{\partial^2 u}{\partial y^2} = 0.
\]
L'opérateur $\Delta = \partial_{xx} + \partial_{yy}$ est le \emph{laplacien}.
\end{definition}

\begin{theoreme}[Holomorphie et harmonicité]\label{thm:holomorphe-harmonique}
Si $f = u + iv$ est holomorphe sur un ouvert $\Omega$, alors $u$ et $v$ sont harmoniques sur $\Omega$.
\end{theoreme}

\begin{proof}
On admettra temporairement que toute fonction holomorphe est de classe $\mathscr{C}^\infty$ (ce résultat fondamental sera démontré au chapitre 5 via la formule intégrale de Cauchy). Les équations de Cauchy-Riemann donnent $u_x = v_y$ et $u_y = -v_x$. En dérivant :
\[
  u_{xx} = v_{yx}, \quad u_{yy} = -v_{xy}.
\]
Comme $f$ est $\mathscr{C}^2$, le théorème de Schwarz donne $v_{yx} = v_{xy}$. Donc :
\[
  \Delta u = u_{xx} + u_{yy} = v_{yx} - v_{xy} = 0.
\]
De même, $v_{xx} = -u_{yx}$ et $v_{yy} = u_{xy}$, d'où $\Delta v = -u_{yx} + u_{xy} = 0$.
\end{proof}

\begin{remarque}\label{rem:laplacien-wirtinger}
Le laplacien s'exprime élégamment à l'aide des opérateurs de Wirtinger :
\[
  \Delta = 4\frac{\partial^2}{\partial z\,\partial\conj{z}}.
\]
En effet, $4\frac{\partial^2}{\partial z\,\partial\conj{z}} = 4\cdot\frac{1}{2}\left(\frac{\partial}{\partial x}-i\frac{\partial}{\partial y}\right)\cdot\frac{1}{2}\left(\frac{\partial}{\partial x}+i\frac{\partial}{\partial y}\right) = \frac{\partial^2}{\partial x^2}+\frac{\partial^2}{\partial y^2} = \Delta$.
\end{remarque}

\begin{definition}[Conjuguée harmonique]\label{def:conjuguee-harmonique}
Soit $u$ une fonction harmonique sur un ouvert $\Omega$. Une fonction harmonique $v$ est une \emph{conjuguée harmonique} de $u$ si $f = u + iv$ est holomorphe sur $\Omega$, c'est-à-dire si $u$ et $v$ satisfont les équations de Cauchy-Riemann.
\end{definition}

\begin{theoreme}[Existence de la conjuguée harmonique]\label{thm:conjuguee-harmonique}
Soit $\Omega \subseteq \C$ un domaine \emph{simplement connexe} et $u : \Omega \to \R$ une fonction harmonique. Alors $u$ admet une conjuguée harmonique $v$, unique à une constante additive près.
\end{theoreme}

\begin{proof}
Fixons un point $z_0 = x_0 + iy_0 \in \Omega$. On cherche $v$ telle que $v_x = -u_y$ et $v_y = u_x$. On définit, pour $(x,y) \in \Omega$ :
\[
  v(x,y) = \int_\gamma \left(-u_y\,\dd x' + u_x\,\dd y'\right),
\]
où $\gamma$ est un chemin quelconque dans $\Omega$ de $(x_0, y_0)$ à $(x,y)$.

Pour que cette intégrale soit bien définie (indépendante du chemin), il faut vérifier la condition d'exactitude. La $1$-forme $\omega = -u_y\,\dd x + u_x\,\dd y$ est fermée car :
\[
  \frac{\partial(-u_y)}{\partial y} = -u_{yy}, \qquad \frac{\partial(u_x)}{\partial x} = u_{xx}.
\]
La condition $\frac{\partial(-u_y)}{\partial y} = \frac{\partial(u_x)}{\partial x}$ se réécrit $-u_{yy} = u_{xx}$, soit $\Delta u = 0$, ce qui est vrai par hypothèse. Comme $\Omega$ est simplement connexe, toute forme fermée est exacte (lemme de Poincaré), donc l'intégrale est indépendante du chemin.

La fonction $v$ ainsi définie satisfait $v_x = -u_y$ et $v_y = u_x$ par construction, donc $f = u + iv$ est holomorphe.

L'unicité à constante près provient du fait que si $v_1$ et $v_2$ sont deux conjuguées harmoniques de $u$, alors $v_1 - v_2$ est une fonction holomorphe réelle (sa partie imaginaire est nulle), donc constante sur le connexe $\Omega$ (par les équations de Cauchy-Riemann : $(v_1 - v_2)_x = 0$ et $(v_1 - v_2)_y = 0$).
\end{proof}

\begin{remarque}\label{rem:connexite-simple-necessaire}
L'hypothèse de connexité simple est essentielle. Sur $\C^*$, la fonction $u(x,y) = \ln(x^2+y^2)^{1/2} = \frac{1}{2}\ln(x^2+y^2)$ est harmonique, mais sa conjuguée harmonique est $v(x,y) = \arctan(y/x)$ (à constante près), qui n'est pas univoque sur $\C^*$. On ne peut pas définir une fonction $v$ continue sur tout $\C^*$.
\end{remarque}

\begin{exemple}[Construction d'une conjuguée harmonique]\label{ex:conjuguee-harmonique}
Soit $u(x,y) = x^2 - y^2 + x$. Vérifions que $u$ est harmonique :
\[
  u_{xx} = 2, \quad u_{yy} = -2, \quad \Delta u = 0. \quad \checkmark
\]
Cherchons $v$ telle que $v_y = u_x = 2x + 1$ et $v_x = -u_y = 2y$. De la première équation :
\[
  v(x,y) = \int (2x+1)\,\dd y = (2x+1)y + \varphi(x),
\]
où $\varphi$ est une fonction de $x$ seul. De la seconde équation :
\[
  v_x = 2y + \varphi'(x) = 2y \implies \varphi'(x) = 0 \implies \varphi(x) = C.
\]
Donc $v(x,y) = 2xy + y + C$. La fonction holomorphe correspondante est
\[
  f(z) = (x^2-y^2+x) + i(2xy+y) + iC = z^2 + z + iC.
\]
\end{exemple}

\begin{proposition}[Propriétés des fonctions harmoniques]\label{prop:harmoniques-proprietes}
Soit $u$ harmonique sur un domaine $\Omega$.
\begin{enumerate}[label=(\roman*)]
  \item $u$ est de classe $\mathscr{C}^\infty$ (car c'est la partie réelle d'une fonction holomorphe, elle-même $\mathscr{C}^\infty$).
  \item $u$ satisfait la propriété de la moyenne : pour tout disque $\conj{D}(a,r) \subseteq \Omega$,
  \[
    u(a) = \frac{1}{2\pi}\int_0^{2\pi} u(a + re^{i\theta})\,\dd\theta.
  \]
  \item $u$ satisfait le principe du maximum : si $u$ atteint son maximum (ou son minimum) en un point intérieur de $\Omega$, alors $u$ est constante.
\end{enumerate}
\end{proposition}

\begin{proof}
Ces résultats seront démontrés rigoureusement au chapitre 5 comme conséquences de la formule intégrale de Cauchy. Donnons cependant l'idée pour le principe du maximum.

Supposons que $u$ atteint son maximum $M$ en $a \in \Omega$. Par la propriété de la moyenne,
\[
  M = u(a) = \frac{1}{2\pi}\int_0^{2\pi} u(a+re^{i\theta})\,\dd\theta.
\]
Comme $u(a+re^{i\theta}) \leq M$ pour tout $\theta$, et que la moyenne de fonctions majorées par $M$ vaut $M$, on doit avoir $u(a+re^{i\theta}) = M$ pour tout $\theta$. En itérant, $u = M$ sur tout disque centré en $a$, puis par connexité, sur tout $\Omega$.
\end{proof}

\begin{exemple}[Application : problème de Dirichlet]\label{ex:dirichlet}
Le \emph{problème de Dirichlet} consiste à trouver une fonction harmonique $u$ sur un domaine $\Omega$ dont les valeurs au bord $\partial\Omega$ sont prescrites. Par exemple, sur le disque $\D$, si $u$ est harmonique et continue sur $\conj{\D}$, la formule intégrale de Poisson (qui sera vue au chapitre 5) donne :
\[
  u(re^{i\theta}) = \frac{1}{2\pi}\int_0^{2\pi} \frac{1 - r^2}{1 - 2r\cos(\theta-\varphi)+r^2}\, u(e^{i\varphi})\,\dd\varphi,
\]
pour $0 \leq r < 1$. Le noyau $P_r(\theta-\varphi) = \frac{1-r^2}{1-2r\cos(\theta-\varphi)+r^2}$ est le \emph{noyau de Poisson}.
\end{exemple}

\section{Conditions de Cauchy-Riemann en coordonnées polaires}\label{sec:CR-polaires}

\begin{proposition}[Cauchy-Riemann en coordonnées polaires]\label{prop:CR-polaires}
Soit $f = u + iv$ holomorphe sur un ouvert $\Omega \subseteq \C^*$. En coordonnées polaires $z = re^{i\theta}$, les équations de Cauchy-Riemann s'écrivent :
\begin{equation}\label{eq:CR-polaires}
  \frac{\partial u}{\partial r} = \frac{1}{r}\frac{\partial v}{\partial \theta}, \qquad
  \frac{\partial v}{\partial r} = -\frac{1}{r}\frac{\partial u}{\partial \theta}.
\end{equation}
La dérivée complexe est
\[
  f'(z) = e^{-i\theta}\left(\frac{\partial u}{\partial r} + i\frac{\partial v}{\partial r}\right).
\]
\end{proposition}

\begin{proof}
On utilise le changement de variables $x = r\cos\theta$, $y = r\sin\theta$, d'où par la règle de la chaîne :
\[
  \frac{\partial u}{\partial r} = u_x \cos\theta + u_y \sin\theta, \qquad
  \frac{\partial u}{\partial \theta} = -u_x r\sin\theta + u_y r\cos\theta.
\]
En utilisant $u_x = v_y$ et $u_y = -v_x$ :
\begin{align*}
  \frac{\partial u}{\partial r} &= v_y\cos\theta - v_x\sin\theta = \frac{1}{r}\left(v_x r\cos\theta \cdot 0 + \cdots\right).
\end{align*}
Calculons plutôt $\frac{1}{r}\frac{\partial v}{\partial\theta}$ :
\[
  \frac{1}{r}\frac{\partial v}{\partial \theta} = \frac{1}{r}\left(-v_x r\sin\theta + v_y r\cos\theta\right) = -v_x\sin\theta + v_y\cos\theta.
\]
Et
\[
  \frac{\partial u}{\partial r} = u_x\cos\theta + u_y\sin\theta = v_y\cos\theta + (-v_x)\sin\theta = \frac{1}{r}\frac{\partial v}{\partial\theta}.
\]
La deuxième équation se vérifie de manière analogue.
\end{proof}

\begin{exemple}[Vérification avec $f(z) = \Log z$]\label{ex:CR-polaires-log}
Le logarithme principal est $\Log z = \ln r + i\theta$ (pour $\theta \in (-\pi,\pi)$). Ici $u = \ln r$ et $v = \theta$. Vérifions :
\[
  u_r = \frac{1}{r}, \quad \frac{1}{r}v_\theta = \frac{1}{r}\cdot 1 = \frac{1}{r}. \quad \checkmark
\]
\[
  v_r = 0, \quad -\frac{1}{r}u_\theta = -\frac{1}{r}\cdot 0 = 0. \quad \checkmark
\]
La dérivée est $f'(z) = e^{-i\theta}(u_r + iv_r) = e^{-i\theta}\cdot\frac{1}{r} = \frac{1}{re^{i\theta}} = \frac{1}{z}$.
\end{exemple}

\section{Exercices}\label{sec:exo-ch2}

\begin{exercice}[$\star$]\label{exo:CR-simple}
Vérifier les conditions de Cauchy-Riemann et calculer la dérivée complexe pour :
\begin{enumerate}[label=(\alph*)]
  \item $f(z) = z^3$.
  \item $f(z) = \frac{1}{z}$ sur $\C^*$.
  \item $f(z) = e^{z^2}$.
\end{enumerate}
\end{exercice}

\begin{exercice}[$\star$]\label{exo:non-holomorphe-CR}
Montrer à l'aide des conditions de Cauchy-Riemann que les fonctions suivantes ne sont holomorphes en aucun point :
\begin{enumerate}[label=(\alph*)]
  \item $f(z) = \re(z)$.
  \item $f(z) = \conj{z}^2$.
  \item $f(z) = z\conj{z}$.
\end{enumerate}
Pour (c), la fonction est-elle $\C$-dérivable en un point isolé ?
\end{exercice}

\begin{exercice}[$\star\star$]\label{exo:conjuguee-harmonique-construction}
Pour chacune des fonctions harmoniques suivantes, trouver une conjuguée harmonique et la fonction holomorphe correspondante :
\begin{enumerate}[label=(\alph*)]
  \item $u(x,y) = e^x\cos y$.
  \item $u(x,y) = x^3 - 3xy^2$.
  \item $u(x,y) = \ln(x^2 + y^2)$ sur $\C^*$ (attention à la connexité simple).
  \item $u(x,y) = \frac{x}{x^2+y^2}$ sur $\C^*$.
\end{enumerate}
\end{exercice}

\begin{exercice}[$\star\star$]\label{exo:harmonique-verification}
Montrer qu'une fonction harmonique de la forme $u(x,y) = ax^2 + bxy + cy^2 + dx + ey + f$ (avec $a,b,c,d,e,f \in \R$) est nécessairement telle que $a + c = 0$. Trouver la conjuguée harmonique et la fonction holomorphe associée.
\end{exercice}

\begin{exercice}[$\star\star$]\label{exo:conformite-exp}
Soit $f(z) = e^z$.
\begin{enumerate}[label=(\alph*)]
  \item Déterminer l'image par $f$ de la bande $S = \{z \in \C : 0 < \im(z) < \pi\}$.
  \item Déterminer l'image par $f$ du rectangle $\{z : 0 < \re(z) < 1,\; 0 < \im(z) < \pi/2\}$.
  \item En quel(s) point(s) $f$ est-elle injective ? Est-elle injective sur $S$ ?
\end{enumerate}
\end{exercice}

\begin{exercice}[$\star\star$]\label{exo:jacobien-holomorphe}
Soit $f = u + iv$ holomorphe sur un ouvert $\Omega$.
\begin{enumerate}[label=(\alph*)]
  \item Montrer que le jacobien de $f$ (vue comme application $\R^2 \to \R^2$) est $J_f = \abs{f'}^2$.
  \item En déduire que $f$ est localement injective en tout point $z_0$ tel que $f'(z_0) \neq 0$ (par le théorème d'inversion locale).
  \item Montrer que si $f$ est injective sur $\Omega$, alors $f'(z) \neq 0$ pour tout $z \in \Omega$.
\end{enumerate}
\end{exercice}

\begin{exercice}[$\star\star$]\label{exo:wirtinger-calcul}
\begin{enumerate}[label=(\alph*)]
  \item Montrer que $\frac{\partial}{\partial\conj{z}}(z^n) = 0$ et $\frac{\partial}{\partial z}(z^n) = nz^{n-1}$ pour tout $n \in \N$.
  \item Montrer que $\frac{\partial}{\partial\conj{z}}(\conj{z}^n) = n\conj{z}^{n-1}$ et $\frac{\partial}{\partial z}(\conj{z}^n) = 0$.
  \item Calculer $\frac{\partial}{\partial z}(\abs{z}^2)$ et $\frac{\partial}{\partial\conj{z}}(\abs{z}^2)$.
  \item Montrer que $\Delta = 4\frac{\partial^2}{\partial z\,\partial\conj{z}}$.
\end{enumerate}
\end{exercice}

\begin{exercice}[$\star\star\star$]\label{exo:CR-polaires-exo}
\begin{enumerate}[label=(\alph*)]
  \item Écrire les conditions de Cauchy-Riemann en coordonnées polaires et les vérifier pour $f(z) = z^n$ ($n \in \Z$).
  \item En déduire le laplacien en coordonnées polaires :
  \[
    \Delta u = \frac{\partial^2 u}{\partial r^2} + \frac{1}{r}\frac{\partial u}{\partial r} + \frac{1}{r^2}\frac{\partial^2 u}{\partial\theta^2}.
  \]
  \item Montrer que les fonctions $u(r,\theta) = r^n\cos(n\theta)$ et $v(r,\theta) = r^n\sin(n\theta)$ sont harmoniques pour tout $n \in \N$.
\end{enumerate}
\end{exercice}

\begin{exercice}[$\star\star\star$]\label{exo:principe-maximum-harmonique}
Soit $u$ harmonique et non constante sur un domaine borné $\Omega$, continue sur $\conj{\Omega}$.
\begin{enumerate}[label=(\alph*)]
  \item En admettant la propriété de la moyenne, montrer que $u$ ne peut pas atteindre son maximum en un point intérieur (principe du maximum).
  \item En déduire que $\max_{\conj{\Omega}} u = \max_{\partial\Omega} u$.
  \item En déduire l'unicité de la solution du problème de Dirichlet : si $u_1$ et $u_2$ sont harmoniques sur $\Omega$, continues sur $\conj{\Omega}$, et $u_1 = u_2$ sur $\partial\Omega$, alors $u_1 = u_2$ sur $\Omega$.
\end{enumerate}
\end{exercice}

\begin{exercice}[$\star\star\star$]\label{exo:holomorphe-constante}
Soit $f$ holomorphe sur un domaine $\Omega$. Montrer que $f$ est constante si l'une des conditions suivantes est satisfaite :
\begin{enumerate}[label=(\alph*)]
  \item $\conj{f}$ est holomorphe sur $\Omega$.
  \item $\abs{f}$ est constante sur $\Omega$.
  \item $\re(f)$ est constante sur $\Omega$.
  \item $\arg(f)$ est constante sur $\Omega$ (et $f$ ne s'annule pas).
  \item $f(\Omega) \subseteq \R$.
  \item $f(\Omega)$ est contenu dans un cercle.
\end{enumerate}
(\emph{Indication :} utiliser les équations de Cauchy-Riemann.)
\end{exercice}

\begin{exercice}[$\star\star\star$]\label{exo:conforme-applications}
\begin{enumerate}[label=(\alph*)]
  \item Montrer que la transformation de Joukowski $f(z) = z + 1/z$ est conforme sur $\C \setminus \{-1, 0, 1\}$. Déterminer l'image du cercle $\abs{z} = R$ pour $R > 1$ et $R = 1$.
  \item Montrer que $f(z) = \frac{1}{2}(z + 1/z)$ envoie l'extérieur du disque unité conformément sur $\C \setminus [-1,1]$.
\end{enumerate}
\end{exercice}

%%%%%%%%%%%%%%%%%%%%%%%%%%%%%%%%%%%%%%%%%%%%%%%%%%%%%%%%%%%%%%%%%%%%
%  CHAPITRE 3 — Fonctions Élémentaires
%%%%%%%%%%%%%%%%%%%%%%%%%%%%%%%%%%%%%%%%%%%%%%%%%%%%%%%%%%%%%%%%%%%%
\chapter{Fonctions Élémentaires}\label{chap:fonctions-elementaires}

\section{Introduction}

Les fonctions élémentaires --- exponentielle, logarithme, puissances, fonctions
trigonométriques et hyperboliques --- constituent le socle de l'analyse complexe.
Leur étude révèle des phénomènes sans analogue réel : la périodicité de
l'exponentielle selon l'axe imaginaire, le caractère multivalué du logarithme et
des puissances complexes, et l'unification profonde entre fonctions
trigonométriques et hyperboliques par les formules d'Euler.

Ce chapitre développe systématiquement ces fonctions, en insistant sur les
propriétés géométriques (image de bandes, de droites), sur la notion de
\emph{branche} d'une fonction multivaluée, et sur les calculs explicites.

%=======================================================================
\section{L'exponentielle complexe}\label{sec:exp}
%=======================================================================

\begin{definition}[Exponentielle complexe]\label{def:exp-complexe}
Pour $z = x + iy \in \C$ (avec $x, y \in \R$), on définit
\[
  \exp(z) = e^z = e^x(\cos y + i \sin y).
\]
De manière équivalente, $\exp$ est définie par la série entière convergente sur $\C$ tout entier :
\[
  \exp(z) = \sum_{n=0}^{\infty} \frac{z^n}{n!}.
\]
\end{definition}

\begin{theoreme}[Propriétés fondamentales de l'exponentielle]\label{thm:proprietes-exp}
La fonction $\exp : \C \to \C$ vérifie les propriétés suivantes.
\begin{enumerate}[label=(\roman*)]
  \item \textbf{Holomorphie.} $\exp$ est holomorphe sur $\C$ tout entier (fonction entière) et $\exp'(z) = \exp(z)$ pour tout $z \in \C$.
  \item \textbf{Propriété fonctionnelle.} Pour tous $z, w \in \C$,
  \[
    e^{z+w} = e^z \cdot e^w.
  \]
  \item \textbf{Non-annulation.} $e^z \neq 0$ pour tout $z \in \C$, et $e^{-z} = (e^z)^{-1}$.
  \item \textbf{Module et argument.} $\abs{e^z} = e^{\re(z)} = e^x$ et $\arg(e^z) = \im(z) + 2k\pi$ pour $k \in \Z$.
  \item \textbf{Périodicité.} $\exp$ est périodique de période $2\pi i$ :
  \[
    e^{z + 2\pi i} = e^z \quad \text{pour tout } z \in \C.
  \]
  De plus, $2\pi i$ est une période fondamentale : si $e^{z+\omega} = e^z$ pour tout $z$, alors $\omega \in 2\pi i \Z$.
  \item \textbf{Surjectivité.} $\exp : \C \to \C \setminus \{0\}$ est surjective.
\end{enumerate}
\end{theoreme}

\begin{proof}
\emph{(i)} La série entière $\sum z^n/n!$ a un rayon de convergence infini (par le critère de d'Alembert : $\abs{z}/(n+1) \to 0$). La dérivée terme à terme donne $\exp'(z) = \sum_{n=1}^{\infty} z^{n-1}/(n-1)! = \exp(z)$.

\medskip

\emph{(ii)} On utilise le produit de Cauchy des séries :
\[
  e^z \cdot e^w = \left(\sum_{n=0}^{\infty} \frac{z^n}{n!}\right)
                  \left(\sum_{m=0}^{\infty} \frac{w^m}{m!}\right)
               = \sum_{k=0}^{\infty} \frac{1}{k!}\sum_{j=0}^{k} \binom{k}{j} z^j w^{k-j}
               = \sum_{k=0}^{\infty} \frac{(z+w)^k}{k!} = e^{z+w},
\]
la dernière égalité provenant du binôme de Newton.

\medskip

\emph{(iii)} De $e^z \cdot e^{-z} = e^0 = 1$, on déduit $e^z \neq 0$ et $e^{-z} = (e^z)^{-1}$.

\medskip

\emph{(iv)} On a $\abs{e^z} = \abs{e^x(\cos y + i\sin y)} = e^x \abs{\cos y + i\sin y} = e^x$, car $\cos^2 y + \sin^2 y = 1$.

\medskip

\emph{(v)} $e^{z+2\pi i} = e^z \cdot e^{2\pi i} = e^z(\cos 2\pi + i\sin 2\pi) = e^z$. Pour la minimalité, si $e^\omega = 1$ avec $\omega = a + ib$, alors $e^a = 1$ donc $a = 0$, et $\cos b + i\sin b = 1$ donc $b \in 2\pi\Z$.

\medskip

\emph{(vi)} Soit $w \in \C \setminus \{0\}$. Écrivons $w = \rho e^{i\theta}$ avec $\rho > 0$ et $\theta \in \R$. Alors $z = \ln\rho + i\theta$ vérifie $e^z = e^{\ln\rho}(\cos\theta + i\sin\theta) = \rho e^{i\theta} = w$.
\end{proof}

\begin{remarque}\label{rem:exp-non-injective}
La périodicité de $\exp$ implique qu'elle n'est \emph{pas} injective sur $\C$.
Plus précisément, $e^{z_1} = e^{z_2}$ si et seulement si $z_1 - z_2 \in 2\pi i\Z$.
En revanche, $\exp$ est injective sur toute bande horizontale de largeur strictement
inférieure à $2\pi$, par exemple sur
$\{z \in \C : -\pi < \im(z) \leq \pi\}$.
\end{remarque}

\subsection{Image de bandes horizontales et verticales}

L'un des aspects les plus éclairants de l'exponentielle complexe est son action
géométrique sur les bandes du plan.

\begin{proposition}[Image de bandes horizontales]\label{prop:exp-bandes-horizontales}
Soit $\alpha < \beta$ avec $\beta - \alpha \leq 2\pi$. Alors $\exp$ envoie la
bande horizontale
\[
  B_{\alpha,\beta} = \{z \in \C : \alpha < \im(z) < \beta\}
\]
sur le secteur angulaire
\[
  S_{\alpha,\beta} = \{w \in \C \setminus \{0\} : \alpha < \arg(w) < \beta\}.
\]
En particulier :
\begin{enumerate}[label=(\alph*)]
  \item La bande $\{-\pi < \im(z) < \pi\}$ est envoyée sur $\C \setminus \R_{\leq 0}$ (le plan coupé le long de la demi-droite réelle négative).
  \item La bande $\{0 < \im(z) < \pi\}$ est envoyée sur le demi-plan supérieur $\{\im(w) > 0\}$.
  \item La bande $\{-\pi/2 < \im(z) < \pi/2\}$ est envoyée sur le demi-plan $\{\re(w) > 0\}$.
\end{enumerate}
\end{proposition}

\begin{proof}
Soit $z = x + iy$ avec $\alpha < y < \beta$. Alors $e^z = e^x e^{iy}$.
Le module $e^x$ parcourt $(0,+\infty)$ quand $x$ parcourt $\R$, et
l'argument $y$ varie dans $(\alpha,\beta)$. Donc l'image est exactement
l'ensemble des $w \neq 0$ de module quelconque et d'argument dans $(\alpha,\beta)$.
Les cas particuliers s'en déduisent immédiatement.
\end{proof}

\begin{proposition}[Image de droites]\label{prop:exp-droites}
\leavevmode
\begin{enumerate}[label=(\roman*)]
  \item \textbf{Droite horizontale} $\im(z) = y_0$ (i.e.\ $z = x + iy_0$, $x \in \R$) : l'image est la demi-droite issue de l'origine, de direction $e^{iy_0}$, privée de l'origine.
  \item \textbf{Droite verticale} $\re(z) = x_0$ (i.e.\ $z = x_0 + iy$, $y \in \R$) : l'image est le cercle de centre $0$ et de rayon $e^{x_0}$, parcouru dans le sens trigonométrique.
\end{enumerate}
\end{proposition}

\begin{proof}
\emph{(i)} Si $z = x + iy_0$, alors $e^z = e^x e^{iy_0}$. Quand $x$ parcourt $\R$, $e^x$ parcourt $(0,+\infty)$, et la direction $e^{iy_0}$ est fixe.

\emph{(ii)} Si $z = x_0 + iy$, alors $e^z = e^{x_0} e^{iy}$. Le module $e^{x_0}$ est fixe, et $e^{iy}$ parcourt le cercle unité quand $y$ parcourt $[0,2\pi)$.
\end{proof}

\begin{center}
\begin{tikzpicture}[>=Stealth, scale=0.85]
  % ── Plan z (gauche) ──
  \begin{scope}[shift={(-5.5,0)}]
    \node[above] at (0,3.3) {\textbf{Plan $z$}};
    \draw[->] (-3,0) -- (3,0) node[right] {$\re$};
    \draw[->] (0,-3) -- (0,3) node[above] {$\im$};

    % Bande horizontale -pi < y < pi
    \fill[blue!10] (-2.8,-2.5) rectangle (2.8,2.5);
    \draw[blue!70!black, thick, dashed] (-2.8,2.5) -- (2.8,2.5)
      node[right, font=\footnotesize] {$y = \pi$};
    \draw[blue!70!black, thick, dashed] (-2.8,-2.5) -- (2.8,-2.5)
      node[right, font=\footnotesize] {$y = -\pi$};

    % Quelques droites horizontales
    \draw[red!70!black, thick] (-2.5,0) -- (2.5,0);
    \node[red!70!black, right, font=\footnotesize] at (2.5,0.2) {$y=0$};

    \draw[orange!80!black, thick] (-2.5,1.25) -- (2.5,1.25);
    \node[orange!80!black, right, font=\footnotesize] at (2.5,1.45) {$y=\frac{\pi}{2}$};

    \draw[green!60!black, thick] (-2.5,-1.25) -- (2.5,-1.25);
    \node[green!60!black, right, font=\footnotesize] at (2.5,-1.05) {$y=-\frac{\pi}{2}$};

    % Quelques droites verticales
    \draw[purple!70!black, thick, densely dotted] (-1.5,-2.5) -- (-1.5,2.5);
    \draw[purple!70!black, thick, densely dotted] (0,-2.5) -- (0,2.5);
    \draw[purple!70!black, thick, densely dotted] (1.5,-2.5) -- (1.5,2.5);
  \end{scope}

  % ── Flèche ──
  \draw[->, very thick, gray] (-1.2,0) -- (0.5,0)
    node[midway, above, font=\small] {$\exp$};

  % ── Plan w (droite) ──
  \begin{scope}[shift={(5,0)}]
    \node[above] at (0,3.3) {\textbf{Plan $w = e^z$}};
    \draw[->] (-3,0) -- (3,0) node[right] {$\re$};
    \draw[->] (0,-3) -- (0,3) node[above] {$\im$};

    % Demi-droites (images des horizontales)
    \draw[red!70!black, thick, ->] (0,0) -- (2.8,0);
    \node[red!70!black, above right, font=\footnotesize] at (2.2,0.1) {$\arg=0$};

    \draw[orange!80!black, thick, ->] (0,0) -- (0,2.8);
    \node[orange!80!black, right, font=\footnotesize] at (0.1,2.3) {$\arg=\frac{\pi}{2}$};

    \draw[green!60!black, thick, ->] (0,0) -- (0,-2.8);
    \node[green!60!black, right, font=\footnotesize] at (0.1,-2.3) {$\arg=-\frac{\pi}{2}$};

    % Cercles (images des verticales)
    \draw[purple!70!black, thick, densely dotted] (0,0) circle (0.5);
    \draw[purple!70!black, thick, densely dotted] (0,0) circle (1.3);
    \draw[purple!70!black, thick, densely dotted] (0,0) circle (2.5);

    % Coupure
    \draw[gray, very thick, decorate, decoration={zigzag, segment length=4, amplitude=1}]
      (-2.8,0) -- (0,0);
    \node[gray, below, font=\footnotesize] at (-1.5,-0.2) {coupure};
  \end{scope}
\end{tikzpicture}
\captionof{figure}{Image par $\exp$ de la bande $\{-\pi < \im(z) < \pi\}$. Les droites horizontales deviennent des demi-droites, les droites verticales deviennent des cercles.}
\end{center}

\begin{exemple}[Résolution de $e^z = w$]\label{ex:resolution-exp}
Résolvons $e^z = -1 + i$.

On écrit $-1 + i$ en forme polaire : $\abs{-1+i} = \sqrt{2}$ et
$\arg(-1+i) = \frac{3\pi}{4} + 2k\pi$ ($k \in \Z$). Donc
\[
  e^z = \sqrt{2}\,e^{i(3\pi/4 + 2k\pi)}.
\]
On cherche $z = x + iy$ tel que $e^x = \sqrt{2}$ et $y = \frac{3\pi}{4} + 2k\pi$.
Ainsi $x = \ln\sqrt{2} = \frac{1}{2}\ln 2$ et
\[
  z = \frac{1}{2}\ln 2 + i\left(\frac{3\pi}{4} + 2k\pi\right), \quad k \in \Z.
\]
\end{exemple}

\begin{exemple}[Image d'un rectangle]\label{ex:exp-rectangle}
Soit $R = \{z \in \C : 0 < \re(z) < 1,\; 0 < \im(z) < \pi/2\}$, le rectangle ouvert.

L'image $\exp(R)$ est le quart de couronne :
\[
  \exp(R) = \{w \in \C : 1 < \abs{w} < e,\; 0 < \arg(w) < \pi/2\}.
\]
En effet, $\re(z) \in (0,1)$ donne $\abs{w} = e^{\re(z)} \in (1,e)$, et $\im(z) \in (0,\pi/2)$ donne $\arg(w) \in (0,\pi/2)$.
\end{exemple}

%=======================================================================
\section{Le logarithme complexe}\label{sec:log}
%=======================================================================

\subsection{L'argument et sa multivaluation}

\begin{definition}[Argument principal]\label{def:arg-principal}
Pour $z \in \C \setminus \{0\}$, l'\emph{argument principal} de $z$, noté $\Arg(z)$,
est l'unique $\theta \in (-\pi, \pi]$ tel que
\[
  z = \abs{z} e^{i\theta}.
\]
L'\emph{argument} (multivalué) de $z$ est l'ensemble
$\arg(z) = \{\Arg(z) + 2k\pi : k \in \Z\}$.
\end{definition}

\begin{remarque}\label{rem:arg-discontinu}
La fonction $\Arg : \C \setminus \{0\} \to (-\pi,\pi]$ est discontinue sur la demi-droite réelle négative $\R_{<0}$. En effet, si $z = -1 + i\varepsilon$ avec $\varepsilon > 0$ petit, $\Arg(z) \approx \pi$, tandis que si $z = -1 - i\varepsilon$, $\Arg(z) \approx -\pi$. En revanche, $\Arg$ est continue sur $\C \setminus \R_{\leq 0}$.
\end{remarque}

\subsection{Définition du logarithme complexe}

\begin{definition}[Logarithme complexe (multivalué)]\label{def:log-multivalué}
Pour $z \in \C \setminus \{0\}$, on définit le \emph{logarithme complexe} (multivalué) de $z$ par
\[
  \log(z) = \ln\abs{z} + i\arg(z) = \{\ln\abs{z} + i(\Arg(z) + 2k\pi) : k \in \Z\}.
\]
Chaque valeur de cet ensemble est appelée une \emph{détermination} du logarithme de $z$.
\end{definition}

\begin{definition}[Logarithme principal]\label{def:log-principal}
Le \emph{logarithme principal} est la fonction
\[
  \Log : \C \setminus \R_{\leq 0} \longrightarrow \C, \quad
  \Log(z) = \ln\abs{z} + i\Arg(z),
\]
où $\Arg(z) \in (-\pi,\pi)$ est l'argument principal de $z$
(on exclut $\R_{\leq 0}$ pour assurer la continuité).
\end{definition}

\begin{remarque}\label{rem:log-domaine}
Le domaine $\C \setminus \R_{\leq 0}$ est obtenu en retirant la demi-droite
$\{x \in \R : x \leq 0\}$ du plan complexe. Cette demi-droite est appelée
\emph{coupure} (\textit{branch cut}) du logarithme principal. C'est un choix
conventionnel : on pourrait choisir n'importe quel rayon issu de l'origine comme coupure.
\end{remarque}

\begin{theoreme}[Propriétés du logarithme principal]\label{thm:proprietes-log}
\leavevmode
\begin{enumerate}[label=(\roman*)]
  \item $\Log$ est holomorphe sur $\C \setminus \R_{\leq 0}$, et
  \[
    \Log'(z) = \frac{1}{z} \quad \text{pour tout } z \in \C \setminus \R_{\leq 0}.
  \]
  \item $\exp(\Log(z)) = z$ pour tout $z \in \C \setminus \R_{\leq 0}$.
  \item $\Log(\exp(z)) = z$ pour tout $z$ tel que $-\pi < \im(z) < \pi$.
  \item $\re(\Log(z)) = \ln\abs{z}$ et $\im(\Log(z)) = \Arg(z) \in (-\pi,\pi)$.
  \item L'image de $\Log$ est la bande $\{w \in \C : -\pi < \im(w) < \pi\}$.
\end{enumerate}
\end{theoreme}

\begin{proof}
\emph{(i)} Écrivons $z = re^{i\theta}$ avec $r > 0$ et $\theta = \Arg(z) \in (-\pi,\pi)$. En coordonnées polaires, $\Log(z) = \ln r + i\theta$. Posons $u(r,\theta) = \ln r$ et $v(r,\theta) = \theta$. Les équations de Cauchy-Riemann en coordonnées polaires s'écrivent
\[
  \frac{\partial u}{\partial r} = \frac{1}{r}\frac{\partial v}{\partial \theta}, \qquad
  \frac{\partial v}{\partial r} = -\frac{1}{r}\frac{\partial u}{\partial \theta}.
\]
On a $\partial u/\partial r = 1/r$, $\partial v/\partial \theta = 1$, $\partial v/\partial r = 0$, $\partial u/\partial \theta = 0$, donc les deux équations sont satisfaites. La dérivée complexe est
\[
  \Log'(z) = e^{-i\theta}\left(\frac{\partial u}{\partial r} + i\frac{\partial v}{\partial r}\right) = e^{-i\theta} \cdot \frac{1}{r} = \frac{1}{re^{i\theta}} = \frac{1}{z}.
\]

\emph{(ii)} $\exp(\Log(z)) = \exp(\ln r + i\theta) = r(\cos\theta + i\sin\theta) = re^{i\theta} = z$.

\emph{(iii)} Si $z = x + iy$ avec $-\pi < y < \pi$, alors $\abs{e^z} = e^x$ et $\Arg(e^z) = y$, donc $\Log(e^z) = x + iy = z$.

\emph{(iv)} et \emph{(v)} découlent directement de la définition.
\end{proof}

\begin{remarque}\label{rem:log-pas-morphisme}
Contrairement au logarithme réel, le logarithme principal ne vérifie \emph{pas}
en général $\Log(z_1 z_2) = \Log(z_1) + \Log(z_2)$. On a seulement
\[
  \Log(z_1 z_2) = \Log(z_1) + \Log(z_2) + 2k\pi i
\]
pour un certain $k \in \{-1, 0, 1\}$. L'entier $k$ corrige le fait que la somme
$\Arg(z_1) + \Arg(z_2)$ peut sortir de l'intervalle $(-\pi,\pi]$.
\end{remarque}

\begin{exemple}[Calculs de $\Log(z)$]\label{ex:calculs-log}
\leavevmode
\begin{enumerate}[label=(\alph*)]
  \item $\Log(1) = \ln 1 + i \cdot 0 = 0$.
  \item $\Log(-1)$ n'est pas défini par $\Log$ (car $-1 \in \R_{\leq 0}$), mais la valeur multivaluée est $\log(-1) = i\pi + 2ki\pi = i(2k+1)\pi$.
  \item $\Log(i) = \ln 1 + i\frac{\pi}{2} = \frac{i\pi}{2}$.
  \item $\Log(-i) = \ln 1 + i\bigl(-\frac{\pi}{2}\bigr) = -\frac{i\pi}{2}$.
  \item $\Log(1+i) = \ln\abs{1+i} + i\Arg(1+i) = \ln\sqrt{2} + \frac{i\pi}{4}
         = \frac{1}{2}\ln 2 + \frac{i\pi}{4}$.
  \item $\Log(-1-i) = \ln\sqrt{2} + i\bigl(-\frac{3\pi}{4}\bigr)
         = \frac{1}{2}\ln 2 - \frac{3i\pi}{4}$.
  \item $\Log(e^{3+4i})$. On a $\abs{e^{3+4i}} = e^3$ et $\Arg(e^{3+4i}) = 4 - 2\pi$ (car $4 > \pi$, on soustrait $2\pi$ pour ramener dans $(-\pi,\pi)$). Donc $\Log(e^{3+4i}) = 3 + i(4 - 2\pi)$.
\end{enumerate}
\end{exemple}

\begin{center}
\begin{tikzpicture}[>=Stealth, scale=0.85]
  % ── Plan z (gauche) ──
  \begin{scope}[shift={(-5.5,0)}]
    \node[above] at (0,3.3) {\textbf{Plan $z$}};
    \draw[->] (-3.2,0) -- (3.2,0) node[right] {$\re$};
    \draw[->] (0,-3.2) -- (0,3.2) node[above] {$\im$};

    % Coupure
    \draw[red!70!black, very thick, decorate,
      decoration={zigzag, segment length=4, amplitude=1}]
      (-3,0) -- (0,0);
    \node[red!70!black, below, font=\footnotesize] at (-1.5,-0.3) {coupure $\R_{\leq 0}$};

    % Point z = 1+i
    \fill[blue!70!black] (1.5,1.5) circle (2pt)
      node[above right, font=\footnotesize] {$1+i$};
    % Point z = i
    \fill[orange!80!black] (0,2) circle (2pt)
      node[above right, font=\footnotesize] {$i$};
    % Point z = -i
    \fill[green!60!black] (0,-2) circle (2pt)
      node[below right, font=\footnotesize] {$-i$};

    % Cercle |z| = r
    \draw[purple!60!black, thick, dashed] (0,0) circle (2);
    \node[purple!60!black, font=\footnotesize] at (1.6,1.2) {$\abs{z}=r$};

    % Rayon
    \draw[gray, dashed] (0,0) -- (2,1.5);
    \node[gray, font=\footnotesize] at (1.3,0.5) {$\theta$};
    \draw[gray] (0.5,0) arc[start angle=0, end angle=36.87, radius=0.5];
  \end{scope}

  % ── Flèche ──
  \draw[->, very thick, gray] (-1.2,0) -- (0.5,0)
    node[midway, above, font=\small] {$\Log$};

  % ── Plan w (droite) ──
  \begin{scope}[shift={(5,0)}]
    \node[above] at (0,3.3) {\textbf{Plan $w = \Log(z)$}};
    \draw[->] (-3.2,0) -- (3.2,0) node[right] {$\re$};
    \draw[->] (0,-3.2) -- (0,3.2) node[above] {$\im$};

    % Bande image
    \fill[blue!8] (-2.8,-2.5) rectangle (2.8,2.5);
    \draw[red!70!black, thick, dashed] (-2.8,2.5) -- (2.8,2.5)
      node[right, font=\footnotesize] {$\im = \pi$};
    \draw[red!70!black, thick, dashed] (-2.8,-2.5) -- (2.8,-2.5)
      node[right, font=\footnotesize] {$\im = -\pi$};

    % Points images
    \fill[blue!70!black] ({ln(sqrt(2))*1.5},{0.625*1.5}) circle (2pt)
      node[above right, font=\footnotesize] {$\frac{1}{2}\ln 2 + \frac{i\pi}{4}$};
    \fill[orange!80!black] (0,{0.5*3.14159*2.5/3.14159}) circle (2pt)
      node[right, font=\footnotesize] {$\frac{i\pi}{2}$};
    \fill[green!60!black] (0,{-0.5*3.14159*2.5/3.14159}) circle (2pt)
      node[right, font=\footnotesize] {$-\frac{i\pi}{2}$};

    % Image du cercle = droite verticale
    \draw[purple!60!black, thick, dashed] ({ln(2)*1.5/1.2},-2.5) -- ({ln(2)*1.5/1.2},2.5);
    \node[purple!60!black, right, font=\footnotesize] at ({ln(2)*1.5/1.2+0.05},2) {$\re=\ln r$};
  \end{scope}
\end{tikzpicture}
\captionof{figure}{Le logarithme principal $\Log$ envoie $\C \setminus \R_{\leq 0}$ bijectivement sur la bande $\{-\pi < \im(w) < \pi\}$. Les cercles centrés à l'origine deviennent des segments verticaux.}
\end{center}

\subsection{Branches du logarithme}

\begin{definition}[Branche du logarithme]\label{def:branche-log}
Soit $\Omega \subset \C \setminus \{0\}$ un ouvert connexe. Une \emph{branche du logarithme} sur $\Omega$ est une fonction continue $\ell : \Omega \to \C$ telle que
\[
  \exp(\ell(z)) = z \quad \text{pour tout } z \in \Omega.
\]
\end{definition}

\begin{proposition}[Caractérisation des branches]\label{prop:branches-log-carac}
Soit $\ell$ une branche du logarithme sur un ouvert connexe $\Omega \subset \C \setminus \{0\}$. Alors :
\begin{enumerate}[label=(\roman*)]
  \item $\ell$ est holomorphe sur $\Omega$ et $\ell'(z) = 1/z$.
  \item Toute autre branche du logarithme sur $\Omega$ est de la forme $\ell(z) + 2k\pi i$ pour un $k \in \Z$ fixé (indépendant de $z$).
  \item Pour tout $z \in \Omega$, $\ell(z) = \ln\abs{z} + i\theta(z)$ où $\theta : \Omega \to \R$ est une détermination continue de l'argument sur $\Omega$.
\end{enumerate}
\end{proposition}

\begin{proof}
\emph{(i)} Puisque $\exp(\ell(z)) = z$, on a $\exp'(\ell(z)) \cdot \ell'(z) = 1$ partout où $\ell$ est dérivable (par le théorème d'inversion locale, l'holomorphie de $\ell$ découle de la continuité de $\ell$ et du fait que $\exp' = \exp$ ne s'annule jamais). Donc $\ell'(z) = 1/\exp(\ell(z)) = 1/z$.

\emph{(ii)} Si $\ell_1$ et $\ell_2$ sont deux branches, alors $\exp(\ell_1(z) - \ell_2(z)) = z/z = 1$, donc $\ell_1(z) - \ell_2(z) \in 2\pi i\Z$ pour tout $z$. Par continuité et connexité de $\Omega$, cette différence est constante.

\emph{(iii)} De $\exp(\ell(z)) = z$, on tire $\abs{\exp(\ell(z))} = \abs{z}$, d'où $e^{\re(\ell(z))} = \abs{z}$, soit $\re(\ell(z)) = \ln\abs{z}$. La partie imaginaire $\theta(z) = \im(\ell(z))$ est alors une détermination continue de l'argument.
\end{proof}

\begin{theoreme}[Existence d'une branche du logarithme]\label{thm:existence-branche-log}
Soit $\Omega \subset \C \setminus \{0\}$ un ouvert simplement connexe. Alors il existe une branche du logarithme sur $\Omega$.
\end{theoreme}

\begin{proof}
Fixons $z_0 \in \Omega$ et choisissons $w_0 \in \C$ tel que $e^{w_0} = z_0$. Puisque $\Omega$ est simplement connexe et $1/z$ est holomorphe sur $\Omega$ (car $0 \notin \Omega$), la fonction $1/z$ admet une primitive sur $\Omega$ (par le théorème de Cauchy pour les domaines simplement connexes). Posons
\[
  \ell(z) = w_0 + \int_{z_0}^{z} \frac{1}{\zeta}\,\dd\zeta,
\]
l'intégrale étant calculée le long de n'importe quel chemin dans $\Omega$ reliant $z_0$ à $z$ (le résultat ne dépend pas du chemin par simple connexité). Alors $\ell'(z) = 1/z$. Considérons $g(z) = z \exp(-\ell(z))$. On a
\[
  g'(z) = \exp(-\ell(z)) + z \cdot (-\ell'(z)) \exp(-\ell(z))
         = \exp(-\ell(z))(1 - z/z) = 0.
\]
Donc $g$ est constante sur $\Omega$ : $g(z) = g(z_0) = z_0 e^{-w_0} = 1$. Ainsi $\exp(\ell(z)) = z$ pour tout $z \in \Omega$.
\end{proof}

\begin{exemple}[Branche du logarithme sur $\C \setminus \R_{\geq 0}$]\label{ex:branche-log-coupe-positive}
On peut définir une branche du logarithme sur $\C \setminus \R_{\geq 0}$ en prenant
\[
  \ell(z) = \ln\abs{z} + i\theta(z), \quad \theta(z) \in (0, 2\pi).
\]
Cette branche est holomorphe sur $\C \setminus \R_{\geq 0}$, et sa partie imaginaire prend ses valeurs dans $(0, 2\pi)$ au lieu de $(-\pi, \pi)$ pour $\Log$.
Par exemple, $\ell(-1) = \ln 1 + i\pi = i\pi$ et $\ell(-i) = \ln 1 + i\frac{3\pi}{2} = \frac{3i\pi}{2}$, alors que $\Log(-i) = -\frac{i\pi}{2}$.
\end{exemple}

\subsection{Intuition sur les surfaces de Riemann}\label{subsec:surface-riemann-log}

Le logarithme complexe est intrinsèquement multivalué : à chaque $z \neq 0$ correspond une infinité de valeurs $\ln\abs{z} + i(\Arg(z) + 2k\pi)$, $k \in \Z$. Pour rendre le logarithme uniforme, Riemann a proposé de remplacer le plan $\C \setminus \{0\}$ par une \emph{surface de Riemann} sur laquelle les différentes déterminations se <<\,raccordent\,>> de manière naturelle.

\begin{center}
\begin{tikzpicture}[>=Stealth, scale=0.9]
  % Feuillet k=-1 (en bas)
  \begin{scope}[shift={(0,-2.5)}, opacity=0.6]
    \fill[green!15] (-2.5,-1) -- (2.5,-1) -- (3,0.5) -- (-2,0.5) -- cycle;
    \draw[green!60!black, thick] (-2.5,-1) -- (2.5,-1) -- (3,0.5) -- (-2,0.5) -- cycle;
    \draw[green!60!black, thick, dashed] (0.25,-0.25) -- (3,0.15);
    \node[green!60!black, font=\footnotesize] at (-1.5,0) {$k = -1$};
    \node[green!60!black, font=\footnotesize, right] at (3,0) {coupure};
  \end{scope}

  % Feuillet k=0 (au milieu)
  \begin{scope}[shift={(0,0)}, opacity=0.8]
    \fill[blue!15] (-2.5,-1) -- (2.5,-1) -- (3,0.5) -- (-2,0.5) -- cycle;
    \draw[blue!70!black, thick] (-2.5,-1) -- (2.5,-1) -- (3,0.5) -- (-2,0.5) -- cycle;
    \draw[blue!70!black, thick, dashed] (0.25,-0.25) -- (3,0.15);
    \node[blue!70!black, font=\footnotesize] at (-1.5,0) {$k = 0$ (principal)};
  \end{scope}

  % Feuillet k=1 (en haut)
  \begin{scope}[shift={(0,2.5)}, opacity=0.6]
    \fill[red!15] (-2.5,-1) -- (2.5,-1) -- (3,0.5) -- (-2,0.5) -- cycle;
    \draw[red!70!black, thick] (-2.5,-1) -- (2.5,-1) -- (3,0.5) -- (-2,0.5) -- cycle;
    \draw[red!70!black, thick, dashed] (0.25,-0.25) -- (3,0.15);
    \node[red!70!black, font=\footnotesize] at (-1.5,0) {$k = 1$};
  \end{scope}

  % Flèches de connexion
  \draw[->, thick, gray, dashed] (2.8,-1.8) to[out=60, in=-60] (2.8,-0.7);
  \draw[->, thick, gray, dashed] (2.8,0.7) to[out=60, in=-60] (2.8,1.8);

  % Points de branchement
  \fill[black] (0.25,-0.25) circle (2pt);
  \fill[black] (0.25,-0.25+2.5) circle (2pt);
  \fill[black] (0.25,-0.25-2.5) circle (2pt);

  % Étiquettes
  \node[right, font=\small] at (3.5,0) {%
    \begin{minipage}{4.5cm}
    Chaque feuillet porte une\\
    détermination de $\log(z)$.\\
    Traverser la coupure fait\\
    passer au feuillet suivant\\
    (ajout de $2\pi i$).
    \end{minipage}};

  % Vdots
  \node at (0,4.2) {$\vdots$};
  \node at (0,-4.2) {$\vdots$};
\end{tikzpicture}
\captionof{figure}{Intuition de la surface de Riemann du logarithme : une infinité de feuillets, chacun portant une branche, raccordés le long de la coupure.}
\end{center}

\begin{remarque}\label{rem:surface-riemann-intuition}
La surface de Riemann de $\log$ est un revêtement universel de $\C \setminus \{0\}$. C'est une surface connexe, simplement connexe, homéomorphe à $\C$, sur laquelle le logarithme devient une bijection holomorphe. Le point $z = 0$ est un \emph{point de branchement} : en tournant autour de $0$, on passe d'un feuillet au suivant. L'étude rigoureuse des surfaces de Riemann dépasse le cadre de ce cours, mais cette vision géométrique éclaire la nature multivaluée du logarithme.
\end{remarque}

%=======================================================================
\section{Puissances complexes}\label{sec:puissances}
%=======================================================================

\begin{definition}[Puissance complexe]\label{def:puissance-complexe}
Pour $z \in \C \setminus \{0\}$ et $\alpha \in \C$, on définit la \emph{puissance complexe} (multivaluée)
\[
  z^\alpha = \exp(\alpha \log z) = \exp\bigl(\alpha(\ln\abs{z} + i\arg(z))\bigr).
\]
La \emph{détermination principale} de $z^\alpha$ est
\[
  z^\alpha_{\mathrm{princ}} = \exp(\alpha \Log z), \quad z \in \C \setminus \R_{\leq 0}.
\]
\end{definition}

\begin{remarque}\label{rem:puissance-multivaluee}
Les différentes valeurs de $z^\alpha$ sont
\[
  z^\alpha = e^{\alpha \Log(z)} \cdot e^{2k\pi i \alpha}, \quad k \in \Z.
\]
\begin{itemize}
  \item Si $\alpha = n \in \Z$, alors $e^{2k\pi i n} = 1$ pour tout $k$, donc $z^n$ est univaluée (comme attendu).
  \item Si $\alpha = p/q \in \Q$ (fraction irréductible, $q \geq 2$), alors $e^{2k\pi i p/q}$ prend exactement $q$ valeurs distinctes. On retrouve les $q$ racines $q$-ièmes.
  \item Si $\alpha \notin \Q$, alors $e^{2k\pi i \alpha}$ prend une infinité de valeurs distinctes.
\end{itemize}
\end{remarque}

\begin{proposition}[Holomorphie de la puissance principale]\label{prop:puissance-holomorphe}
Pour tout $\alpha \in \C$, la détermination principale $z \mapsto \exp(\alpha \Log z)$ est holomorphe sur $\C \setminus \R_{\leq 0}$, et
\[
  \frac{\dd}{\dd z}\bigl[\exp(\alpha \Log z)\bigr] = \frac{\alpha}{z}\exp(\alpha \Log z) = \alpha z^{\alpha - 1},
\]
où $z^{\alpha-1}$ désigne la détermination principale.
\end{proposition}

\begin{proof}
C'est la composée de fonctions holomorphes : $z \mapsto \Log z \mapsto \alpha \Log z \mapsto \exp(\alpha \Log z)$. Par la règle de la chaîne,
\[
  \frac{\dd}{\dd z}\exp(\alpha \Log z) = \exp(\alpha \Log z) \cdot \alpha \cdot \frac{1}{z} = \frac{\alpha}{z}\exp(\alpha \Log z).
\]
Or $\exp(\alpha \Log z)/z = \exp(\alpha \Log z) \cdot \exp(-\Log z) = \exp((\alpha-1)\Log z) = z^{\alpha-1}$.
\end{proof}

\begin{exemple}[Calcul de $i^i$]\label{ex:i-puissance-i}
On cherche toutes les valeurs de $i^i$. On a
\[
  i^i = \exp(i \log i) = \exp\!\left(i\left(\ln 1 + i\left(\frac{\pi}{2} + 2k\pi\right)\right)\right)
     = \exp\!\left(-\frac{\pi}{2} - 2k\pi\right), \quad k \in \Z.
\]
Toutes les valeurs sont donc \emph{réelles et positives}. La valeur principale (correspondant à $k = 0$) est
\[
  i^i_{\mathrm{princ}} = e^{-\pi/2} \approx 0{,}2079.
\]
C'est un résultat remarquable d'Euler : $i^i$ est un nombre réel.
\end{exemple}

\begin{exemple}[Calcul de $(-1)^{\sqrt{2}}$]\label{ex:moins-un-puissance-sqrt2}
Les valeurs de $(-1)^{\sqrt{2}}$ sont
\[
  (-1)^{\sqrt{2}} = \exp(\sqrt{2}\,\log(-1)) = \exp\bigl(\sqrt{2}\,i(2k+1)\pi\bigr)
  = e^{i(2k+1)\sqrt{2}\,\pi}, \quad k \in \Z.
\]
Puisque $\sqrt{2}$ est irrationnel, les nombres $(2k+1)\sqrt{2}$ sont tous distincts modulo $2$, donc les valeurs de $(-1)^{\sqrt{2}}$ forment un sous-ensemble dense du cercle unité $\T$.
\end{exemple}

\begin{exemple}[Racine carrée principale]\label{ex:racine-carree}
La détermination principale de $z^{1/2}$ est
\[
  z^{1/2} = \exp\!\left(\frac{1}{2}\Log z\right)
           = \exp\!\left(\frac{1}{2}\ln\abs{z} + \frac{i}{2}\Arg(z)\right)
           = \sqrt{\abs{z}}\,e^{i\Arg(z)/2}.
\]
C'est une fonction holomorphe sur $\C \setminus \R_{\leq 0}$, dont la dérivée est $\frac{1}{2}z^{-1/2}$.

Par exemple, $(-4)^{1/2}$ au sens de la valeur principale n'est pas défini (car $-4 \in \R_{\leq 0}$), mais au sens multivalué on a $(-4)^{1/2} = \pm 2i$.
\end{exemple}

\begin{exemple}[Fonction $z^{1/3}$]\label{ex:racine-cubique}
Les trois valeurs de $z^{1/3}$ pour $z = -8$ sont
\[
  (-8)^{1/3} = \exp\!\left(\frac{1}{3}\log(-8)\right)
             = \exp\!\left(\frac{1}{3}\bigl(\ln 8 + i(2k+1)\pi\bigr)\right)
             = 2\exp\!\left(\frac{i(2k+1)\pi}{3}\right), \quad k = 0, 1, 2.
\]
Cela donne :
\begin{itemize}
  \item $k=0$ : $2e^{i\pi/3} = 2\bigl(\frac{1}{2} + \frac{i\sqrt{3}}{2}\bigr) = 1 + i\sqrt{3}$,
  \item $k=1$ : $2e^{i\pi} = -2$,
  \item $k=2$ : $2e^{i5\pi/3} = 1 - i\sqrt{3}$.
\end{itemize}
La racine cubique réelle $-2$ est bien l'une des trois déterminations.
\end{exemple}

%=======================================================================
\section{Fonctions trigonométriques complexes}\label{sec:trig}
%=======================================================================

\subsection{Définitions via l'exponentielle}

\begin{definition}[Cosinus et sinus complexes]\label{def:cos-sin-complexes}
Pour $z \in \C$, on définit
\[
  \cos z = \frac{e^{iz} + e^{-iz}}{2}, \qquad
  \sin z = \frac{e^{iz} - e^{-iz}}{2i}.
\]
\end{definition}

\begin{definition}[Tangente et cotangente complexes]\label{def:tan-cot-complexes}
Pour $z \in \C$ tel que $\cos z \neq 0$ (resp.\ $\sin z \neq 0$),
\[
  \tan z = \frac{\sin z}{\cos z}, \qquad
  \cot z = \frac{\cos z}{\sin z}.
\]
\end{definition}

\begin{theoreme}[Propriétés des fonctions trigonométriques complexes]\label{thm:proprietes-trig}
\leavevmode
\begin{enumerate}[label=(\roman*)]
  \item \textbf{Holomorphie.} $\cos$ et $\sin$ sont des fonctions entières, avec
  \[
    \cos'(z) = -\sin(z), \qquad \sin'(z) = \cos(z).
  \]
  \item \textbf{Formule d'Euler.} Pour tout $z \in \C$,
  \[
    e^{iz} = \cos z + i\sin z.
  \]
  \item \textbf{Identité fondamentale.} $\cos^2 z + \sin^2 z = 1$.
  \item \textbf{Périodicité.} $\cos$ et $\sin$ sont $2\pi$-périodiques.
  \item \textbf{Zéros.} $\cos z = 0 \iff z = \frac{\pi}{2} + k\pi$ ($k \in \Z$) et $\sin z = 0 \iff z = k\pi$ ($k \in \Z$).
  \item \textbf{Formules d'addition.} Pour tous $z, w \in \C$,
  \begin{align*}
    \cos(z+w) &= \cos z \cos w - \sin z \sin w, \\
    \sin(z+w) &= \sin z \cos w + \cos z \sin w.
  \end{align*}
  \item \textbf{Non-bornitude.} Contrairement au cas réel, $\cos$ et $\sin$ ne sont \emph{pas} bornées sur $\C$.
\end{enumerate}
\end{theoreme}

\begin{proof}
\emph{(i)} La dérivation de $\cos z = (e^{iz} + e^{-iz})/2$ donne
$\cos'(z) = (ie^{iz} - ie^{-iz})/2 = -\sin z$. De même pour $\sin$.

\emph{(ii)} $\cos z + i\sin z = \frac{e^{iz}+e^{-iz}}{2} + i\cdot\frac{e^{iz}-e^{-iz}}{2i} = \frac{e^{iz}+e^{-iz}}{2} + \frac{e^{iz}-e^{-iz}}{2} = e^{iz}$.

\emph{(iii)} $\cos^2 z + \sin^2 z = \frac{(e^{iz}+e^{-iz})^2}{4} + \frac{(e^{iz}-e^{-iz})^2}{-4} = \frac{e^{2iz}+2+e^{-2iz}}{4} + \frac{-(e^{2iz}-2+e^{-2iz})}{4} = \frac{4}{4} = 1$.

\emph{(iv)} Découle de la $2\pi i$-périodicité de $\exp$ : $\cos(z+2\pi) = \frac{e^{i(z+2\pi)}+e^{-i(z+2\pi)}}{2} = \frac{e^{iz}+e^{-iz}}{2} = \cos z$.

\emph{(v)} $\sin z = 0 \iff e^{iz} = e^{-iz} \iff e^{2iz} = 1 \iff 2iz \in 2\pi i\Z \iff z \in \pi\Z$. Pour $\cos z = 0$, on utilise $\cos z = \sin(z + \pi/2)$.

\emph{(vi)} Calcul direct à partir des exponentielles et de la propriété fonctionnelle.

\emph{(vii)} On a $\sin(iy) = \frac{e^{-y}-e^y}{2i} = \frac{i(e^y - e^{-y})}{2} = i\sinh y$. Quand $y \to +\infty$, $\abs{\sin(iy)} = \sinh y \to +\infty$.
\end{proof}

\begin{proposition}[Parties réelle et imaginaire]\label{prop:trig-re-im}
Pour $z = x + iy$ avec $x, y \in \R$ :
\begin{align}
  \cos(x+iy) &= \cos x \cosh y - i \sin x \sinh y, \label{eq:cos-re-im} \\
  \sin(x+iy) &= \sin x \cosh y + i \cos x \sinh y. \label{eq:sin-re-im}
\end{align}
En conséquence :
\[
  \abs{\cos z}^2 = \cos^2 x + \sinh^2 y, \qquad
  \abs{\sin z}^2 = \sin^2 x + \sinh^2 y.
\]
\end{proposition}

\begin{proof}
Par les formules d'addition :
\[
  \cos(x+iy) = \cos x \cos(iy) - \sin x \sin(iy).
\]
Or $\cos(iy) = \frac{e^{-y}+e^y}{2} = \cosh y$ et $\sin(iy) = \frac{e^{-y}-e^y}{2i} = i\sinh y$. D'où le résultat pour le cosinus ; le cas du sinus est analogue.

Pour les modules, $\abs{\cos z}^2 = \cos^2 x \cosh^2 y + \sin^2 x \sinh^2 y = \cos^2 x(1+\sinh^2 y) + \sin^2 x \sinh^2 y = \cos^2 x + \sinh^2 y$.
\end{proof}

\begin{remarque}\label{rem:trig-non-bornee}
La formule $\abs{\sin z}^2 = \sin^2 x + \sinh^2 y$ montre que sur la droite $x = \pi/2$ (i.e.\ $z = \pi/2 + iy$), $\abs{\cos z} = \sinh\abs{y}$, qui croît exponentiellement. Les fonctions trigonométriques complexes sont non bornées et prennent toutes les valeurs complexes (sauf peut-être une, par le petit théorème de Picard).
\end{remarque}

%=======================================================================
\section{Fonctions hyperboliques complexes}\label{sec:hyp}
%=======================================================================

\begin{definition}[Cosinus et sinus hyperboliques]\label{def:cosh-sinh}
Pour $z \in \C$, on définit
\[
  \cosh z = \frac{e^z + e^{-z}}{2}, \qquad
  \sinh z = \frac{e^z - e^{-z}}{2}.
\]
De même, $\tanh z = \sinh z / \cosh z$ pour $\cosh z \neq 0$.
\end{definition}

\begin{theoreme}[Propriétés des fonctions hyperboliques]\label{thm:proprietes-hyp}
\leavevmode
\begin{enumerate}[label=(\roman*)]
  \item $\cosh$ et $\sinh$ sont des fonctions entières, avec $\cosh'(z) = \sinh(z)$ et $\sinh'(z) = \cosh(z)$.
  \item $\cosh^2 z - \sinh^2 z = 1$.
  \item $\cosh$ et $\sinh$ sont $2\pi i$-périodiques.
  \item $\cosh z = 0 \iff z = i(\pi/2 + k\pi)$ et $\sinh z = 0 \iff z = ik\pi$, $k \in \Z$.
\end{enumerate}
\end{theoreme}

\begin{proof}
Les preuves sont analogues à celles du 
ef{thm:proprietes-trig}, en utilisant les définitions exponentielles.

\emph{(ii)} $\cosh^2 z - \sinh^2 z = \frac{(e^z+e^{-z})^2 - (e^z-e^{-z})^2}{4} = \frac{4}{4} = 1$.

\emph{(iv)} $\sinh z = 0 \iff e^z = e^{-z} \iff e^{2z} = 1 \iff z \in i\pi\Z$.
\end{proof}

%=======================================================================
\section{Lien entre fonctions trigonométriques et hyperboliques}\label{sec:lien-trig-hyp}
%=======================================================================

\begin{theoreme}[Relations fondamentales]\label{thm:relations-trig-hyp}
Pour tout $z \in \C$ :
\begin{align}
  \cos(iz) &= \cosh(z), \label{eq:cos-iz} \\
  \sin(iz) &= i\sinh(z), \label{eq:sin-iz} \\
  \cosh(iz) &= \cos(z), \label{eq:cosh-iz} \\
  \sinh(iz) &= i\sin(z). \label{eq:sinh-iz}
\end{align}
\end{theoreme}

\begin{proof}
Calcul direct :
\[
  \cos(iz) = \frac{e^{i(iz)} + e^{-i(iz)}}{2} = \frac{e^{-z} + e^{z}}{2} = \cosh z.
\]
\[
  \sin(iz) = \frac{e^{i(iz)} - e^{-i(iz)}}{2i} = \frac{e^{-z} - e^z}{2i}
           = \frac{-(e^z - e^{-z})}{2i} = \frac{i(e^z - e^{-z})}{2} = i\sinh z.
\]
Les deux autres s'obtiennent en substituant $iz$ par $z$ dans les premières.
\end{proof}

\begin{remarque}\label{rem:unification-trig-hyp}
Ces relations montrent que la distinction entre fonctions trigonométriques et
hyperboliques disparaît dans le cadre complexe : les deux familles ne sont que
des <<\,rotations\,>> l'une de l'autre dans le plan complexe. C'est l'une des
simplifications conceptuelles majeures apportées par l'analyse complexe.
\end{remarque}

\begin{corollaire}[Formules d'Euler généralisées]\label{cor:euler-generalisees}
Pour tout $z \in \C$ :
\[
  e^z = \cosh z + \sinh z, \qquad
  e^{-z} = \cosh z - \sinh z.
\]
Plus généralement, pour $z = x + iy$ :
\[
  e^z = e^x\cos y + ie^x\sin y = (\cosh x + \sinh x)(\cos y + i\sin y).
\]
\end{corollaire}

\begin{proof}
La première ligne découle de $\cosh z + \sinh z = \frac{e^z+e^{-z}}{2} + \frac{e^z-e^{-z}}{2} = e^z$. La seconde est la définition de $\exp$.
\end{proof}

%=======================================================================
\section{Tableau récapitulatif}\label{sec:tableau-recap}
%=======================================================================

Le tableau suivant résume les principales propriétés des fonctions élémentaires.

\begin{center}
\renewcommand{\arraystretch}{1.6}
\begin{tabular}{@{}lcccc@{}}
\toprule
\textbf{Fonction} & \textbf{Domaine} & \textbf{Période} & \textbf{Zéros} & \textbf{Dérivée} \\
\midrule
$e^z$ & $\C$ & $2\pi i$ & aucun & $e^z$ \\
$\Log z$ & $\C \setminus \R_{\leq 0}$ & --- & $z = 1$ & $1/z$ \\
$z^\alpha$ (princ.) & $\C \setminus \R_{\leq 0}$ & --- & aucun ($\alpha \neq 0$) & $\alpha z^{\alpha-1}$ \\
$\cos z$ & $\C$ & $2\pi$ & $\frac{\pi}{2} + k\pi$ & $-\sin z$ \\
$\sin z$ & $\C$ & $2\pi$ & $k\pi$ & $\cos z$ \\
$\cosh z$ & $\C$ & $2\pi i$ & $i(\frac{\pi}{2}+k\pi)$ & $\sinh z$ \\
$\sinh z$ & $\C$ & $2\pi i$ & $ik\pi$ & $\cosh z$ \\
\bottomrule
\end{tabular}
\end{center}

%=======================================================================
\section{Exemples détaillés}\label{sec:exemples-detaillés}
%=======================================================================

\begin{exemple}[Détermination de branches sur un domaine à trou]\label{ex:branche-domaine-trou}
Soit $\Omega = \{z \in \C : 1 < \abs{z} < 3\}$ la couronne ouverte centrée à l'origine.
Ce domaine n'est \emph{pas} simplement connexe (le cercle $\abs{z} = 2$ n'est pas
contractible dans $\Omega$). En conséquence, il n'existe pas de branche du logarithme
sur $\Omega$ tout entier.

En effet, supposons qu'une telle branche $\ell$ existe. Alors $\ell'(z) = 1/z$ et
\[
  \int_{\abs{z}=2} \frac{1}{z}\,\dd z = \int_{\abs{z}=2} \ell'(z)\,\dd z = \ell(z)\Big|_{\text{tour complet}} = 0,
\]
car $\ell$ est univaluée. Or, par calcul direct,
$\int_{\abs{z}=2} \frac{\dd z}{z} = 2\pi i \neq 0$. Contradiction.

En revanche, si l'on coupe la couronne le long d'un rayon, par exemple
$\Omega' = \Omega \setminus \R_{\leq 0}$, le domaine obtenu est simplement connexe et
admet des branches du logarithme.
\end{exemple}

\begin{exemple}[Calcul de $(1+i)^{3-i}$]\label{ex:puissance-complexe-detaillee}
Calculons la valeur principale de $(1+i)^{3-i}$.

\textbf{Étape 1 : $\Log(1+i)$.}
$\abs{1+i} = \sqrt{2}$, $\Arg(1+i) = \pi/4$. Donc
$\Log(1+i) = \frac{1}{2}\ln 2 + \frac{i\pi}{4}$.

\textbf{Étape 2 : $\alpha \Log z$.}
$(3-i)\Log(1+i) = (3-i)\left(\frac{1}{2}\ln 2 + \frac{i\pi}{4}\right)$.

Développons :
\begin{align*}
  (3-i)\left(\frac{\ln 2}{2} + \frac{i\pi}{4}\right)
  &= \frac{3\ln 2}{2} + \frac{3i\pi}{4} - \frac{i\ln 2}{2} - \frac{i^2\pi}{4} \\
  &= \left(\frac{3\ln 2}{2} + \frac{\pi}{4}\right)
     + i\left(\frac{3\pi}{4} - \frac{\ln 2}{2}\right).
\end{align*}

\textbf{Étape 3 : Exponentielle.}
\[
  (1+i)^{3-i}_{\mathrm{princ}} = \exp\!\left[\left(\frac{3\ln 2}{2} + \frac{\pi}{4}\right)
  + i\left(\frac{3\pi}{4} - \frac{\ln 2}{2}\right)\right].
\]
Le module est
$\exp\!\left(\frac{3\ln 2}{2} + \frac{\pi}{4}\right) = 2^{3/2}\,e^{\pi/4} \approx 6{,}168$.

L'argument est $\frac{3\pi}{4} - \frac{\ln 2}{2} \approx 2{,}009$ radians.
\end{exemple}

\begin{exemple}[Résolution de $\cos z = 2$]\label{ex:cos-egal-2}
Résolvons $\cos z = 2$. On a
\[
  \frac{e^{iz}+e^{-iz}}{2} = 2, \quad\text{soit}\quad e^{iz} + e^{-iz} = 4.
\]
En posant $w = e^{iz}$, on obtient $w + w^{-1} = 4$, soit $w^2 - 4w + 1 = 0$.
Le discriminant est $16 - 4 = 12$, d'où
\[
  w = \frac{4 \pm 2\sqrt{3}}{2} = 2 \pm \sqrt{3}.
\]
Les deux valeurs sont réelles positives. Alors $e^{iz} = 2 \pm \sqrt{3}$, d'où
$iz = \ln(2 \pm \sqrt{3}) + 2k\pi i$, et finalement
\[
  z = 2k\pi - i\ln(2 \pm \sqrt{3}), \quad k \in \Z.
\]
Puisque $2 - \sqrt{3} = 1/(2+\sqrt{3})$, on a $\ln(2-\sqrt{3}) = -\ln(2+\sqrt{3})$,
et les solutions se réduisent à
\[
  z = 2k\pi \pm i\ln(2+\sqrt{3}), \quad k \in \Z.
\]
Remarquons que toutes ces solutions sont non réelles, confirmant que $\cos$ prend
des valeurs en dehors de $[-1,1]$ pour des arguments complexes.
\end{exemple}

\begin{exemple}[Résolution de $\sinh z = i$]\label{ex:sinh-egal-i}
On résout $\sinh z = i$, c'est-à-dire $\frac{e^z - e^{-z}}{2} = i$. En posant $w = e^z$, on obtient $w - w^{-1} = 2i$, soit $w^2 - 2iw - 1 = 0$. Le discriminant est $-4 + 4 = 0$, donc $w = i$ (racine double). Alors $e^z = i$, d'où
\[
  z = \log(i) = i\left(\frac{\pi}{2} + 2k\pi\right), \quad k \in \Z.
\]
\end{exemple}

%=======================================================================
\section{Transformations géométriques}\label{sec:transformations-geometriques}
%=======================================================================

\begin{proposition}[Image du demi-plan par $\sin$]\label{prop:sin-demi-plan}
La fonction $\sin$ envoie la bande $\{z \in \C : -\pi/2 < \re(z) < \pi/2\}$
bijectivement sur $\C \setminus ((-\infty,-1] \cup [1,+\infty))$.
\end{proposition}

\begin{proof}
Posons $z = x + iy$ avec $-\pi/2 < x < \pi/2$ et $y \in \R$. Par la

ef{prop:trig-re-im},
\[
  \sin z = \sin x \cosh y + i\cos x \sinh y.
\]
Soit $w = u + iv$ l'image. Alors $u = \sin x \cosh y$ et $v = \cos x \sinh y$.

Pour $y \neq 0$, on peut écrire
\[
  \frac{u^2}{\cosh^2 y} + \frac{v^2}{\sinh^2 y}
  = \sin^2 x + \cos^2 x = 1,
\]
c'est-à-dire que les images des droites horizontales $\im(z) = y_0$ (avec $y_0 \neq 0$, $-\pi/2 < \re(z) < \pi/2$) sont des demi-ellipses.

De même, pour $x$ fixé avec $-\pi/2 < x < \pi/2$ et $x \neq 0$,
\[
  \frac{u^2}{\sin^2 x} - \frac{v^2}{\cos^2 x} = \cosh^2 y - \sinh^2 y = 1,
\]
donc les images des droites verticales sont des branches d'hyperboles.

L'injectivité se vérifie par un argument d'unicité, et la surjectivité par résolution
de $\sin z = w$ pour $w \notin (-\infty,-1] \cup [1,+\infty)$.
\end{proof}

\begin{center}
\begin{tikzpicture}[>=Stealth, scale=0.78]
  % ── Plan z ──
  \begin{scope}[shift={(-5.5,0)}]
    \node[above] at (0,3.5) {\textbf{Plan $z$}};
    \draw[->] (-3,0) -- (3,0) node[right] {$\re$};
    \draw[->] (0,-3.2) -- (0,3.2) node[above] {$\im$};

    % Bande
    \fill[cyan!10] (-1.57,-3) rectangle (1.57,3);
    \draw[cyan!70!black, thick, dashed] (-1.57,-3) -- (-1.57,3)
      node[above, font=\footnotesize] {$-\frac{\pi}{2}$};
    \draw[cyan!70!black, thick, dashed] (1.57,-3) -- (1.57,3)
      node[above, font=\footnotesize] {$\frac{\pi}{2}$};

    % Droites horizontales
    \foreach \y/\col in {-2/green!60!black, -1/orange!80!black, 0/red!70!black,
                          1/blue!70!black, 2/purple!70!black} {
      \draw[\col, thick] (-1.4,\y) -- (1.4,\y);
    }

    % Droites verticales
    \foreach \x/\col in {-1/gray, -0.5/gray!60, 0.5/gray!60, 1/gray} {
      \draw[\col, thick, densely dotted] (\x,-2.8) -- (\x,2.8);
    }
  \end{scope}

  % ── Flèche ──
  \draw[->, very thick, gray] (-1.2,0) -- (0.5,0)
    node[midway, above, font=\small] {$\sin$};

  % ── Plan w ──
  \begin{scope}[shift={(5.5,0)}]
    \node[above] at (0,3.5) {\textbf{Plan $w = \sin z$}};
    \draw[->] (-3.2,0) -- (3.2,0) node[right] {$\re$};
    \draw[->] (0,-3.2) -- (0,3.2) node[above] {$\im$};

    % Coupures
    \draw[cyan!70!black, very thick, decorate,
      decoration={zigzag, segment length=4, amplitude=1}]
      (-3,0) -- (-1,0);
    \draw[cyan!70!black, very thick, decorate,
      decoration={zigzag, segment length=4, amplitude=1}]
      (1,0) -- (3,0);

    % Image de y=0 : segment [-1,1]
    \draw[red!70!black, very thick] (-1,0) -- (1,0);

    % Ellipses (images des horizontales)
    \draw[blue!70!black, thick]
      (0,0) ellipse ({cosh(1)*sin(70)*(1/1)} and {sinh(1)*cos(70)*(1/0.5)});

    \draw[purple!70!black, thick]
      (0,0) ellipse ({cosh(2)*sin(70)*1/1} and {sinh(2)*cos(70)*1/0.5});

    % Hyperbole (image d'une verticale) - approximation
    \draw[gray, thick, densely dotted, domain=-2.5:2.5, samples=50]
      plot ({sin(30)*cosh(\x)}, {cos(30)*sinh(\x)});
    \draw[gray, thick, densely dotted, domain=-2.5:2.5, samples=50]
      plot ({-sin(30)*cosh(\x)}, {-cos(30)*sinh(\x)});
  \end{scope}
\end{tikzpicture}
\captionof{figure}{Image de la bande $\{-\pi/2 < \re(z) < \pi/2\}$ par $\sin$. Les droites horizontales deviennent des ellipses, les droites verticales des hyperboles.}
\end{center}

%=======================================================================
\section{Compléments sur les branches}\label{sec:complements-branches}
%=======================================================================

\begin{definition}[Branche d'une puissance]\label{def:branche-puissance}
Soit $\Omega \subset \C \setminus \{0\}$ un ouvert connexe et $\alpha \in \C$. Une \emph{branche de $z^\alpha$} sur $\Omega$ est une fonction continue $f : \Omega \to \C$ telle qu'il existe une branche du logarithme $\ell$ sur $\Omega$ avec
\[
  f(z) = \exp(\alpha\,\ell(z)) \quad \text{pour tout } z \in \Omega.
\]
\end{definition}

\begin{proposition}[Branches de la racine carrée]\label{prop:branches-racine}
Soit $\Omega \subset \C \setminus \{0\}$ un ouvert simplement connexe. Alors il existe exactement deux branches de $z^{1/2}$ sur $\Omega$, et elles sont opposées l'une de l'autre.
\end{proposition}

\begin{proof}
Par le 
ef{thm:existence-branche-log}, il existe une branche du logarithme $\ell$ sur $\Omega$. Les branches de $z^{1/2}$ sont $f_k(z) = \exp(\frac{1}{2}(\ell(z) + 2k\pi i))$ pour $k \in \Z$. Or $\exp(\frac{1}{2} \cdot 2k\pi i) = \exp(k\pi i) = (-1)^k$. Donc $f_k = f_0$ si $k$ est pair, et $f_k = -f_0$ si $k$ est impair. Il n'y a que deux branches distinctes : $f_0$ et $-f_0$.
\end{proof}

\begin{exemple}[Branche de $\sqrt{z^2-1}$]\label{ex:branche-sqrt-z2-1}
Considérons la fonction multivaluée $\sqrt{z^2 - 1} = \sqrt{(z-1)(z+1)}$. Les points de branchement sont $z = 1$ et $z = -1$. On peut choisir comme coupure le segment $[-1,1]$ sur l'axe réel. Sur le domaine $\Omega = \C \setminus [-1,1]$, qui est simplement connexe, il existe deux branches de $\sqrt{z^2-1}$, holomorphes sur $\Omega$.

La branche principale est déterminée par la condition que $\sqrt{z^2-1} > 0$ pour $z > 1$ réel. On peut l'écrire
\[
  \sqrt{z^2-1} = \exp\!\left(\frac{1}{2}\Log(z^2-1)\right),
\]
mais cette formule n'est valable que pour $z^2-1 \notin \R_{\leq 0}$. Une écriture plus correcte est
\[
  \sqrt{z^2-1} = \exp\!\left(\frac{1}{2}\bigl(\ell_1(z-1) + \ell_2(z+1)\bigr)\right),
\]
où $\ell_1$ et $\ell_2$ sont des branches du logarithme convenablement choisies sur $\Omega$, compatibles avec la coupure $[-1,1]$.
\end{exemple}

%=======================================================================
\section{Exercices}\label{sec:exercices-ch3}
%=======================================================================

\begin{exercice}[Calculs élémentaires]\label{exo:calculs-exp-log}
Calculer :
\begin{enumerate}[label=(\alph*)]
  \item $e^{2+i\pi/3}$, $e^{-1+3\pi i}$, $e^{i\pi/6}$.
  \item $\Log(2i)$, $\Log(-3)$, $\Log(1-i)$.
  \item Toutes les valeurs de $\log(-2i)$.
\end{enumerate}
\end{exercice}

\begin{exercice}[Équations exponentielles]\label{exo:equations-exp}
Résoudre dans $\C$ :
\begin{enumerate}[label=(\alph*)]
  \item $e^z = 1 + i\sqrt{3}$.
  \item $e^z = -e$.
  \item $e^{2z} - 3e^z + 2 = 0$.
\end{enumerate}
\end{exercice}

\begin{exercice}[Puissances complexes]\label{exo:puissances}
\begin{enumerate}[label=(\alph*)]
  \item Calculer toutes les valeurs de $2^i$. Sont-elles réelles ?
  \item Déterminer la valeur principale de $(1-i)^{1+i}$.
  \item Montrer que $i^i = e^{-\pi/2}$ est la valeur principale et que toutes les valeurs sont réelles.
  \item Trouver toutes les valeurs de $(-1)^i$.
\end{enumerate}
\end{exercice}

\begin{exercice}[Équations trigonométriques]\label{exo:equations-trig}
Résoudre dans $\C$ :
\begin{enumerate}[label=(\alph*)]
  \item $\sin z = 3$.
  \item $\cos z = i$.
  \item $\cosh z = 0$.
  \item $\tan z = 2i$.
\end{enumerate}
\end{exercice}

\begin{exercice}[Images géométriques]\label{exo:images-geometriques}
\begin{enumerate}[label=(\alph*)]
  \item Déterminer l'image par $\exp$ de la bande $\{z : 0 < \im(z) < \pi/4\}$.
  \item Déterminer l'image par $\exp$ du rectangle $\{z : -1 < \re(z) < 1,\; \pi/2 < \im(z) < \pi\}$.
  \item Quelle est l'image du demi-plan $\re(z) > 0$ par $z \mapsto \Log z$ ?
  \item Déterminer l'image du disque $\abs{z-1} < 1$ par $\Log$.
\end{enumerate}
\end{exercice}

\begin{exercice}[Identités]\label{exo:identites}
Démontrer les identités suivantes pour tous $z, w \in \C$ :
\begin{enumerate}[label=(\alph*)]
  \item $\cos^2 z = \frac{1+\cos 2z}{2}$, \quad $\sin^2 z = \frac{1-\cos 2z}{2}$.
  \item $\cosh^2 z = \frac{1+\cosh 2z}{2}$, \quad $\sinh^2 z = \frac{\cosh 2z - 1}{2}$.
  \item $\sin(iz) = i\sinh z$ et $\cos(iz) = \cosh z$.
  \item $\abs{\sin z}^2 + \abs{\cos z}^2 = \frac{1}{2}(\cosh 2y + 1)$ où $y = \im(z)$.
\end{enumerate}
\end{exercice}

\begin{exercice}[Propriété du logarithme]\label{exo:log-propriete}
\begin{enumerate}[label=(\alph*)]
  \item Trouver des $z_1, z_2 \in \C$ tels que $\Log(z_1 z_2) \neq \Log(z_1) + \Log(z_2)$.
  \item Montrer que $\Log(z_1 z_2) = \Log(z_1) + \Log(z_2) + 2k\pi i$ pour un unique $k \in \{-1, 0, 1\}$.
  \item La formule $\Log(z^2) = 2\Log(z)$ est-elle toujours vraie ? Trouver un contre-exemple.
\end{enumerate}
\end{exercice}

\begin{exercice}[Branches du logarithme]\label{exo:branches-log}
\begin{enumerate}[label=(\alph*)]
  \item Montrer qu'il n'existe pas de branche continue du logarithme sur $\C \setminus \{0\}$.
  \item Soit $\Omega = \C \setminus \R_{\geq 0}$. Construire explicitement une branche du logarithme sur $\Omega$ dont la partie imaginaire prend ses valeurs dans $(0, 2\pi)$.
  \item Soit $\Omega = \C \setminus i\R_{\geq 0}$ (on retire la demi-droite imaginaire positive). Construire une branche du logarithme sur $\Omega$.
\end{enumerate}
\end{exercice}

\begin{exercice}[Non-existence de branche sur une couronne]\label{exo:branche-couronne}
Soit $A = \{z \in \C : r < \abs{z} < R\}$ avec $0 < r < R$.
\begin{enumerate}[label=(\alph*)]
  \item Montrer qu'il n'existe pas de branche du logarithme sur $A$.
  \item En déduire qu'il n'existe pas de branche de $z^{1/n}$ ($n \geq 2$) sur $A$.
  \item Montrer qu'en revanche, il existe une branche holomorphe de $z^2$ sur $A$ (triviale : $z \mapsto z^2$).
\end{enumerate}
\end{exercice}

\begin{exercice}[Calcul de $(\sqrt{3}+i)^{10}$]\label{exo:puissance-entiere}
\begin{enumerate}[label=(\alph*)]
  \item Écrire $\sqrt{3}+i$ sous forme exponentielle.
  \item En déduire $(\sqrt{3}+i)^{10}$ sans utiliser le binôme de Newton.
  \item Retrouver le résultat en utilisant le logarithme principal.
\end{enumerate}
\end{exercice}

\begin{exercice}[Application conforme de l'exponentielle]\label{exo:conforme-exp}
\begin{enumerate}[label=(\alph*)]
  \item Montrer que $\exp$ est une application conforme (holomorphe et injective) de la bande $S = \{z : 0 < \im(z) < \pi\}$ sur le demi-plan supérieur $\mathbb{H} = \{w : \im(w) > 0\}$.
  \item En déduire que $\Log$ est une application conforme de $\mathbb{H}$ sur $S$.
  \item Déterminer les images par $\exp$ des côtés du rectangle $\{z : a < \re(z) < b,\; 0 < \im(z) < \pi\}$.
\end{enumerate}
\end{exercice}

\begin{exercice}[Fonction $z^z$]\label{exo:z-puissance-z}
On considère $f(z) = z^z = \exp(z\Log z)$, définie sur $\C \setminus \R_{\leq 0}$ par la branche principale.
\begin{enumerate}[label=(\alph*)]
  \item Calculer $f'(z)$.
  \item Déterminer $f(i)$, $f(-i)$ (valeurs principales seulement, en écrivant soigneusement les étapes).
  \item Montrer que $\abs{f(e^{i\theta})} = e^{-\theta}$ pour $\theta \in (-\pi,\pi)$.
\end{enumerate}
\end{exercice}
%%%%%%%%%%%%%%%%%%%%%%%%%%%%%%%%%%%%%%%%%%%%%%%%%%%%%%%%%%%%%%%%%%%%
%  CHAPITRE 4 — Intégration Complexe — Théorème de Cauchy
%%%%%%%%%%%%%%%%%%%%%%%%%%%%%%%%%%%%%%%%%%%%%%%%%%%%%%%%%%%%%%%%%%%%
\chapter{Intégration Complexe --- Théorème de Cauchy}\label{chap:integration-cauchy}

\section{Introduction}

L'intégration complexe constitue le cœur de l'analyse complexe. Alors qu'en
analyse réelle l'intégrale d'une fonction sur un intervalle ne dépend que des
bornes, dans le plan complexe l'intégrale dépend \emph{a priori} du chemin
suivi entre deux points. Le théorème de Cauchy affirme que, sous des
hypothèses convenables d'holomorphie et de topologie du domaine, cette
dépendance disparaît : l'intégrale d'une fonction holomorphe le long d'un
lacet est nulle. Ce résultat, d'apparence simple, a des conséquences
profondes que nous développerons dans les chapitres suivants.

Ce chapitre s'organise comme suit. Nous commençons par la notion de chemin
et d'intégrale curviligne complexe (\S\ref{sec:chemins}--\S\ref{sec:integrale-curviligne}),
puis nous établissons les résultats fondamentaux : le lemme de Goursat
(\S\ref{sec:goursat}), le théorème de Cauchy sous ses différentes formes
(\S\ref{sec:cauchy-etoile}--\S\ref{sec:cauchy-homotopie}), la notion d'indice
(\S\ref{sec:indice}), et enfin la version générale du théorème de Cauchy
(\S\ref{sec:cauchy-general}).

%% ════════════════════════════════════════════════════════════════════
\section{Chemins dans le plan complexe}\label{sec:chemins}
%% ════════════════════════════════════════════════════════════════════

\begin{definition}[Chemin]\label{def:chemin}
Un \textbf{chemin}\index{chemin} (ou \textbf{arc paramétré}) dans $\C$ est une
application continue $\gamma \colon [a,b] \to \C$, où $[a,b]$ est un intervalle
fermé borné de $\R$. On appelle $\gamma(a)$ le \textbf{point initial} et
$\gamma(b)$ le \textbf{point terminal} du chemin. L'\textbf{image} (ou
\textbf{trace}) de $\gamma$ est l'ensemble
\[
  \gamma^* = \gamma([a,b]) = \{\gamma(t) : t \in [a,b]\} \subset \C.
\]
\end{definition}

\begin{definition}[Chemin lisse, lisse par morceaux]\label{def:chemin-lisse}
Un chemin $\gamma \colon [a,b] \to \C$ est dit :
\begin{enumerate}[label=\textup{(\roman*)}]
  \item \textbf{lisse}\index{chemin!lisse} (ou \textbf{régulier}) s'il est de
    classe $\mathcal{C}^1$ sur $[a,b]$ et si $\gamma'(t) \neq 0$ pour tout
    $t \in [a,b]$ ;
  \item \textbf{lisse par morceaux}\index{chemin!lisse par morceaux} s'il
    existe une subdivision $a = t_0 < t_1 < \cdots < t_n = b$ telle que la
    restriction $\gamma|_{[t_{k-1},t_k]}$ soit lisse pour chaque
    $k \in \{1,\ldots,n\}$.
\end{enumerate}
\end{definition}

\begin{definition}[Lacet]\label{def:lacet}
Un chemin $\gamma \colon [a,b] \to \C$ est un \textbf{lacet}\index{lacet}
(ou \textbf{courbe fermée}) si $\gamma(a) = \gamma(b)$.
\end{definition}

\begin{definition}[Lacet simple]\label{def:lacet-simple}
Un lacet $\gamma \colon [a,b] \to \C$ est dit \textbf{simple}\index{lacet!simple}
(ou \textbf{courbe de Jordan}) s'il est injectif sur $[a,b)$, c'est-à-dire si
$\gamma(t_1) = \gamma(t_2)$ avec $t_1, t_2 \in [a,b)$ implique $t_1 = t_2$.
\end{definition}

\begin{definition}[Chemin inverse, concaténation]\label{def:chemin-operations}
Soit $\gamma \colon [a,b] \to \C$ un chemin.
\begin{enumerate}[label=\textup{(\roman*)}]
  \item Le \textbf{chemin inverse}\index{chemin!inverse} $\gamma^-$ est défini par
    \[
      \gamma^-(t) = \gamma(a + b - t), \quad t \in [a,b].
    \]
    Il parcourt la même trace que $\gamma$ mais en sens opposé.
  \item Si $\gamma_1 \colon [a,b] \to \C$ et $\gamma_2 \colon [b,c] \to \C$
    sont deux chemins tels que $\gamma_1(b) = \gamma_2(b)$, la
    \textbf{concaténation}\index{concaténation} $\gamma_1 + \gamma_2 \colon
    [a,c] \to \C$ est définie par
    \[
      (\gamma_1 + \gamma_2)(t) =
      \begin{cases}
        \gamma_1(t) & \text{si } t \in [a,b], \\
        \gamma_2(t) & \text{si } t \in [b,c].
      \end{cases}
    \]
\end{enumerate}
\end{definition}

\begin{exemple}[Cercle de centre $z_0$ et rayon $r$]\label{ex:cercle-chemin}
Le chemin canonique du cercle de centre $z_0 \in \C$ et de rayon $r > 0$,
parcouru dans le sens trigonométrique, est
\[
  \gamma(t) = z_0 + r e^{it}, \quad t \in [0, 2\pi].
\]
C'est un lacet simple lisse. Sa trace est le cercle
$\{z \in \C : \abs{z - z_0} = r\}$.
\end{exemple}

\begin{exemple}[Segment de droite]\label{ex:segment-chemin}
Le segment de $z_1$ à $z_2$ ($z_1 \neq z_2$) est paramétré par
\[
  \gamma(t) = (1-t)z_1 + t z_2 = z_1 + t(z_2 - z_1), \quad t \in [0,1].
\]
On note parfois ce chemin $[z_1, z_2]$.
\end{exemple}

\begin{center}
\begin{tikzpicture}[scale=1.0, >=Stealth]
  % Cercle
  \begin{scope}[xshift=-4cm]
    \draw[->] (-1.8,0) -- (1.8,0);
    \draw[->] (0,-1.8) -- (0,1.8);
    \draw[blue!70!black, very thick, decoration={markings,
      mark=at position 0.15 with {\arrow{>}},
      mark=at position 0.4 with {\arrow{>}},
      mark=at position 0.65 with {\arrow{>}},
      mark=at position 0.9 with {\arrow{>}}},
      postaction={decorate}] (0,0) circle (1.2);
    \fill[blue!70!black] (1.2,0) circle (2pt);
    \node[blue!70!black, below right] at (1.2,0) {$z_0 + r$};
    \fill (0,0) circle (1.5pt) node[below left] {$z_0$};
    \node[below] at (0,-2) {Cercle $\gamma(t) = z_0 + re^{it}$};
  \end{scope}
  % Segment
  \begin{scope}[xshift=3cm]
    \draw[->] (-1.5,0) -- (2.5,0);
    \draw[->] (0,-1) -- (0,2);
    \draw[red!70!black, very thick, decoration={markings,
      mark=at position 0.5 with {\arrow{>}}},
      postaction={decorate}] (-0.5,-0.3) -- (1.8,1.5);
    \fill[red!70!black] (-0.5,-0.3) circle (2pt) node[below left] {$z_1$};
    \fill[red!70!black] (1.8,1.5) circle (2pt) node[above right] {$z_2$};
    \node[below] at (0.5,-1) {Segment $[z_1,z_2]$};
  \end{scope}
  % Chemin lisse par morceaux
  \begin{scope}[xshift=9.5cm]
    \draw[->] (-1.5,0) -- (2.5,0);
    \draw[->] (0,-1) -- (0,2);
    \draw[green!50!black, very thick, decoration={markings,
      mark=at position 0.25 with {\arrow{>}},
      mark=at position 0.6 with {\arrow{>}},
      mark=at position 0.85 with {\arrow{>}}},
      postaction={decorate}]
      (-1,0.5) -- (0.5,1.5) -- (1.5,0.3) -- (2,1.2);
    \fill[green!50!black] (-1,0.5) circle (2pt) node[left] {$z_1$};
    \fill[green!50!black] (2,1.2) circle (2pt) node[right] {$z_n$};
    \foreach \p in {(0.5,1.5),(1.5,0.3)} {
      \fill[green!50!black] \p circle (2pt);
    }
    \node[below] at (0.5,-1) {Chemin lisse par morceaux};
  \end{scope}
\end{tikzpicture}
\captionof{figure}{Exemples de chemins dans le plan complexe.}\label{fig:exemples-chemins}
\end{center}

\begin{definition}[Longueur d'un chemin]\label{def:longueur-chemin}
La \textbf{longueur}\index{longueur d'un chemin} d'un chemin lisse par morceaux
$\gamma \colon [a,b] \to \C$ est
\[
  \ell(\gamma) = \int_a^b \abs{\gamma'(t)} \, \dd t.
\]
\end{definition}

%% ════════════════════════════════════════════════════════════════════
\section{Intégrale curviligne complexe}\label{sec:integrale-curviligne}
%% ════════════════════════════════════════════════════════════════════

\begin{definition}[Intégrale le long d'un chemin]\label{def:integrale-curviligne}
Soit $\gamma \colon [a,b] \to \C$ un chemin lisse par morceaux et soit
$f \colon \gamma^* \to \C$ une fonction continue. L'\textbf{intégrale
curviligne}\index{intégrale curviligne} de $f$ le long de $\gamma$ est
\[
  \int_\gamma f(z) \, \dd z
  = \int_a^b f\bigl(\gamma(t)\bigr) \, \gamma'(t) \, \dd t.
\]
Si $\gamma$ est seulement lisse par morceaux (avec subdivision
$a = t_0 < \cdots < t_n = b$), on pose
\[
  \int_\gamma f(z) \, \dd z
  = \sum_{k=1}^{n} \int_{t_{k-1}}^{t_k}
    f\bigl(\gamma(t)\bigr) \, \gamma'(t) \, \dd t.
\]
\end{definition}

\begin{remarque}\label{rem:integrale-reparametrage}
L'intégrale $\int_\gamma f(z)\,\dd z$ ne dépend pas du paramétrage choisi
pour le chemin, tant que l'orientation est préservée. Si $\varphi \colon
[c,d] \to [a,b]$ est un difféomorphisme croissant de classe $\mathcal{C}^1$,
et $\tilde{\gamma} = \gamma \circ \varphi$, alors
$\int_{\tilde{\gamma}} f(z)\,\dd z = \int_\gamma f(z)\,\dd z$.
\end{remarque}

\begin{proposition}[Propriétés de l'intégrale curviligne]\label{prop:proprietes-integrale}
Soient $\gamma$ un chemin lisse par morceaux et $f, g$ continues sur $\gamma^*$.
\begin{enumerate}[label=\textup{(\roman*)}]
  \item \textbf{Linéarité.} Pour tous $\alpha, \beta \in \C$,
    \[
      \int_\gamma \bigl[\alpha f(z) + \beta g(z)\bigr] \, \dd z
      = \alpha \int_\gamma f(z)\,\dd z + \beta \int_\gamma g(z)\,\dd z.
    \]
  \item \textbf{Chemin inverse.}
    $\displaystyle \int_{\gamma^-} f(z)\,\dd z = -\int_\gamma f(z)\,\dd z$.
  \item \textbf{Concaténation.} Si $\gamma = \gamma_1 + \gamma_2$, alors
    \[
      \int_\gamma f(z)\,\dd z
      = \int_{\gamma_1} f(z)\,\dd z + \int_{\gamma_2} f(z)\,\dd z.
    \]
\end{enumerate}
\end{proposition}

\begin{proof}
\textbf{(i)} Résulte directement de la linéarité de l'intégrale réelle.

\textbf{(ii)} Posons $\gamma \colon [a,b] \to \C$ et $\gamma^-(t) =
\gamma(a+b-t)$. Alors $(\gamma^-)'(t) = -\gamma'(a+b-t)$ et
\begin{align*}
  \int_{\gamma^-} f(z)\,\dd z
  &= \int_a^b f\bigl(\gamma(a+b-t)\bigr)\,\bigl(-\gamma'(a+b-t)\bigr)\,\dd t \\
  &= -\int_a^b f\bigl(\gamma(a+b-t)\bigr)\,\gamma'(a+b-t)\,\dd t.
\end{align*}
Le changement de variable $u = a+b-t$ ($\dd u = -\dd t$) donne
\[
  = -\int_b^a f\bigl(\gamma(u)\bigr)\,\gamma'(u)\,(-\dd u)
  = -\int_a^b f\bigl(\gamma(u)\bigr)\,\gamma'(u)\,\dd u
  = -\int_\gamma f(z)\,\dd z.
\]

\textbf{(iii)} Découle immédiatement de l'additivité de l'intégrale sur les
sous-intervalles.
\end{proof}

\begin{theoreme}[Inégalité ML]\label{thm:inegalite-ML}
Soit $\gamma$ un chemin lisse par morceaux de longueur $\ell(\gamma)$ et
$f$ continue sur $\gamma^*$. Posons $M = \sup_{z \in \gamma^*} \abs{f(z)}$.
Alors
\[
  \abs{\int_\gamma f(z)\,\dd z} \leq M \cdot \ell(\gamma).
\]
\end{theoreme}

\begin{proof}
On a
\begin{align*}
  \abs{\int_\gamma f(z)\,\dd z}
  &= \abs{\int_a^b f(\gamma(t))\,\gamma'(t)\,\dd t}
  \leq \int_a^b \abs{f(\gamma(t))} \cdot \abs{\gamma'(t)} \, \dd t \\
  &\leq M \int_a^b \abs{\gamma'(t)} \, \dd t
  = M \cdot \ell(\gamma). \qedhere
\end{align*}
\end{proof}

\begin{remarque}\label{rem:ML-utile}
L'inégalité ML est l'un des outils les plus utilisés en analyse complexe pour
majorer des intégrales. Elle est particulièrement utile pour montrer que
certaines intégrales tendent vers zéro lorsque le rayon d'un cercle tend
vers l'infini (ou vers zéro).
\end{remarque}

%% ════════════════════════════════════════════════════════════════════
\section{Exemples fondamentaux}\label{sec:exemples-fondamentaux}
%% ════════════════════════════════════════════════════════════════════

Les calculs suivants sont à la base de toute la théorie.

\begin{proposition}[Intégrale de $z^n$ sur un cercle]\label{prop:integrale-zn-cercle}
Soit $\gamma(t) = z_0 + re^{it}$, $t \in [0, 2\pi]$, le cercle de centre
$z_0$ et de rayon $r > 0$ parcouru dans le sens positif. Pour tout
$n \in \Z$,
\[
  \int_\gamma (z - z_0)^n \, \dd z =
  \begin{cases}
    2\pi i & \text{si } n = -1, \\
    0      & \text{si } n \neq -1.
  \end{cases}
\]
\end{proposition}

\begin{proof}
On paramétrise par $z = z_0 + re^{it}$, $\dd z = ire^{it}\,\dd t$, et
$(z - z_0)^n = r^n e^{int}$. Donc
\[
  \int_\gamma (z-z_0)^n\,\dd z
  = \int_0^{2\pi} r^n e^{int} \cdot ire^{it}\,\dd t
  = ir^{n+1} \int_0^{2\pi} e^{i(n+1)t}\,\dd t.
\]
\textbf{Cas $n = -1$.} On obtient $ir^0 \int_0^{2\pi} e^{0}\,\dd t
= i \cdot 2\pi = 2\pi i$.

\textbf{Cas $n \neq -1$.} On a $n+1 \neq 0$, et
\[
  \int_0^{2\pi} e^{i(n+1)t}\,\dd t
  = \left[\frac{e^{i(n+1)t}}{i(n+1)}\right]_0^{2\pi}
  = \frac{e^{2\pi i(n+1)} - 1}{i(n+1)} = \frac{1-1}{i(n+1)} = 0,
\]
car $e^{2\pi i(n+1)} = 1$ pour $n+1 \in \Z$.
\end{proof}

\begin{remarque}\label{rem:cas-n-moins-un-crucial}
Le cas $n = -1$ est le résultat crucial : c'est la seule puissance de
$(z-z_0)$ dont l'intégrale sur un cercle autour de $z_0$ est non nulle. Ce
fait sous-tend toute la théorie des résidus (chapitre~9).
\end{remarque}

\begin{exemple}[Intégrale de $\conj{z}$ sur le cercle unité]\label{ex:integrale-zbar}
Soit $\gamma(t) = e^{it}$, $t \in [0, 2\pi]$, le cercle unité. On calcule
\[
  \int_\gamma \conj{z} \, \dd z
  = \int_0^{2\pi} \overline{e^{it}} \cdot ie^{it} \, \dd t
  = \int_0^{2\pi} e^{-it} \cdot ie^{it} \, \dd t
  = i \int_0^{2\pi} \dd t = 2\pi i.
\]
On observe que $\conj{z}$ n'est pas holomorphe, et pourtant cette intégrale
est non nulle --- ce qui est cohérent car le théorème de Cauchy ne s'applique
pas aux fonctions non holomorphes.
\end{exemple}

\begin{exemple}[Intégrale de $\re(z)$ sur le cercle unité]\label{ex:integrale-rez}
Avec le même $\gamma$, puisque $\re(z) = \tfrac{1}{2}(z + \conj{z})$,
\[
  \int_\gamma \re(z) \, \dd z
  = \frac{1}{2}\int_\gamma z\,\dd z + \frac{1}{2}\int_\gamma \conj{z}\,\dd z
  = \frac{1}{2} \cdot 0 + \frac{1}{2} \cdot 2\pi i = \pi i.
\]
L'intégrale de $z\,\dd z$ sur le cercle unité vaut $0$ car $z^1$
admet la primitive $z^2/2$.
\end{exemple}

%% ════════════════════════════════════════════════════════════════════
\section{Primitives et indépendance du chemin}\label{sec:primitives}
%% ════════════════════════════════════════════════════════════════════

\begin{definition}[Primitive]\label{def:primitive}
Soit $U \subset \C$ un ouvert et $f \colon U \to \C$ continue. Une
\textbf{primitive}\index{primitive} de $f$ sur $U$ est une fonction holomorphe
$F \colon U \to \C$ telle que $F'(z) = f(z)$ pour tout $z \in U$.
\end{definition}

\begin{theoreme}[Primitive et indépendance du chemin]\label{thm:primitive-indep}
Soit $U \subset \C$ un ouvert connexe et $f \colon U \to \C$ continue.
Les assertions suivantes sont équivalentes :
\begin{enumerate}[label=\textup{(\roman*)}]
  \item $f$ admet une primitive sur $U$.
  \item Pour tout lacet lisse par morceaux $\gamma$ dans $U$,
    $\displaystyle \int_\gamma f(z)\,\dd z = 0$.
  \item L'intégrale $\int_\gamma f(z)\,\dd z$ ne dépend que des extrémités
    de $\gamma$ (et non du chemin lisse par morceaux $\gamma$ dans $U$).
\end{enumerate}
De plus, si $F$ est une primitive de $f$, alors pour tout chemin lisse par
morceaux $\gamma$ de $z_1$ à $z_2$ dans $U$,
\[
  \int_\gamma f(z)\,\dd z = F(z_2) - F(z_1).
\]
\end{theoreme}

\begin{proof}
\textbf{(i) $\Rightarrow$ (iii) et formule.}
Supposons que $F' = f$ sur $U$. Si $\gamma \colon [a,b] \to U$ est lisse,
la règle de la chaîne donne
\[
  \frac{\dd}{\dd t}\bigl[F(\gamma(t))\bigr] = F'(\gamma(t))\,\gamma'(t)
  = f(\gamma(t))\,\gamma'(t).
\]
Donc
\[
  \int_\gamma f(z)\,\dd z
  = \int_a^b f(\gamma(t))\,\gamma'(t)\,\dd t
  = \bigl[F(\gamma(t))\bigr]_a^b
  = F(\gamma(b)) - F(\gamma(a))
  = F(z_2) - F(z_1).
\]
Le cas lisse par morceaux s'obtient par sommation télescopique. L'intégrale ne
dépend que de $z_1 = \gamma(a)$ et $z_2 = \gamma(b)$.

\textbf{(iii) $\Rightarrow$ (ii).} Si $\gamma$ est un lacet, $z_1 = z_2$,
donc l'intégrale vaut $F(z_1) - F(z_1) = 0$.

\textbf{(ii) $\Rightarrow$ (i).} Fixons $z_0 \in U$. Puisque $U$ est ouvert
connexe, pour tout $z \in U$ il existe un chemin lisse par morceaux
$\gamma_z$ de $z_0$ à $z$ dans $U$. Posons
\[
  F(z) = \int_{\gamma_z} f(\zeta)\,\dd\zeta.
\]
La condition (ii) garantit que $F$ est bien définie (indépendante du chemin
choisi). Montrons que $F'(z) = f(z)$.

Soit $z \in U$. Puisque $U$ est ouvert, il existe $\delta > 0$ tel que
$D(z,\delta) \subset U$. Pour $h \in \C$ avec $0 < \abs{h} < \delta$,
le segment $[z, z+h]$ est dans $U$, et
\[
  F(z+h) - F(z) = \int_{[z,z+h]} f(\zeta)\,\dd\zeta.
\]
Or, en paramétrant $\zeta = z + th$, $t \in [0,1]$,
\[
  \int_{[z,z+h]} f(\zeta)\,\dd\zeta = h \int_0^1 f(z+th)\,\dd t.
\]
Donc
\[
  \frac{F(z+h) - F(z)}{h} - f(z)
  = \int_0^1 \bigl[f(z+th) - f(z)\bigr]\,\dd t.
\]
Par continuité de $f$, pour tout $\varepsilon > 0$ il existe
$\delta' \leq \delta$ tel que $\abs{h} < \delta'$ implique
$\abs{f(z+th) - f(z)} < \varepsilon$ pour tout $t \in [0,1]$. Ainsi
\[
  \abs{\frac{F(z+h) - F(z)}{h} - f(z)} \leq \varepsilon,
\]
ce qui prouve que $F'(z) = f(z)$.
\end{proof}

\begin{corollaire}\label{cor:zn-primitive}
Pour $n \in \Z$, $n \neq -1$, la fonction $z^n$ admet la primitive
$\dfrac{z^{n+1}}{n+1}$ sur $\C$ (si $n \geq 0$) ou sur $\C^*$ (si
$n \leq -2$). En particulier, $\displaystyle \int_\gamma z^n\,\dd z = 0$
pour tout lacet $\gamma$ dans le domaine correspondant.

En revanche, $1/z$ n'admet \emph{pas} de primitive sur $\C^*$.
\end{corollaire}

\begin{proof}
Pour $n \neq -1$, on vérifie que $\frac{\dd}{\dd z}\frac{z^{n+1}}{n+1}
= z^n$. Le théorème~\ref{thm:primitive-indep} s'applique. Pour $1/z$, si
une primitive existait sur $\C^*$, l'intégrale de $1/z$ sur tout lacet dans
$\C^*$ serait nulle. Or, la proposition~\ref{prop:integrale-zn-cercle}
montre que $\int_\gamma (1/z)\,\dd z = 2\pi i \neq 0$ pour $\gamma$ le
cercle unité.
\end{proof}

%% ════════════════════════════════════════════════════════════════════
\section{Lemme de Goursat}\label{sec:goursat}
%% ════════════════════════════════════════════════════════════════════

Le lemme de Goursat est un résultat technique fondamental : il établit que
l'intégrale d'une fonction holomorphe sur le bord d'un triangle est nulle,
\emph{sans} supposer la continuité de la dérivée. Ceci est remarquable car
la preuve du théorème de Cauchy via la formule de Green nécessite cette
continuité.

\begin{lemme}[Goursat]\label{lem:goursat}
Soit $U \subset \C$ un ouvert et $f \colon U \to \C$ holomorphe. Soit
$T \subset U$ un triangle fermé (c'est-à-dire l'enveloppe convexe de trois
points non alignés), avec $T \subset U$. Si $\partial T$ désigne le bord
de $T$ parcouru dans le sens positif, alors
\[
  \int_{\partial T} f(z) \, \dd z = 0.
\]
\end{lemme}

\begin{proof}
La preuve repose sur un argument de \textbf{bissection itérée}.

\medskip

\noindent\textbf{Étape 1 : Construction par subdivision.}
Notons $T^{(0)} = T$ et $I_0 = \int_{\partial T^{(0)}} f(z)\,\dd z$.
En reliant les milieux des trois côtés de $T^{(0)}$, on obtient quatre
sous-triangles $T_1, T_2, T_3, T_4$ (voir figure~\ref{fig:goursat}).

\begin{center}
\begin{tikzpicture}[scale=2.2, >=Stealth]
  % Triangle principal
  \coordinate (A) at (0,0);
  \coordinate (B) at (3,0);
  \coordinate (C) at (1.2,2.2);
  \coordinate (MAB) at ($(A)!0.5!(B)$);
  \coordinate (MBC) at ($(B)!0.5!(C)$);
  \coordinate (MAC) at ($(A)!0.5!(C)$);

  % Remplissage
  \fill[blue!8] (A) -- (B) -- (C) -- cycle;
  \fill[red!15] (MAB) -- (MBC) -- (MAC) -- cycle;

  % Triangle extérieur
  \draw[thick] (A) -- (B) -- (C) -- cycle;
  % Subdivisions
  \draw[dashed, gray] (MAB) -- (MBC) -- (MAC) -- cycle;

  % Flèches d'orientation sur le bord extérieur
  \draw[blue!70!black, very thick, decoration={markings,
    mark=at position 0.5 with {\arrow{>}}},
    postaction={decorate}] (A) -- (B);
  \draw[blue!70!black, very thick, decoration={markings,
    mark=at position 0.5 with {\arrow{>}}},
    postaction={decorate}] (B) -- (C);
  \draw[blue!70!black, very thick, decoration={markings,
    mark=at position 0.5 with {\arrow{>}}},
    postaction={decorate}] (C) -- (A);

  % Labels
  \node[below left] at (A) {$A$};
  \node[below right] at (B) {$B$};
  \node[above] at (C) {$C$};
  \node[below] at (MAB) {$M_{AB}$};
  \node[right] at (MBC) {$M_{BC}$};
  \node[left] at (MAC) {$M_{AC}$};

  % Labels des sous-triangles
  \node at (0.7,0.5) {$T_1$};
  \node at (2,0.5) {$T_2$};
  \node at (1.1,1.5) {$T_3$};
  \node[red!70!black] at (1.4,0.85) {$T_4$};

  % Annotation
  \node[right] at (3.3,1) {%
    \begin{minipage}{4cm}
      \small Les intégrales sur les côtés intérieurs s'annulent
      deux à deux (parcours opposés).

      $\displaystyle I_0 = \sum_{j=1}^{4} \int_{\partial T_j} f(z)\,\dd z$
    \end{minipage}};
\end{tikzpicture}
\captionof{figure}{Subdivision d'un triangle dans la preuve du lemme de
Goursat.}\label{fig:goursat}
\end{center}

Les intégrales sur les côtés intérieurs (ceux introduits par la subdivision)
s'annulent deux à deux car chaque côté intérieur est parcouru une fois dans
chaque sens. Donc
\[
  I_0 = \int_{\partial T^{(0)}} f(z)\,\dd z
  = \sum_{j=1}^{4} \int_{\partial T_j} f(z)\,\dd z.
\]
Par l'inégalité triangulaire, il existe $j_0 \in \{1,2,3,4\}$ tel que
\[
  \abs{\int_{\partial T_{j_0}} f(z)\,\dd z}
  \geq \frac{1}{4}\abs{I_0}.
\]
On pose $T^{(1)} = T_{j_0}$.

\medskip

\noindent\textbf{Étape 2 : Itération.}
En itérant ce procédé, on construit une suite de triangles emboîtés
$T^{(0)} \supset T^{(1)} \supset T^{(2)} \supset \cdots$ tels que, pour
tout $k \geq 0$ :
\begin{enumerate}[label=(\alph*)]
  \item $\displaystyle \abs{\int_{\partial T^{(k)}} f(z)\,\dd z}
    \geq \frac{\abs{I_0}}{4^k}$,
  \item $\ell(\partial T^{(k)}) = \dfrac{\ell(\partial T^{(0)})}{2^k}$,
  \item $\operatorname{diam}(T^{(k)}) = \dfrac{\operatorname{diam}(T^{(0)})}{2^k}$.
\end{enumerate}

\medskip

\noindent\textbf{Étape 3 : Point limite.}
Par le théorème des fermés emboîtés, $\displaystyle \bigcap_{k=0}^{\infty}
T^{(k)} = \{z_0\}$ pour un certain $z_0 \in T \subset U$. Puisque $f$ est
holomorphe en $z_0$, on peut écrire
\[
  f(z) = f(z_0) + f'(z_0)(z-z_0) + (z-z_0)\,\varphi(z),
\]
où $\varphi(z) \to 0$ quand $z \to z_0$.

\medskip

\noindent\textbf{Étape 4 : Majoration.}
La fonction $z \mapsto f(z_0) + f'(z_0)(z-z_0)$ est un polynôme de degré
$\leq 1$ en $z$, donc elle admet une primitive. Par le
théorème~\ref{thm:primitive-indep},
\[
  \int_{\partial T^{(k)}} \bigl[f(z_0) + f'(z_0)(z-z_0)\bigr]\,\dd z = 0.
\]
Donc
\[
  \int_{\partial T^{(k)}} f(z)\,\dd z
  = \int_{\partial T^{(k)}} (z-z_0)\,\varphi(z)\,\dd z.
\]
Pour $z \in \partial T^{(k)}$, on a $\abs{z - z_0} \leq
\operatorname{diam}(T^{(k)}) = d_0/2^k$ où $d_0 = \operatorname{diam}(T^{(0)})$.
Soit $\varepsilon > 0$. Il existe $K$ tel que pour $k \geq K$ et
$z \in T^{(k)}$, on a $\abs{\varphi(z)} < \varepsilon$. Par l'inégalité ML,
\[
  \abs{\int_{\partial T^{(k)}} f(z)\,\dd z}
  \leq \frac{d_0}{2^k} \cdot \varepsilon \cdot \frac{L_0}{2^k}
  = \frac{\varepsilon \, d_0 \, L_0}{4^k},
\]
où $L_0 = \ell(\partial T^{(0)})$. Or, d'après (a),
$\abs{I_0}/4^k \leq \abs{\int_{\partial T^{(k)}} f(z)\,\dd z}$. Donc
\[
  \abs{I_0} \leq \varepsilon \, d_0 \, L_0.
\]
Comme ceci vaut pour tout $\varepsilon > 0$, on conclut que $I_0 = 0$.
\end{proof}

\begin{remarque}\label{rem:goursat-vs-cauchy}
La force du lemme de Goursat est de ne supposer que l'existence de $f'(z)$,
sans supposer la continuité de $f'$. On sait, grâce au théorème de
Looman--Menchov, que cette continuité est en fait automatique, mais la preuve
de Goursat est élémentaire et directe.
\end{remarque}

%% ════════════════════════════════════════════════════════════════════
\section{Théorème de Cauchy pour les domaines étoilés}\label{sec:cauchy-etoile}
%% ════════════════════════════════════════════════════════════════════

\begin{definition}[Domaine étoilé]\label{def:domaine-etoile}
Un ouvert $U \subset \C$ est \textbf{étoilé}\index{domaine étoilé} par rapport
à un point $z_0 \in U$ si, pour tout $z \in U$, le segment $[z_0, z]$ est
contenu dans $U$. On dit simplement que $U$ est étoilé s'il est étoilé par
rapport à au moins un de ses points.
\end{definition}

\begin{center}
\begin{tikzpicture}[scale=0.9, >=Stealth]
  % Domaine étoilé
  \begin{scope}[xshift=-4.5cm]
    \fill[blue!10] plot[smooth cycle, tension=0.7]
      coordinates {(-1.5,-1) (0,-1.5) (2,-0.8) (2.2,0.5) (1.5,1.8) (0,2) (-1.5,1.2) (-2,0)};
    \draw[blue!70!black, thick] plot[smooth cycle, tension=0.7]
      coordinates {(-1.5,-1) (0,-1.5) (2,-0.8) (2.2,0.5) (1.5,1.8) (0,2) (-1.5,1.2) (-2,0)};
    \fill[red!70!black] (0,0.3) circle (2pt) node[below right] {$z_0$};
    % Quelques segments
    \foreach \a/\b in {-1.2/-0.7, 1.5/-0.5, 1.8/0.3, 1.2/1.5, -0.3/1.7, -1.3/0.8} {
      \draw[red!50!black, dashed, thin] (0,0.3) -- (\a,\b);
      \fill[red!50!black] (\a,\b) circle (1.5pt);
    }
    \node[below] at (0,-2) {Domaine étoilé};
  \end{scope}
  % Domaine convexe
  \begin{scope}[xshift=4cm]
    \fill[green!10] plot[smooth cycle, tension=0.85]
      coordinates {(-1.5,-0.8) (0.5,-1.2) (2,-0.3) (1.8,1.2) (0,1.8) (-1.5,0.8)};
    \draw[green!50!black, thick] plot[smooth cycle, tension=0.85]
      coordinates {(-1.5,-0.8) (0.5,-1.2) (2,-0.3) (1.8,1.2) (0,1.8) (-1.5,0.8)};
    \fill[red!70!black] (0.2,0.3) circle (2pt) node[below right] {$z_0$};
    \node[below] at (0.2,-2) {Domaine convexe (étoilé};
    \node[below] at (0.2,-2.5) {par rapport à tout point)};
  \end{scope}
\end{tikzpicture}
\captionof{figure}{Domaine étoilé et domaine convexe.}\label{fig:domaine-etoile}
\end{center}

\begin{theoreme}[Cauchy pour les domaines étoilés]\label{thm:cauchy-etoile}
Soit $U \subset \C$ un ouvert étoilé et $f \colon U \to \C$ holomorphe. Alors :
\begin{enumerate}[label=\textup{(\roman*)}]
  \item $f$ admet une primitive $F$ sur $U$.
  \item Pour tout lacet lisse par morceaux $\gamma$ dans $U$,
    $\displaystyle \int_\gamma f(z)\,\dd z = 0$.
\end{enumerate}
\end{theoreme}

\begin{proof}
Soit $z_0 \in U$ un point par rapport auquel $U$ est étoilé. Pour tout
$z \in U$, le segment $[z_0, z]$ est dans $U$. On définit
\[
  F(z) = \int_{[z_0, z]} f(\zeta)\,\dd\zeta
  = \int_0^1 f\bigl(z_0 + t(z-z_0)\bigr)(z-z_0)\,\dd t.
\]

\textbf{Montrons que $F'(z) = f(z)$.} Soit $z \in U$ et $h \in \C$ assez petit.
Le triangle de sommets $z_0, z, z+h$ est contenu dans $U$ (car $U$ est étoilé
par rapport à $z_0$, les segments $[z_0,z]$, $[z_0,z+h]$ sont dans $U$, et
par convexité locale le segment $[z,z+h]$ est aussi dans $U$ pour $\abs{h}$
assez petit). Par le lemme de Goursat~\ref{lem:goursat},
\[
  \int_{[z_0,z]} f + \int_{[z,z+h]} f + \int_{[z+h,z_0]} f = 0.
\]
Donc
\[
  F(z+h) - F(z) = \int_{[z_0,z+h]} f - \int_{[z_0,z]} f
  = \int_{[z,z+h]} f.
\]
En paramétrant le segment par $\zeta = z + th$, $t \in [0,1]$,
\[
  \frac{F(z+h)-F(z)}{h} = \int_0^1 f(z+th)\,\dd t \xrightarrow[h \to 0]{} f(z)
\]
par continuité de $f$. Donc $F$ est holomorphe et $F' = f$.

L'assertion (ii) découle alors du théorème~\ref{thm:primitive-indep}.
\end{proof}

%% ════════════════════════════════════════════════════════════════════
\section{Théorème de Cauchy pour les domaines simplement connexes}\label{sec:cauchy-simplement-connexe}
%% ════════════════════════════════════════════════════════════════════

\begin{definition}[Domaine simplement connexe]\label{def:simplement-connexe}
Un ouvert connexe $U \subset \C$ est \textbf{simplement connexe}%
\index{simplement connexe} si tout lacet continu dans $U$ est homotope à un
point dans $U$, c'est-à-dire si $\pi_1(U, z_0) = \{e\}$ pour tout (ou un)
$z_0 \in U$.

De manière équivalente (pour les ouverts de $\C$), $U$ est simplement connexe
si et seulement si $\C \setminus U$ est connexe dans $\C \cup \{\infty\}$
(la sphère de Riemann).
\end{definition}

\begin{center}
\begin{tikzpicture}[scale=0.85, >=Stealth]
  % Simplement connexe
  \begin{scope}[xshift=-5cm]
    \fill[green!12] plot[smooth cycle, tension=0.7]
      coordinates {(-2,-1.2) (0,-1.8) (2.2,-0.6) (2,1) (0.5,2) (-1,1.8) (-2.2,0.5)};
    \draw[green!50!black, thick] plot[smooth cycle, tension=0.7]
      coordinates {(-2,-1.2) (0,-1.8) (2.2,-0.6) (2,1) (0.5,2) (-1,1.8) (-2.2,0.5)};
    \node[green!50!black, font=\large\bfseries] at (0,0) {$U$};
    \node[below] at (0,-2.3) {Simplement connexe};
    \node[below, font=\small] at (0,-2.8) {(pas de trou)};
  \end{scope}
  % Multiplement connexe
  \begin{scope}[xshift=4cm]
    \fill[orange!12] plot[smooth cycle, tension=0.7]
      coordinates {(-2.5,-1.5) (0,-2) (2.5,-1) (2.5,1) (1,2.2) (-1,2) (-2.5,0.5)};
    \fill[white] (0.8,0) circle (0.6);
    \fill[white] (-0.8,0.5) circle (0.45);
    \draw[orange!70!black, thick] plot[smooth cycle, tension=0.7]
      coordinates {(-2.5,-1.5) (0,-2) (2.5,-1) (2.5,1) (1,2.2) (-1,2) (-2.5,0.5)};
    \draw[orange!70!black, thick] (0.8,0) circle (0.6);
    \draw[orange!70!black, thick] (-0.8,0.5) circle (0.45);
    \node[orange!70!black, font=\large\bfseries] at (-0.2,-1) {$V$};
    \node[below] at (0,-2.5) {Multiplement connexe};
    \node[below, font=\small] at (0,-3) {(deux trous)};
    % Lacet non contractible
    \draw[red!70!black, thick, decoration={markings,
      mark=at position 0.3 with {\arrow{>}},
      mark=at position 0.8 with {\arrow{>}}},
      postaction={decorate}] (0.8,0) circle (1.0);
  \end{scope}
\end{tikzpicture}
\captionof{figure}{Domaines simplement et multiplement connexes.}\label{fig:simple-multiple}
\end{center}

\begin{theoreme}[Cauchy pour les domaines simplement connexes]\label{thm:cauchy-simple-connexe}
Soit $U \subset \C$ un ouvert simplement connexe et $f \colon U \to \C$
holomorphe. Alors :
\begin{enumerate}[label=\textup{(\roman*)}]
  \item $f$ admet une primitive sur $U$.
  \item Pour tout lacet lisse par morceaux $\gamma$ dans $U$,
    $\displaystyle \int_\gamma f(z)\,\dd z = 0$.
\end{enumerate}
\end{theoreme}

Nous reportons la preuve complète à la section~\ref{sec:cauchy-homotopie},
où elle sera obtenue comme conséquence de la version homotopique du théorème
de Cauchy.

\begin{exemple}\label{ex:cauchy-exemples-domaines}
\leavevmode
\begin{enumerate}[label=\textup{(\arabic*)}]
  \item $\C$, tout disque $D(z_0, r)$, tout demi-plan sont simplement connexes.
    Le théorème de Cauchy s'y applique.
  \item $\C^* = \C \setminus \{0\}$ n'est \emph{pas} simplement connexe :
    le cercle unité est un lacet non contractible dans $\C^*$.
  \item $\C \setminus \R^-$ (le plan privé du demi-axe réel négatif) est
    simplement connexe. C'est le domaine naturel pour la détermination
    principale du logarithme.
\end{enumerate}
\end{exemple}

%% ════════════════════════════════════════════════════════════════════
\section{Preuve du théorème de Cauchy via la formule de Green}\label{sec:cauchy-green}
%% ════════════════════════════════════════════════════════════════════

La preuve qui suit utilise la formule de Green (Stokes en dimension 2) et
suppose la continuité de $f'$. On montrera plus tard (comme conséquence
de la formule intégrale de Cauchy) que toute fonction holomorphe a une dérivée
continue, de sorte que cette hypothèse supplémentaire est finalement
superflue.

\begin{theoreme}[Cauchy via Green]\label{thm:cauchy-green}
Soit $U \subset \C$ un ouvert borné dont le bord $\partial U$ est une courbe
de Jordan lisse par morceaux (ou une réunion finie de telles courbes), orienté
positivement. Si $f$ est holomorphe sur un ouvert contenant $\conj{U}$ et si
$f'$ est continue, alors
\[
  \oint_{\partial U} f(z)\,\dd z = 0.
\]
\end{theoreme}

\begin{proof}
Écrivons $f = u + iv$ et $z = x + iy$, de sorte que
$f(z)\,\dd z = (u+iv)(\dd x + i\,\dd y)$. En développant,
\[
  f(z)\,\dd z = (u\,\dd x - v\,\dd y) + i(v\,\dd x + u\,\dd y).
\]
Par la formule de Green (théorème de Stokes en dimension 2), pour un domaine
$D$ à bord lisse par morceaux orienté positivement,
\[
  \oint_{\partial D} (P\,\dd x + Q\,\dd y)
  = \iint_D \left(\frac{\partial Q}{\partial x}
    - \frac{\partial P}{\partial y}\right) \dd x\,\dd y.
\]
\textbf{Partie réelle.} Avec $P = u$, $Q = -v$ :
\[
  \oint_{\partial U} (u\,\dd x - v\,\dd y)
  = \iint_U \left(-\frac{\partial v}{\partial x}
    - \frac{\partial u}{\partial y}\right) \dd x\,\dd y.
\]
Par les équations de Cauchy--Riemann ($u_x = v_y$ et $u_y = -v_x$), on a
$-v_x - u_y = -(-u_y) - u_y = u_y - u_y = 0$.

\textbf{Partie imaginaire.} Avec $P = v$, $Q = u$ :
\[
  \oint_{\partial U} (v\,\dd x + u\,\dd y)
  = \iint_U \left(\frac{\partial u}{\partial x}
    - \frac{\partial v}{\partial y}\right) \dd x\,\dd y
  = \iint_U (u_x - v_y)\,\dd x\,\dd y = 0.
\]
Les parties réelle et imaginaire de l'intégrale sont toutes deux nulles, d'où
$\oint_{\partial U} f(z)\,\dd z = 0$.
\end{proof}

\begin{remarque}\label{rem:green-hypothese}
Cette preuve a l'avantage de la clarté conceptuelle : le théorème de Cauchy
est une conséquence directe des équations de Cauchy--Riemann et du théorème
de Stokes. Cependant, elle requiert l'hypothèse supplémentaire que $f'$ soit
continue, ce que la preuve via le lemme de Goursat évite.
\end{remarque}

%% ════════════════════════════════════════════════════════════════════
\section{Version homotopique du théorème de Cauchy}\label{sec:cauchy-homotopie}
%% ════════════════════════════════════════════════════════════════════

\begin{definition}[Homotopie de chemins]\label{def:homotopie}
Soit $U \subset \C$ un ouvert. Deux chemins $\gamma_0, \gamma_1 \colon
[a,b] \to U$ ayant les mêmes extrémités ($\gamma_0(a) = \gamma_1(a)$ et
$\gamma_0(b) = \gamma_1(b)$) sont dits \textbf{homotopes à extrémités
fixes}\index{homotopie} dans $U$ s'il existe une application continue
$H \colon [a,b] \times [0,1] \to U$ telle que
\begin{align*}
  H(t,0) &= \gamma_0(t) \quad &&\text{pour tout } t \in [a,b], \\
  H(t,1) &= \gamma_1(t) \quad &&\text{pour tout } t \in [a,b], \\
  H(a,s) &= \gamma_0(a) = \gamma_1(a) \quad &&\text{pour tout } s \in [0,1], \\
  H(b,s) &= \gamma_0(b) = \gamma_1(b) \quad &&\text{pour tout } s \in [0,1].
\end{align*}
L'application $H$ est appelée une \textbf{homotopie} de $\gamma_0$ vers
$\gamma_1$. On note $\gamma_0 \simeq \gamma_1$ dans $U$.
\end{definition}

\begin{definition}[Lacets homotopes]\label{def:lacets-homotopes}
Deux lacets $\gamma_0, \gamma_1$ basés en $z_0 \in U$ sont homotopes (en
tant que lacets) s'il existe une homotopie $H$ comme ci-dessus avec
$H(a,s) = H(b,s) = z_0$ pour tout $s$.

Un lacet est dit \textbf{contractile}\index{lacet!contractile} (ou
\textbf{homotope à un point}) s'il est homotope au lacet constant $z_0$.
\end{definition}

\begin{center}
\begin{tikzpicture}[scale=1.0, >=Stealth]
  % Domaine
  \fill[gray!8] plot[smooth cycle, tension=0.7]
    coordinates {(-3,-2) (0,-2.5) (3.5,-1.5) (3.5,1.5) (1,2.5) (-2,2.5) (-3.5,0.5)};
  \draw[gray!50, thick] plot[smooth cycle, tension=0.7]
    coordinates {(-3,-2) (0,-2.5) (3.5,-1.5) (3.5,1.5) (1,2.5) (-2,2.5) (-3.5,0.5)};

  % Trou
  \fill[white] (0.5,0) circle (0.5);
  \draw[gray!50, thick] (0.5,0) circle (0.5);
  \node at (0.5,0) {\small trou};

  % Deux chemins homotopes
  \coordinate (A) at (-2.5,-1);
  \coordinate (B) at (2.5,1);
  \fill[black] (A) circle (2pt) node[left] {$z_1$};
  \fill[black] (B) circle (2pt) node[right] {$z_2$};

  \draw[blue!70!black, very thick, decoration={markings,
    mark=at position 0.5 with {\arrow{>}}},
    postaction={decorate}]
    (A) .. controls (-1,1.5) and (1,2) .. (B);
  \node[blue!70!black, above] at (-0.3,1.8) {$\gamma_0$};

  \draw[red!70!black, very thick, decoration={markings,
    mark=at position 0.5 with {\arrow{>}}},
    postaction={decorate}]
    (A) .. controls (-1,-1.8) and (2,-1.5) .. (B);
  \node[red!70!black, below] at (0.5,-1.8) {$\gamma_1$};

  % Homotopie indication
  \draw[green!50!black, dashed, thin, decoration={markings,
    mark=at position 0.5 with {\arrow{>}}},
    postaction={decorate}]
    (A) .. controls (-0.5,0.5) and (1.5,0.8) .. (B);

  \node[right, font=\small] at (3.8,0) {%
    $\gamma_0 \simeq \gamma_1$ dans $U$};
\end{tikzpicture}
\captionof{figure}{Homotopie de chemins --- $\gamma_0$ peut être continûment
déformé en $\gamma_1$ sans sortir de $U$ ni traverser le trou.}\label{fig:homotopie}
\end{center}

\begin{theoreme}[Cauchy --- version homotopique]\label{thm:cauchy-homotopie}
Soit $U \subset \C$ un ouvert et $f \colon U \to \C$ holomorphe.
\begin{enumerate}[label=\textup{(\roman*)}]
  \item Si $\gamma_0$ et $\gamma_1$ sont deux chemins lisses par morceaux
    homotopes à extrémités fixes dans $U$, alors
    \[
      \int_{\gamma_0} f(z)\,\dd z = \int_{\gamma_1} f(z)\,\dd z.
    \]
  \item En particulier, si $\gamma$ est un lacet lisse par morceaux
    contractile dans $U$, alors $\displaystyle \int_\gamma f(z)\,\dd z = 0$.
\end{enumerate}
\end{theoreme}

\begin{proof}
Nous donnons les grandes lignes de la preuve ; les détails techniques
d'approximation lisse par morceaux sont admis.

\medskip

\noindent\textbf{Idée clé.}
Soit $H \colon [a,b] \times [0,1] \to U$ l'homotopie. Le compact
$K = H([a,b] \times [0,1]) \subset U$ est à distance strictement positive du
complémentaire de $U$. On subdivise le rectangle $[a,b] \times [0,1]$ en sous-rectangles
suffisamment petits pour que l'image de chacun soit contenue dans un disque
ouvert inclus dans $U$.

Pour chaque paire de chemins intermédiaires adjacents $\gamma_s$ et
$\gamma_{s+\Delta s}$ (correspondant à deux valeurs consécutives du paramètre
d'homotopie), la différence des intégrales se décompose en une somme
d'intégrales sur de petits lacets, chacun contenu dans un disque (donc étoilé).
Par le théorème de Cauchy pour les domaines étoilés
(\ref{thm:cauchy-etoile}), chacune de ces intégrales est nulle.

Plus précisément, on subdivise $[a,b] = \bigcup [t_{j-1}, t_j]$ et
$[0,1] = \bigcup [s_{k-1}, s_k]$ en sous-intervalles assez fins. Pour
chaque rectangle $[t_{j-1}, t_j] \times [s_{k-1}, s_k]$, l'image
$H([t_{j-1}, t_j] \times [s_{k-1}, s_k])$ est contenue dans un disque
$D_{jk} \subset U$. En reliant les images des sommets par des segments, on
construit un quadrilatère dans $D_{jk}$. L'intégrale de $f$ le long du bord
de ce quadrilatère est nulle (car $D_{jk}$ est étoilé). En sommant sur tous
les sous-rectangles, les contributions des côtés intérieurs s'annulent, et il
ne reste que la différence entre les intégrales le long de $\gamma_0$ et
$\gamma_1$ (à une approximation polygonale près qui converge).

L'assertion (ii) en découle immédiatement en prenant $\gamma_1$ constant.
\end{proof}

\begin{corollaire}[Preuve du théorème~\ref{thm:cauchy-simple-connexe}]\label{cor:preuve-cauchy-sc}
Si $U$ est simplement connexe, tout lacet dans $U$ est contractile. Le
théorème~\ref{thm:cauchy-homotopie}(ii) donne donc
$\int_\gamma f(z)\,\dd z = 0$ pour tout lacet $\gamma$ et toute $f$
holomorphe sur $U$. L'existence d'une primitive résulte alors du
théorème~\ref{thm:primitive-indep}.
\end{corollaire}

%% ════════════════════════════════════════════════════════════════════
\section{Indice d'un point par rapport à un lacet}\label{sec:indice}
%% ════════════════════════════════════════════════════════════════════

\begin{definition}[Indice (nombre d'enroulement)]\label{def:indice}
Soit $\gamma$ un lacet lisse par morceaux dans $\C$ et $z_0 \notin \gamma^*$.
L'\textbf{indice}\index{indice (nombre d'enroulement)} (ou \textbf{nombre
d'enroulement}) de $\gamma$ par rapport à $z_0$ est
\[
  \ind(\gamma, z_0) = \frac{1}{2\pi i} \int_\gamma \frac{\dd z}{z - z_0}.
\]
\end{definition}

\begin{theoreme}[Intégralité de l'indice]\label{thm:indice-entier}
Pour tout lacet lisse par morceaux $\gamma$ et tout $z_0 \notin \gamma^*$,
l'indice $\ind(\gamma, z_0)$ est un entier.
\end{theoreme}

\begin{proof}
Soit $\gamma \colon [a,b] \to \C$ avec $\gamma(a) = \gamma(b)$. Considérons
la fonction
\[
  \varphi(t) = \int_a^t \frac{\gamma'(s)}{\gamma(s) - z_0}\,\dd s,
  \quad t \in [a,b].
\]
On a $\varphi(a) = 0$ et $\varphi(b) = \int_\gamma \frac{\dd z}{z-z_0}
= 2\pi i \, \ind(\gamma, z_0)$.

Définissons $\psi(t) = (\gamma(t) - z_0)\,e^{-\varphi(t)}$.
Pour $t$ en un point de dérivabilité de $\gamma$,
\begin{align*}
  \psi'(t)
  &= \gamma'(t)\,e^{-\varphi(t)}
    + (\gamma(t)-z_0)\,(-\varphi'(t))\,e^{-\varphi(t)} \\
  &= \gamma'(t)\,e^{-\varphi(t)}
    - (\gamma(t)-z_0) \cdot \frac{\gamma'(t)}{\gamma(t)-z_0} \cdot e^{-\varphi(t)} \\
  &= \gamma'(t)\,e^{-\varphi(t)} - \gamma'(t)\,e^{-\varphi(t)} = 0.
\end{align*}
(Sur chaque morceau lisse, $\psi$ est constante.) Puisque $\gamma$ est lisse
par morceaux et $\psi$ est continue, $\psi$ est constante sur $[a,b]$ entier.
Donc $\psi(b) = \psi(a)$, c'est-à-dire
\[
  (\gamma(b)-z_0)\,e^{-\varphi(b)} = (\gamma(a)-z_0)\,e^{-\varphi(a)}.
\]
Or $\gamma(a) = \gamma(b)$ et $\varphi(a) = 0$, donc $e^{-\varphi(b)} = 1$.
Ceci signifie $\varphi(b) \in 2\pi i\,\Z$, d'où
$\ind(\gamma, z_0) = \varphi(b)/(2\pi i) \in \Z$.
\end{proof}

\begin{proposition}[Propriétés de l'indice]\label{prop:proprietes-indice}
Soit $\gamma$ un lacet lisse par morceaux dans $\C$.
\begin{enumerate}[label=\textup{(\roman*)}]
  \item \textbf{Constance sur les composantes.} La fonction
    $z_0 \mapsto \ind(\gamma, z_0)$ est continue (et à valeurs entières, donc
    constante) sur chaque composante connexe de $\C \setminus \gamma^*$.
  \item \textbf{Valeur à l'infini.} $\ind(\gamma, z_0) = 0$ pour tout
    $z_0$ dans la composante connexe non bornée de $\C \setminus \gamma^*$.
  \item \textbf{Chemin inverse.} $\ind(\gamma^-, z_0) = -\ind(\gamma, z_0)$.
  \item \textbf{Additivité.} Si $\gamma = \gamma_1 + \gamma_2$ est une
    concaténation de lacets, alors
    $\ind(\gamma, z_0) = \ind(\gamma_1, z_0) + \ind(\gamma_2, z_0)$.
\end{enumerate}
\end{proposition}

\begin{proof}
\textbf{(i)} Soit $z_0 \notin \gamma^*$. La fonction $z \mapsto
\frac{1}{2\pi i} \int_\gamma \frac{\dd\zeta}{\zeta - z}$ est continue en $z$
(on peut dériver sous le signe intégral puisque $\gamma^*$ est compact et
$z \notin \gamma^*$). Étant continue et à valeurs dans $\Z$, elle est
localement constante, donc constante sur chaque composante connexe.

\textbf{(ii)} Si $\abs{z_0}$ est assez grand (plus grand que $\max_{\gamma^*}
\abs{\zeta}$), on peut majorer :
\[
  \abs{\ind(\gamma, z_0)}
  = \frac{1}{2\pi}\abs{\int_\gamma \frac{\dd\zeta}{\zeta - z_0}}
  \leq \frac{1}{2\pi} \cdot
    \frac{1}{\abs{z_0} - \max\abs{\zeta}} \cdot \ell(\gamma)
  \xrightarrow[\abs{z_0}\to\infty]{} 0.
\]
Comme $\ind(\gamma, z_0) \in \Z$, il vaut $0$ pour $\abs{z_0}$ grand, donc
sur toute la composante non bornée par (i).

\textbf{(iii)} et \textbf{(iv)} résultent directement des propriétés de
l'intégrale curviligne (proposition~\ref{prop:proprietes-integrale}).
\end{proof}

\begin{exemple}[Indice du cercle]\label{ex:indice-cercle}
Soit $\gamma(t) = z_0 + re^{it}$, $t \in [0, 2\pi]$. Pour $z_1$ à
l'intérieur du cercle ($\abs{z_1 - z_0} < r$),
\[
  \ind(\gamma, z_1) = \frac{1}{2\pi i}\int_\gamma \frac{\dd z}{z - z_1}.
\]
On montre que $\ind(\gamma, z_1) = 1$ (car $z_1$ est dans la composante
bornée de $\C \setminus \gamma^*$, et $\ind(\gamma, z_0) = 1$ par le calcul
de la proposition~\ref{prop:integrale-zn-cercle}).
Pour $z_1$ à l'extérieur ($\abs{z_1 - z_0} > r$), $\ind(\gamma, z_1) = 0$.
\end{exemple}

\begin{center}
\begin{tikzpicture}[scale=1.1, >=Stealth]
  \draw[->] (-3.5,0) -- (3.5,0) node[right] {$\re$};
  \draw[->] (0,-2.5) -- (0,2.5) node[above] {$\im$};

  % Cercle
  \draw[blue!70!black, very thick, decoration={markings,
    mark=at position 0.12 with {\arrow{>}},
    mark=at position 0.37 with {\arrow{>}},
    mark=at position 0.62 with {\arrow{>}},
    mark=at position 0.87 with {\arrow{>}}},
    postaction={decorate}] (0,0) circle (1.8);

  % Points
  \fill[red!70!black] (0,0) circle (2pt);
  \node[red!70!black, below right] at (0.1,0) {$z_0$};
  \node[red!70!black] at (0,0.7) {$\ind = 1$};

  \fill[green!50!black] (2.8,0.5) circle (2pt);
  \node[green!50!black, right] at (2.8,0.5) {$z_1$};
  \node[green!50!black] at (2.8,-0.3) {$\ind = 0$};

  % Enroulement double
  \begin{scope}[xshift=8cm]
    \draw[->] (-2.5,0) -- (2.5,0) node[right] {$\re$};
    \draw[->] (0,-2.5) -- (0,2.5) node[above] {$\im$};

    % Double cercle (2 tours)
    \draw[purple!70!black, very thick, decoration={markings,
      mark=at position 0.06 with {\arrow{>}},
      mark=at position 0.31 with {\arrow{>}},
      mark=at position 0.56 with {\arrow{>}},
      mark=at position 0.81 with {\arrow{>}}},
      postaction={decorate}]
      (0:1.5) arc (0:720:1.5);

    \fill[red!70!black] (0,0) circle (2pt);
    \node[red!70!black] at (0,0.5) {$\ind = 2$};
    \node[below] at (0,-2.8) {Lacet enroulé 2 fois};
  \end{scope}
\end{tikzpicture}
\captionof{figure}{Nombre d'enroulement : le lacet simple encercle $z_0$ une
fois ($\ind = 1$) et ne tourne pas autour de $z_1$ ($\ind = 0$). À droite,
un lacet parcourant deux fois le cercle donne $\ind = 2$.}\label{fig:indice}
\end{center}

%% ════════════════════════════════════════════════════════════════════
\section{Version générale du théorème de Cauchy}\label{sec:cauchy-general}
%% ════════════════════════════════════════════════════════════════════

La version générale du théorème de Cauchy fait intervenir l'indice et ne
requiert aucune hypothèse de simple connexité.

\begin{theoreme}[Cauchy --- version générale]\label{thm:cauchy-general}
Soit $U \subset \C$ un ouvert et $f \colon U \to \C$ holomorphe. Soit
$\gamma$ un lacet lisse par morceaux dans $U$ tel que
\[
  \ind(\gamma, z_0) = 0 \quad \text{pour tout } z_0 \in \C \setminus U.
\]
Alors
\[
  \int_\gamma f(z)\,\dd z = 0.
\]
\end{theoreme}

\begin{remarque}\label{rem:cauchy-general-interpretation}
La condition $\ind(\gamma, z_0) = 0$ pour tout $z_0 \notin U$ signifie que le
lacet $\gamma$ ne tourne autour d'aucun point extérieur à $U$. C'est une
condition \emph{homologique} sur $\gamma$ : on dit que $\gamma$ est
\emph{homologue à zéro} dans $U$\index{homologue à zéro}, et on note
$\gamma \sim 0$ dans $U$.

Le théorème affirme donc : si $f$ est holomorphe sur $U$ et $\gamma \sim 0$
dans $U$, alors $\int_\gamma f = 0$.
\end{remarque}

Nous admettons la preuve complète du théorème~\ref{thm:cauchy-general}
(qui sera établie au chapitre~5 comme conséquence de la formule intégrale
de Cauchy), et nous en développons plutôt les applications.

\begin{corollaire}[Déformation de contours]\label{cor:deformation-contours}
Soit $f$ holomorphe sur un ouvert $U$ et soient $\gamma_0, \gamma_1$ deux
lacets lisses par morceaux dans $U$ tels que
\[
  \ind(\gamma_0, z) = \ind(\gamma_1, z) \quad \text{pour tout }
  z \in \C \setminus U.
\]
Alors $\displaystyle \int_{\gamma_0} f(z)\,\dd z = \int_{\gamma_1} f(z)\,\dd z$.
\end{corollaire}

\begin{proof}
Le lacet $\gamma = \gamma_0 + \gamma_1^-$ vérifie $\ind(\gamma, z) =
\ind(\gamma_0, z) - \ind(\gamma_1, z) = 0$ pour tout $z \notin U$.
Par le théorème~\ref{thm:cauchy-general},
$0 = \int_\gamma f = \int_{\gamma_0} f - \int_{\gamma_1} f$.
\end{proof}

\begin{exemple}[Contournement d'une singularité]\label{ex:contournement}
Soit $f$ holomorphe sur $\C \setminus \{0\}$. Soient $C_1$ le cercle
$\abs{z} = 1$ et $C_2$ le cercle $\abs{z} = 2$, tous deux parcourus dans le
sens positif. Alors $\ind(C_1, 0) = \ind(C_2, 0) = 1$ et
$\ind(C_j, z) = 0$ pour $z \notin \C \setminus \{0\}$ (c'est-à-dire pour
$z = 0$, mais l'hypothèse est sur $z \notin U = \C\setminus\{0\}$, et
$0 \notin U$). Puisque les indices sur $\C \setminus U = \{0\}$ sont égaux
($= 1$), on a $\int_{C_1} f = \int_{C_2} f$.
\end{exemple}

\begin{theoreme}[Cauchy pour un domaine à bord multiple]\label{thm:cauchy-bord-multiple}
Soit $U$ un ouvert borné dont le bord est constitué de $n+1$ courbes de
Jordan lisses par morceaux : une courbe extérieure $\Gamma_0$ parcourue dans
le sens positif et des courbes intérieures $\Gamma_1, \ldots, \Gamma_n$
parcourues dans le sens négatif (horaire). Si $f$ est holomorphe sur un ouvert
contenant $\conj{U}$, alors
\[
  \int_{\Gamma_0} f(z)\,\dd z
  + \sum_{k=1}^{n} \int_{\Gamma_k} f(z)\,\dd z = 0,
\]
c'est-à-dire
\[
  \int_{\Gamma_0} f(z)\,\dd z
  = \sum_{k=1}^{n} \int_{\Gamma_k^+} f(z)\,\dd z,
\]
où $\Gamma_k^+$ désigne $\Gamma_k$ parcourue dans le sens positif.
\end{theoreme}

\begin{center}
\begin{tikzpicture}[scale=1.0, >=Stealth]
  % Domaine extérieur
  \fill[blue!8] plot[smooth cycle, tension=0.7]
    coordinates {(-3,-2) (0,-2.5) (3.5,-1.5) (3.5,2) (0.5,2.5) (-2.5,2) (-3.5,0)};
  % Trous
  \fill[white] (-1,0) circle (0.7);
  \fill[white] (1.5,0.3) circle (0.55);

  % Bord extérieur
  \draw[blue!70!black, very thick, decoration={markings,
    mark=at position 0.1 with {\arrow{>}},
    mark=at position 0.35 with {\arrow{>}},
    mark=at position 0.6 with {\arrow{>}},
    mark=at position 0.85 with {\arrow{>}}},
    postaction={decorate}]
    plot[smooth cycle, tension=0.7]
    coordinates {(-3,-2) (0,-2.5) (3.5,-1.5) (3.5,2) (0.5,2.5) (-2.5,2) (-3.5,0)};
  \node[blue!70!black] at (-2.8,2.3) {$\Gamma_0$ (sens $+$)};

  % Trou 1
  \draw[red!70!black, very thick, decoration={markings,
    mark=at position 0.25 with {\arrow{>}},
    mark=at position 0.75 with {\arrow{>}}},
    postaction={decorate}] (-1,0) circle (0.7);
  \node[red!70!black] at (-1,-1.1) {$\Gamma_1$ (sens $-$)};

  % Trou 2
  \draw[red!70!black, very thick, decoration={markings,
    mark=at position 0.25 with {\arrow{>}},
    mark=at position 0.75 with {\arrow{>}}},
    postaction={decorate}] (1.5,0.3) circle (0.55);
  \node[red!70!black] at (1.5,-0.65) {$\Gamma_2$ (sens $-$)};

  % Label domaine
  \node[font=\large] at (0.3,1.5) {$U$};

  % Singularités
  \fill (-1,0) circle (1.5pt);
  \fill (1.5,0.3) circle (1.5pt);
\end{tikzpicture}
\captionof{figure}{Domaine multiplement connexe avec bord orienté :
$\Gamma_0$ dans le sens trigonométrique, $\Gamma_1, \Gamma_2$ dans le sens
horaire.}\label{fig:bord-multiple}
\end{center}

%% ════════════════════════════════════════════════════════════════════
\section{Applications et exemples}\label{sec:applications-integration}
%% ════════════════════════════════════════════════════════════════════

\begin{exemple}[Calcul de $\displaystyle\int_\gamma \frac{\dd z}{z^2+1}$]\label{ex:calcul-z2+1}
Soit $\gamma$ le cercle $\abs{z} = 2$ parcouru dans le sens positif. La
fonction $f(z) = \frac{1}{z^2+1} = \frac{1}{(z-i)(z+i)}$ a des pôles en
$z = i$ et $z = -i$, tous deux à l'intérieur du cercle.

Par décomposition en éléments simples :
\[
  \frac{1}{z^2+1} = \frac{1}{2i}\left(\frac{1}{z-i} - \frac{1}{z+i}\right).
\]
Donc, en utilisant la proposition~\ref{prop:integrale-zn-cercle} et la
linéarité,
\[
  \int_\gamma \frac{\dd z}{z^2+1}
  = \frac{1}{2i}\left(\int_\gamma \frac{\dd z}{z-i}
    - \int_\gamma \frac{\dd z}{z+i}\right)
  = \frac{1}{2i}(2\pi i - 2\pi i) = 0.
\]
\end{exemple}

\begin{exemple}[Calcul de $\displaystyle\int_\gamma \frac{z}{z^2+1}\,\dd z$]\label{ex:calcul-z-sur-z2+1}
Avec le même cercle $\gamma$, $\abs{z} = 2$ :
\[
  \frac{z}{z^2+1} = \frac{z}{(z-i)(z+i)}
  = \frac{1}{2}\left(\frac{1}{z-i} + \frac{1}{z+i}\right).
\]
Donc
\[
  \int_\gamma \frac{z}{z^2+1}\,\dd z
  = \frac{1}{2}(2\pi i + 2\pi i) = 2\pi i.
\]
\end{exemple}

\begin{exemple}[Invariance par déformation]\label{ex:deformation}
Soit $f(z) = \frac{1}{z}$, holomorphe sur $\C^*$. Soit $\gamma$ un lacet
simple lisse par morceaux dans $\C^*$ tel que $\ind(\gamma, 0) = 1$.
Par le corollaire~\ref{cor:deformation-contours}, l'intégrale
$\int_\gamma \frac{\dd z}{z}$ est égale à l'intégrale sur le cercle unité
(qui a le même indice par rapport à $0$), donc
\[
  \int_\gamma \frac{\dd z}{z} = 2\pi i.
\]
Ce résultat illustre la puissance de la déformation de contours : peu importe
la forme du lacet, seul l'indice compte.
\end{exemple}

\begin{proposition}[Intégrale de $\frac{1}{z-a}$ le long d'un lacet]\label{prop:integrale-1-sur-z-a}
Soit $\gamma$ un lacet lisse par morceaux et $a \notin \gamma^*$. Alors
\[
  \int_\gamma \frac{\dd z}{z - a} = 2\pi i \, \ind(\gamma, a).
\]
\end{proposition}

\begin{proof}
C'est la définition même de l'indice (définition~\ref{def:indice}).
\end{proof}

%% ════════════════════════════════════════════════════════════════════
\section{Exercices}\label{sec:exercices-ch4}
%% ════════════════════════════════════════════════════════════════════

\begin{exercice}[Calculs directs d'intégrales curvilignes]\label{exo:calculs-directs}
Calculer les intégrales suivantes.
\begin{enumerate}[label=\textup{(\alph*)}]
  \item $\displaystyle\int_\gamma z^2\,\dd z$ où $\gamma$ est le segment de
    $0$ à $1+i$.
  \item $\displaystyle\int_\gamma \abs{z}^2\,\dd z$ où $\gamma$ est le segment
    de $0$ à $1+i$.
  \item $\displaystyle\int_\gamma \abs{z}^2\,\dd z$ où $\gamma$ est le chemin
    constitué du segment de $0$ à $1$ suivi du segment de $1$ à $1+i$.
  \item $\displaystyle\int_\gamma e^z\,\dd z$ où $\gamma$ est un chemin
    quelconque de $0$ à $\pi i$.
\end{enumerate}
\end{exercice}

\noindent\textit{Solution.}\enspace
\textbf{(a)} $z^2$ admet la primitive $z^3/3$, donc
$\int_\gamma z^2\,\dd z = (1+i)^3/3 - 0 = (1+3i+3i^2+i^3)/3 =
(1+3i-3-i)/3 = (-2+2i)/3$.

\textbf{(b)} Paramétrons $\gamma(t) = t(1+i)$, $t \in [0,1]$. Alors
$\gamma'(t) = 1+i$ et $\abs{\gamma(t)}^2 = t^2 \cdot 2$. Donc
\[
  \int_\gamma \abs{z}^2\,\dd z = \int_0^1 2t^2 \cdot (1+i)\,\dd t
  = (1+i) \cdot \frac{2}{3} = \frac{2(1+i)}{3}.
\]

\textbf{(c)} Sur le premier segment, $\gamma_1(t) = t$, $t \in [0,1]$,
$\abs{z}^2 = t^2$, $\gamma_1' = 1$ :
$\int_{\gamma_1} = \int_0^1 t^2\,\dd t = 1/3$. Sur le second,
$\gamma_2(t) = 1+it$, $t \in [0,1]$, $\abs{z}^2 = 1+t^2$, $\gamma_2' = i$ :
$\int_{\gamma_2} = i\int_0^1(1+t^2)\,\dd t = i(1+1/3) = 4i/3$. Total :
$1/3 + 4i/3$.

On constate que les réponses (b) et (c) diffèrent : $\abs{z}^2$ n'est pas
holomorphe, l'intégrale \emph{dépend} du chemin.

\textbf{(d)} $e^z$ admet la primitive $e^z$, donc
$\int_\gamma e^z\,\dd z = e^{\pi i} - e^0 = -1 - 1 = -2$. \hfill$\square$

\bigskip

\begin{exercice}[Inégalité ML en pratique]\label{exo:ML-pratique}
Soit $\gamma_R$ le demi-cercle supérieur $\gamma_R(t) = Re^{it}$,
$t \in [0, \pi]$, avec $R > 1$.
\begin{enumerate}[label=\textup{(\alph*)}]
  \item Montrer que $\displaystyle\abs{\int_{\gamma_R}
    \frac{\dd z}{z^2+1}} \leq \frac{\pi R}{R^2 - 1}$.
  \item En déduire que cette intégrale tend vers $0$ quand $R \to +\infty$.
\end{enumerate}
\end{exercice}

\noindent\textit{Solution.}\enspace
\textbf{(a)} Sur $\gamma_R$, $\abs{z} = R$ donc $\abs{z^2+1} \geq
\abs{z}^2 - 1 = R^2 - 1$ (inégalité triangulaire inverse). D'où
$\abs{1/(z^2+1)} \leq 1/(R^2-1)$. La longueur de $\gamma_R$ est
$\pi R$. Par l'inégalité ML :
\[
  \abs{\int_{\gamma_R} \frac{\dd z}{z^2+1}} \leq \frac{1}{R^2-1} \cdot \pi R
  = \frac{\pi R}{R^2-1}.
\]

\textbf{(b)} $\frac{\pi R}{R^2-1} = \frac{\pi}{R - 1/R} \to 0$ quand
$R \to +\infty$. \hfill$\square$

\bigskip

\begin{exercice}[Primitive de $1/z$ sur un demi-plan]\label{exo:primitive-1z}
Soit $U = \{z \in \C : \re(z) > 0\}$ le demi-plan droit.
\begin{enumerate}[label=\textup{(\alph*)}]
  \item Montrer que $U$ est simplement connexe.
  \item Montrer que $F(z) = \Log(z)$ (détermination principale du logarithme)
    est une primitive de $1/z$ sur $U$.
  \item Calculer $\displaystyle\int_\gamma \frac{\dd z}{z}$ où $\gamma$ est
    un chemin quelconque de $1$ à $i$ dans $U$.
\end{enumerate}
\end{exercice}

\noindent\textit{Solution.}\enspace
\textbf{(a)} Le demi-plan $U$ est convexe (pour $z_1, z_2 \in U$ et
$t \in [0,1]$, $\re((1-t)z_1+tz_2) = (1-t)\re(z_1)+t\re(z_2) > 0$). Tout
convexe est étoilé, et tout étoilé connexe est simplement connexe.

\textbf{(b)} $\Log(z) = \ln\abs{z} + i\Arg(z)$ est holomorphe sur
$\C \setminus \R^- \supset U$, et $\Log'(z) = 1/z$.

\textbf{(c)} Par le théorème~\ref{thm:primitive-indep},
\[
  \int_\gamma \frac{\dd z}{z} = \Log(i) - \Log(1)
  = (\ln 1 + i\pi/2) - 0 = \frac{\pi i}{2}. \quad\square
\]

\bigskip

\begin{exercice}[Lemme de Goursat et régularité]\label{exo:goursat-application}
Soit $f \colon \C \to \C$ continue, holomorphe sur $\C \setminus \{0\}$.
Montrer que $f$ est holomorphe sur $\C$ tout entier.

\noindent\textit{Indication.} Montrer que $\int_{\partial T} f = 0$ pour tout
triangle $T$ (distinguer selon que $0$ est sommet, sur un côté, ou à
l'intérieur de $T$), puis utiliser le théorème de Morera (chapitre~5).
\end{exercice}

\bigskip

\begin{exercice}[Calcul de l'indice]\label{exo:calcul-indice}
Soit $\gamma$ le lacet rectangulaire de sommets $-2-i$, $2-i$, $2+i$, $-2+i$
parcouru dans le sens positif.
\begin{enumerate}[label=\textup{(\alph*)}]
  \item Déterminer $\ind(\gamma, 0)$, $\ind(\gamma, 3)$ et $\ind(\gamma, i)$.
  \item Calculer $\displaystyle\int_\gamma \frac{\dd z}{z}$ et
    $\displaystyle\int_\gamma \frac{\dd z}{z-3}$.
\end{enumerate}
\end{exercice}

\noindent\textit{Solution.}\enspace
\textbf{(a)} Les points $0$ et $i$ sont à l'intérieur du rectangle
(car $\abs{\re(0)} < 2$ et $\abs{\im(0)} < 1$, et $\im(i) = 1$ est sur le
bord, donc exclu ; en fait $i$ est sur le bord $\gamma^*$, l'indice n'y est
pas défini).

Corrigeons : $i$ est un sommet du rectangle, donc $i \in \gamma^*$ et
$\ind(\gamma, i)$ n'est pas défini.

Pour $0$ : $0$ est à l'intérieur du rectangle, donc $\ind(\gamma, 0) = 1$.
Pour $3$ : $3$ est à l'extérieur ($\re(3) = 3 > 2$), donc
$\ind(\gamma, 3) = 0$.

\textbf{(b)} Par la proposition~\ref{prop:integrale-1-sur-z-a},
$\int_\gamma \frac{\dd z}{z} = 2\pi i \cdot 1 = 2\pi i$ et
$\int_\gamma \frac{\dd z}{z-3} = 2\pi i \cdot 0 = 0$. \hfill$\square$

\bigskip

\begin{exercice}[Théorème de Cauchy et intégrales réelles]\label{exo:cauchy-integrale-reelle}
En intégrant $f(z) = e^{-z^2}$ le long du contour rectangulaire de sommets
$-R$, $R$, $R+ia$, $-R+ia$ (avec $a > 0$ fixé et $R > 0$), montrer que
\[
  \int_{-\infty}^{+\infty} e^{-t^2}\cos(2at)\,\dd t = \sqrt{\pi}\,e^{-a^2}.
\]
\noindent\textit{Indication.} On admettra que
$\int_{-\infty}^{+\infty} e^{-t^2}\,\dd t = \sqrt{\pi}$.
\end{exercice}

\noindent\textit{Solution.}\enspace
La fonction $f(z) = e^{-z^2}$ est entière. Le contour rectangulaire $\gamma_R$
est un lacet dans $\C$ (simplement connexe), donc par le théorème de Cauchy,
$\oint_{\gamma_R} e^{-z^2}\,\dd z = 0$.

Le contour se décompose en quatre segments. Sur le côté inférieur ($z = t$,
$t \in [-R,R]$), l'intégrale vaut $\int_{-R}^R e^{-t^2}\,\dd t$.

Sur le côté supérieur ($z = t + ia$, $t$ de $R$ à $-R$), l'intégrale vaut
$-\int_{-R}^R e^{-(t+ia)^2}\,\dd t$. Or
$e^{-(t+ia)^2} = e^{-(t^2+2iat-a^2)} = e^{a^2-t^2}\,e^{-2iat}$.

Sur les côtés verticaux ($z = \pm R + is$, $s \in [0,a]$), on majore :
$\abs{e^{-(R+is)^2}} = e^{-(R^2-s^2)} \leq e^{-(R^2-a^2)} \to 0$ quand
$R \to \infty$. Donc ces intégrales tendent vers $0$.

En passant à la limite $R \to \infty$ :
\[
  0 = \int_{-\infty}^{+\infty} e^{-t^2}\,\dd t
    - e^{a^2}\int_{-\infty}^{+\infty} e^{-t^2}e^{-2iat}\,\dd t.
\]
Donc
\[
  \int_{-\infty}^{+\infty} e^{-t^2}e^{-2iat}\,\dd t = e^{-a^2}\sqrt{\pi}.
\]
En prenant la partie réelle :
$\int_{-\infty}^{+\infty} e^{-t^2}\cos(2at)\,\dd t = \sqrt{\pi}\,e^{-a^2}$.
\hfill$\square$

\bigskip

\begin{exercice}[Lacets et indices composés]\label{exo:lacets-composes}
Soit $\gamma$ le lacet en forme de « huit » défini comme suit : $\gamma$
parcourt d'abord le cercle $\abs{z-1} = 1/2$ dans le sens positif, puis le
cercle $\abs{z+1} = 1/2$ dans le sens négatif.
\begin{enumerate}[label=\textup{(\alph*)}]
  \item Déterminer $\ind(\gamma, 1)$, $\ind(\gamma, -1)$ et $\ind(\gamma, 0)$.
  \item Calculer $\displaystyle\int_\gamma \frac{\dd z}{z^2-1}$.
\end{enumerate}
\end{exercice}

\noindent\textit{Solution.}\enspace
\textbf{(a)} Par additivité de l'indice :
\begin{itemize}
  \item $\ind(\gamma, 1) = \ind(C_1^+, 1) + \ind(C_2^-, 1) = 1 + 0 = 1$
    (car $1$ est à l'intérieur de $C_1$ et à l'extérieur de $C_2$).
  \item $\ind(\gamma, -1) = \ind(C_1^+, -1) + \ind(C_2^-, -1) = 0 + (-1) = -1$
    (car $-1$ est à l'extérieur de $C_1$ et à l'intérieur de $C_2$,
    parcourue en sens négatif).
  \item $\ind(\gamma, 0) = 0 + 0 = 0$ (car $0$ est extérieur aux deux petits
    cercles).
\end{itemize}

\textbf{(b)} On a $\frac{1}{z^2-1} = \frac{1}{(z-1)(z+1)} =
\frac{1}{2}\left(\frac{1}{z-1} - \frac{1}{z+1}\right)$. Donc
\begin{align*}
  \int_\gamma \frac{\dd z}{z^2-1}
  &= \frac{1}{2}\left(\int_\gamma \frac{\dd z}{z-1}
    - \int_\gamma \frac{\dd z}{z+1}\right) \\
  &= \frac{1}{2}\bigl(2\pi i\,\ind(\gamma,1) - 2\pi i\,\ind(\gamma,-1)\bigr)
  = \frac{1}{2}(2\pi i - (-2\pi i)) = 2\pi i. \quad\square
\end{align*}

\bigskip

\begin{exercice}[Indice et topologie]\label{exo:indice-topologie}
Soit $\gamma$ un lacet lisse par morceaux dans $\C \setminus \{0,1\}$ tel que
$\ind(\gamma, 0) = 3$ et $\ind(\gamma, 1) = -2$.
\begin{enumerate}[label=\textup{(\alph*)}]
  \item Calculer $\displaystyle\int_\gamma \frac{\dd z}{z}$ et
    $\displaystyle\int_\gamma \frac{\dd z}{z-1}$.
  \item Calculer $\displaystyle\int_\gamma \frac{2z-1}{z(z-1)}\,\dd z$.
  \item Montrer que $\gamma$ n'est pas contractile dans
    $\C \setminus \{0,1\}$.
\end{enumerate}
\end{exercice}

\noindent\textit{Solution.}\enspace
\textbf{(a)} $\int_\gamma \frac{\dd z}{z} = 2\pi i \cdot 3 = 6\pi i$ et
$\int_\gamma \frac{\dd z}{z-1} = 2\pi i \cdot (-2) = -4\pi i$.

\textbf{(b)} $\frac{2z-1}{z(z-1)} = \frac{1}{z} + \frac{1}{z-1}$
(vérification : $\frac{z-1+z}{z(z-1)} = \frac{2z-1}{z(z-1)}$). Donc
\[
  \int_\gamma \frac{2z-1}{z(z-1)}\,\dd z = 6\pi i + (-4\pi i) = 2\pi i.
\]

\textbf{(c)} Si $\gamma$ était contractile dans $\C\setminus\{0,1\}$, on
aurait $\ind(\gamma, z) = 0$ pour tout $z \notin \C\setminus\{0,1\}$, en
particulier $\ind(\gamma, 0) = 0$ et $\ind(\gamma, 1) = 0$. Contradiction.
\hfill$\square$

\bigskip

\begin{exercice}[Formule de Cauchy pour l'indice et logarithme]\label{exo:indice-log}
Soit $\gamma \colon [0,1] \to \C^*$ un lacet lisse par morceaux. Montrer que
\[
  \ind(\gamma, 0) = \frac{1}{2\pi}\bigl[\arg\bigl(\gamma(t)\bigr)\bigr]_0^1,
\]
où $\arg(\gamma(t))$ désigne toute détermination continue de l'argument le
long de $\gamma$.

\noindent\textit{Indication.} Si $\theta(t)$ est une détermination continue
de l'argument de $\gamma(t)$, écrire $\gamma(t) = \abs{\gamma(t)}e^{i\theta(t)}$,
calculer $\gamma'(t)/\gamma(t)$, et intégrer.
\end{exercice}

\noindent\textit{Solution.}\enspace
Posons $r(t) = \abs{\gamma(t)}$ et $\theta(t)$ une détermination continue
de l'argument. Alors $\gamma(t) = r(t)e^{i\theta(t)}$ et
\[
  \frac{\gamma'(t)}{\gamma(t)} = \frac{r'(t)}{r(t)} + i\theta'(t).
\]
En intégrant,
\[
  \int_0^1 \frac{\gamma'(t)}{\gamma(t)}\,\dd t
  = \bigl[\ln r(t)\bigr]_0^1 + i\bigl[\theta(t)\bigr]_0^1.
\]
Comme $\gamma$ est un lacet, $r(1) = r(0)$, donc $[\ln r]_0^1 = 0$. Ainsi
$\int_\gamma \frac{\dd z}{z} = i(\theta(1)-\theta(0))$ et
\[
  \ind(\gamma, 0) = \frac{1}{2\pi i}\cdot i(\theta(1)-\theta(0))
  = \frac{\theta(1)-\theta(0)}{2\pi} = \frac{[\arg(\gamma(t))]_0^1}{2\pi}.
  \quad\square
\]

\bigskip

\begin{exercice}[Inégalité de la moyenne pour l'intégrale]\label{exo:inegalite-moyenne}
Soit $f$ holomorphe sur un ouvert contenant le disque fermé
$\overline{D}(0,R)$.
\begin{enumerate}[label=\textup{(\alph*)}]
  \item En utilisant le théorème de Cauchy, montrer que
    $\displaystyle\int_0^{2\pi} f(Re^{i\theta})\,\dd\theta = 2\pi f(0)$.

    \noindent\textit{Indication.} Poser $\gamma(t) = Re^{it}$ et relier
    $\int_\gamma \frac{f(z)}{z}\,\dd z$ à l'intégrale ci-dessus.
  \item En déduire que $\displaystyle\abs{f(0)} \leq
    \frac{1}{2\pi}\int_0^{2\pi}\abs{f(Re^{i\theta})}\,\dd\theta$.
\end{enumerate}
\end{exercice}

\noindent\textit{Solution.}\enspace
\textbf{(a)} Posons $\gamma(t) = Re^{it}$. Par le résultat de la
proposition~\ref{prop:integrale-zn-cercle} appliqué à la fonction
$f(z)/z$ (qui n'est holomorphe que si $f(0) = 0$, mais nous anticipons ici la
formule intégrale de Cauchy ou nous raisonnons directement) :

On calcule $\int_\gamma \frac{f(z)}{z}\,\dd z = \int_0^{2\pi}
\frac{f(Re^{it})}{Re^{it}}\cdot iRe^{it}\,\dd t = i\int_0^{2\pi}
f(Re^{it})\,\dd t$.

Ce résultat sera identifié à $2\pi i\,f(0)$ par la formule intégrale de
Cauchy (chapitre~5). En anticipant, $\int_0^{2\pi}f(Re^{i\theta})\,\dd\theta
= 2\pi f(0)$.

\textbf{(b)} Par l'inégalité triangulaire pour les intégrales :
\[
  2\pi\abs{f(0)} = \abs{\int_0^{2\pi}f(Re^{i\theta})\,\dd\theta}
  \leq \int_0^{2\pi}\abs{f(Re^{i\theta})}\,\dd\theta. \quad\square
\]

\bigskip

\begin{exercice}[Un lemme utile pour les résidus]\label{exo:lemme-residus}
Soit $f$ holomorphe sur $D(a, R) \setminus \{a\}$ avec un pôle simple en $a$.
Soit $\gamma_\varepsilon$ l'arc de cercle $\gamma_\varepsilon(t) = a +
\varepsilon e^{it}$, $t \in [\alpha, \beta]$, avec $0 < \varepsilon < R$.
\begin{enumerate}[label=\textup{(\alph*)}]
  \item Montrer que si $f(z) = \frac{c_{-1}}{z-a} + g(z)$ où $g$ est holomorphe
    sur $D(a,R)$, alors
    \[
      \lim_{\varepsilon \to 0} \int_{\gamma_\varepsilon} f(z)\,\dd z
      = i(\beta - \alpha)\,c_{-1}.
    \]
  \item En déduire le cas particulier $\alpha = 0$, $\beta = 2\pi$
    (cercle complet).
\end{enumerate}
\end{exercice}

\noindent\textit{Solution.}\enspace
\textbf{(a)} On a $\int_{\gamma_\varepsilon} f(z)\,\dd z =
c_{-1}\int_{\gamma_\varepsilon}\frac{\dd z}{z-a} +
\int_{\gamma_\varepsilon}g(z)\,\dd z$.

Pour le premier terme : $\int_{\gamma_\varepsilon}\frac{\dd z}{z-a} =
\int_\alpha^\beta \frac{i\varepsilon e^{it}}{\varepsilon e^{it}}\,\dd t =
i(\beta-\alpha)$.

Pour le second : $g$ est continue (donc bornée) sur $\overline{D}(a,
\varepsilon_0)$ pour $\varepsilon_0$ petit, disons $\abs{g(z)} \leq M$.
Par l'inégalité ML, $\abs{\int_{\gamma_\varepsilon}g(z)\,\dd z} \leq
M \cdot \varepsilon(\beta-\alpha) \to 0$.

Donc $\lim_{\varepsilon\to 0}\int_{\gamma_\varepsilon}f(z)\,\dd z =
i(\beta-\alpha)\,c_{-1}$.

\textbf{(b)} Pour $\alpha = 0$, $\beta = 2\pi$, on retrouve
$\int_{\gamma_\varepsilon}f(z)\,\dd z \to 2\pi i\,c_{-1} = 2\pi i\,\Res(f,a)$.
\hfill$\square$
%%%%%%%%%%%%%%%%%%%%%%%%%%%%%%%%%%%%%%%%%%%%%%%%%%%%%%%%%%%%%%%%%%%%
%  CHAPITRE 5 — Formule Intégrale de Cauchy et Conséquences
%%%%%%%%%%%%%%%%%%%%%%%%%%%%%%%%%%%%%%%%%%%%%%%%%%%%%%%%%%%%%%%%%%%%
\chapter{Formule Intégrale de Cauchy et Conséquences}\label{chap:cauchy-integral}

\section{Introduction}

La formule intégrale de Cauchy constitue l'un des résultats les plus profonds
et les plus féconds de l'analyse complexe. Elle exprime la valeur d'une fonction
holomorphe en un point intérieur à un contour fermé comme une intégrale curviligne
sur ce contour. De ce résultat fondamental découlent des conséquences remarquables :
formules pour les dérivées d'ordre supérieur, inégalités de Cauchy, théorème de
Liouville, développement en série de Taylor, et les principes du maximum et du minimum.

\section{Formule intégrale de Cauchy pour un disque}

\begin{theoreme}[Formule intégrale de Cauchy pour un disque]\label{thm:cauchy-integral-disk}
Soit $f$ holomorphe sur un ouvert $U \subset \C$ contenant le disque fermé
$\overline{D}(z_0,R) = \{z \in \C : \abs{z-z_0} \leq R\}$. Alors pour tout
$z \in D(z_0,R)$, on a
\[
  f(z) = \frac{1}{2\pi i} \oint_{\abs{\zeta - z_0} = R}
         \frac{f(\zeta)}{\zeta - z} \, \dd\zeta,
\]
où le cercle $\abs{\zeta - z_0} = R$ est parcouru dans le sens positif
(trigonométrique).
\end{theoreme}

\begin{proof}
Soit $z \in D(z_0,R)$ fixé. Posons $r = \abs{z - z_0} < R$ et notons
$C_R$ le cercle $\abs{\zeta - z_0} = R$ parcouru dans le sens positif,
et $C_\varepsilon$ le cercle $\abs{\zeta - z} = \varepsilon$ parcouru dans
le sens positif, avec $\varepsilon > 0$ assez petit pour que
$\overline{D}(z,\varepsilon) \subset D(z_0,R)$.

\medskip

\noindent\textbf{Étape 1 : Réduction par le théorème de Cauchy.}
La fonction $g(\zeta) = \dfrac{f(\zeta)}{\zeta - z}$ est holomorphe sur
$U \setminus \{z\}$. Par le théorème de Cauchy-Goursat appliqué à la
couronne $\{w : \varepsilon \leq \abs{w - z} \leq R'\}$ (avec $R'$
convenablement choisi et en reliant les deux cercles par des segments),
on obtient
\[
  \oint_{C_R} \frac{f(\zeta)}{\zeta - z} \, \dd\zeta
  = \oint_{C_\varepsilon} \frac{f(\zeta)}{\zeta - z} \, \dd\zeta.
\]

Plus précisément, on considère le domaine simplement connexe obtenu en
retirant une fente radiale de la couronne. Les intégrales le long des
deux bords de la fente s'annulent car elles sont parcourues en sens
opposés, ce qui donne l'égalité ci-dessus.

\medskip

\noindent\textbf{Étape 2 : Évaluation sur le petit cercle.}
On paramétrise $C_\varepsilon$ par $\zeta = z + \varepsilon e^{i\theta}$,
$\theta \in [0, 2\pi]$. Alors $\dd\zeta = i\varepsilon e^{i\theta}\,\dd\theta$
et
\[
  \oint_{C_\varepsilon} \frac{f(\zeta)}{\zeta - z}\,\dd\zeta
  = \int_0^{2\pi} \frac{f(z + \varepsilon e^{i\theta})}
    {\varepsilon e^{i\theta}} \cdot i\varepsilon e^{i\theta}\,\dd\theta
  = i \int_0^{2\pi} f(z + \varepsilon e^{i\theta})\,\dd\theta.
\]

Puisque $f$ est continue en $z$, pour tout $\delta > 0$ il existe
$\varepsilon_0 > 0$ tel que pour $\varepsilon < \varepsilon_0$,
\[
  \abs{f(z + \varepsilon e^{i\theta}) - f(z)} < \delta
  \quad \text{pour tout } \theta \in [0, 2\pi].
\]
Par conséquent,
\[
  \abs{\oint_{C_\varepsilon} \frac{f(\zeta)}{\zeta - z}\,\dd\zeta
       - 2\pi i \, f(z)}
  = \abs{i\int_0^{2\pi}\bigl[f(z + \varepsilon e^{i\theta}) - f(z)\bigr]
    \,\dd\theta}
  \leq 2\pi \delta.
\]

Comme le membre de gauche est l'intégrale sur $C_R$ (qui ne dépend pas
de $\varepsilon$) moins $2\pi i\,f(z)$, et comme $\delta > 0$ est
arbitraire, on conclut
\[
  \oint_{C_R} \frac{f(\zeta)}{\zeta - z}\,\dd\zeta = 2\pi i\,f(z),
\]
ce qui achève la preuve.
\end{proof}

\begin{center}
\begin{tikzpicture}[scale=1.2, >=Stealth]
  % Grand cercle C_R
  \draw[thick, blue!70!black, postaction={decorate,
    decoration={markings,
      mark=at position 0.15 with {\arrow{>}},
      mark=at position 0.4 with {\arrow{>}},
      mark=at position 0.65 with {\arrow{>}},
      mark=at position 0.9 with {\arrow{>}}
    }}] (0,0) circle (2.5);
  % Petit cercle C_epsilon
  \draw[thick, red!70!black, postaction={decorate,
    decoration={markings,
      mark=at position 0.25 with {\arrow{>}},
      mark=at position 0.75 with {\arrow{>}}
    }}] (0.8,0.4) circle (0.5);
  % Points
  \fill (0,0) circle (1.5pt) node[below left] {$z_0$};
  \fill (0.8,0.4) circle (1.5pt) node[below right] {$z$};
  % Labels
  \node[blue!70!black] at (2.1,1.8) {$C_R$};
  \node[red!70!black] at (1.7,0.9) {$C_\varepsilon$};
  % Rayon
  \draw[dashed, gray] (0,0) -- node[above, font=\small] {$R$} (2.5*0.707, 2.5*0.707);
  \draw[dashed, gray] (0.8,0.4) -- node[right, font=\small] {$\varepsilon$} (0.8+0.5*0.707, 0.4+0.5*0.707);
\end{tikzpicture}
\captionof{figure}{Contour d'intégration pour la formule de Cauchy : le grand
cercle $C_R$ et le petit cercle $C_\varepsilon$ autour de $z$.}
\end{center}

\section{Formule intégrale de Cauchy --- version générale}

\begin{theoreme}[Formule intégrale de Cauchy --- version générale]\label{thm:cauchy-integral-general}
Soit $U \subset \C$ un ouvert, $f \colon U \to \C$ holomorphe, et
$\gamma$ un cycle dans $U$ tel que $\ind(\gamma, w) = 0$ pour tout
$w \notin U$. Alors pour tout $z \in U \setminus \gamma^*$,
\[
  \ind(\gamma, z) \cdot f(z) = \frac{1}{2\pi i}
  \oint_\gamma \frac{f(\zeta)}{\zeta - z}\,\dd\zeta,
\]
où $\gamma^*$ désigne l'image de $\gamma$.
\end{theoreme}

\begin{remarque}[Rôle de l'indice]\label{rem:role-indice}
Lorsque $\gamma$ est un cercle parcouru une fois dans le sens positif
et $z$ est à l'intérieur, $\ind(\gamma, z) = 1$ et l'on retrouve la
formule du théorème~\ref{thm:cauchy-integral-disk}. Pour un point
extérieur, $\ind(\gamma, z) = 0$ et l'intégrale est nulle.
\end{remarque}

\begin{exemple}[Calcul direct par la formule de Cauchy]\label{ex:cauchy-calc-1}
Calculons $\displaystyle I = \oint_{\abs{z}=2} \frac{e^z}{z-1}\,\dd z$.

La fonction $f(z) = e^z$ est entière, et $z_0 = 1$ est à l'intérieur
du cercle $\abs{z} = 2$. Par la formule intégrale de Cauchy,
\[
  I = 2\pi i \cdot f(1) = 2\pi i \cdot e.
\]
\end{exemple}

\begin{exemple}[Point extérieur au contour]\label{ex:cauchy-calc-2}
Calculons $\displaystyle J = \oint_{\abs{z}=1} \frac{\sin z}{z - 3}\,\dd z$.

Le point $z_0 = 3$ est à l'\emph{extérieur} du cercle $\abs{z}=1$,
donc la fonction $\zeta \mapsto \dfrac{\sin \zeta}{\zeta - 3}$ est
holomorphe sur un ouvert contenant $\overline{D}(0,1)$. Par le
théorème de Cauchy-Goursat, $J = 0$.
\end{exemple}

\section{Formule de Cauchy pour les dérivées}

\begin{theoreme}[Formule de Cauchy pour les dérivées]\label{thm:cauchy-derivatives}
Soit $f$ holomorphe sur un ouvert $U$ et soit $\gamma$ un cercle positivement
orienté dans $U$ dont l'intérieur est contenu dans $U$. Alors $f$ est
indéfiniment dérivable sur l'intérieur de $\gamma$, et pour tout entier
$n \geq 0$ et tout $z$ à l'intérieur de $\gamma$,
\[
  f^{(n)}(z) = \frac{n!}{2\pi i} \oint_\gamma
  \frac{f(\zeta)}{(\zeta - z)^{n+1}}\,\dd\zeta.
\]
\end{theoreme}

\begin{proof}
Nous démontrons le résultat par récurrence sur $n$.

\medskip
\noindent\textbf{Cas $n = 0$.}
C'est la formule intégrale de Cauchy (théorème~\ref{thm:cauchy-integral-disk}).

\medskip
\noindent\textbf{Cas $n = 1$ (passage au quotient différentiel).}
Soit $z$ à l'intérieur de $\gamma$ et $h \in \C^*$ assez petit pour
que $z + h$ soit aussi à l'intérieur. Par la formule de Cauchy,
\begin{align*}
  \frac{f(z+h) - f(z)}{h}
  &= \frac{1}{2\pi i} \oint_\gamma f(\zeta)
     \cdot \frac{1}{h}\left(\frac{1}{\zeta - z - h}
     - \frac{1}{\zeta - z}\right)\dd\zeta \\
  &= \frac{1}{2\pi i} \oint_\gamma
     \frac{f(\zeta)}{(\zeta - z - h)(\zeta - z)}\,\dd\zeta.
\end{align*}

Notons $d = \min_{\zeta \in \gamma^*} \abs{\zeta - z} > 0$. Pour
$\abs{h} < d/2$, on a $\abs{\zeta - z - h} \geq d/2$ pour tout
$\zeta \in \gamma^*$, et donc
\[
  \frac{f(\zeta)}{(\zeta - z - h)(\zeta - z)}
  - \frac{f(\zeta)}{(\zeta - z)^2}
  = \frac{h \cdot f(\zeta)}{(\zeta - z)^2(\zeta - z - h)}.
\]
En posant $M = \max_{\zeta \in \gamma^*} \abs{f(\zeta)}$ et
$L = \text{longueur}(\gamma)$, on obtient la majoration
\[
  \abs{\frac{f(z+h) - f(z)}{h}
  - \frac{1}{2\pi i}\oint_\gamma \frac{f(\zeta)}{(\zeta-z)^2}\,\dd\zeta}
  \leq \frac{1}{2\pi} \cdot \frac{\abs{h} M L}{d^2 \cdot (d/2)}
  = \frac{M L \abs{h}}{\pi d^3} \xrightarrow[h \to 0]{} 0.
\]
Ainsi $f'(z) = \dfrac{1}{2\pi i}\displaystyle\oint_\gamma
\dfrac{f(\zeta)}{(\zeta - z)^2}\,\dd\zeta$.

\medskip
\noindent\textbf{Cas général (récurrence).}
Le passage de $n$ à $n+1$ suit exactement le même schéma. On forme
le quotient différentiel
\[
  \frac{f^{(n)}(z+h) - f^{(n)}(z)}{h}
  = \frac{n!}{2\pi i}\oint_\gamma f(\zeta) \cdot \frac{1}{h}
    \left[\frac{1}{(\zeta - z - h)^{n+1}}
    - \frac{1}{(\zeta - z)^{n+1}}\right]\dd\zeta,
\]
et l'on montre par une majoration analogue que cette expression tend vers
\[
  \frac{(n+1)!}{2\pi i}\oint_\gamma
  \frac{f(\zeta)}{(\zeta - z)^{n+2}}\,\dd\zeta
\]
lorsque $h \to 0$. Le facteur $(n+1)$ provient de l'identité algébrique
\[
  \lim_{h \to 0} \frac{1}{h}\left[\frac{1}{(\zeta - z - h)^{n+1}}
  - \frac{1}{(\zeta - z)^{n+1}}\right]
  = \frac{n+1}{(\zeta - z)^{n+2}},
\]
et la convergence uniforme en $\zeta \in \gamma^*$ permet d'intervertir
limite et intégrale.
\end{proof}

\begin{corollaire}[Holomorphie implique analyticité]\label{cor:holo-implies-analytic}
Toute fonction holomorphe est indéfiniment dérivable. En particulier,
si $f$ est holomorphe sur un ouvert $U$, alors $f'$, $f''$, \ldots\ sont
toutes holomorphes sur $U$.
\end{corollaire}

\begin{proof}
Le théorème~\ref{thm:cauchy-derivatives} montre que $f^{(n)}$ existe
pour tout $n \geq 0$ en tout point de $U$. Chaque $f^{(n)}$ s'exprime
par une intégrale de Cauchy et est donc continue, puis holomorphe.
\end{proof}

\begin{exemple}[Application de la formule des dérivées]\label{ex:cauchy-deriv-calc}
Calculons $\displaystyle I = \oint_{\abs{z}=1} \frac{e^z}{z^4}\,\dd z$.

On reconnaît la formule de Cauchy pour la dérivée avec $f(\zeta) = e^\zeta$,
$z_0 = 0$, et $n = 3$. Ainsi
\[
  I = \frac{2\pi i}{3!}\,f'''(0) = \frac{2\pi i}{6}\cdot e^0 = \frac{\pi i}{3}.
\]
\end{exemple}

\section{Théorème de Morera}

\begin{theoreme}[Morera]\label{thm:morera}
Soit $f \colon U \to \C$ continue sur un ouvert $U \subset \C$. Si
\[
  \oint_T f(z)\,\dd z = 0
\]
pour tout triangle $T$ dont l'intérieur est contenu dans $U$, alors $f$
est holomorphe sur $U$.
\end{theoreme}

\begin{proof}
Il suffit de montrer que $f$ est holomorphe au voisinage de chaque point
de $U$. Soit $z_0 \in U$ et $r > 0$ tel que $D(z_0, r) \subset U$.
Définissons $F \colon D(z_0, r) \to \C$ par
\[
  F(z) = \int_{[z_0, z]} f(\zeta)\,\dd\zeta,
\]
où $[z_0, z]$ désigne le segment de $z_0$ à $z$. L'hypothèse sur les
intégrales triangulaires assure que cette intégrale ne dépend pas du
chemin (dans le disque), et donc $F$ est bien définie.

Montrons que $F'(z) = f(z)$ pour tout $z \in D(z_0, r)$. Pour $h$
assez petit,
\[
  F(z+h) - F(z) = \int_{[z, z+h]} f(\zeta)\,\dd\zeta,
\]
d'où
\[
  \frac{F(z+h) - F(z)}{h} - f(z)
  = \frac{1}{h}\int_{[z,z+h]} \bigl[f(\zeta) - f(z)\bigr]\,\dd\zeta.
\]
Par continuité de $f$, pour tout $\varepsilon > 0$, il existe $\delta > 0$
tel que $\abs{f(\zeta) - f(z)} < \varepsilon$ pour $\abs{\zeta - z} < \delta$.
Si $\abs{h} < \delta$,
\[
  \abs{\frac{F(z+h)-F(z)}{h} - f(z)}
  \leq \frac{1}{\abs{h}} \cdot \abs{h} \cdot \varepsilon = \varepsilon.
\]
Donc $F'(z) = f(z)$, et $F$ est holomorphe sur $D(z_0, r)$. Par le
corollaire~\ref{cor:holo-implies-analytic}, $f = F'$ est aussi holomorphe
sur $D(z_0,r)$.
\end{proof}

\begin{remarque}[Réciproque partielle]\label{rem:morera-reciproque}
Le théorème de Morera est une réciproque partielle du théorème de
Cauchy-Goursat : Cauchy-Goursat dit que l'holomorphie implique la
nullité des intégrales le long de contours fermés ; Morera dit que
la nullité des intégrales triangulaires (avec la continuité) implique
l'holomorphie.
\end{remarque}

\section{Inégalités de Cauchy}

\begin{theoreme}[Inégalités de Cauchy]\label{thm:cauchy-inequalities}
Soit $f$ holomorphe sur un ouvert contenant $\overline{D}(z_0, R)$. Alors
pour tout $n \geq 0$,
\[
  \abs{f^{(n)}(z_0)} \leq \frac{n!\,M_R}{R^n},
\]
où $M_R = \max_{\abs{z - z_0} = R} \abs{f(z)}$.
\end{theoreme}

\begin{proof}
Par la formule de Cauchy pour les dérivées
(théorème~\ref{thm:cauchy-derivatives}), en paramétrant le cercle
$C_R : \zeta = z_0 + Re^{i\theta}$,
\begin{align*}
  \abs{f^{(n)}(z_0)}
  &= \abs{\frac{n!}{2\pi i}\oint_{C_R}
     \frac{f(\zeta)}{(\zeta - z_0)^{n+1}}\,\dd\zeta} \\
  &\leq \frac{n!}{2\pi} \cdot \frac{M_R}{R^{n+1}} \cdot 2\pi R
  = \frac{n!\,M_R}{R^n}.
\end{align*}
On a utilisé l'inégalité ML : $\abs{\oint_\gamma g(\zeta)\,\dd\zeta}
\leq \max_{\gamma^*}\abs{g} \cdot \text{longueur}(\gamma) = \dfrac{M_R}{R^{n+1}}
\cdot 2\pi R$.
\end{proof}

\section{Théorème de Liouville}

\begin{theoreme}[Liouville]\label{thm:liouville}
Toute fonction entière bornée est constante. Autrement dit, si
$f \colon \C \to \C$ est holomorphe et s'il existe $M > 0$ tel que
$\abs{f(z)} \leq M$ pour tout $z \in \C$, alors $f$ est constante.
\end{theoreme}

\begin{proof}
Soit $z_0 \in \C$ quelconque. Pour tout $R > 0$, l'inégalité de Cauchy
pour $n = 1$ donne
\[
  \abs{f'(z_0)} \leq \frac{M}{R}.
\]
En faisant tendre $R \to +\infty$, on obtient $f'(z_0) = 0$. Ceci étant
vrai pour tout $z_0 \in \C$, la fonction $f$ est de dérivée nulle sur
$\C$ (connexe), donc constante.
\end{proof}

\begin{theoreme}[Théorème fondamental de l'algèbre]\label{thm:fta}
Tout polynôme non constant $P \in \C[z]$ admet au moins une racine dans $\C$.
\end{theoreme}

\begin{proof}
Supposons par l'absurde que $P(z) \neq 0$ pour tout $z \in \C$.
Alors $g(z) = \dfrac{1}{P(z)}$ est entière. Écrivons
$P(z) = a_n z^n + \cdots + a_0$ avec $a_n \neq 0$ et $n \geq 1$. Pour
$\abs{z}$ assez grand,
\[
  \abs{P(z)} = \abs{a_n}\abs{z}^n
  \abs{1 + \frac{a_{n-1}}{a_n z} + \cdots + \frac{a_0}{a_n z^n}}
  \geq \frac{\abs{a_n}}{2}\abs{z}^n,
\]
car le terme entre barres de valeur absolue tend vers $1$. Plus
précisément, il existe $R_0 > 0$ tel que pour $\abs{z} \geq R_0$,
$\abs{P(z)} \geq \frac{\abs{a_n}}{2}\abs{z}^n$.

Donc pour $\abs{z} \geq R_0$,
$\abs{g(z)} \leq \dfrac{2}{\abs{a_n}\abs{z}^n} \leq \dfrac{2}{\abs{a_n}R_0^n}$.

Sur le compact $\overline{D}(0, R_0)$, $g$ est continue (car $P$ ne
s'annule pas) donc bornée. Ainsi $g$ est bornée sur $\C$ tout entier.
Par le théorème de Liouville, $g$ est constante, ce qui contredit
$\deg P \geq 1$.
\end{proof}

\begin{corollaire}[Factorisation complète]\label{cor:factorisation-complete}
Tout polynôme $P \in \C[z]$ de degré $n \geq 1$ admet exactement $n$
racines dans $\C$ (comptées avec multiplicité) :
\[
  P(z) = a_n(z - z_1)^{m_1}(z - z_2)^{m_2}\cdots(z - z_k)^{m_k},
  \quad m_1 + m_2 + \cdots + m_k = n.
\]
\end{corollaire}

\section{Théorème de la moyenne}

\begin{theoreme}[Propriété de la valeur moyenne]\label{thm:mean-value}
Soit $f$ holomorphe sur un ouvert contenant $\overline{D}(z_0,R)$. Alors
\[
  f(z_0) = \frac{1}{2\pi}\int_0^{2\pi} f(z_0 + Re^{i\theta})\,\dd\theta.
\]
\end{theoreme}

\begin{proof}
On paramétrise le cercle $\abs{\zeta - z_0} = R$ par
$\zeta = z_0 + Re^{i\theta}$, $\dd\zeta = iRe^{i\theta}\,\dd\theta$.
Par la formule intégrale de Cauchy,
\[
  f(z_0) = \frac{1}{2\pi i}\oint_{\abs{\zeta - z_0}=R}
  \frac{f(\zeta)}{\zeta - z_0}\,\dd\zeta
  = \frac{1}{2\pi i}\int_0^{2\pi}
  \frac{f(z_0 + Re^{i\theta})}{Re^{i\theta}}\cdot iRe^{i\theta}\,\dd\theta
  = \frac{1}{2\pi}\int_0^{2\pi}f(z_0 + Re^{i\theta})\,\dd\theta.
\]
\end{proof}

\begin{remarque}[Interprétation]\label{rem:mean-value-interp}
La valeur d'une fonction holomorphe en un point est la moyenne de ses
valeurs sur tout cercle centré en ce point et contenu dans le domaine
d'holomorphie. C'est une propriété caractéristique des fonctions
harmoniques (dont les parties réelle et imaginaire d'une fonction
holomorphe sont des exemples).
\end{remarque}

\section{Principe du maximum}

\begin{theoreme}[Principe du maximum]\label{thm:max-principle}
Soit $U \subset \C$ un ouvert connexe et $f \colon U \to \C$ holomorphe.
Si $\abs{f}$ atteint son maximum en un point de $U$, alors $f$ est
constante sur $U$.
\end{theoreme}

\begin{proof}
Supposons qu'il existe $z_0 \in U$ tel que $\abs{f(z_0)} \geq \abs{f(z)}$
pour tout $z \in U$. Posons $M = \abs{f(z_0)}$.

\medskip
\noindent\textbf{Cas $M = 0$.} Alors $f \equiv 0$ sur $U$.

\medskip
\noindent\textbf{Cas $M > 0$.} Soit $R > 0$ tel que
$\overline{D}(z_0,R) \subset U$. Par la propriété de la valeur
moyenne (théorème~\ref{thm:mean-value}),
\[
  M = \abs{f(z_0)} = \abs{\frac{1}{2\pi}\int_0^{2\pi}
  f(z_0 + Re^{i\theta})\,\dd\theta}
  \leq \frac{1}{2\pi}\int_0^{2\pi}
  \abs{f(z_0 + Re^{i\theta})}\,\dd\theta \leq M.
\]
La dernière inégalité vient de $\abs{f} \leq M$ sur $U$. On a donc
égalité partout, ce qui nécessite
\[
  \abs{f(z_0 + Re^{i\theta})} = M \quad \text{pour tout } \theta \in [0,2\pi].
\]

En effet, si $\abs{f(z_0 + R e^{i\theta_0})} < M$ pour un certain
$\theta_0$, par continuité, il existerait un arc sur lequel
$\abs{f} < M - \varepsilon$, ce qui donnerait une inégalité stricte,
contradiction.

Ainsi $\abs{f}$ est constante (égale à $M$) sur tout cercle de centre
$z_0$ contenu dans $U$, donc sur tout le disque $D(z_0, R)$. Écrivons
$f(z) = Me^{i\varphi(z)}$ sur ce disque ; une fonction holomorphe de
module constant est constante (car $\abs{f}^2 = u^2 + v^2$ est constant,
et les équations de Cauchy-Riemann entraînent $f' = 0$).

Donc $f$ est constante sur $D(z_0,R)$. Par le principe du prolongement
analytique (ou par connexité de $U$ et le fait que l'ensemble des points
où $f$ est constante est à la fois ouvert et fermé dans $U$), $f$ est
constante sur $U$.
\end{proof}

\begin{corollaire}[Principe du maximum pour les domaines bornés]\label{cor:max-principle-bounded}
Soit $U$ un ouvert borné et connexe, et $f$ continue sur $\overline{U}$
et holomorphe sur $U$. Alors
\[
  \max_{z \in \overline{U}} \abs{f(z)} = \max_{z \in \partial U} \abs{f(z)}.
\]
En d'autres termes, le maximum de $\abs{f}$ sur l'adhérence est atteint
sur le bord.
\end{corollaire}

\begin{proof}
Puisque $\overline{U}$ est compact et $\abs{f}$ est continue, le maximum
est atteint en un certain $z_1 \in \overline{U}$. Si $z_1 \in U$, le
théorème~\ref{thm:max-principle} implique que $f$ est constante, et le
maximum est atteint partout, en particulier sur $\partial U$. Si
$z_1 \in \partial U$, c'est immédiat.
\end{proof}

\begin{theoreme}[Principe du minimum]\label{thm:min-principle}
Soit $U \subset \C$ un ouvert connexe et $f \colon U \to \C$ holomorphe
et non identiquement nulle. Si $\abs{f}$ atteint son minimum en un point
de $U$, alors $f$ est constante.
\end{theoreme}

\begin{proof}
Supposons que $\abs{f}$ atteint son minimum en $z_0 \in U$.

Si $f(z_0) = 0$, alors le minimum de $\abs{f}$ est $0$, ce qui signifie
que $f$ s'annule en $z_0$. Dans ce cas, le résultat ne donne pas
directement la constance (il faut que $f$ ne s'annule pas). Mais si
$f(z_0) \neq 0$, on considère $g = 1/f$, qui est holomorphe sur $U$
(car $f$ ne s'annule pas --- en effet, si $f(z_1) = 0$ pour un certain
$z_1 \in U$, on aurait $\abs{f(z_1)} = 0 < \abs{f(z_0)}$, contredisant
la minimalité en $z_0$). Alors $\abs{g}$ atteint son maximum en $z_0$,
et par le principe du maximum, $g$ est constante, donc $f$ l'est aussi.
\end{proof}

\begin{center}
\begin{tikzpicture}[scale=1.0, >=Stealth]
  % Domaine
  \draw[thick, fill=blue!5] plot[smooth cycle, tension=0.7]
    coordinates {(-2.5,0) (-1.8,1.5) (0,2.2) (1.8,1.6) (2.5,0.3)
                 (2,-1.3) (0.5,-2) (-1,-1.8) (-2.3,-1)};
  \node at (0,0) {$U$};
  % Point intérieur
  \fill[red] (0.5,0.8) circle (2pt) node[above right] {$z_0$};
  % Bord
  \node[blue!70!black] at (2.8,1) {$\partial U$};
  % Annotation
  \node[below, font=\small] at (0,-2.5) {$\max_{\overline{U}}\abs{f}$
    est atteint sur $\partial U$};
\end{tikzpicture}
\captionof{figure}{Principe du maximum : sur un domaine borné, le
maximum de $\abs{f}$ ne peut être atteint qu'au bord.}
\end{center}

\section{Développement en série de Taylor}

\begin{theoreme}[Développement de Taylor]\label{thm:taylor-series}
Soit $f$ holomorphe sur un ouvert $U \subset \C$ et $z_0 \in U$.
Soit $R = d(z_0, \C \setminus U)$ le rayon du plus grand disque ouvert
de centre $z_0$ contenu dans $U$ (avec $R = +\infty$ si $U = \C$).
Alors $f$ admet un développement en série entière convergente sur
$D(z_0,R)$ :
\[
  f(z) = \sum_{n=0}^{\infty} a_n (z - z_0)^n, \quad \abs{z - z_0} < R,
\]
avec $a_n = \dfrac{f^{(n)}(z_0)}{n!}$. La convergence est normale
(donc uniforme) sur tout compact de $D(z_0,R)$.
\end{theoreme}

\begin{proof}
Soit $z \in D(z_0, R)$ fixé. Choisissons $r$ tel que
$\abs{z - z_0} < r < R$ et notons $C_r$ le cercle $\abs{\zeta - z_0} = r$
parcouru positivement. Par la formule intégrale de Cauchy,
\[
  f(z) = \frac{1}{2\pi i}\oint_{C_r}
  \frac{f(\zeta)}{\zeta - z}\,\dd\zeta.
\]

La clé est de développer le noyau de Cauchy en série géométrique. On
écrit
\[
  \frac{1}{\zeta - z}
  = \frac{1}{(\zeta - z_0) - (z - z_0)}
  = \frac{1}{\zeta - z_0}\cdot
    \frac{1}{1 - \dfrac{z - z_0}{\zeta - z_0}}.
\]
Pour $\zeta \in C_r$, on a $\abs{\zeta - z_0} = r$ et
$\abs{z - z_0} < r$, donc
$\abs{\dfrac{z-z_0}{\zeta - z_0}} = \dfrac{\abs{z-z_0}}{r} < 1$.
La série géométrique converge :
\[
  \frac{1}{\zeta - z}
  = \sum_{n=0}^{\infty} \frac{(z - z_0)^n}{(\zeta - z_0)^{n+1}}.
\]
La convergence est uniforme en $\zeta \in C_r$ car
$\abs{\dfrac{z-z_0}{\zeta - z_0}} \leq \dfrac{\abs{z-z_0}}{r} < 1$
est indépendant de $\zeta$. On peut donc intervertir somme et intégrale :
\begin{align*}
  f(z) &= \frac{1}{2\pi i}\oint_{C_r}
  \sum_{n=0}^{\infty}\frac{f(\zeta)\,(z-z_0)^n}{(\zeta - z_0)^{n+1}}
  \,\dd\zeta \\
  &= \sum_{n=0}^{\infty}\left[\frac{1}{2\pi i}\oint_{C_r}
  \frac{f(\zeta)}{(\zeta - z_0)^{n+1}}\,\dd\zeta\right](z-z_0)^n.
\end{align*}

Par la formule de Cauchy pour les dérivées
(théorème~\ref{thm:cauchy-derivatives}),
\[
  \frac{1}{2\pi i}\oint_{C_r}\frac{f(\zeta)}{(\zeta-z_0)^{n+1}}\,\dd\zeta
  = \frac{f^{(n)}(z_0)}{n!} = a_n.
\]
On obtient bien $f(z) = \displaystyle\sum_{n=0}^\infty a_n(z-z_0)^n$.

Pour la convergence normale, soit $K \subset D(z_0,R)$ compact. Il
existe $\rho < R$ tel que $K \subset \overline{D}(z_0,\rho)$. Choisissons
$r$ avec $\rho < r < R$. Pour $z \in K$,
\[
  \abs{a_n(z-z_0)^n} \leq \frac{M_r}{r^n}\rho^n
  = M_r\left(\frac{\rho}{r}\right)^n,
\]
où $M_r = \max_{\abs{\zeta - z_0}=r}\abs{f(\zeta)}$. Puisque
$\rho/r < 1$, la série $\sum M_r(\rho/r)^n$ converge, ce qui prouve
la convergence normale sur $K$.
\end{proof}

\begin{center}
\begin{tikzpicture}[scale=0.85, >=Stealth]
  % Domaine U (forme irrégulière)
  \draw[thick, blue!50!black, fill=blue!5] plot[smooth cycle, tension=0.65]
    coordinates {(-3.5,0.5) (-2.5,2.8) (-0.5,3.5) (2,3) (3.5,1.5)
                 (4,-0.5) (3,-2.5) (1,-3.5) (-1.5,-3) (-3,-1.5)};
  \node[blue!50!black] at (3.2,2.8) {$U$};
  % Point z_0
  \fill (0,0) circle (2pt) node[below left] {$z_0$};
  % Disque de convergence
  \draw[thick, red!70!black, dashed] (0,0) circle (2.5);
  \node[red!70!black] at (1.8,-2.0) {$D(z_0,R)$};
  % Rayon R
  \draw[<->, red!70!black] (0,0) -- node[above, font=\small] {$R$} (2.5*0.866, 2.5*0.5);
  % Point du bord le plus proche
  \fill[red] (-2.5*0.866+0.15, -2.5*0.5-0.45) circle (0pt);
  % Annotation singularité
  \node[font=\small, red!70!black] at (-2.2, -2.8) {Point de $\partial U$};
  \node[font=\small, red!70!black] at (-2.2, -3.2) {le plus proche};
  \draw[->, red!70!black, thin] (-2.2,-2.6) -- (-2.5,-1.6);
\end{tikzpicture}
\captionof{figure}{Le rayon de convergence de la série de Taylor est au
moins la distance de $z_0$ au bord de $U$.}
\end{center}

\begin{exemple}[Série de Taylor de $1/(1-z)$]\label{ex:taylor-1/(1-z)}
La fonction $f(z) = \dfrac{1}{1-z}$ est holomorphe sur $\C \setminus \{1\}$.
Autour de $z_0 = 0$, le plus grand disque contenu dans le domaine est
$D(0,1)$. On a
\[
  \frac{1}{1-z} = \sum_{n=0}^{\infty} z^n, \quad \abs{z} < 1.
\]
Le rayon de convergence est $R = 1 = d(0, \{1\}) = d(0, \C\setminus U)$.
\end{exemple}

\begin{exemple}[Série de Taylor de $\Log(1+z)$]\label{ex:taylor-log}
La fonction $f(z) = \Log(1+z)$ est holomorphe sur $\C \setminus (-\infty, -1]$.
Autour de $z_0 = 0$, le rayon de convergence est $R = 1$ (la singularité
la plus proche est $z = -1$). On a $f(0) = 0$, $f'(z) = \dfrac{1}{1+z}$,
$f^{(n)}(z) = \dfrac{(-1)^{n-1}(n-1)!}{(1+z)^n}$ pour $n \geq 1$, donc
$a_n = \dfrac{f^{(n)}(0)}{n!} = \dfrac{(-1)^{n-1}}{n}$. Ainsi
\[
  \Log(1+z) = \sum_{n=1}^{\infty}\frac{(-1)^{n-1}}{n}z^n
  = z - \frac{z^2}{2} + \frac{z^3}{3} - \cdots, \quad \abs{z} < 1.
\]
\end{exemple}

\section{Exercices}

\begin{exercice}[Calculs par la formule de Cauchy]\label{exo:cauchy-calculs}
Calculer les intégrales suivantes à l'aide de la formule intégrale de
Cauchy ou de ses variantes pour les dérivées.
\begin{enumerate}[label=\textbf{\alph*)}]
  \item $\displaystyle\oint_{\abs{z}=3} \frac{\cos z}{z - \pi}\,\dd z$
  \item $\displaystyle\oint_{\abs{z}=1} \frac{z^2 + 1}{z(z-2)}\,\dd z$
  \item $\displaystyle\oint_{\abs{z}=2} \frac{e^{2z}}{(z-1)^3}\,\dd z$
  \item $\displaystyle\oint_{\abs{z}=1} \frac{\sin z}{z^5}\,\dd z$
  \item $\displaystyle\oint_{\abs{z}=4}
        \frac{z}{(z-1)(z-3)}\,\dd z$
\end{enumerate}
\end{exercice}

\begin{exercice}[Théorème de Liouville et applications]\label{exo:liouville}
\begin{enumerate}[label=\textbf{\alph*)}]
  \item Soit $f$ entière telle que $\abs{f(z)} \leq A + B\abs{z}^n$
        pour tout $z \in \C$, où $A, B > 0$ et $n \in \N$. Montrer
        que $f$ est un polynôme de degré au plus $n$.
  \item Soit $f$ entière et injective. Montrer que $f(z) = az + b$
        avec $a \neq 0$.
        \emph{Indication :} étudier le comportement de $f$ au voisinage
        de l'infini.
  \item Soit $f$ entière telle que $\re(f(z)) \leq M$ pour tout
        $z \in \C$. Montrer que $f$ est constante.
        \emph{Indication :} considérer $g(z) = e^{f(z)}$.
\end{enumerate}
\end{exercice}

\begin{exercice}[Principe du maximum]\label{exo:max-principle}
\begin{enumerate}[label=\textbf{\alph*)}]
  \item Soit $f$ holomorphe sur $\D = D(0,1)$, continue sur
        $\overline{\D}$, avec $\abs{f(z)} = 1$ pour tout $z \in \T$.
        Montrer que $f$ est constante ou que $f$ admet au moins un zéro
        dans $\D$.
        \emph{Indication :} appliquer le principe du minimum à $f$ si
        $f$ ne s'annule pas.
  \item Soit $f$ holomorphe sur $\D$ avec $\abs{f(z)} < 1$ pour tout
        $z \in \D$ et $f(0) = 0$. Montrer que $\abs{f(z)} \leq \abs{z}$
        pour tout $z \in \D$ (lemme de Schwarz).
        \emph{Indication :} appliquer le principe du maximum à
        $g(z) = f(z)/z$.
  \item Soit $f$ holomorphe et non constante sur un ouvert connexe $U$.
        Montrer que $\abs{f}$ ne peut atteindre de minimum local strict
        en un point de $U$ où $f$ ne s'annule pas.
\end{enumerate}
\end{exercice}

\begin{exercice}[Propriété de la valeur moyenne]\label{exo:mean-value}
\begin{enumerate}[label=\textbf{\alph*)}]
  \item Soit $f(z) = z^2 + 3z + 1$. Vérifier directement la propriété
        de la valeur moyenne pour $z_0 = 0$ et $R = 1$.
  \item Soit $u(x,y)$ une fonction harmonique sur un ouvert
        $U \subset \R^2$. En utilisant le fait que $u$ est localement
        la partie réelle d'une fonction holomorphe, montrer que $u$
        satisfait aussi la propriété de la valeur moyenne.
\end{enumerate}
\end{exercice}

\begin{exercice}[Morera et limites uniformes]\label{exo:morera-limits}
Soit $(f_n)$ une suite de fonctions holomorphes sur un ouvert $U$
qui converge uniformément sur tout compact de $U$ vers une fonction
$f \colon U \to \C$.
\begin{enumerate}[label=\textbf{\alph*)}]
  \item Montrer que $f$ est continue sur $U$.
  \item En utilisant le théorème de Morera, montrer que $f$ est
        holomorphe sur $U$.
  \item Montrer que $f_n' \to f'$ uniformément sur tout compact de $U$.
\end{enumerate}
\end{exercice}

\begin{exercice}[Série de Taylor et rayon de convergence]\label{exo:taylor-radius}
Déterminer la série de Taylor autour de $z_0 = 0$ et son rayon de
convergence pour chacune des fonctions suivantes.
\begin{enumerate}[label=\textbf{\alph*)}]
  \item $f(z) = \dfrac{1}{1+z^2}$
  \item $f(z) = \dfrac{z}{(1-z)(1-2z)}$
  \item $f(z) = \Log\!\left(\dfrac{1+z}{1-z}\right)$
  \item $f(z) = e^{\sin z}$
  \emph{(calculer les premiers termes)}
\end{enumerate}
\end{exercice}

\begin{exercice}[Inégalités de Cauchy et applications]\label{exo:cauchy-ineq}
\begin{enumerate}[label=\textbf{\alph*)}]
  \item Soit $f$ entière avec $\abs{f(z)} \leq e^{\abs{z}}$ pour tout
        $z \in \C$. Montrer que $\abs{f^{(n)}(0)} \leq \left(\dfrac{e}{n}\right)^n n!$
        pour tout $n \geq 1$, en optimisant le choix de $R$ dans
        l'inégalité de Cauchy.
  \item Soit $f$ holomorphe sur $D(0,R)$ avec $f(0) = 0$ et
        $\abs{f(z)} \leq M$ pour $\abs{z} < R$. Montrer que
        $\abs{f(z)} \leq \dfrac{M\abs{z}}{R}$ pour $\abs{z} < R$.
  \item (\emph{Inégalité de Cauchy généralisée}) Soit $f$ holomorphe
        sur un ouvert contenant $\overline{D}(z_0,R)$ et soit
        $z \in D(z_0,R)$. Montrer que
        \[
          \abs{f^{(n)}(z)} \leq \frac{n!\,M_R}{(R - \abs{z - z_0})^n},
        \]
        où $M_R = \max_{\abs{\zeta - z_0}=R}\abs{f(\zeta)}$.
\end{enumerate}
\end{exercice}

\begin{exercice}[Formule de Cauchy et intégrales réelles]\label{exo:cauchy-real}
En utilisant la formule intégrale de Cauchy, calculer
\[
  \int_0^{2\pi} \frac{\dd\theta}{a + \cos\theta}, \quad a > 1.
\]
\emph{Indication :} poser $z = e^{i\theta}$ et exprimer $\cos\theta$
en fonction de $z$ et $\conj{z} = 1/z$ sur $\T$.
\end{exercice}


%%%%%%%%%%%%%%%%%%%%%%%%%%%%%%%%%%%%%%%%%%%%%%%%%%%%%%%%%%%%%%%%%%%%
%  CHAPITRE 6 — Séries de Laurent et Singularités Isolées
%%%%%%%%%%%%%%%%%%%%%%%%%%%%%%%%%%%%%%%%%%%%%%%%%%%%%%%%%%%%%%%%%%%%
\chapter{Séries de Laurent et Singularités Isolées}\label{chap:laurent}

\section{Introduction}

Lorsqu'une fonction holomorphe présente des singularités isolées, le
développement en série de Taylor ne suffit plus à la décrire au voisinage
de ces points. La série de Laurent, qui autorise les puissances négatives,
offre un outil adapté. Ce chapitre développe la théorie des séries de
Laurent, classifie les singularités isolées, et aboutit au théorème de
Casorati-Weierstrass.

\section{Séries de Laurent : existence et unicité}

\begin{definition}[Couronne]\label{def:annulus}
Soient $z_0 \in \C$ et $0 \leq r_1 < r_2 \leq +\infty$. La
\emph{couronne} (ou anneau) ouverte de centre $z_0$ et de rayons
$r_1, r_2$ est
\[
  A(z_0; r_1, r_2) = \{z \in \C : r_1 < \abs{z - z_0} < r_2\}.
\]
\end{definition}

\begin{center}
\begin{tikzpicture}[scale=0.9, >=Stealth]
  % Couronne
  \fill[blue!8] (0,0) circle (2.8);
  \fill[white] (0,0) circle (1.2);
  \draw[thick, blue!60!black] (0,0) circle (2.8);
  \draw[thick, blue!60!black] (0,0) circle (1.2);
  % Centre
  \fill (0,0) circle (1.5pt) node[below] {$z_0$};
  % Rayons
  \draw[<->, red!70!black] (0,0) -- node[below, font=\small] {$r_1$} (1.2,0);
  \draw[<->, red!70!black] (0,0) -- node[above, font=\small, pos=0.7] {$r_2$}
    (2.8*0.707, 2.8*0.707);
  % Label
  \node at (2.0, -0.4) {$A(z_0;r_1,r_2)$};
\end{tikzpicture}
\captionof{figure}{Couronne ouverte $A(z_0; r_1, r_2)$.}
\end{center}

\begin{definition}[Série de Laurent]\label{def:laurent-series}
Une \emph{série de Laurent} centrée en $z_0$ est une série de la forme
\[
  \sum_{n=-\infty}^{+\infty} a_n (z - z_0)^n
  = \underbrace{\sum_{n=0}^{\infty} a_n (z-z_0)^n}_{\text{partie régulière}}
  + \underbrace{\sum_{n=1}^{\infty} a_{-n} (z-z_0)^{-n}}_{\text{partie principale}}.
\]
On dit qu'elle converge si les deux sous-séries convergent séparément.
\end{definition}

\begin{theoreme}[Développement de Laurent]\label{thm:laurent}
Soit $f$ holomorphe sur la couronne $A(z_0; r_1, r_2)$ avec
$0 \leq r_1 < r_2 \leq +\infty$. Alors $f$ admet un unique
développement en série de Laurent
\[
  f(z) = \sum_{n=-\infty}^{+\infty} a_n(z - z_0)^n,
  \quad r_1 < \abs{z - z_0} < r_2,
\]
où les coefficients sont donnés par
\[
  a_n = \frac{1}{2\pi i}\oint_{C_\rho}
  \frac{f(\zeta)}{(\zeta - z_0)^{n+1}}\,\dd\zeta,
  \quad n \in \Z,
\]
pour tout $\rho \in (r_1, r_2)$, et $C_\rho$ est le cercle
$\abs{\zeta - z_0} = \rho$ parcouru positivement. La convergence est
normale sur toute couronne fermée $A(z_0; \rho_1, \rho_2)$ avec
$r_1 < \rho_1 \leq \rho_2 < r_2$.
\end{theoreme}

\begin{proof}
Soit $z \in A(z_0; r_1, r_2)$ fixé. Choisissons $\rho_1, \rho_2$
tels que $r_1 < \rho_1 < \abs{z - z_0} < \rho_2 < r_2$. Notons
$C_1 = C_{\rho_1}$ et $C_2 = C_{\rho_2}$ les cercles de centre $z_0$
et de rayons $\rho_1$ et $\rho_2$ respectivement, parcourus dans le
sens positif.

\medskip
\noindent\textbf{Étape 1 : Décomposition par la formule de Cauchy.}
Par le théorème de Cauchy appliqué à la couronne (en joignant les deux
cercles par des fentes), on obtient
\[
  f(z) = \frac{1}{2\pi i}\oint_{C_2}\frac{f(\zeta)}{\zeta - z}\,\dd\zeta
  - \frac{1}{2\pi i}\oint_{C_1}\frac{f(\zeta)}{\zeta - z}\,\dd\zeta.
\]
Le signe moins devant la deuxième intégrale provient du fait que $C_1$
doit être parcouru dans le sens négatif dans le contour bordant la
couronne.

\begin{center}
\begin{tikzpicture}[scale=0.85, >=Stealth]
  % Couronne
  \fill[blue!5] (0,0) circle (2.8);
  \fill[white] (0,0) circle (1.2);
  % Cercle extérieur C_2
  \draw[thick, blue!70!black, postaction={decorate,
    decoration={markings,
      mark=at position 0.12 with {\arrow{>}},
      mark=at position 0.37 with {\arrow{>}},
      mark=at position 0.62 with {\arrow{>}},
      mark=at position 0.87 with {\arrow{>}}
    }}] (0,0) circle (2.8);
  % Cercle intérieur C_1
  \draw[thick, red!70!black, postaction={decorate,
    decoration={markings,
      mark=at position 0.15 with {\arrow{>}},
      mark=at position 0.65 with {\arrow{>}}
    }}] (0,0) circle (1.2);
  % Points
  \fill (0,0) circle (1.5pt) node[below left] {$z_0$};
  \fill (1.8, 0.6) circle (1.5pt) node[right] {$z$};
  % Labels
  \node[blue!70!black] at (2.4, 2.2) {$C_2$};
  \node[red!70!black] at (0.4, 1.1) {$C_1$};
  % Fentes
  \draw[thick, green!50!black, dashed] (1.2, 0) -- (2.8, 0);
  \node[green!50!black, font=\small] at (2.0, -0.3) {fente};
\end{tikzpicture}
\captionof{figure}{Contour pour la preuve du théorème de Laurent :
les deux cercles $C_1$ et $C_2$ reliés par une fente.}
\end{center}

\medskip
\noindent\textbf{Étape 2 : Développement de l'intégrale sur $C_2$.}
Pour $\zeta \in C_2$, on a $\abs{z - z_0} < \rho_2 = \abs{\zeta - z_0}$,
donc comme dans la preuve de la série de Taylor,
\[
  \frac{1}{\zeta - z} = \sum_{n=0}^{\infty}
  \frac{(z-z_0)^n}{(\zeta - z_0)^{n+1}},
\]
avec convergence uniforme en $\zeta \in C_2$. En intégrant terme à terme,
\[
  \frac{1}{2\pi i}\oint_{C_2}\frac{f(\zeta)}{\zeta - z}\,\dd\zeta
  = \sum_{n=0}^{\infty}\left[\frac{1}{2\pi i}\oint_{C_2}
  \frac{f(\zeta)}{(\zeta - z_0)^{n+1}}\,\dd\zeta\right](z-z_0)^n.
\]
C'est la partie régulière de la série de Laurent.

\medskip
\noindent\textbf{Étape 3 : Développement de l'intégrale sur $C_1$.}
Pour $\zeta \in C_1$, on a $\abs{\zeta - z_0} = \rho_1 < \abs{z - z_0}$,
donc cette fois on développe autrement :
\begin{align*}
  \frac{1}{\zeta - z}
  &= \frac{-1}{(z - z_0) - (\zeta - z_0)}
  = \frac{-1}{z - z_0}\cdot\frac{1}{1 - \dfrac{\zeta - z_0}{z - z_0}} \\
  &= -\sum_{m=0}^{\infty}\frac{(\zeta - z_0)^m}{(z - z_0)^{m+1}},
\end{align*}
avec convergence uniforme en $\zeta \in C_1$. En intégrant :
\begin{align*}
  -\frac{1}{2\pi i}\oint_{C_1}\frac{f(\zeta)}{\zeta - z}\,\dd\zeta
  &= \sum_{m=0}^{\infty}\left[\frac{1}{2\pi i}\oint_{C_1}
  f(\zeta)(\zeta - z_0)^m\,\dd\zeta\right]\frac{1}{(z-z_0)^{m+1}}.
\end{align*}
En posant $n = -(m+1)$, c'est-à-dire $m = -n-1$, on obtient
\[
  -\frac{1}{2\pi i}\oint_{C_1}\frac{f(\zeta)}{\zeta - z}\,\dd\zeta
  = \sum_{n=-\infty}^{-1}\left[\frac{1}{2\pi i}\oint_{C_1}
  \frac{f(\zeta)}{(\zeta - z_0)^{n+1}}\,\dd\zeta\right](z-z_0)^n.
\]
C'est la partie principale.

\medskip
\noindent\textbf{Étape 4 : Indépendance du rayon.}
Par le théorème de Cauchy, les intégrales
$\dfrac{1}{2\pi i}\oint_{C_\rho}\dfrac{f(\zeta)}{(\zeta - z_0)^{n+1}}\,\dd\zeta$
ne dépendent pas de $\rho \in (r_1, r_2)$, car l'intégrande est
holomorphe dans la couronne. On peut donc écrire les coefficients
avec un seul rayon $\rho$ quelconque.

\medskip
\noindent\textbf{Unicité.}
Supposons que $f(z) = \sum_{n \in \Z} b_n(z - z_0)^n$ sur la couronne.
Multiplions par $(z - z_0)^{-k-1}$ et intégrons sur $C_\rho$ :
\[
  \frac{1}{2\pi i}\oint_{C_\rho}\frac{f(\zeta)}{(\zeta-z_0)^{k+1}}\,\dd\zeta
  = \sum_{n \in \Z} b_n \cdot \frac{1}{2\pi i}\oint_{C_\rho}
  (\zeta - z_0)^{n-k-1}\,\dd\zeta.
\]
Or $\dfrac{1}{2\pi i}\oint_{C_\rho}(\zeta - z_0)^{m}\,\dd\zeta
= \delta_{m,-1}$ (vaut $1$ si $m = -1$, et $0$ sinon). Donc le membre
de droite vaut $b_k$, ce qui donne $b_k = a_k$ pour tout $k \in \Z$.
\end{proof}

\section{Exemples de développements de Laurent}

\begin{exemple}[Laurent de $\dfrac{1}{z(z-1)}$]\label{ex:laurent-1/z(z-1)}
Développons $f(z) = \dfrac{1}{z(z-1)}$ en série de Laurent autour de
$z_0 = 0$. La fonction a des singularités en $z = 0$ et $z = 1$.
Il y a deux couronnes à considérer.

\medskip
\noindent\textbf{Région $0 < \abs{z} < 1$ :}
On décompose en éléments simples :
\[
  \frac{1}{z(z-1)} = \frac{-1}{z} + \frac{1}{z-1}
  = -\frac{1}{z} - \frac{1}{1-z}.
\]
Pour $\abs{z} < 1$, $\dfrac{1}{1-z} = \sum_{n=0}^\infty z^n$, donc
\[
  f(z) = -\frac{1}{z} - \sum_{n=0}^{\infty} z^n
  = -\frac{1}{z} - 1 - z - z^2 - \cdots
  = \sum_{n=-1}^{\infty} a_n z^n,
\]
avec $a_{-1} = -1$ et $a_n = -1$ pour $n \geq 0$.

\medskip
\noindent\textbf{Région $\abs{z} > 1$ :}
On écrit $\dfrac{1}{z-1} = \dfrac{1}{z}\cdot\dfrac{1}{1-1/z}
= \dfrac{1}{z}\sum_{n=0}^\infty z^{-n} = \sum_{n=0}^\infty z^{-(n+1)}$
pour $\abs{z} > 1$. Donc
\[
  f(z) = -\frac{1}{z} + \sum_{n=0}^{\infty} z^{-(n+1)}
  = -\frac{1}{z} + \frac{1}{z} + \frac{1}{z^2} + \frac{1}{z^3} + \cdots
  = \sum_{n=2}^{\infty} z^{-n}.
\]
\end{exemple}

\begin{exemple}[Laurent de $e^{1/z}$]\label{ex:laurent-exp1/z}
La fonction $f(z) = e^{1/z}$ est holomorphe sur $\C^* = \C\setminus\{0\}$,
c'est-à-dire sur la couronne $A(0; 0, +\infty)$. En remplaçant $z$
par $1/z$ dans le développement de l'exponentielle, on obtient
\[
  e^{1/z} = \sum_{n=0}^{\infty} \frac{1}{n!\,z^n}
  = 1 + \frac{1}{z} + \frac{1}{2!\,z^2} + \frac{1}{3!\,z^3} + \cdots
\]
La partie principale a une infinité de termes non nuls.
\end{exemple}

\begin{exemple}[Laurent de $\dfrac{z^2-2z+5}{(z-2)(z^2+1)}$]\label{ex:laurent-partial}
Développons $f(z) = \dfrac{z^2-2z+5}{(z-2)(z^2+1)}$ en série de Laurent
autour de $z_0 = 2$. La décomposition en éléments simples donne
\[
  f(z) = \frac{1}{z-2} + \frac{1}{z^2+1}.
\]
Le premier terme est déjà sous forme de Laurent en $z_0 = 2$. Pour le
second, on pose $w = z - 2$ et on écrit
\[
  \frac{1}{z^2+1} = \frac{1}{(w+2)^2 + 1} = \frac{1}{w^2+4w+5}
  = \frac{1}{5}\cdot\frac{1}{1 + \frac{4w+w^2}{5}}.
\]
Pour $\abs{w}$ assez petit (en fait $\abs{w} < \sqrt{5}$, la distance
de $2$ aux pôles $\pm i$), on développe en série géométrique :
\[
  \frac{1}{z^2+1} = \frac{1}{5}\sum_{n=0}^\infty
  \left(-\frac{4w+w^2}{5}\right)^n,
\]
ce qui donne un développement en puissances positives de $(z-2)$.
Ainsi, dans la couronne $0 < \abs{z-2} < \sqrt{5}$, la série de Laurent
est
\[
  f(z) = \frac{1}{z-2} + \frac{1}{5} - \frac{4}{25}(z-2) + \cdots
\]
\end{exemple}

\begin{exemple}[Laurent de $\dfrac{\sin z}{z^3}$]\label{ex:laurent-sinz/z3}
Pour $z \neq 0$,
\[
  \frac{\sin z}{z^3} = \frac{1}{z^3}\left(z - \frac{z^3}{3!}
  + \frac{z^5}{5!} - \cdots\right)
  = \frac{1}{z^2} - \frac{1}{6} + \frac{z^2}{120} - \cdots
\]
La partie principale est $z^{-2}$, c'est un pôle d'ordre $2$ en $z = 0$.
\end{exemple}

\section{Classification des singularités isolées}

\begin{definition}[Singularité isolée]\label{def:isolated-sing}
Soit $f$ holomorphe sur un ensemble de la forme $D(z_0, R)\setminus\{z_0\}$
pour un certain $R > 0$. On dit que $z_0$ est une \emph{singularité isolée}
de $f$.
\end{definition}

\begin{definition}[Types de singularités]\label{def:sing-types}
Soit $z_0$ une singularité isolée de $f$ et
$f(z) = \sum_{n=-\infty}^{+\infty} a_n(z-z_0)^n$ son développement
de Laurent dans $D(z_0,R)\setminus\{z_0\}$. On distingue trois cas :
\begin{enumerate}[label=\textbf{(\roman*)}]
  \item \textbf{Singularité éliminable} (ou effaçable) : $a_n = 0$
        pour tout $n < 0$. Dans ce cas, $f$ se prolonge en une fonction
        holomorphe sur $D(z_0,R)$ en posant $f(z_0) = a_0$.
  \item \textbf{Pôle d'ordre $m$} (avec $m \geq 1$) : $a_{-m} \neq 0$
        et $a_n = 0$ pour tout $n < -m$. La partie principale est
        $\sum_{n=1}^{m} a_{-n}(z-z_0)^{-n}$.
  \item \textbf{Singularité essentielle} : $a_n \neq 0$ pour une
        infinité de valeurs de $n < 0$. La partie principale a une
        infinité de termes.
\end{enumerate}
\end{definition}

\begin{remarque}[Terminologie]\label{rem:sing-terminology}
Un pôle d'ordre $1$ est appelé \emph{pôle simple}. Un pôle d'ordre $2$
est un \emph{pôle double}, etc. Le coefficient $a_{-m}$ (pour un pôle
d'ordre $m$) est parfois appelé le \emph{coefficient dominant} de la
partie principale.
\end{remarque}

\section{Critère de Riemann pour les singularités éliminables}

\begin{theoreme}[Critère de Riemann]\label{thm:riemann-removable}
Soit $z_0$ une singularité isolée de $f$. Alors $z_0$ est une singularité
éliminable si et seulement si $f$ est bornée au voisinage de $z_0$,
c'est-à-dire s'il existe $M, \delta > 0$ tels que $\abs{f(z)} \leq M$
pour $0 < \abs{z - z_0} < \delta$.
\end{theoreme}

\begin{proof}
\noindent$(\Rightarrow)$ Si $z_0$ est éliminable, $f$ se prolonge en une
fonction holomorphe sur $D(z_0,R)$, qui est continue donc bornée au
voisinage de $z_0$.

\medskip
\noindent$(\Leftarrow)$ Supposons $\abs{f(z)} \leq M$ pour
$0 < \abs{z-z_0} < \delta$. Montrons que $a_n = 0$ pour tout $n < 0$.

Pour $n < 0$ et $0 < \rho < \delta$,
\[
  \abs{a_n} = \abs{\frac{1}{2\pi i}\oint_{\abs{\zeta - z_0}=\rho}
  \frac{f(\zeta)}{(\zeta - z_0)^{n+1}}\,\dd\zeta}
  \leq \frac{1}{2\pi}\cdot\frac{M}{\rho^{n+1}}\cdot 2\pi\rho
  = \frac{M}{\rho^n} = M\rho^{-n}.
\]
Puisque $n < 0$, on a $-n > 0$, et en faisant $\rho \to 0^+$,
$\abs{a_n} \leq M\rho^{-n} \to 0$. Donc $a_n = 0$ pour tout $n < 0$.
\end{proof}

\begin{corollaire}[Critère par la limite]\label{cor:removable-limit}
La singularité $z_0$ est éliminable si et seulement si
$\lim_{z \to z_0}(z - z_0)f(z) = 0$.
\end{corollaire}

\begin{proof}
Si $z_0$ est éliminable, $f$ est bornée au voisinage de $z_0$, donc
$(z-z_0)f(z) \to 0$. Réciproquement, si $(z-z_0)f(z) \to 0$, la
fonction $g(z) = (z-z_0)f(z)$ (étendue par $g(z_0) = 0$) est holomorphe
en $z_0$ avec $g(z_0) = 0$. Donc $g(z) = (z-z_0)h(z)$ avec $h$
holomorphe en $z_0$, et $f(z) = h(z)$ pour $z \neq z_0$, ce qui montre
que $z_0$ est éliminable.
\end{proof}

\begin{exemple}[Singularité éliminable de $\dfrac{\sin z}{z}$]\label{ex:sinz/z-removable}
La fonction $f(z) = \dfrac{\sin z}{z}$ a une singularité en $z = 0$.
Puisque $\sin z = z - \dfrac{z^3}{6} + \cdots$, on a
$f(z) = 1 - \dfrac{z^2}{6} + \cdots$. Tous les coefficients de Laurent
négatifs sont nuls : la singularité est éliminable. On pose $f(0) = 1$.
\end{exemple}

\section{Ordre d'un pôle et partie principale}

\begin{proposition}[Caractérisation des pôles]\label{prop:pole-char}
Soit $z_0$ une singularité isolée de $f$. Les conditions suivantes
sont équivalentes :
\begin{enumerate}[label=\textbf{(\roman*)}]
  \item $z_0$ est un pôle d'ordre $m$ de $f$.
  \item $\lim_{z \to z_0}\abs{f(z)} = +\infty$ et $(z-z_0)^m f(z)$
        admet une limite finie non nulle en $z_0$.
  \item Il existe une fonction $g$ holomorphe au voisinage de $z_0$
        avec $g(z_0) \neq 0$ telle que
        $f(z) = \dfrac{g(z)}{(z-z_0)^m}$ au voisinage de $z_0$.
  \item La fonction $1/f$ a un zéro d'ordre $m$ en $z_0$.
\end{enumerate}
\end{proposition}

\begin{proof}
\noindent\textbf{(i) $\Rightarrow$ (ii).} Si $z_0$ est un pôle d'ordre $m$,
le développement de Laurent est
$f(z) = \dfrac{a_{-m}}{(z-z_0)^m} + \cdots + \dfrac{a_{-1}}{z-z_0}
+ a_0 + a_1(z-z_0) + \cdots$
avec $a_{-m} \neq 0$. Donc
$(z-z_0)^m f(z) = a_{-m} + a_{-m+1}(z-z_0) + \cdots \to a_{-m} \neq 0$.
Et $\abs{f(z)} \to +\infty$ car le terme dominant est $a_{-m}(z-z_0)^{-m}$.

\medskip
\noindent\textbf{(ii) $\Rightarrow$ (iii).} On pose $g(z) = (z-z_0)^m f(z)$
pour $z \neq z_0$ et $g(z_0) = \lim_{z \to z_0}(z-z_0)^m f(z) \neq 0$.
La fonction $g$ est holomorphe sur $D(z_0,R)\setminus\{z_0\}$ et bornée
au voisinage de $z_0$, donc elle se prolonge holomorphiquement par le
critère de Riemann.

\medskip
\noindent\textbf{(iii) $\Rightarrow$ (iv).} On a $1/f(z) = (z-z_0)^m/g(z)$
au voisinage de $z_0$, et $g(z_0) \neq 0$, donc $1/f$ a un zéro d'ordre
$m$ en $z_0$.

\medskip
\noindent\textbf{(iv) $\Rightarrow$ (i).} Si $1/f$ a un zéro d'ordre $m$
en $z_0$, alors $1/f(z) = (z-z_0)^m h(z)$ avec $h$ holomorphe et
$h(z_0) \neq 0$. Donc $f(z) = 1/[(z-z_0)^m h(z)]$ et $1/h$ est
holomorphe au voisinage de $z_0$ avec $1/h(z_0) \neq 0$. Posons
$g = 1/h$; alors $f(z) = g(z)/(z-z_0)^m$ et le développement de Taylor
de $g$ donne le développement de Laurent de $f$ avec $a_{-m} = g(z_0) \neq 0$.
\end{proof}

\begin{definition}[Partie principale]\label{def:principal-part}
Soit $z_0$ un pôle d'ordre $m$ de $f$. La \emph{partie principale}
de $f$ en $z_0$ est
\[
  P(z) = \sum_{n=1}^{m}\frac{a_{-n}}{(z - z_0)^n}
  = \frac{a_{-m}}{(z-z_0)^m} + \cdots + \frac{a_{-1}}{z-z_0}.
\]
Le coefficient $a_{-1}$ est appelé le \emph{résidu} de $f$ en $z_0$,
noté $\Res(f, z_0)$ ou $\Res_{z_0} f$.
\end{definition}

\begin{exemple}[Pôles de $\dfrac{1}{\sin z}$]\label{ex:poles-1/sinz}
La fonction $f(z) = \dfrac{1}{\sin z}$ a des pôles aux points
$z_k = k\pi$, $k \in \Z$. Puisque $\sin z$ a un zéro simple en
chaque $z_k$ (car $\sin'(k\pi) = \cos(k\pi) = (-1)^k \neq 0$),
$f$ a un pôle simple en chaque $z_k$.

Pour $z_0 = 0$, la partie principale est $\dfrac{1}{z}$, car
$\sin z = z - z^3/6 + \cdots$, donc
$\dfrac{1}{\sin z} = \dfrac{1}{z}\cdot\dfrac{1}{1 - z^2/6 + \cdots}
= \dfrac{1}{z}(1 + z^2/6 + \cdots) = \dfrac{1}{z} + \dfrac{z}{6} + \cdots$.
Le résidu en $0$ est $\Res(f, 0) = 1$.
\end{exemple}

\section{Théorème de Casorati-Weierstrass}

\begin{theoreme}[Casorati-Weierstrass]\label{thm:casorati-weierstrass}
Soit $z_0$ une singularité essentielle de $f$. Alors pour tout
$\varepsilon > 0$, l'image $f(D(z_0,\varepsilon)\setminus\{z_0\})$
est dense dans $\C$. En d'autres termes, pour tout $w \in \C$ et
tout $\delta > 0$, il existe $z$ avec $0 < \abs{z - z_0} < \varepsilon$
tel que $\abs{f(z) - w} < \delta$.
\end{theoreme}

\begin{proof}
Raisonnons par contraposée. Supposons qu'il existe $\varepsilon > 0$
tel que $f(D(z_0,\varepsilon)\setminus\{z_0\})$ n'est pas dense dans
$\C$. Alors il existe $w_0 \in \C$ et $\delta > 0$ tels que
\[
  \abs{f(z) - w_0} \geq \delta \quad \text{pour tout }
  z \in D(z_0,\varepsilon)\setminus\{z_0\}.
\]
La fonction $g(z) = \dfrac{1}{f(z) - w_0}$ est alors holomorphe sur
$D(z_0,\varepsilon)\setminus\{z_0\}$ et satisfait $\abs{g(z)} \leq 1/\delta$.
Par le critère de Riemann (théorème~\ref{thm:riemann-removable}), $z_0$
est une singularité éliminable de $g$.

\medskip
\noindent\textbf{Cas 1 :} $g(z_0) \neq 0$ (après prolongement). Alors
$f(z) = w_0 + 1/g(z)$ est holomorphe en $z_0$, donc $z_0$ est une
singularité éliminable de $f$.

\medskip
\noindent\textbf{Cas 2 :} $g(z_0) = 0$. Alors $g$ a un zéro d'ordre
$m \geq 1$ en $z_0$, et $f(z) = w_0 + 1/g(z)$ a un pôle d'ordre $m$
en $z_0$.

\medskip
Dans les deux cas, $z_0$ n'est pas une singularité essentielle. Par
contraposée, si $z_0$ est une singularité essentielle, alors l'image
est dense.
\end{proof}

\begin{remarque}[Théorème de Picard]\label{rem:picard}
Le théorème de Casorati-Weierstrass est considérablement renforcé par
le \emph{grand théorème de Picard}, qui affirme que $f$ prend en fait
toute valeur de $\C$ à l'exception d'au plus une (une infinité de fois)
dans tout voisinage épointé de $z_0$. Par exemple, $e^{1/z}$ prend
toute valeur sauf $0$ au voisinage de $0$.
\end{remarque}

\begin{exemple}[Comportement de $e^{1/z}$ près de $0$]\label{ex:casorati-exp1/z}
La fonction $f(z) = e^{1/z}$ a une singularité essentielle en $z = 0$
(car la série de Laurent $\sum_{n=0}^\infty \frac{1}{n!\,z^n}$ a une
infinité de termes en puissances négatives).

Par le théorème de Casorati-Weierstrass, l'image de tout voisinage
épointé de $0$ est dense dans $\C$. Vérifions-le directement : soit
$w_0 \in \C^*$. On cherche $z$ proche de $0$ tel que $e^{1/z} = w_0$,
c'est-à-dire $1/z = \Log w_0 + 2k\pi i$ pour $k \in \Z$. Donc
$z = \dfrac{1}{\Log w_0 + 2k\pi i}$, et $z \to 0$ quand
$k \to \pm\infty$. Ainsi $f$ prend la valeur $w_0$ dans tout voisinage
de $0$. La seule valeur non atteinte est $w_0 = 0$ (car $e^w \neq 0$
pour tout $w$).
\end{exemple}

\section{Développements de Laurent : exemples supplémentaires}

\begin{exemple}[Laurent de $\dfrac{1}{z^2+1}$ autour de $z_0 = i$]\label{ex:laurent-1/(z2+1)}
On a $z^2+1 = (z-i)(z+i)$. Autour de $z_0 = i$, posons $w = z - i$ :
\[
  \frac{1}{z^2+1} = \frac{1}{w(w+2i)} = \frac{1}{2iw}\cdot\frac{1}{1+w/(2i)}.
\]
Pour $0 < \abs{w} < 2$ (soit $0 < \abs{z-i} < 2$), on développe :
\[
  \frac{1}{1+w/(2i)} = \sum_{n=0}^{\infty}\left(-\frac{w}{2i}\right)^n
  = \sum_{n=0}^{\infty}\frac{(-1)^n}{(2i)^n}w^n.
\]
Donc
\[
  \frac{1}{z^2+1} = \frac{1}{2i}\sum_{n=0}^{\infty}
  \frac{(-1)^n}{(2i)^n}(z-i)^{n-1}
  = \frac{1}{2i(z-i)} + \sum_{n=0}^{\infty}
  \frac{(-1)^{n+1}}{(2i)^{n+2}}(z-i)^n.
\]
Il y a un pôle simple en $z = i$ avec résidu $\Res(f, i) = \dfrac{1}{2i}$.
\end{exemple}

\begin{exemple}[Laurent de $z\,e^{1/(z-1)}$ autour de $z_0 = 1$]\label{ex:laurent-ze1/(z-1)}
Posons $w = z - 1$, de sorte que $z = w + 1$. Alors
\begin{align*}
  z\,e^{1/(z-1)} &= (w+1)e^{1/w}
  = (w+1)\sum_{n=0}^{\infty}\frac{1}{n!\,w^n} \\
  &= \sum_{n=0}^{\infty}\frac{w}{n!\,w^n}
     + \sum_{n=0}^{\infty}\frac{1}{n!\,w^n} \\
  &= \sum_{n=0}^{\infty}\frac{1}{n!\,w^{n-1}}
     + \sum_{n=0}^{\infty}\frac{1}{n!\,w^n}.
\end{align*}
En posant $m = n - 1$ dans la première somme :
\[
  z\,e^{1/(z-1)} = \sum_{m=-1}^{\infty}\frac{1}{(m+1)!\,w^m}
  + \sum_{n=0}^{\infty}\frac{1}{n!\,w^n}
  = \frac{1}{w^{-1}}\,\text{[...]}
\]
Plus explicitement, regroupons les termes. Pour $m \leq -1$ (c'est-à-dire
les puissances négatives), la première somme contribue
$\dfrac{1}{(m+1)!w^m}$ pour $m \geq -1$, c'est-à-dire un seul terme
$\dfrac{1}{0!\,w^{-1}} = w$ (pour $m = -1$), ce qui n'est pas une
puissance négative. Recalculons plus soigneusement :
\begin{align*}
  (w+1)e^{1/w} &= we^{1/w} + e^{1/w} \\
  &= w\left(1 + \frac{1}{w} + \frac{1}{2w^2} + \frac{1}{6w^3}+\cdots\right)
   + \left(1 + \frac{1}{w} + \frac{1}{2w^2} + \frac{1}{6w^3}+\cdots\right) \\
  &= \left(w + 1 + \frac{1}{2w} + \frac{1}{6w^2} + \cdots\right)
   + \left(1 + \frac{1}{w} + \frac{1}{2w^2} + \frac{1}{6w^3} + \cdots\right) \\
  &= w + 2 + \frac{3}{2w} + \frac{2}{3w^2} + \cdots
\end{align*}
En revenant à $z - 1$ :
\[
  z\,e^{1/(z-1)} = (z-1) + 2 + \frac{3}{2(z-1)} + \frac{2}{3(z-1)^2} + \cdots
\]
La partie principale a une infinité de termes, donc $z_0 = 1$ est une
singularité essentielle. Le résidu est $a_{-1} = \dfrac{3}{2}$.
\end{exemple}

\begin{exemple}[Laurent de $\cos\!\left(\dfrac{1}{z}\right)$]\label{ex:laurent-cos1/z}
Pour $z \neq 0$, en remplaçant $z$ par $1/z$ dans le développement de
$\cos$ :
\[
  \cos\!\left(\frac{1}{z}\right) = \sum_{n=0}^{\infty}
  \frac{(-1)^n}{(2n)!}\cdot\frac{1}{z^{2n}}
  = 1 - \frac{1}{2z^2} + \frac{1}{24z^4} - \cdots
\]
Singularité essentielle en $z = 0$. Le résidu est $a_{-1} = 0$
(seules les puissances paires apparaissent).
\end{exemple}

\section{Singularité à l'infini}

\begin{definition}[Singularité à l'infini]\label{def:sing-infinity}
Soit $f$ holomorphe sur $\{z \in \C : \abs{z} > R\}$ pour un certain
$R \geq 0$. On dit que $f$ a une singularité à l'infini. On étudie
la nature de cette singularité en considérant la fonction
$g(w) = f(1/w)$, qui est holomorphe sur $D(0, 1/R)\setminus\{0\}$.
\begin{enumerate}[label=\textbf{(\roman*)}]
  \item $\infty$ est une \emph{singularité éliminable} de $f$ si $0$
        est une singularité éliminable de $g$.
  \item $\infty$ est un \emph{pôle d'ordre $m$} de $f$ si $0$ est un
        pôle d'ordre $m$ de $g$.
  \item $\infty$ est une \emph{singularité essentielle} de $f$ si $0$
        est une singularité essentielle de $g$.
\end{enumerate}
\end{definition}

\begin{exemple}[Singularités à l'infini]\label{ex:sing-infty}
\begin{enumerate}[label=\textbf{\alph*)}]
  \item $f(z) = \dfrac{1}{z^2+1}$ : on a $g(w) = \dfrac{w^2}{1+w^2}$,
        qui a une singularité éliminable en $w = 0$ (avec $g(0) = 0$).
        Donc $\infty$ est une singularité éliminable de $f$.
  \item $f(z) = z^3 + 2z$ : on a $g(w) = \dfrac{1}{w^3} + \dfrac{2}{w}$,
        qui a un pôle d'ordre $3$ en $w = 0$. Donc $\infty$ est un pôle
        d'ordre $3$ de $f$ (comme tout polynôme de degré $3$).
  \item $f(z) = e^z$ : on a $g(w) = e^{1/w}$, qui a une singularité
        essentielle en $w = 0$. Donc $\infty$ est une singularité
        essentielle de $e^z$.
\end{enumerate}
\end{exemple}

\begin{proposition}[Laurent et singularité à l'infini]\label{prop:laurent-infty}
Soit $f$ holomorphe sur $\{z : \abs{z} > R\}$. Le développement de
Laurent de $f$ en $z_0 = 0$ dans la couronne $\abs{z} > R$ est
\[
  f(z) = \sum_{n=-\infty}^{+\infty} a_n z^n.
\]
Alors :
\begin{itemize}
  \item $\infty$ est éliminable $\iff$ $a_n = 0$ pour tout $n > 0$
        $\iff$ $f$ est bornée pour $\abs{z}$ grand.
  \item $\infty$ est un pôle d'ordre $m$ $\iff$ $a_m \neq 0$ et
        $a_n = 0$ pour $n > m$.
  \item $\infty$ est une singularité essentielle $\iff$ $a_n \neq 0$
        pour une infinité de $n > 0$.
\end{itemize}
\end{proposition}

\begin{proof}
Il suffit de remarquer que le développement de Laurent de
$g(w) = f(1/w)$ en $w = 0$ est
$g(w) = \sum_{n=-\infty}^{+\infty} a_n w^{-n}
= \sum_{k=-\infty}^{+\infty} a_{-k} w^k$.
Les puissances négatives de $w$ (partie principale de $g$) correspondent
aux puissances positives de $z$.
\end{proof}

\section{Résumé : tableau des singularités}

\begin{center}
\renewcommand{\arraystretch}{1.6}
\begin{tabular}{@{}lccc@{}}
\toprule
& \textbf{Éliminable} & \textbf{Pôle d'ordre $m$} & \textbf{Essentielle} \\
\midrule
Partie principale & nulle & finie ($m$ termes) & infinie \\
$\lim_{z \to z_0} f(z)$ & existe (finie) & $\infty$ & n'existe pas \\
$\lim_{z \to z_0}(z-z_0)^m f(z)$ & $0$ ($m \geq 1$)
  & fini $\neq 0$ & n'existe pas \\
Comportement de $\abs{f}$ & borné & $\to +\infty$
  & oscillation \\
Image de $D'(z_0,\varepsilon)$ & $\approx f(z_0)$
  & $\approx \infty$ & dense dans $\C$ \\
\bottomrule
\end{tabular}
\captionof{table}{Résumé comparatif des trois types de singularités isolées.}
\end{center}

\begin{center}
\begin{tikzpicture}[scale=1.1, >=Stealth]
  % --- Singularité éliminable ---
  \begin{scope}[xshift=-4.5cm]
    \draw[->] (-1.3,0) -- (1.3,0);
    \draw[->] (0,-0.5) -- (0,1.8);
    \draw[thick, blue!70!black, domain=-1.1:1.1, samples=80]
      plot (\x, {sin(3*\x r)/(3*\x r+0.001)+0.8});
    \draw[fill=white, draw=blue!70!black] (0,1.1) circle (2pt);
    \fill[blue!70!black] (0,1.1) circle (0pt);
    \node[below, font=\small\bfseries] at (0,-0.7) {Éliminable};
    \node[font=\tiny] at (0,-1.0) {$\frac{\sin z}{z}$};
  \end{scope}
  % --- Pôle ---
  \begin{scope}
    \draw[->] (-1.3,0) -- (1.3,0);
    \draw[->] (0,-0.8) -- (0,1.8);
    \draw[thick, red!70!black, domain=0.15:1.1, samples=40]
      plot (\x, {0.15/(\x*\x)});
    \draw[thick, red!70!black, domain=-1.1:-0.15, samples=40]
      plot (\x, {0.15/(\x*\x)});
    \draw[dashed, gray] (0,0) -- (0,1.7);
    \node[below, font=\small\bfseries] at (0,-1.0) {Pôle};
    \node[font=\tiny] at (0,-1.3) {$\frac{1}{z^2}$};
  \end{scope}
  % --- Essentielle ---
  \begin{scope}[xshift=4.5cm]
    \draw[->] (-1.3,0) -- (1.3,0);
    \draw[->] (0,-1.2) -- (0,1.8);
    \draw[thick, green!50!black, domain=0.12:1.1, samples=100]
      plot (\x, {exp(1/\x)*0.05*sin(20/\x r)});
    \draw[thick, green!50!black, domain=-1.1:-0.12, samples=100]
      plot (\x, {exp(-1/\x)*0.05*sin(20/\x r)});
    \node[below, font=\small\bfseries] at (0,-1.3) {Essentielle};
    \node[font=\tiny] at (0,-1.6) {$e^{1/z}\sin(1/z)$};
  \end{scope}
\end{tikzpicture}
\captionof{figure}{Comportement schématique des trois types de singularités.}
\end{center}

\section{Exercices}

\begin{exercice}[Développements de Laurent]\label{exo:laurent-dev}
Déterminer le développement de Laurent des fonctions suivantes dans
les régions indiquées.
\begin{enumerate}[label=\textbf{\alph*)}]
  \item $f(z) = \dfrac{1}{z(z-1)(z-2)}$ dans les couronnes
        $0 < \abs{z} < 1$, $1 < \abs{z} < 2$, et $\abs{z} > 2$.
  \item $f(z) = \dfrac{e^z}{(z-1)^2}$ dans la couronne
        $0 < \abs{z-1} < +\infty$.
  \item $f(z) = z^2\sin\!\left(\dfrac{1}{z}\right)$ dans $\abs{z} > 0$.
  \item $f(z) = \dfrac{1}{e^z - 1}$ dans $0 < \abs{z} < 2\pi$
        (calculer les premiers termes).
\end{enumerate}
\end{exercice}

\begin{exercice}[Classification des singularités]\label{exo:classify-sing}
Déterminer et classifier les singularités isolées des fonctions suivantes.
\begin{enumerate}[label=\textbf{\alph*)}]
  \item $f(z) = \dfrac{z^3 - 1}{z - 1}$
  \item $f(z) = \dfrac{1}{\cos z - 1}$
  \item $f(z) = \dfrac{e^z - 1}{z^2}$
  \item $f(z) = \dfrac{1}{z^2\sin z}$
  \item $f(z) = e^{1/(z^2+1)}$
  \item $f(z) = \dfrac{z}{\sin^2 z}$
\end{enumerate}
\end{exercice}

\begin{exercice}[Singularités éliminables]\label{exo:removable}
\begin{enumerate}[label=\textbf{\alph*)}]
  \item Soit $f$ holomorphe sur $D(z_0,R)\setminus\{z_0\}$ et continue
        sur $D(z_0,R)$. Montrer que $z_0$ est une singularité éliminable.
  \item Soit $f$ holomorphe sur $D(z_0,R)\setminus\{z_0\}$. Montrer
        que $z_0$ est une singularité éliminable si et seulement si
        $\lim_{z \to z_0}(z - z_0)f(z) = 0$.
  \item Soit $f$ holomorphe et bornée sur le demi-plan
        $\{\re(z) > 0\}\setminus\{1\}$. Montrer que $f$ se prolonge
        holomorphiquement en $z = 1$.
\end{enumerate}
\end{exercice}

\begin{exercice}[Pôles et résidus]\label{exo:poles-residues}
Déterminer l'ordre des pôles et les résidus des fonctions suivantes.
\begin{enumerate}[label=\textbf{\alph*)}]
  \item $f(z) = \dfrac{z}{(z^2+1)^2}$ aux pôles $z = \pm i$.
  \item $f(z) = \dfrac{1}{\sin^2 z}$ au pôle $z = 0$.
  \item $f(z) = \dfrac{e^z}{z^2(z-1)}$ aux pôles $z = 0$ et $z = 1$.
  \item $f(z) = \dfrac{z^3}{(z-1)^2(z+2)}$ aux pôles $z = 1$ et $z = -2$.
\end{enumerate}
\end{exercice}

\begin{exercice}[Casorati-Weierstrass]\label{exo:casorati}
\begin{enumerate}[label=\textbf{\alph*)}]
  \item Montrer directement (sans utiliser le théorème de
        Casorati-Weierstrass) que pour tout $w_0 \in \C$ et tout
        $\varepsilon > 0$, il existe $z$ avec $0 < \abs{z} < \varepsilon$
        tel que $\abs{\sin(1/z) - w_0} < \varepsilon$.
  \item Soit $f$ holomorphe sur $D(0,1)\setminus\{0\}$ avec une
        singularité essentielle en $0$. Montrer qu'il existe une suite
        $(z_n)$ dans $D(0,1)\setminus\{0\}$ telle que $z_n \to 0$
        et $\abs{f(z_n)} \to +\infty$.
  \item Construire un exemple de fonction ayant une singularité
        essentielle en $z = 0$ et telle que $\re(f(z)) > 0$ pour
        $z$ dans un certain voisinage épointé de $0$, ou montrer
        qu'un tel exemple n'existe pas.
\end{enumerate}
\end{exercice}

\begin{exercice}[Singularité à l'infini]\label{exo:sing-infty}
Déterminer la nature de la singularité à l'infini des fonctions suivantes.
\begin{enumerate}[label=\textbf{\alph*)}]
  \item $f(z) = \dfrac{z^4 + z}{z^2 + 1}$
  \item $f(z) = e^{1/z}$
  \item $f(z) = \dfrac{z^2}{z^2+1}$
  \item $f(z) = \tan z$
  \item $f(z) = z\,e^z$
\end{enumerate}
\end{exercice}

\begin{exercice}[Principe de Casorati-Weierstrass et topologie]\label{exo:casorati-topo}
Soit $f$ une fonction entière non polynomiale.
\begin{enumerate}[label=\textbf{\alph*)}]
  \item Montrer que $f$ a une singularité essentielle à l'infini.
  \item En déduire que pour tout $w \in \C$ et tout $R > 0$, il existe
        $z$ avec $\abs{z} > R$ tel que $\abs{f(z) - w} < 1$.
        \emph{Indication :} appliquer Casorati-Weierstrass à $g(w) = f(1/w)$.
  \item En déduire le \emph{petit théorème de Picard} : l'image d'une
        fonction entière non constante est $\C$ ou $\C$ privé d'un point.
        \emph{(Cet exercice ne demande pas une preuve complète de Picard,
        mais une discussion de ce que Casorati-Weierstrass implique.)}
\end{enumerate}
\end{exercice}

\begin{exercice}[Unicité du développement de Laurent]\label{exo:laurent-uniqueness}
\begin{enumerate}[label=\textbf{\alph*)}]
  \item Soit $f(z) = \displaystyle\sum_{n=-\infty}^{+\infty} a_n z^n$
        convergente dans une couronne $A(0; r_1, r_2)$. Montrer que
        $a_n = \dfrac{1}{2\pi}\displaystyle\int_0^{2\pi}
        f(\rho e^{i\theta})e^{-in\theta}\,\dd\theta \cdot \rho^{-n}$
        pour tout $\rho \in (r_1, r_2)$.
  \item En déduire que si $f(z) = \displaystyle\sum_{n=-\infty}^{+\infty}
        a_n z^n = \displaystyle\sum_{n=-\infty}^{+\infty} b_n z^n$ dans
        la même couronne, alors $a_n = b_n$ pour tout $n$.
  \item Utiliser l'unicité pour trouver le développement de Laurent de
        $f(z) = \dfrac{1}{z-1} - \dfrac{1}{z-2}$ dans $1 < \abs{z} < 2$
        sans calculer les intégrales de Cauchy directement.
\end{enumerate}
\end{exercice}

\begin{exercice}[Développements de Laurent par manipulation]\label{exo:laurent-manip}
Sans calculer les intégrales de Cauchy, déterminer les développements de
Laurent indiqués en utilisant des développements connus et des
manipulations algébriques.
\begin{enumerate}[label=\textbf{\alph*)}]
  \item $\dfrac{z+1}{z-1}$ dans $\abs{z} > 1$.
        \emph{Indication :} écrire $\dfrac{z+1}{z-1} = 1 + \dfrac{2}{z-1}$.
  \item $\dfrac{1}{z^2-3z+2}$ dans $1 < \abs{z} < 2$.
        \emph{Indication :} décomposer en éléments simples.
  \item $\Log\!\left(\dfrac{z-1}{z-2}\right)$ dans $\abs{z} > 2$.
        \emph{Indication :} écrire $\Log(1-1/z) - \Log(1-2/z)$ et
        utiliser le développement de $\Log(1-w)$.
  \item $\dfrac{z}{(z-1)(z-2)(z-3)}$ dans $2 < \abs{z} < 3$.
\end{enumerate}
\end{exercice}
%%%%%%%%%%%%%%%%%%%%%%%%%%%%%%%%%%%%%%%%%%%%%%%%%%%%%%%%%%%%%%%%%%%%
%  Chapitre 7 : Théorème des Résidus et Applications
%%%%%%%%%%%%%%%%%%%%%%%%%%%%%%%%%%%%%%%%%%%%%%%%%%%%%%%%%%%%%%%%%%%%
\chapter{Théorème des Résidus et Applications}\label{chap:residus}

\minitoc
\vspace{1em}

Le théorème des résidus constitue l'un des résultats les plus puissants de l'analyse
complexe. Il ramène le calcul d'intégrales curvilignes à un problème purement
algébrique~: la détermination de coefficients dans un développement en série de Laurent.
Ce chapitre développe la théorie des résidus et démontre sa remarquable efficacité
pour le calcul d'intégrales réelles qui résistent aux méthodes classiques.

%% ═══════════════════════════════════════════════════════════════════
\section{Résidus~: définition et calcul}\label{sec:residus-def}
%% ═══════════════════════════════════════════════════════════════════

\subsection{Définition du résidu}

Soit $f$ une fonction holomorphe sur un ouvert $U \subset \C$ privé d'un point
isolé $z_0$. Au voisinage de $z_0$, la fonction $f$ admet un développement en
série de Laurent~:
\[
  f(z) = \sum_{n=-\infty}^{+\infty} a_n (z - z_0)^n, \qquad 0 < \abs{z - z_0} < R.
\]

\begin{definition}[Résidu en une singularité isolée]\label{def:residu}
Soit $f$ holomorphe sur un voisinage épointé de $z_0$. Le \textbf{résidu} de $f$
en $z_0$ est le coefficient $a_{-1}$ du développement de Laurent de $f$ au
voisinage de $z_0$~:
\[
  \Res(f, z_0) = a_{-1} = \frac{1}{2\pi i} \oint_{\gamma} f(z)\, \dd z,
\]
où $\gamma$ est un cercle de centre $z_0$, de rayon suffisamment petit, parcouru
dans le sens direct.
\end{definition}

\begin{remarque}\label{rem:residu-intrinseque}
Le résidu est une notion intrinsèque~: il ne dépend pas du choix du cercle $\gamma$,
pourvu que celui-ci soit contenu dans le domaine d'holomorphie et n'entoure
que la singularité $z_0$. Plus généralement, $\gamma$ peut être remplacé par
tout lacet simple entourant $z_0$ une fois dans le sens positif.
\end{remarque}

\subsection{Formules de calcul des résidus}

\begin{proposition}[Résidu en un pôle simple]\label{prop:residu-pole-simple}
Si $f$ possède un pôle simple en $z_0$, alors~:
\[
  \Res(f, z_0) = \lim_{z \to z_0} (z - z_0) f(z).
\]
\end{proposition}

\begin{proof}
Au voisinage de $z_0$, on a $f(z) = \frac{a_{-1}}{z - z_0} + a_0 + a_1(z-z_0) + \cdots$.
En multipliant par $(z - z_0)$~:
\[
  (z - z_0) f(z) = a_{-1} + a_0 (z-z_0) + a_1 (z-z_0)^2 + \cdots
\]
et le passage à la limite $z \to z_0$ donne $a_{-1} = \Res(f, z_0)$.
\end{proof}

\begin{proposition}[Résidu par quotient]\label{prop:residu-quotient}
Si $f(z) = \frac{g(z)}{h(z)}$ où $g$ et $h$ sont holomorphes au voisinage de $z_0$,
avec $g(z_0) \neq 0$, $h(z_0) = 0$ et $h'(z_0) \neq 0$, alors~:
\[
  \Res(f, z_0) = \frac{g(z_0)}{h'(z_0)}.
\]
\end{proposition}

\begin{proof}
On écrit~:
\[
  \Res(f, z_0) = \lim_{z \to z_0} (z - z_0) \frac{g(z)}{h(z)}
  = \lim_{z \to z_0} g(z) \cdot \frac{z - z_0}{h(z) - h(z_0)}
  = \frac{g(z_0)}{h'(z_0)}. \qedhere
\]
\end{proof}

\begin{proposition}[Résidu en un pôle d'ordre $n$]\label{prop:residu-pole-ordre-n}
Si $f$ possède un pôle d'ordre $n \geq 1$ en $z_0$, alors~:
\[
  \Res(f, z_0) = \frac{1}{(n-1)!} \lim_{z \to z_0}
  \frac{\dd^{n-1}}{\dd z^{n-1}} \bigl[ (z-z_0)^n f(z) \bigr].
\]
\end{proposition}

\begin{proof}
La fonction $\varphi(z) = (z-z_0)^n f(z)$ est holomorphe au voisinage de $z_0$
avec $\varphi(z_0) \neq 0$. Son développement de Taylor est~:
\[
  \varphi(z) = \sum_{k=0}^{\infty} \frac{\varphi^{(k)}(z_0)}{k!} (z - z_0)^k.
\]
Or $f(z) = \frac{\varphi(z)}{(z-z_0)^n}$, donc~:
\[
  f(z) = \sum_{k=0}^{\infty} \frac{\varphi^{(k)}(z_0)}{k!} (z - z_0)^{k-n}.
\]
Le coefficient de $(z-z_0)^{-1}$ correspond à $k = n-1$~:
\[
  \Res(f, z_0) = \frac{\varphi^{(n-1)}(z_0)}{(n-1)!}. \qedhere
\]
\end{proof}

\begin{remarque}\label{rem:pole-double}
Pour un pôle double ($n=2$), la formule se réduit à~:
\[
  \Res(f, z_0) = \lim_{z \to z_0} \frac{\dd}{\dd z} \bigl[ (z-z_0)^2 f(z) \bigr].
\]
\end{remarque}

\begin{exemple}[Calculs de résidus élémentaires]\label{ex:residus-elementaires}
\begin{enumerate}[label=\textbf{\arabic*.}]
  \item $f(z) = \frac{e^z}{z^2 + 1}$. Les pôles sont $z = \pm i$, simples.
  \[
    \Res(f, i) = \frac{e^i}{2i}, \qquad
    \Res(f, -i) = \frac{e^{-i}}{-2i}.
  \]

  \item $f(z) = \frac{1}{z^2(z-1)}$. Pôle double en $0$, pôle simple en $1$.
  \[
    \Res(f, 1) = \lim_{z \to 1} \frac{1}{z^2} = 1, \qquad
    \Res(f, 0) = \lim_{z \to 0} \frac{\dd}{\dd z}\!\left[\frac{1}{z-1}\right]
    = \frac{-1}{1} = -1.
  \]

  \item $f(z) = \frac{\cos z}{z^3}$. On développe~:
  \[
    f(z) = \frac{1}{z^3} - \frac{1}{2z} + \frac{z}{24} - \cdots,
    \quad \text{donc } \Res(f, 0) = -\frac{1}{2}.
  \]
\end{enumerate}
\end{exemple}

%% ═══════════════════════════════════════════════════════════════════
\section{Le théorème des résidus}\label{sec:thm-residus}
%% ═══════════════════════════════════════════════════════════════════

\begin{theoreme}[Théorème des résidus]\label{thm:residus}
Soit $U \subset \C$ un ouvert, $\{z_1, \ldots, z_N\}$ un ensemble fini de points
de $U$, et $f \colon U \setminus \{z_1, \ldots, z_N\} \to \C$ holomorphe.
Si $\gamma$ est un cycle dans $U \setminus \{z_1, \ldots, z_N\}$
homologue à zéro dans $U$, alors~:
\[
  \frac{1}{2\pi i} \oint_{\gamma} f(z)\, \dd z
  = \sum_{k=1}^{N} \ind(\gamma, z_k)\, \Res(f, z_k).
\]
\end{theoreme}

\begin{proof}
Pour chaque $k$, soit $P_k(z) = \sum_{n=1}^{\infty} a_{-n}^{(k)} (z - z_k)^{-n}$
la partie principale du développement de Laurent de $f$ en $z_k$.
Posons $g(z) = f(z) - \sum_{k=1}^{N} P_k(z)$. En chaque $z_k$, la singularité
de $f$ est exactement compensée par $P_k$, et les autres $P_j$ ($j \neq k$) sont
holomorphes en $z_k$. Donc $g$ se prolonge holomorphiquement à $U$ entier.

Puisque $\gamma$ est homologue à zéro dans $U$, le théorème de Cauchy donne
$\oint_{\gamma} g(z)\, \dd z = 0$. Donc~:
\[
  \oint_{\gamma} f(z)\, \dd z = \sum_{k=1}^{N} \oint_{\gamma} P_k(z)\, \dd z.
\]

Par convergence uniforme sur $\gamma$~:
\[
  \oint_{\gamma} P_k(z)\, \dd z
  = \sum_{n=1}^{\infty} a_{-n}^{(k)} \oint_{\gamma} (z - z_k)^{-n}\, \dd z.
\]
Pour $n \geq 2$, $(z - z_k)^{-n}$ admet une primitive sur $\C \setminus \{z_k\}$,
donc l'intégrale est nulle. Pour $n = 1$~:
\[
  \oint_{\gamma} \frac{\dd z}{z - z_k} = 2\pi i\, \ind(\gamma, z_k).
\]
Ainsi $\oint_{\gamma} P_k(z)\, \dd z = 2\pi i\, \Res(f, z_k)\, \ind(\gamma, z_k)$,
et en sommant on obtient le résultat.
\end{proof}

\begin{corollaire}[Version pour les lacets simples]\label{cor:residus-lacet-simple}
Si $\gamma$ est un lacet simple parcouru dans le sens direct entourant exactement
les singularités $z_1, \ldots, z_m$, alors~:
\[
  \oint_{\gamma} f(z)\, \dd z = 2\pi i \sum_{k=1}^{m} \Res(f, z_k).
\]
\end{corollaire}

\begin{exemple}\label{ex:application-residus-1}
Calculons $I = \oint_{\abs{z}=2} \frac{e^z}{z(z-1)^2}\, \dd z$.

Les singularités $z = 0$ (pôle simple) et $z = 1$ (pôle double) sont dans $\abs{z}<2$.

\textbf{Résidu en $z = 0$}~:
$\Res(f, 0) = \lim_{z \to 0} \frac{e^z}{(z-1)^2} = 1$.

\textbf{Résidu en $z = 1$}~:
$\Res(f, 1) = \lim_{z \to 1} \frac{\dd}{\dd z}\!\left[\frac{e^z}{z}\right]
= \lim_{z \to 1} \frac{ze^z - e^z}{z^2} = 0$.

Donc $I = 2\pi i$.
\end{exemple}

%% ═══════════════════════════════════════════════════════════════════
\section{Applications au calcul d'intégrales réelles}\label{sec:integrales-reelles}
%% ═══════════════════════════════════════════════════════════════════

La puissance du théorème des résidus réside dans sa capacité à évaluer des intégrales
réelles en les reliant à des intégrales complexes sur des contours fermés
judicieusement choisis.

\subsection{Intégrales de fractions rationnelles sur $\R$}\label{subsec:int-rationnelles}

\begin{theoreme}[Intégrales rationnelles sur $\R$]\label{thm:int-rationnelle}
Soit $R(x) = \frac{P(x)}{Q(x)}$ une fraction rationnelle telle que~:
\begin{enumerate}[label=(\roman*)]
  \item $Q$ n'a pas de zéro réel~;
  \item $\deg Q \geq \deg P + 2$.
\end{enumerate}
Alors~:
\[
  \int_{-\infty}^{+\infty} R(x)\, \dd x
  = 2\pi i \sum_{\substack{z_k \text{ pôle de } R \\ \im(z_k) > 0}} \Res(R, z_k).
\]
\end{theoreme}

\begin{proof}
On considère le contour $\Gamma_r$ formé du segment $[-r, r]$ et du demi-cercle
$C_r = \{re^{i\theta} : \theta \in [0, \pi]\}$ parcouru dans le sens direct.

\begin{center}
\begin{tikzpicture}[scale=1.2, >=stealth]
  % Axes
  \draw[->] (-3.5,0) -- (3.5,0) node[right] {$\re$};
  \draw[->] (0,-0.5) -- (0,3) node[above] {$\im$};
  % Contour
  \draw[thick, blue, ->-/.style={decoration={markings, mark=at position #1 with {\arrow{>}}}, postaction={decorate}},
    ->-=0.3, ->-=0.7]
    (-2.8,0) -- (2.8,0);
  \draw[thick, blue, ->-/.style={decoration={markings, mark=at position #1 with {\arrow{>}}}, postaction={decorate}},
    ->-=0.5]
    (2.8,0) arc[start angle=0, end angle=180, radius=2.8];
  % Labels
  \node[below] at (-2.8,0) {$-r$};
  \node[below] at (2.8,0) {$r$};
  \node[blue, above right] at (2.2,1.8) {$C_r$};
  % Poles
  \filldraw[red] (0.8,1.2) circle (1.5pt) node[right] {$z_k$};
  \filldraw[red] (-1,0.8) circle (1.5pt) node[left] {$z_j$};
  \node at (1.5,-0.7) {$\Gamma_r$};
\end{tikzpicture}
\end{center}

Pour $r$ assez grand, tous les pôles de $R$ dans le demi-plan supérieur sont
à l'intérieur de $\Gamma_r$. Par le théorème des résidus~:
\[
  \int_{-r}^{r} R(x)\, \dd x + \int_{C_r} R(z)\, \dd z
  = 2\pi i \sum_{\substack{z_k \\ \im(z_k) > 0}} \Res(R, z_k).
\]

Il suffit de montrer que $\int_{C_r} R(z)\, \dd z \to 0$ quand $r \to +\infty$.
Sur $C_r$, $\abs{z} = r$ et la condition $\deg Q \geq \deg P + 2$ entraîne
$\abs{R(z)} = O(1/r^2)$. Donc~:
\[
  \left|\int_{C_r} R(z)\, \dd z\right|
  \leq \pi r \cdot \sup_{z \in C_r} \abs{R(z)}
  = O\!\left(\frac{1}{r}\right) \xrightarrow[r \to \infty]{} 0. \qedhere
\]
\end{proof}

\begin{exemple}[Intégrale de $1/(1+x^2)^2$]\label{ex:int-1-plus-x2-carre}
Calculons $\displaystyle I = \int_{-\infty}^{+\infty} \frac{\dd x}{(1+x^2)^2}$.

La fraction $R(z) = \frac{1}{(1+z^2)^2} = \frac{1}{(z-i)^2(z+i)^2}$
a un pôle double en $z = i$ (dans le demi-plan supérieur) et en $z = -i$.

\textbf{Résidu en $z = i$}~:
\[
  \Res(R, i) = \lim_{z \to i} \frac{\dd}{\dd z}\!\left[\frac{1}{(z+i)^2}\right]
  = \lim_{z \to i} \frac{-2}{(z+i)^3} = \frac{-2}{(2i)^3} = \frac{-2}{-8i} = \frac{1}{4i}.
\]

Donc~:
\[
  I = 2\pi i \cdot \frac{1}{4i} = \frac{\pi}{2}.
\]
\end{exemple}

\begin{exemple}[Intégrale de $1/(1+x^4)$]\label{ex:int-1-plus-x4}
Calculons $\displaystyle I = \int_{-\infty}^{+\infty} \frac{\dd x}{1 + x^4}$.

Les pôles de $\frac{1}{1+z^4}$ sont les racines quatrièmes de $-1$~:
$z_k = e^{i(2k+1)\pi/4}$ pour $k = 0, 1, 2, 3$. Ceux dans le demi-plan
supérieur sont~:
\[
  z_0 = e^{i\pi/4} = \frac{1+i}{\sqrt{2}}, \qquad
  z_1 = e^{3i\pi/4} = \frac{-1+i}{\sqrt{2}}.
\]

Tous sont des pôles simples. Avec $h(z) = 1 + z^4$, $h'(z) = 4z^3$~:
\[
  \Res(R, z_k) = \frac{1}{4z_k^3} = \frac{1}{4} \cdot z_k^{-3} = \frac{z_k}{4z_k^4} = \frac{z_k}{4(-1)} = -\frac{z_k}{4}.
\]

Donc~:
\[
  I = 2\pi i \left(-\frac{z_0}{4} - \frac{z_1}{4}\right)
  = -\frac{\pi i}{2}(z_0 + z_1)
  = -\frac{\pi i}{2} \cdot \frac{(1+i)+(-1+i)}{\sqrt{2}}
  = -\frac{\pi i}{2} \cdot \frac{2i}{\sqrt{2}}
  = -\frac{\pi i^2 \sqrt{2}}{1} \cdot\frac{1}{1}
  = \frac{\pi}{\sqrt{2}} = \frac{\pi\sqrt{2}}{2}.
\]
\end{exemple}

\subsection{Intégrales trigonométriques}\label{subsec:int-trig}

\begin{theoreme}[Intégrales trigonométriques]\label{thm:int-trig}
Soit $R(u,v)$ une fraction rationnelle en deux variables, sans pôle sur le cercle
unité (paramétré par $(\cos\theta, \sin\theta)$). Alors~:
\[
  \int_0^{2\pi} R(\cos\theta, \sin\theta)\, \dd\theta
  = \oint_{\abs{z}=1} R\!\left(\frac{z+z^{-1}}{2},\, \frac{z-z^{-1}}{2i}\right)
  \frac{\dd z}{iz}.
\]
\end{theoreme}

\begin{proof}
On paramètre le cercle unité par $z = e^{i\theta}$, $\theta \in [0, 2\pi]$.
Alors $\dd z = ie^{i\theta}\,\dd\theta = iz\,\dd\theta$, soit
$\dd\theta = \frac{\dd z}{iz}$. De plus~:
\[
  \cos\theta = \frac{e^{i\theta}+e^{-i\theta}}{2} = \frac{z+z^{-1}}{2}, \qquad
  \sin\theta = \frac{e^{i\theta}-e^{-i\theta}}{2i} = \frac{z-z^{-1}}{2i}.
\]
La substitution transforme l'intégrale en une intégrale curviligne sur $\T$,
évaluable par le théorème des résidus.
\end{proof}

\begin{exemple}[Intégrale trigonométrique classique]\label{ex:int-trig-1}
Calculons $\displaystyle I = \int_0^{2\pi} \frac{\dd\theta}{2 + \cos\theta}$.

Avec $z = e^{i\theta}$, $\cos\theta = \frac{z + z^{-1}}{2}$~:
\[
  I = \oint_{\T} \frac{1}{2 + \frac{z+z^{-1}}{2}} \cdot \frac{\dd z}{iz}
  = \oint_{\T} \frac{2\, \dd z}{iz(4 + z + z^{-1})}
  = \oint_{\T} \frac{2\, \dd z}{i(z^2 + 4z + 1)}.
\]

Les racines de $z^2 + 4z + 1 = 0$ sont $z = -2 \pm \sqrt{3}$.
Seule $z_0 = -2 + \sqrt{3}$ vérifie $\abs{z_0} < 1$ (car $\abs{-2+\sqrt{3}} \approx 0{,}27$).

C'est un pôle simple~:
\[
  \Res\!\left(\frac{2}{i(z^2+4z+1)},\, z_0\right)
  = \frac{2}{i \cdot 2z_0 + 4i} = \frac{2}{i(2(-2+\sqrt{3})+4)} = \frac{2}{2i\sqrt{3}} = \frac{1}{i\sqrt{3}}.
\]

Donc~:
\[
  I = 2\pi i \cdot \frac{1}{i\sqrt{3}} = \frac{2\pi}{\sqrt{3}} = \frac{2\pi\sqrt{3}}{3}.
\]
\end{exemple}

\begin{exemple}[Intégrale avec $\cos^2$]\label{ex:int-trig-2}
Calculons $\displaystyle I = \int_0^{2\pi} \frac{\dd\theta}{a + b\cos^2\theta}$,
avec $a > 0$, $b > 0$.

On utilise $\cos^2\theta = \frac{1+\cos 2\theta}{2}$, mais il est plus direct
d'employer la substitution $z = e^{i\theta}$ avec
$\cos^2\theta = \left(\frac{z+z^{-1}}{2}\right)^2 = \frac{z^2 + 2 + z^{-2}}{4}$~:
\[
  I = \oint_{\T} \frac{1}{a + b\,\frac{z^2+2+z^{-2}}{4}} \cdot \frac{\dd z}{iz}
  = \oint_{\T} \frac{4\,\dd z}{iz(4a+2b+b(z^2+z^{-2}))}
  = \oint_{\T} \frac{4\,\dd z}{i(bz^4 + (4a+2b)z^2 + b) / z}.
\]

En posant $w = z^2$, l'intégrale se simplifie. Le dénominateur en $z$ est
$bz^4 + (4a+2b)z^2 + b$, dont les racines en $w = z^2$ sont~:
\[
  w = \frac{-(4a+2b) \pm \sqrt{(4a+2b)^2 - 4b^2}}{2b}
  = \frac{-(2a+b) \pm 2\sqrt{a^2+ab}}{b}.
\]
Pour $a = b = 1$~: $w = \frac{-3 \pm 2\sqrt{2}}{1}$, d'où $w_1 = -3 + 2\sqrt{2} \approx -0{,}17$,
$w_2 = -3 - 2\sqrt{2} \approx -5{,}83$. Les pôles en $z$ avec $\abs{z} < 1$
donnent, après calcul des résidus~:
\[
  I = \frac{2\pi}{\sqrt{a(a+b)}}.
\]
\end{exemple}

\subsection{Lemme de Jordan et intégrales oscillantes}\label{subsec:jordan}

Pour traiter des intégrales de la forme $\int_{-\infty}^{+\infty} R(x) e^{iax}\, \dd x$
avec $a > 0$, la condition $\deg Q \geq \deg P + 2$ peut être affaiblie grâce au
lemme de Jordan.

\begin{lemme}[Lemme de Jordan]\label{lem:jordan}
Soit $C_r = \{re^{i\theta} : \theta \in [0,\pi]\}$ le demi-cercle supérieur de
rayon $r$, et soit $a > 0$. Si $M_r = \sup_{\theta \in [0,\pi]} \abs{g(re^{i\theta})}$
et $M_r \to 0$ quand $r \to +\infty$, alors~:
\[
  \lim_{r \to +\infty} \int_{C_r} g(z) e^{iaz}\, \dd z = 0.
\]
\end{lemme}

\begin{proof}
On paramètre $z = re^{i\theta}$, $\dd z = ire^{i\theta}\,\dd\theta$~:
\[
  \left|\int_{C_r} g(z) e^{iaz}\, \dd z\right|
  \leq \int_0^{\pi} \abs{g(re^{i\theta})} \cdot e^{-ar\sin\theta} \cdot r\, \dd\theta
  \leq M_r \cdot r \int_0^{\pi} e^{-ar\sin\theta}\, \dd\theta.
\]

L'inégalité de Jordan $\sin\theta \geq \frac{2\theta}{\pi}$ pour
$\theta \in [0, \pi/2]$ donne~:
\[
  r \int_0^{\pi} e^{-ar\sin\theta}\, \dd\theta
  = 2r \int_0^{\pi/2} e^{-ar\sin\theta}\, \dd\theta
  \leq 2r \int_0^{\pi/2} e^{-2ar\theta/\pi}\, \dd\theta
  = 2r \cdot \frac{\pi}{2ar}\left(1 - e^{-ar}\right)
  = \frac{\pi}{a}\left(1 - e^{-ar}\right) \leq \frac{\pi}{a}.
\]

Donc $\left|\int_{C_r} g(z) e^{iaz}\, \dd z\right| \leq \frac{\pi M_r}{a} \to 0$.
\end{proof}

\begin{center}
\begin{tikzpicture}[scale=1.1, >=stealth]
  \draw[->] (-3.5,0) -- (3.5,0) node[right] {$\re$};
  \draw[->] (0,-0.5) -- (0,3.2) node[above] {$\im$};
  \draw[thick, blue, decoration={markings, mark=at position 0.25 with {\arrow{>}},
    mark=at position 0.75 with {\arrow{>}}}, postaction={decorate}]
    (-2.8,0) -- (2.8,0);
  \draw[thick, blue, decoration={markings, mark=at position 0.5 with {\arrow{>}}},
    postaction={decorate}]
    (2.8,0) arc[start angle=0, end angle=180, radius=2.8];
  \node[below] at (-2.8,0) {$-R$};
  \node[below] at (2.8,0) {$R$};
  \node[blue] at (2.4,2.2) {$C_R$};
  % Shading to show exponential decay
  \fill[blue!8] (2.8,0) arc[start angle=0, end angle=180, radius=2.8] -- cycle;
  \node at (0,1.2) {\small $e^{iaz}$ décroît};
  \node at (0,0.7) {\small car $\im(az) = a r\sin\theta > 0$};
\end{tikzpicture}
\end{center}

\begin{theoreme}[Intégrales de type Fourier]\label{thm:int-fourier}
Soit $R(x) = \frac{P(x)}{Q(x)}$ une fraction rationnelle avec $\deg Q \geq \deg P + 1$
et $Q$ sans zéro réel, et soit $a > 0$. Alors~:
\[
  \int_{-\infty}^{+\infty} R(x) e^{iax}\, \dd x
  = 2\pi i \sum_{\substack{z_k \text{ pôle de } R \\ \im(z_k) > 0}} \Res(R(z) e^{iaz}, z_k).
\]
\end{theoreme}

\begin{proof}
On utilise le même contour en demi-cercle. Le lemme de Jordan avec
$g(z) = R(z)$ assure que $M_r \to 0$ dès que $\deg Q \geq \deg P + 1$,
et l'intégrale sur $C_r$ tend vers zéro.
\end{proof}

\begin{exemple}[Intégrale de Fourier classique]\label{ex:fourier-1}
Calculons $\displaystyle I = \int_{-\infty}^{+\infty} \frac{\cos x}{x^2 + 1}\, \dd x$.

On écrit $I = \re\!\left(\int_{-\infty}^{+\infty} \frac{e^{ix}}{x^2+1}\, \dd x\right)$.
La fonction $\frac{e^{iz}}{z^2+1}$ a un pôle simple en $z = i$ dans le demi-plan supérieur~:
\[
  \Res\!\left(\frac{e^{iz}}{z^2+1}, i\right) = \frac{e^{i \cdot i}}{2i} = \frac{e^{-1}}{2i}.
\]

Donc $\int_{-\infty}^{+\infty} \frac{e^{ix}}{x^2+1}\, \dd x = 2\pi i \cdot \frac{e^{-1}}{2i} = \frac{\pi}{e}$, et~:
\[
  I = \re\!\left(\frac{\pi}{e}\right) = \frac{\pi}{e}.
\]
\end{exemple}

\begin{exemple}[Intégrale de $x\sin x/(x^2+a^2)$]\label{ex:fourier-2}
Calculons $\displaystyle I = \int_{-\infty}^{+\infty} \frac{x\sin x}{x^2 + a^2}\, \dd x$
pour $a > 0$.

On a $I = \im\!\left(\int_{-\infty}^{+\infty} \frac{xe^{ix}}{x^2+a^2}\,\dd x\right)$.
Le pôle dans le demi-plan supérieur est $z = ia$~:
\[
  \Res\!\left(\frac{ze^{iz}}{z^2+a^2}, ia\right)
  = \frac{ia \cdot e^{-a}}{2ia} = \frac{e^{-a}}{2}.
\]

Donc $\int_{-\infty}^{+\infty} \frac{xe^{ix}}{x^2+a^2}\,\dd x = 2\pi i \cdot \frac{e^{-a}}{2} = \pi i e^{-a}$, et~:
\[
  I = \im(\pi i e^{-a}) = \pi e^{-a}.
\]
\end{exemple}

\subsection{Intégrales avec coupure logarithmique}\label{subsec:int-log}

Pour calculer des intégrales du type $\int_0^{+\infty} \frac{x^{\alpha-1}}{Q(x)}\,\dd x$
avec $0 < \alpha < 1$ et $Q$ un polynôme sans zéro sur $\R_+$, on utilise un
contour en \og trou de serrure \fg{} (ou \emph{keyhole contour}).

\begin{center}
\begin{tikzpicture}[scale=1.0, >=stealth]
  \draw[->] (-3.5,0) -- (3.8,0) node[right] {$\re$};
  \draw[->] (0,-3.2) -- (0,3.2) node[above] {$\im$};

  % Grand cercle
  \draw[thick, blue, decoration={markings,
    mark=at position 0.15 with {\arrow{>}},
    mark=at position 0.5 with {\arrow{>}},
    mark=at position 0.85 with {\arrow{>}}},
    postaction={decorate}]
    (3,0.15) arc[start angle=3, end angle=357, radius=3];

  % Petit cercle
  \draw[thick, red, decoration={markings,
    mark=at position 0.5 with {\arrow{<}}},
    postaction={decorate}]
    (0.4,-0.15) arc[start angle=-20, end angle=340, radius=0.4];

  % Lignes sur l'axe réel
  \draw[thick, blue, decoration={markings, mark=at position 0.5 with {\arrow{>}}},
    postaction={decorate}]
    (0.45,0.12) -- (2.95,0.12);
  \draw[thick, blue, decoration={markings, mark=at position 0.5 with {\arrow{<}}},
    postaction={decorate}]
    (0.45,-0.12) -- (2.95,-0.12);

  % Labels
  \node[blue] at (2.5,2.5) {$C_R$};
  \node[red] at (-0.7,0.5) {$C_\varepsilon$};
  \node at (1.7,0.45) {\small $\arg z = 0$};
  \node at (1.7,-0.5) {\small $\arg z = 2\pi$};
  \node at (0,-3.7) {Contour en trou de serrure};

  % Branch cut wavy line
  \draw[thick, decorate, decoration={zigzag, segment length=4, amplitude=1}]
    (0.05,0) -- (-0.5,0);

  % Poles
  \filldraw[red] (1,1.5) circle (1.5pt);
  \filldraw[red] (-1.5,1) circle (1.5pt);
  \filldraw[red] (1,-1.5) circle (1.5pt);
  \filldraw[red] (-1.5,-1) circle (1.5pt);
\end{tikzpicture}
\end{center}

\begin{theoreme}[Intégrales de type Mellin]\label{thm:int-mellin}
Soit $Q$ un polynôme de degré $d \geq 1$ sans zéro sur $[0, +\infty)$, et soit
$0 < \alpha < d$ un réel non entier. Alors, en choisissant la détermination
$z^{\alpha-1} = \abs{z}^{\alpha-1} e^{i(\alpha-1)\arg z}$ avec $\arg z \in [0, 2\pi)$~:
\[
  \int_0^{+\infty} \frac{x^{\alpha-1}}{Q(x)}\, \dd x
  = \frac{2\pi i}{1 - e^{2\pi i \alpha}}
  \sum_{k} \Res\!\left(\frac{z^{\alpha-1}}{Q(z)},\, z_k\right),
\]
où la somme porte sur tous les pôles $z_k$ de $1/Q$ dans
$\C \setminus [0, +\infty)$.
\end{theoreme}

\begin{proof}
On intègre $f(z) = \frac{z^{\alpha-1}}{Q(z)}$ sur le contour en trou de serrure
formé de~:
\begin{itemize}
  \item le segment supérieur de $\varepsilon$ à $R$ (avec $\arg z = 0$)~;
  \item le grand cercle $C_R$ de rayon $R$ (sens direct)~;
  \item le segment inférieur de $R$ à $\varepsilon$ (avec $\arg z = 2\pi$)~;
  \item le petit cercle $C_\varepsilon$ de rayon $\varepsilon$ (sens indirect).
\end{itemize}

Par le théorème des résidus~:
\[
  \int_\varepsilon^R \frac{x^{\alpha-1}}{Q(x)}\,\dd x
  + \int_{C_R} f(z)\,\dd z
  + \int_R^\varepsilon \frac{(xe^{2\pi i})^{\alpha-1}}{Q(x)}\,\dd x
  + \int_{C_\varepsilon} f(z)\,\dd z
  = 2\pi i \sum_k \Res(f, z_k).
\]

Sur le bord inférieur, $z = xe^{2\pi i}$ donne $z^{\alpha-1} = x^{\alpha-1} e^{2\pi i(\alpha-1)}$.
L'intégrale sur ce bord vaut~:
\[
  -e^{2\pi i(\alpha-1)} \int_\varepsilon^R \frac{x^{\alpha-1}}{Q(x)}\,\dd x.
\]

Les intégrales sur $C_R$ et $C_\varepsilon$ tendent vers $0$ quand $R \to +\infty$
et $\varepsilon \to 0$ (grâce à la condition $0 < \alpha < d$). Il reste~:
\[
  (1 - e^{2\pi i(\alpha-1)}) \int_0^{+\infty} \frac{x^{\alpha-1}}{Q(x)}\,\dd x
  = 2\pi i \sum_k \Res(f, z_k).
\]
Or $1 - e^{2\pi i(\alpha-1)} = 1 - e^{-2\pi i} e^{2\pi i\alpha} = 1 - e^{2\pi i\alpha}$
(car $e^{-2\pi i} = 1$), d'où le résultat.
\end{proof}

\begin{exemple}[Intégrale de $1/(1+x^3)$]\label{ex:int-mellin-1}
Calculons $\displaystyle I = \int_0^{+\infty} \frac{\dd x}{1 + x^3}$.

On pose $\alpha = 1$ (attention, $\alpha - 1 = 0$, c'est un cas limite).
On utilise plutôt un contour sectoriel. Considérons le secteur d'angle $2\pi/3$~:

\begin{center}
\begin{tikzpicture}[scale=1.0, >=stealth]
  \draw[->] (-1.5,0) -- (3.5,0) node[right] {$\re$};
  \draw[->] (0,-0.5) -- (0,3) node[above] {$\im$};
  % Contour
  \draw[thick, blue, decoration={markings, mark=at position 0.5 with {\arrow{>}}},
    postaction={decorate}]
    (0,0) -- (3,0);
  \draw[thick, blue, decoration={markings, mark=at position 0.5 with {\arrow{>}}},
    postaction={decorate}]
    (3,0) arc[start angle=0, end angle=120, radius=3];
  \draw[thick, blue, decoration={markings, mark=at position 0.5 with {\arrow{>}}},
    postaction={decorate}]
    ({3*cos(120)},{3*sin(120)}) -- (0,0);
  % Labels
  \node[below] at (3,0) {$R$};
  \node[above left] at ({3*cos(120)},{3*sin(120)}) {$Re^{2i\pi/3}$};
  \node at (1,0.3) {\small $\theta=0$};
  \node at ({1.2*cos(100)},{1.2*sin(100)}) {\small $\frac{2\pi}{3}$};
  % Pole
  \filldraw[red] ({cos(60)},{sin(60)}) circle (1.5pt)
    node[right] {$e^{i\pi/3}$};
\end{tikzpicture}
\end{center}

Sur le rayon $\arg z = 2\pi/3$, on pose $z = te^{2i\pi/3}$ avec $t$ de $R$ à $0$~:
\[
  \int_{\text{rayon}} \frac{\dd z}{1+z^3}
  = \int_R^0 \frac{e^{2i\pi/3}\,\dd t}{1+t^3 e^{2i\pi}} = -e^{2i\pi/3} \int_0^R \frac{\dd t}{1+t^3}.
\]

Le seul pôle dans le secteur est $z_0 = e^{i\pi/3}$, pôle simple~:
\[
  \Res\!\left(\frac{1}{1+z^3}, e^{i\pi/3}\right)
  = \frac{1}{3z_0^2} = \frac{1}{3e^{2i\pi/3}}.
\]

Par le théorème des résidus (l'intégrale sur l'arc tend vers $0$)~:
\[
  (1 - e^{2i\pi/3}) I = 2\pi i \cdot \frac{1}{3e^{2i\pi/3}}.
\]

Calculons $1 - e^{2i\pi/3} = 1 - \left(-\frac{1}{2} + \frac{i\sqrt{3}}{2}\right)
= \frac{3}{2} - \frac{i\sqrt{3}}{2}$. Son module est $\sqrt{3}$, et~:
\[
  I = \frac{2\pi i}{3e^{2i\pi/3}\left(\frac{3}{2} - \frac{i\sqrt{3}}{2}\right)}
  = \frac{2\pi i}{3 \cdot e^{2i\pi/3} \cdot \sqrt{3}\, e^{-i\pi/6}}
  = \frac{2\pi i}{3\sqrt{3}\, e^{i\pi/2}}
  = \frac{2\pi i}{3\sqrt{3} \cdot i}
  = \frac{2\pi}{3\sqrt{3}} = \frac{2\pi\sqrt{3}}{9}.
\]
\end{exemple}

\begin{exemple}[Intégrale de $x^{s-1}/(1+x)$]\label{ex:int-mellin-2}
Calculons $\displaystyle I = \int_0^{+\infty} \frac{x^{s-1}}{1+x}\, \dd x$
pour $0 < s < 1$.

On applique le théorème~\ref{thm:int-mellin} avec $Q(x) = 1+x$, $\alpha = s$.
Le seul pôle de $1/(1+z)$ est $z = -1 = e^{i\pi}$ (dans la coupure
$\arg z \in [0, 2\pi)$). On a~:
\[
  \Res\!\left(\frac{z^{s-1}}{1+z},\, e^{i\pi}\right)
  = \left. z^{s-1}\right|_{z=e^{i\pi}} = e^{i\pi(s-1)} = -e^{i\pi s}.
\]

Attendons~: $z^{s-1} = e^{i\pi(s-1)}$ car $\arg(-1) = \pi$. Donc~:
\[
  I = \frac{2\pi i}{1 - e^{2\pi is}} \cdot (-e^{i\pi s})
  = \frac{-2\pi i\, e^{i\pi s}}{1 - e^{2\pi is}}
  = \frac{2\pi i\, e^{i\pi s}}{e^{2\pi is} - 1}
  = \frac{2\pi i}{e^{i\pi s} - e^{-i\pi s}}
  = \frac{2\pi i}{2i\sin(\pi s)}
  = \frac{\pi}{\sin(\pi s)}.
\]

On retrouve le \textbf{complément d'Euler}~: $\Gamma(s)\Gamma(1-s) = \frac{\pi}{\sin(\pi s)}$.
\end{exemple}

%% ═══════════════════════════════════════════════════════════════════
\section{Principe de l'argument et théorème de Rouché}\label{sec:argument-rouche}
%% ═══════════════════════════════════════════════════════════════════

\begin{theoreme}[Principe de l'argument]\label{thm:principe-argument}
Soit $f$ méromorphe sur un ouvert $U$, avec des zéros $a_1, \ldots, a_p$
(de multiplicités $m_1, \ldots, m_p$) et des pôles $b_1, \ldots, b_q$
(d'ordres $n_1, \ldots, n_q$). Si $\gamma$ est un lacet simple dans $U$
ne passant par aucun zéro ni pôle, parcouru dans le sens direct, alors~:
\[
  \frac{1}{2\pi i} \oint_{\gamma} \frac{f'(z)}{f(z)}\, \dd z
  = \sum_{j=1}^{p} m_j \cdot \ind(\gamma, a_j) - \sum_{k=1}^{q} n_k \cdot \ind(\gamma, b_k).
\]
En particulier, si $\gamma$ entoure tous les zéros et pôles~:
\[
  \frac{1}{2\pi i} \oint_{\gamma} \frac{f'(z)}{f(z)}\, \dd z = Z - P,
\]
où $Z = \sum m_j$ est le nombre de zéros et $P = \sum n_k$ le nombre de pôles,
comptés avec multiplicité.
\end{theoreme}

\begin{proof}
Au voisinage d'un zéro $a_j$ de multiplicité $m_j$, on écrit
$f(z) = (z-a_j)^{m_j} g(z)$ avec $g$ holomorphe et $g(a_j) \neq 0$.
Alors~:
\[
  \frac{f'(z)}{f(z)} = \frac{m_j}{z - a_j} + \frac{g'(z)}{g(z)}.
\]
Donc $\frac{f'}{f}$ a un pôle simple en $a_j$ de résidu $m_j$.

De même, au voisinage d'un pôle $b_k$ d'ordre $n_k$,
$f(z) = (z-b_k)^{-n_k} h(z)$ avec $h(b_k) \neq 0$, et
$\frac{f'(z)}{f(z)} = \frac{-n_k}{z-b_k} + \frac{h'(z)}{h(z)}$.
Le résidu en $b_k$ est $-n_k$.

Le théorème des résidus donne alors le résultat.
\end{proof}

\begin{theoreme}[Théorème de Rouché]\label{thm:rouche}
Soient $f$ et $g$ holomorphes sur un ouvert contenant le lacet simple $\gamma$
et son intérieur. Si~:
\[
  \abs{g(z)} < \abs{f(z)} \qquad \text{pour tout } z \in \gamma,
\]
alors $f$ et $f + g$ ont le même nombre de zéros (comptés avec multiplicité)
à l'intérieur de $\gamma$.
\end{theoreme}

\begin{proof}
Pour $t \in [0,1]$, posons $F_t = f + tg$. Sur $\gamma$, $\abs{tg(z)} \leq \abs{g(z)} < \abs{f(z)}$,
donc $F_t$ ne s'annule pas sur $\gamma$. La quantité
\[
  N(t) = \frac{1}{2\pi i} \oint_\gamma \frac{F_t'(z)}{F_t(z)}\,\dd z
\]
est entière et continue en $t$, donc constante. Ainsi $N(0) = N(1)$, ce qui
signifie que $f = F_0$ et $f + g = F_1$ ont le même nombre de zéros.
\end{proof}

\begin{exemple}[Application du théorème de Rouché]\label{ex:rouche}
Montrons que $p(z) = z^5 + 3z + 1$ a exactement un zéro dans $\abs{z} < 1$.

Posons $f(z) = 3z$ et $g(z) = z^5 + 1$. Sur $\abs{z} = 1$~:
\[
  \abs{f(z)} = 3, \qquad \abs{g(z)} \leq \abs{z}^5 + 1 = 2 < 3 = \abs{f(z)}.
\]
Par Rouché, $p = f + g$ a le même nombre de zéros que $f(z) = 3z$ dans $\abs{z} < 1$,
à savoir exactement un zéro.
\end{exemple}

%% ═══════════════════════════════════════════════════════════════════
\section{Exercices}\label{sec:exo-residus}
%% ═══════════════════════════════════════════════════════════════════

\begin{exercice}[Calcul de résidus]\label{exo:calcul-residus}
Calculer les résidus des fonctions suivantes aux points indiqués~:
\begin{enumerate}[label=\alph*)]
  \item $\dfrac{z^2}{(z^2+1)(z-2)}$ en $z = i$, $z = -i$, $z = 2$~;
  \item $\dfrac{e^z}{z^2(z+1)^3}$ en $z = 0$ et $z = -1$~;
  \item $\dfrac{\sin z}{z^4}$ en $z = 0$~;
  \item $\dfrac{z}{e^z - 1}$ en $z = 2k\pi i$ pour $k \in \Z^*$.
\end{enumerate}
\end{exercice}

\begin{exercice}[Intégrales rationnelles]\label{exo:int-rationnelles}
Calculer à l'aide du théorème des résidus~:
\begin{enumerate}[label=\alph*)]
  \item $\displaystyle \int_{-\infty}^{+\infty} \frac{\dd x}{(x^2+1)(x^2+4)}$~;
  \item $\displaystyle \int_{-\infty}^{+\infty} \frac{x^2\,\dd x}{(x^2+a^2)^3}$
    pour $a > 0$~;
  \item $\displaystyle \int_{-\infty}^{+\infty} \frac{\dd x}{(1+x^2)^n}$
    pour $n \geq 1$ entier.
\end{enumerate}
\end{exercice}

\begin{exercice}[Intégrales trigonométriques]\label{exo:int-trig}
Calculer~:
\begin{enumerate}[label=\alph*)]
  \item $\displaystyle \int_0^{2\pi} \frac{\dd\theta}{5 + 4\cos\theta}$~;
  \item $\displaystyle \int_0^{2\pi} \frac{\cos(2\theta)}{5 - 4\cos\theta}\, \dd\theta$~;
  \item $\displaystyle \int_0^{\pi} \frac{\dd\theta}{(a + \cos\theta)^2}$ pour $a > 1$.
\end{enumerate}
\end{exercice}

\begin{exercice}[Intégrales de Fourier]\label{exo:int-fourier}
Calculer à l'aide du lemme de Jordan~:
\begin{enumerate}[label=\alph*)]
  \item $\displaystyle \int_{-\infty}^{+\infty} \frac{\sin x}{x}\, \dd x$
    \emph{(intégrale de Dirichlet --- utiliser un contour avec indentation)}~;
  \item $\displaystyle \int_{-\infty}^{+\infty} \frac{x\cos x}{(x^2+1)^2}\, \dd x$~;
  \item $\displaystyle \int_0^{+\infty} \frac{\cos(ax)}{x^2+b^2}\, \dd x$
    pour $a > 0$, $b > 0$.
\end{enumerate}
\end{exercice}

\begin{exercice}[Intégrale de Fresnel]\label{exo:fresnel}
En intégrant $e^{-z^2}$ sur le contour formé du segment $[0, R]$, de l'arc
$Re^{i\theta}$ ($\theta \in [0, \pi/4]$) et du segment de $Re^{i\pi/4}$ à $0$,
montrer que~:
\[
  \int_0^{+\infty} \cos(x^2)\,\dd x = \int_0^{+\infty} \sin(x^2)\,\dd x
  = \frac{1}{2}\sqrt{\frac{\pi}{2}}.
\]
\end{exercice}

\begin{exercice}[Intégrales avec coupure]\label{exo:int-coupure}
Calculer~:
\begin{enumerate}[label=\alph*)]
  \item $\displaystyle \int_0^{+\infty} \frac{\sqrt{x}}{x^2 + 1}\, \dd x$~;
  \item $\displaystyle \int_0^{+\infty} \frac{\ln x}{(x^2+1)^2}\, \dd x$~;
  \item $\displaystyle \int_0^{+\infty} \frac{x^{1/3}}{1+x^2}\, \dd x$.
\end{enumerate}
\end{exercice}

\begin{exercice}[Rouché et localisation de zéros]\label{exo:rouche}
\begin{enumerate}[label=\alph*)]
  \item Montrer que $z^7 - 5z^3 + 12$ n'a aucun zéro dans $\abs{z} < 1$
    et a exactement trois zéros dans $\abs{z} < 2$.
  \item Montrer que l'équation $e^z = 3z^2$ a exactement deux solutions dans
    $\abs{z} < 1$.
  \item Soit $p(z) = z^n + a_{n-1}z^{n-1} + \cdots + a_0$ un polynôme de degré $n$.
    Montrer que pour $R$ assez grand, $p$ a exactement $n$ zéros dans $\abs{z} < R$.
    (Retrouver le théorème fondamental de l'algèbre.)
\end{enumerate}
\end{exercice}

\begin{exercice}[Somme de séries par les résidus]\label{exo:somme-series}
En intégrant $\frac{\pi\cot(\pi z)}{z^2}$ sur un carré de côté $2N+1$
centré à l'origine, montrer que~:
\[
  \sum_{n=1}^{+\infty} \frac{1}{n^2} = \frac{\pi^2}{6}.
\]
\emph{Indication~: les résidus de $\pi\cot(\pi z)/z^2$ en $z = n \neq 0$ valent $1/n^2$.}
\end{exercice}

%%%%%%%%%%%%%%%%%%%%%%%%%%%%%%%%%%%%%%%%%%%%%%%%%%%%%%%%%%%%%%%%%%%%
%  Chapitre 8 : Transformations Conformes
%%%%%%%%%%%%%%%%%%%%%%%%%%%%%%%%%%%%%%%%%%%%%%%%%%%%%%%%%%%%%%%%%%%%
\chapter{Transformations Conformes}\label{chap:conforme}

\minitoc
\vspace{1em}

Les transformations conformes --- applications holomorphes à dérivée non nulle ---
préservent les angles et les formes infinitésimales. Elles jouent un rôle central
tant en mathématiques pures (uniformisation des surfaces de Riemann) qu'en physique
mathématique (mécanique des fluides, électrostatique). Ce chapitre étudie les
propriétés fondamentales des applications conformes, avec un accent sur les
transformations de Möbius et leurs applications géométriques.

%% ═══════════════════════════════════════════════════════════════════
\section{Applications conformes~: généralités}\label{sec:conforme-gen}
%% ═══════════════════════════════════════════════════════════════════

\subsection{Définition et conservation des angles}

\begin{definition}[Application conforme]\label{def:conforme}
Soit $U \subset \C$ un ouvert. Une application $f \colon U \to \C$
est dite \textbf{conforme} si elle est holomorphe et $f'(z) \neq 0$ pour tout
$z \in U$. On dit aussi que $f$ est un \textbf{biholomorphisme} (ou application
conforme) de $U$ sur $V = f(U)$ si de plus $f$ est bijective.
\end{definition}

\begin{proposition}[Conservation des angles]\label{prop:conservation-angles}
Une application conforme $f$ préserve les angles orientés entre courbes.
Plus précisément, si $\gamma_1$ et $\gamma_2$ sont deux courbes régulières
se coupant en $z_0$ avec un angle $\alpha$, alors $f \circ \gamma_1$ et
$f \circ \gamma_2$ se coupent en $f(z_0)$ avec le même angle $\alpha$.
\end{proposition}

\begin{proof}
Soient $\gamma_1(t)$ et $\gamma_2(t)$ avec $\gamma_1(0) = \gamma_2(0) = z_0$.
L'angle entre les deux courbes en $z_0$ est
$\alpha = \arg(\gamma_2'(0)) - \arg(\gamma_1'(0))$.

Les courbes images ont pour vecteurs tangents~:
\[
  (f \circ \gamma_j)'(0) = f'(z_0) \cdot \gamma_j'(0), \qquad j = 1, 2.
\]

L'angle entre les courbes images est~:
\begin{align*}
  \arg\bigl(f'(z_0)\gamma_2'(0)\bigr) - \arg\bigl(f'(z_0)\gamma_1'(0)\bigr)
  &= \bigl[\arg f'(z_0) + \arg\gamma_2'(0)\bigr]
    - \bigl[\arg f'(z_0) + \arg\gamma_1'(0)\bigr] \\
  &= \arg\gamma_2'(0) - \arg\gamma_1'(0) = \alpha.
\end{align*}
L'angle est préservé car $f'(z_0) \neq 0$ assure que la multiplication par
$f'(z_0)$ est une similitude directe non dégénérée.
\end{proof}

\begin{remarque}\label{rem:conforme-geom}
Géométriquement, la dérivée $f'(z_0)$ agit par multiplication~:
\begin{itemize}
  \item $\abs{f'(z_0)}$ est le \emph{facteur d'échelle} local (dilatation)~;
  \item $\arg f'(z_0)$ est l'\emph{angle de rotation} local.
\end{itemize}
Ainsi une application conforme est, infinitésimalement, une similitude directe~:
elle dilate et tourne, sans cisailler ni refléter.
\end{remarque}

\begin{center}
\begin{tikzpicture}[scale=0.9, >=stealth]
  % Domaine source
  \begin{scope}[shift={(-4,0)}]
    \draw[thick, gray!50] (-1.5,-1.5) grid[step=0.5] (1.5,1.5);
    \draw[->] (-2,0) -- (2,0);
    \draw[->] (0,-2) -- (0,2);
    \draw[thick, blue] (0.3,0) -- (0.3,1);
    \draw[thick, red] (0.3,0) -- (1.2,0);
    \draw[thick, green!60!black] (0.3,0.05) arc[start angle=10, end angle=80, radius=0.4];
    \node[green!60!black] at (0.75,0.5) {\small $\frac{\pi}{2}$};
    \node[below] at (0,-2.3) {Domaine $U$};
  \end{scope}
  % Flèche
  \draw[->, thick, decorate, decoration={snake, amplitude=2, segment length=10}]
    (-1.3,0) -- (0.7,0) node[midway, above] {$f$};
  % Domaine image
  \begin{scope}[shift={(4,0)}]
    % Grille déformée (schématique)
    \draw[thick, gray!50]
      plot[smooth] coordinates {(-1.5,-1.2) (-0.8,-0.6) (0,0) (0.8,0.8) (1.5,1.5)};
    \draw[thick, gray!50]
      plot[smooth] coordinates {(-1.2,-1.5) (-0.5,-0.7) (0.2,0.1) (1,1) (1.6,1.6)};
    \draw[thick, gray!50]
      plot[smooth] coordinates {(-1.8,-0.8) (-1,-0.3) (-0.2,0.3) (0.6,1.1) (1.2,1.8)};
    \draw[->] (-2,0) -- (2,0);
    \draw[->] (0,-2) -- (0,2);
    \draw[thick, blue] plot[smooth] coordinates {(0.5,0.2) (0.55,0.5) (0.45,0.9) (0.3,1.3)};
    \draw[thick, red] plot[smooth] coordinates {(0.5,0.2) (0.9,0.15) (1.3,0.25)};
    \draw[thick, green!60!black] (0.7,0.3) arc[start angle=5, end angle=95, radius=0.3];
    \node[green!60!black] at (1.0,0.65) {\small $\frac{\pi}{2}$};
    \node[below] at (0,-2.3) {Image $f(U)$};
  \end{scope}
\end{tikzpicture}
\end{center}

\subsection{Exemples fondamentaux}

\begin{exemple}[Transformations affines]\label{ex:affines}
Les applications conformes les plus simples sont~:
\begin{enumerate}[label=\textbf{\arabic*.}]
  \item \textbf{Translation}~: $f(z) = z + b$, $b \in \C$. Déplace le plan sans déformation.
  \item \textbf{Rotation}~: $f(z) = e^{i\theta} z$, $\theta \in \R$. Rotation
    d'angle $\theta$ autour de l'origine.
  \item \textbf{Homothétie}~: $f(z) = \lambda z$, $\lambda > 0$. Dilatation de
    rapport $\lambda$.
  \item \textbf{Similitude directe}~: $f(z) = az + b$, $a \neq 0$. Composition
    d'une rotation, d'une homothétie et d'une translation.
\end{enumerate}
Toutes ces applications sont conformes sur $\C$ entier (et bijectives).
\end{exemple}

\begin{exemple}[La fonction puissance]\label{ex:puissance}
L'application $f(z) = z^n$ ($n \geq 2$ entier) est holomorphe sur $\C$ mais
n'est pas conforme en $z = 0$ (car $f'(0) = 0$). Elle multiplie les angles
par $n$ en l'origine~: le secteur $\{0 < \arg z < \pi/n\}$ est envoyé sur
le demi-plan supérieur.
\end{exemple}

\begin{exemple}[La fonction exponentielle]\label{ex:exp-conforme}
L'application $f(z) = e^z$ est conforme sur $\C$ entier (car $f'(z) = e^z \neq 0$).
Elle envoie~:
\begin{itemize}
  \item la bande $\{0 < \im(z) < \pi\}$ sur le demi-plan supérieur~;
  \item les droites horizontales $\im(z) = c$ sur les demi-droites d'angle $c$~;
  \item les droites verticales $\re(z) = a$ sur les cercles $\abs{w} = e^a$.
\end{itemize}
\end{exemple}

\begin{center}
\begin{tikzpicture}[scale=0.85, >=stealth]
  % Bande source
  \begin{scope}[shift={(-4.5,0)}]
    \fill[blue!8] (-2,0) rectangle (2,2.2);
    \draw[->] (-2.5,0) -- (2.5,0) node[right] {$\re$};
    \draw[->] (0,-0.5) -- (0,2.8) node[above] {$\im$};
    \draw[thick, blue, dashed] (-2,0) -- (2,0) node[right] {\small $\im z=0$};
    \draw[thick, red, dashed] (-2,2.2) -- (2,2.2) node[right] {\small $\im z=\pi$};
    % Grille horizontale
    \foreach \y in {0.55,1.1,1.65} {
      \draw[gray!40] (-2,\y) -- (2,\y);
    }
    % Grille verticale
    \foreach \x in {-1.5,-1,-0.5,0.5,1,1.5} {
      \draw[gray!40] (\x,0) -- (\x,2.2);
    }
    \node[below] at (0,-0.8) {Bande $0 < \im(z) < \pi$};
  \end{scope}
  % Flèche
  \draw[->, thick] (-1.5,1) -- (-0.3,1) node[midway, above] {$e^z$};
  % Demi-plan image
  \begin{scope}[shift={(3.5,0)}]
    \fill[blue!8] (-2.5,0) -- (2.5,0) -- (2.5,2.8) -- (-2.5,2.8) -- cycle;
    \draw[->] (-2.8,0) -- (2.8,0) node[right] {$\re$};
    \draw[->] (0,-0.5) -- (0,2.8) node[above] {$\im$};
    % Cercles (images des verticales)
    \foreach \r in {0.4, 0.8, 1.3, 2.0} {
      \draw[gray!40] (\r,0) arc[start angle=0, end angle=180, radius=\r];
    }
    % Rayons (images des horizontales)
    \foreach \a in {30,60,90,120,150} {
      \draw[gray!40] (0,0) -- ({2.3*cos(\a)},{2.3*sin(\a)});
    }
    \draw[thick, blue] (-2.5,0) -- (2.5,0);
    \node[below] at (0,-0.8) {Demi-plan $\im(w) > 0$};
  \end{scope}
\end{tikzpicture}
\end{center}

\begin{exemple}[L'inversion $z \mapsto 1/z$]\label{ex:inversion}
L'application $f(z) = 1/z$ est un biholomorphisme de $\C^* = \C \setminus \{0\}$
sur lui-même. Elle est son propre inverse ($f \circ f = \mathrm{id}$).
Elle envoie~:
\begin{itemize}
  \item les cercles ne passant pas par $0$ sur des cercles~;
  \item les cercles passant par $0$ sur des droites (et réciproquement)~;
  \item le disque unité $\D$ sur l'extérieur du disque unité, et vice versa.
\end{itemize}
\end{exemple}

%% ═══════════════════════════════════════════════════════════════════
\section{Transformations de Möbius}\label{sec:mobius}
%% ═══════════════════════════════════════════════════════════════════

\subsection{Définition et propriétés algébriques}

\begin{definition}[Transformation de Möbius]\label{def:mobius}
Une \textbf{transformation de Möbius} (ou \emph{homographie}) est une application
de la forme~:
\[
  T(z) = \frac{az + b}{cz + d}, \qquad a, b, c, d \in \C, \quad ad - bc \neq 0.
\]
Elle définit une bijection de la sphère de Riemann $\widehat{\C} = \C \cup \{\infty\}$
sur elle-même, avec les conventions $T(-d/c) = \infty$ et $T(\infty) = a/c$
(si $c \neq 0$), ou $T(\infty) = \infty$ (si $c = 0$).
\end{definition}

\begin{remarque}\label{rem:mobius-matrice}
On associe à $T$ la matrice $M_T = \begin{pmatrix} a & b \\ c & d \end{pmatrix}
\in \mathrm{GL}_2(\C)$. Deux matrices $M$ et $\lambda M$ ($\lambda \neq 0$)
définissent la même transformation. Le groupe des transformations de Möbius
est donc isomorphe au groupe projectif $\mathrm{PGL}_2(\C) = \mathrm{GL}_2(\C) / \C^*$.
La composition correspond au produit matriciel~:
\[
  T_1 \circ T_2 \longleftrightarrow M_{T_1} \cdot M_{T_2}.
\]
L'inverse de $T(z) = \frac{az+b}{cz+d}$ est $T^{-1}(z) = \frac{dz - b}{-cz + a}$.
\end{remarque}

\begin{proposition}[Décomposition des transformations de Möbius]\label{prop:mobius-decomp}
Toute transformation de Möbius est la composée de translations, rotations-homothéties
et d'au plus une inversion $z \mapsto 1/z$.
\end{proposition}

\begin{proof}
Si $c = 0$, alors $T(z) = \frac{a}{d}z + \frac{b}{d}$, qui est une similitude directe.

Si $c \neq 0$, on effectue la division euclidienne~:
\[
  T(z) = \frac{az + b}{cz + d} = \frac{a}{c} - \frac{ad - bc}{c^2} \cdot \frac{1}{z + d/c}.
\]
C'est la composée~: $z \mapsto z + d/c$ (translation), puis $z \mapsto 1/z$ (inversion),
puis $z \mapsto -\frac{ad-bc}{c^2} z$ (similitude), puis $z \mapsto z + a/c$ (translation).
\end{proof}

\begin{theoreme}[Transformations de Möbius et cercles généralisés]\label{thm:mobius-cercles}
Les transformations de Möbius envoient les cercles généralisés (cercles et droites
dans $\widehat{\C}$) sur des cercles généralisés.
\end{theoreme}

\begin{proof}
Par la proposition~\ref{prop:mobius-decomp}, il suffit de le vérifier pour les
translations, les similitudes et l'inversion $z \mapsto 1/z$.

Pour les translations et similitudes, c'est évident (les cercles et droites sont
préservés par ces applications).

Pour l'inversion, un cercle généralisé est l'ensemble des points vérifiant
$A(x^2+y^2) + Bx + Cy + D = 0$ avec $A, B, C, D \in \R$, $(B^2+C^2 > 4AD)$.
En coordonnées complexes, cela s'écrit~:
\[
  A z\conj{z} + \re(Ez) + D = 0, \qquad E = B - iC.
\]

Si $w = 1/z$, alors $z = 1/w$ et l'équation devient~:
\[
  \frac{A}{w\conj{w}} + \re\!\left(\frac{E}{w}\right) + D = 0
  \quad \Longleftrightarrow \quad
  D w\conj{w} + \re(\conj{E}w) + A = 0,
\]
qui est à nouveau l'équation d'un cercle généralisé.
\end{proof}

\subsection{Birapport et action triple transitive}

\begin{definition}[Birapport]\label{def:birapport}
Le \textbf{birapport} de quatre points distincts $z_1, z_2, z_3, z_4 \in \widehat{\C}$
est~:
\[
  (z_1, z_2; z_3, z_4) = \frac{(z_1 - z_3)(z_2 - z_4)}{(z_1 - z_4)(z_2 - z_3)}.
\]
Si l'un des points est $\infty$, on prend la limite naturelle, par exemple
$(\infty, z_2; z_3, z_4) = \frac{z_2 - z_4}{z_2 - z_3}$.
\end{definition}

\begin{theoreme}[Invariance du birapport]\label{thm:birapport-invariant}
Les transformations de Möbius préservent le birapport~: si $T$ est une
transformation de Möbius, alors~:
\[
  (T(z_1), T(z_2); T(z_3), T(z_4)) = (z_1, z_2; z_3, z_4).
\]
\end{theoreme}

\begin{proof}
Il suffit de vérifier pour les générateurs (translations, similitudes, inversion).
Pour une translation $T(z) = z + b$~: chaque différence $T(z_i) - T(z_j) = z_i - z_j$
est inchangée, donc le birapport est préservé.
Pour une similitude $T(z) = az$~: chaque différence est multipliée par $a$, et
les facteurs $a$ se simplifient dans le rapport.
Pour l'inversion $T(z) = 1/z$~:
$T(z_i) - T(z_j) = \frac{z_j - z_i}{z_i z_j}$, et les dénominateurs se simplifient
dans le birapport.
\end{proof}

\begin{theoreme}[Triple transitivité]\label{thm:triple-transitivite}
Pour tout triplet de points distincts $(z_1, z_2, z_3)$ et tout triplet
$(w_1, w_2, w_3)$ dans $\widehat{\C}$, il existe une unique transformation de
Möbius $T$ telle que $T(z_k) = w_k$ pour $k = 1, 2, 3$.
\end{theoreme}

\begin{proof}
\textbf{Existence.} La transformation $S(z) = (z, z_1; z_2, z_3)$ envoie
$z_1 \mapsto 0$, $z_2 \mapsto 1$, $z_3 \mapsto \infty$. De même, soit
$U(w) = (w, w_1; w_2, w_3)$, qui envoie $w_1 \mapsto 0$, $w_2 \mapsto 1$, $w_3 \mapsto \infty$.
Alors $T = U^{-1} \circ S$ convient.

\textbf{Unicité.} Si $T_1$ et $T_2$ conviennent, alors $T_2^{-1} \circ T_1$
fixe $z_1, z_2, z_3$, donc fixe trois points distincts. Or une transformation
de Möbius non triviale a au plus deux points fixes (car $\frac{az+b}{cz+d} = z$
conduit à une équation de degré au plus $2$). Donc $T_2^{-1} \circ T_1 = \mathrm{id}$.
\end{proof}

\subsection{Points fixes et classification}

\begin{definition}[Classification des transformations de Möbius]\label{def:mobius-classif}
Soit $T \neq \mathrm{id}$ une transformation de Möbius. Les points fixes de $T$ sont
les solutions de $T(z) = z$, soit $cz^2 + (d-a)z - b = 0$. On distingue~:
\begin{enumerate}[label=\textbf{\arabic*.}]
  \item \textbf{Parabolique}~: un seul point fixe (discriminant nul). La matrice
    associée a une seule valeur propre. Exemples~: les translations.
  \item \textbf{Elliptique}~: deux points fixes, avec multiplicateur
    $\lambda = \frac{T'(\text{pt fixe}_1)}{1}$ de module $\abs{\lambda} = 1$.
    Exemples~: les rotations.
  \item \textbf{Hyperbolique}~: deux points fixes, avec multiplicateur réel
    positif $\lambda > 0$, $\lambda \neq 1$. Exemples~: les homothéties.
  \item \textbf{Loxodromique}~: deux points fixes, avec multiplicateur
    $\lambda \notin \R_+$ et $\abs{\lambda} \neq 1$. Cas général.
\end{enumerate}
\end{definition}

\begin{remarque}\label{rem:conjugaison-mobius}
Toute transformation de Möbius avec deux points fixes $p, q$ est conjuguée
à $w \mapsto \lambda w$ par la transformation qui envoie $p \mapsto 0$
et $q \mapsto \infty$. Le nombre $\lambda$ est le \emph{multiplicateur}.
\end{remarque}

\begin{center}
\begin{tikzpicture}[scale=0.75, >=stealth]
  % Elliptique
  \begin{scope}[shift={(-5.5,0)}]
    \draw[gray!30] (0,0) circle (1.5);
    \draw[gray!30] (0,0) circle (1);
    \draw[gray!30] (0,0) circle (0.5);
    \filldraw (0,0) circle (2pt) node[below=3pt] {$p$};
    \node at (2,0) {$\bullet$}; \node[right] at (2,0) {$q$};
    \draw[->, blue, thick] (0.9,0.65) arc[start angle=35, end angle=85, radius=1.1];
    \draw[->, blue, thick] (1.3,0.9) arc[start angle=35, end angle=85, radius=1.6];
    \node[below] at (0,-2) {\textbf{Elliptique}};
  \end{scope}
  % Hyperbolique
  \begin{scope}[shift={(0,0)}]
    \filldraw (-1.5,0) circle (2pt) node[below=3pt] {$p$};
    \filldraw (1.5,0) circle (2pt) node[below=3pt] {$q$};
    \draw[->, blue, thick] (-1,0.05) -- (-0.2,0.05);
    \draw[->, blue, thick] (0.2,0.05) -- (1,0.05);
    \draw[->, blue, thick] (-1,-0.05) .. controls (-0.5,0.5) .. (0,0.6)
      .. controls (0.5,0.5) .. (1,-0.05);
    \draw[->, blue, thick] (-1,-0.1) .. controls (-0.5,-0.6) .. (0,-0.7)
      .. controls (0.5,-0.6) .. (1,-0.1);
    \node[below] at (0,-2) {\textbf{Hyperbolique}};
  \end{scope}
  % Loxodromique
  \begin{scope}[shift={(5.5,0)}]
    \filldraw (0,0) circle (2pt) node[below left=2pt] {$p$};
    \draw[->, blue, thick]
      plot[smooth, domain=0:700, samples=100]
      ({(0.5+0.002*\x)*cos(\x)},{(0.5+0.002*\x)*sin(\x)});
    \node[below] at (0,-2) {\textbf{Loxodromique}};
  \end{scope}
\end{tikzpicture}
\end{center}

%% ═══════════════════════════════════════════════════════════════════
\section{La transformation de Cayley}\label{sec:cayley}
%% ═══════════════════════════════════════════════════════════════════

\begin{definition}[Transformation de Cayley]\label{def:cayley}
La \textbf{transformation de Cayley} est la transformation de Möbius~:
\[
  \varphi(z) = \frac{z - i}{z + i}.
\]
\end{definition}

\begin{theoreme}[Propriétés de la transformation de Cayley]\label{thm:cayley}
La transformation de Cayley réalise un biholomorphisme~:
\[
  \varphi \colon \{\im(z) > 0\} \xrightarrow{\;\sim\;} \D = \{\abs{w} < 1\}.
\]
Son inverse est $\varphi^{-1}(w) = i\,\frac{1+w}{1-w}$.
De plus~:
\begin{enumerate}[label=(\roman*)]
  \item $\varphi$ envoie l'axe réel $\R$ sur le cercle unité $\T$ (privé de $1$)~;
  \item $\varphi(0) = -1$, $\varphi(\infty) = 1$, $\varphi(i) = 0$~;
  \item $\varphi(-1) = \frac{-1-i}{-1+i} = i$, $\varphi(1) = \frac{1-i}{1+i} = -i$.
\end{enumerate}
\end{theoreme}

\begin{proof}
Pour $z$ avec $\im(z) > 0$, montrons que $\abs{\varphi(z)} < 1$. Posons $z = x + iy$
avec $y > 0$~:
\[
  \abs{\varphi(z)}^2 = \frac{\abs{z-i}^2}{\abs{z+i}^2}
  = \frac{x^2 + (y-1)^2}{x^2 + (y+1)^2}
  = \frac{x^2 + y^2 - 2y + 1}{x^2 + y^2 + 2y + 1} < 1,
\]
car le numérateur est strictement inférieur au dénominateur (la différence est $-4y < 0$).

Pour $z \in \R$ (i.e.\ $y = 0$), $\abs{\varphi(z)}^2 = \frac{x^2+1}{x^2+1} = 1$,
donc $\varphi(\R) \subset \T$.

L'inverse se calcule en résolvant $w = \frac{z-i}{z+i}$ en $z$~:
$w(z+i) = z - i$, soit $z(w-1) = -i(w+1)$, d'où $z = i\frac{1+w}{1-w}$.
\end{proof}

\begin{center}
\begin{tikzpicture}[scale=1.1, >=stealth]
  % Demi-plan
  \begin{scope}[shift={(-4,0)}]
    \fill[blue!8] (-2.5,0) rectangle (2.5,2.5);
    \draw[->] (-2.8,0) -- (2.8,0) node[right] {$\re$};
    \draw[->] (0,-0.5) -- (0,2.8) node[above] {$\im$};
    \draw[very thick, red] (-2.5,0) -- (2.5,0);
    \filldraw[blue] (0,1) circle (2pt) node[right] {$i$};
    \filldraw (0,0) circle (1.5pt) node[below right] {$0$};
    \filldraw (-1,0) circle (1.5pt) node[below] {$-1$};
    \filldraw (1,0) circle (1.5pt) node[below] {$1$};
    \node at (0,1.8) {$\mathbb{H}$};
    \node[below] at (0,-0.8) {Demi-plan supérieur};
  \end{scope}
  % Flèche
  \draw[->, thick] (-0.8,1) -- (0.8,1) node[midway, above] {$\varphi$};
  \draw[<-, thick] (-0.8,0.5) -- (0.8,0.5) node[midway, below] {$\varphi^{-1}$};
  % Disque
  \begin{scope}[shift={(4,0)}]
    \fill[blue!8] (0,1) circle (1.5);
    \draw[->] (-2.2,0) -- (2.2,0) node[right] {$\re$};
    \draw[->] (0,-0.8) -- (0,2.8) node[above] {$\im$};
    \draw[very thick, red] (0,1) circle (1.5);
    \filldraw[blue] (0,1) circle (2pt) node[above right] {$0$};
    \filldraw (-1.5,1) circle (1.5pt) node[left] {$-1$};
    \filldraw (0,2.5) circle (1.5pt) node[above right] {$i$};
    \filldraw (0,-0.5) circle (1.5pt) node[below right] {$-i$};
    \node at (0.6,1.5) {$\D$};
    \node[below] at (0,-1.1) {Disque unité};
  \end{scope}
\end{tikzpicture}
\end{center}

\begin{proposition}[Automorphismes du disque]\label{prop:auto-disque}
Tout automorphisme (biholomorphisme sur lui-même) du disque unité $\D$ est de la
forme~:
\[
  f(z) = e^{i\theta} \frac{z - a}{1 - \conj{a}z}, \qquad \theta \in \R,\; a \in \D.
\]
\end{proposition}

\begin{proof}
Soit $f \colon \D \to \D$ un automorphisme, et posons $a = f^{-1}(0)$.
L'automorphisme de Blaschke $B_a(z) = \frac{z-a}{1-\conj{a}z}$ vérifie $B_a(a) = 0$.
Donc $g = f \circ B_a^{-1}$ est un automorphisme de $\D$ fixant l'origine.

Par le lemme de Schwarz, $\abs{g(z)} \leq \abs{z}$ pour tout $z \in \D$.
En appliquant le même raisonnement à $g^{-1}$~: $\abs{g^{-1}(w)} \leq \abs{w}$,
soit $\abs{z} \leq \abs{g(z)}$. Donc $\abs{g(z)} = \abs{z}$, et le cas
d'égalité du lemme de Schwarz donne $g(z) = e^{i\theta}z$.
Ainsi $f = g \circ B_a$, ce qui donne la forme annoncée.
\end{proof}

\begin{corollaire}[Automorphismes du demi-plan]\label{cor:auto-demiplan}
Tout automorphisme du demi-plan supérieur $\mathbb{H} = \{\im(z) > 0\}$
est de la forme~:
\[
  T(z) = \frac{az + b}{cz + d}, \qquad a,b,c,d \in \R, \quad ad - bc > 0.
\]
Le groupe $\mathrm{Aut}(\mathbb{H})$ est isomorphe à $\mathrm{PSL}_2(\R)$.
\end{corollaire}

%% ═══════════════════════════════════════════════════════════════════
\section{Le théorème de Riemann}\label{sec:riemann-mapping}
%% ═══════════════════════════════════════════════════════════════════

\begin{theoreme}[Théorème de représentation conforme de Riemann]\label{thm:riemann-mapping}
Soit $U \subsetneq \C$ un ouvert simplement connexe non vide, distinct de $\C$.
Alors il existe un biholomorphisme $f \colon U \to \D$. De plus, pour tout
$z_0 \in U$, il existe un unique tel biholomorphisme $f$ vérifiant
$f(z_0) = 0$ et $f'(z_0) > 0$.
\end{theoreme}

\begin{remarque}\label{rem:riemann-mapping}
Ce théorème fondamental affirme que tous les ouverts simplement connexes
(différents de $\C$) sont conformément équivalents. Sa preuve complète,
due à Koebe (1912) et reposant sur le théorème de Montel, dépasse le cadre
de ce cours. Nous en donnons les grandes lignes~:
\begin{enumerate}[label=\textbf{Étape \arabic*.}]
  \item On considère la famille $\mathcal{F}$ de toutes les injections
    holomorphes $f \colon U \to \D$ avec $f(z_0) = 0$.
  \item On montre que $\mathcal{F}$ est non vide (en utilisant l'hypothèse
    $U \neq \C$ et l'existence d'une racine carrée holomorphe).
  \item On prend $f \in \mathcal{F}$ maximisant $\abs{f'(z_0)}$ (l'existence
    est assurée par le théorème de Montel~: $\mathcal{F}$ est une famille
    normale).
  \item On montre que ce $f$ est surjectif, sinon on pourrait construire
    un élément de $\mathcal{F}$ avec une dérivée strictement plus grande.
\end{enumerate}
\end{remarque}

\begin{remarque}\label{rem:riemann-C}
L'hypothèse $U \neq \C$ est essentielle~: $\C$ n'est pas conformément
équivalent à $\D$. En effet, toute application holomorphe bornée sur
$\C$ est constante (théorème de Liouville).
\end{remarque}

%% ═══════════════════════════════════════════════════════════════════
\section{Applications conformes classiques}\label{sec:conforme-classiques}
%% ═══════════════════════════════════════════════════════════════════

Le théorème de Riemann garantit l'existence de biholomorphismes, mais ne les
construit pas explicitement. Voici un catalogue de transformations conformes
classiques, indispensables en pratique.

\subsection{La fonction de Joukowski}

\begin{definition}[Transformation de Joukowski]\label{def:joukowski}
La \textbf{transformation de Joukowski} est~:
\[
  J(z) = \frac{1}{2}\left(z + \frac{1}{z}\right).
\]
\end{definition}

\begin{proposition}\label{prop:joukowski}
La transformation de Joukowski~:
\begin{enumerate}[label=(\roman*)]
  \item est conforme sur $\C \setminus \{-1, 0, 1\}$ (les points critiques sont $\pm 1$)~;
  \item envoie le cercle $\abs{z} = r > 1$ sur l'ellipse
    $\frac{u^2}{\left(\frac{r+r^{-1}}{2}\right)^2}
    + \frac{v^2}{\left(\frac{r-r^{-1}}{2}\right)^2} = 1$~;
  \item réalise un biholomorphisme de $\{\abs{z} > 1\}$ sur
    $\C \setminus [-1, 1]$.
\end{enumerate}
\end{proposition}

\begin{proof}
On a $J'(z) = \frac{1}{2}(1 - z^{-2})$, qui s'annule en $z = \pm 1$.
Pour $z = re^{i\theta}$ avec $r > 1$~:
\[
  J(re^{i\theta}) = \frac{r + r^{-1}}{2}\cos\theta + i\frac{r - r^{-1}}{2}\sin\theta,
\]
ce qui décrit une ellipse de demi-axes $a = \frac{r+r^{-1}}{2}$ et $b = \frac{r-r^{-1}}{2}$.
Pour la bijectivité sur $\{\abs{z}>1\}$~: si $J(z_1) = J(z_2)$ alors $z_1 + 1/z_1 = z_2 + 1/z_2$,
donc $(z_1 - z_2)(1 - 1/(z_1 z_2)) = 0$, soit $z_1 = z_2$ ou $z_1 z_2 = 1$.
Comme $\abs{z_1}, \abs{z_2} > 1$, la seconde possibilité est exclue.
\end{proof}

\begin{center}
\begin{tikzpicture}[scale=1.0, >=stealth]
  % Source: cercles concentriques
  \begin{scope}[shift={(-4,0)}]
    \draw[->] (-2.5,0) -- (2.5,0) node[right] {$\re$};
    \draw[->] (0,-2.5) -- (0,2.5) node[above] {$\im$};
    \draw[thick, red] (0,0) circle (1);
    \foreach \r in {1.2, 1.5, 2.0} {
      \draw[blue!60] (0,0) circle (\r);
    }
    \foreach \a in {0,30,...,330} {
      \draw[gray!40] (0,0) -- ({2.2*cos(\a)},{2.2*sin(\a)});
    }
    \node[below] at (0,-2.8) {$\abs{z} \geq 1$};
  \end{scope}
  % Flèche
  \draw[->, thick] (-1,0) -- (0.5,0) node[midway, above] {$J$};
  % Image: ellipses
  \begin{scope}[shift={(4,0)}]
    \draw[->] (-2.8,0) -- (2.8,0) node[right] {$\re$};
    \draw[->] (0,-2.2) -- (0,2.2) node[above] {$\im$};
    \draw[thick, red] (-1,0) -- (1,0);  % segment [-1,1]
    \draw[blue!60] (0,0) ellipse (1.25 and 0.42);
    \draw[blue!60] (0,0) ellipse (1.58 and 0.83);
    \draw[blue!60] (0,0) ellipse (2.25 and 1.75);
    \node[below] at (0,-2.8) {$\C \setminus [-1,1]$};
    \node[red, below] at (0.5,0) {\small $[-1,1]$};
  \end{scope}
\end{tikzpicture}
\end{center}

\subsection{Bandes, secteurs et demi-plans}

\begin{exemple}[Secteur angulaire vers demi-plan]\label{ex:secteur-demiplan}
Le secteur $S_\alpha = \{z \in \C : 0 < \arg z < \alpha\}$ (avec $0 < \alpha \leq 2\pi$)
est envoyé sur le demi-plan supérieur par~:
\[
  f(z) = z^{\pi/\alpha}.
\]
En effet, si $z = re^{i\theta}$ avec $0 < \theta < \alpha$, alors
$f(z) = r^{\pi/\alpha} e^{i\pi\theta/\alpha}$ et $0 < \pi\theta/\alpha < \pi$,
donc $\im(f(z)) > 0$.
\end{exemple}

\begin{exemple}[Bande vers demi-plan]\label{ex:bande-demiplan}
La bande horizontale $B = \{z : 0 < \im(z) < \pi\}$ est envoyée sur le
demi-plan supérieur par $f(z) = e^z$ (cf.\ exemple~\ref{ex:exp-conforme}).
\end{exemple}

\begin{exemple}[Bande verticale vers disque]\label{ex:bande-disque}
La bande verticale $V = \{z : -\pi/2 < \re(z) < \pi/2\}$ est envoyée sur le disque par
la composée~: d'abord $z \mapsto iz$ (rotation de $\pi/2$, transformant $V$ en la bande
horizontale $\{0 < \im(w) < \pi/2\}$), puis $w \mapsto e^{2w}$ (envoi sur le
demi-plan supérieur, en ajustant), puis la transformation de Cayley. Plus directement~:
\[
  f(z) = \frac{e^z - 1}{e^z + 1} = \tanh(z/2)
\]
envoie la bande $\{-\pi < \im(z) < \pi\}$ sur $\C \setminus ((-\infty, -1] \cup [1, +\infty))$.
Avec un ajustement, $g(z) = \tanh(z)$ réalise une application conforme de la bande
$\{-\pi/4 < \im(z) < \pi/4\}$ sur un domaine du disque.
\end{exemple}

\begin{exemple}[Lentille]\label{ex:lentille}
Une \textbf{lentille} est l'intersection de deux disques. Considérons la lentille
délimitée par deux arcs de cercle passant par $-1$ et $1$ avec des angles intérieurs
$\alpha$ en ces points. La transformation~:
\[
  f(z) = \left(\frac{1+z}{1-z}\right)^{\pi/\alpha}
\]
envoie cette lentille sur le demi-plan supérieur. En effet, $g(z) = \frac{1+z}{1-z}$
envoie $-1 \mapsto 0$, $1 \mapsto \infty$ et transforme la lentille en un secteur
d'angle $\alpha$, puis $w \mapsto w^{\pi/\alpha}$ ouvre ce secteur en demi-plan.
\end{exemple}

\subsection{Tableau récapitulatif}

\begin{center}
\renewcommand{\arraystretch}{1.6}
\begin{tabular}{lll}
\toprule
\textbf{Domaine source} & \textbf{Domaine image} & \textbf{Application} \\
\midrule
Demi-plan $\mathbb{H}$ & Disque $\D$ & $\varphi(z) = \frac{z-i}{z+i}$ \\[4pt]
Secteur $0 < \arg z < \alpha$ & Demi-plan $\mathbb{H}$ & $z^{\pi/\alpha}$ \\[4pt]
Bande $0 < \im z < \pi$ & Demi-plan $\mathbb{H}$ & $e^z$ \\[4pt]
Bande $0 < \im z < h$ & Demi-plan $\mathbb{H}$ & $e^{\pi z/h}$ \\[4pt]
$\C \setminus (-\infty, 0]$ & Bande $-\pi < \im w < \pi$ & $\Log z$ \\[4pt]
$\{\abs{z} > 1\}$ & $\C \setminus [-1,1]$ & $\frac{1}{2}(z+1/z)$ \\[4pt]
$\D$ & $\D$ & $e^{i\theta}\frac{z-a}{1-\conj{a}z}$ \\
\bottomrule
\end{tabular}
\end{center}

%% ═══════════════════════════════════════════════════════════════════
\section{Exercices}\label{sec:exo-conforme}
%% ═══════════════════════════════════════════════════════════════════

\begin{exercice}[Transformations de Möbius élémentaires]\label{exo:mobius-elem}
\begin{enumerate}[label=\alph*)]
  \item Trouver la transformation de Möbius qui envoie $0 \mapsto 1$,
    $1 \mapsto i$, $\infty \mapsto -1$.
  \item Trouver la transformation de Möbius qui envoie le cercle $\abs{z} = 1$
    sur l'axe réel, avec $1 \mapsto 0$, $i \mapsto 1$, $-1 \mapsto \infty$.
  \item Montrer que la transformation $T(z) = \frac{z-1}{z+1}$ envoie le
    demi-plan $\re(z) > 0$ sur le disque $\abs{w} < 1$.
\end{enumerate}
\end{exercice}

\begin{exercice}[Birapport]\label{exo:birapport}
\begin{enumerate}[label=\alph*)]
  \item Calculer le birapport $(0, 1; 2, \infty)$.
  \item Montrer que quatre points $z_1, z_2, z_3, z_4$ sont cocycliques
    (sur un même cercle généralisé) si et seulement si leur birapport est réel.
  \item Soit $T$ une transformation de Möbius fixant $0$ et $\infty$.
    Montrer que $T(z) = \lambda z$ pour un certain $\lambda \in \C^*$.
\end{enumerate}
\end{exercice}

\begin{exercice}[Points fixes et classification]\label{exo:classif-mobius}
Classifier les transformations de Möbius suivantes (parabolique, elliptique,
hyperbolique, loxodromique) et déterminer leurs points fixes~:
\begin{enumerate}[label=\alph*)]
  \item $T(z) = \frac{2z+1}{z+1}$~;
  \item $T(z) = \frac{3z-2}{z}$~;
  \item $T(z) = \frac{(1+i)z}{z+i}$~;
  \item $T(z) = z + 1$.
\end{enumerate}
\end{exercice}

\begin{exercice}[Automorphismes du disque]\label{exo:auto-disque}
\begin{enumerate}[label=\alph*)]
  \item Montrer que $B_a(z) = \frac{z-a}{1-\conj{a}z}$ envoie $\D$ sur $\D$
    et $\T$ sur $\T$ pour tout $a \in \D$.
  \item Calculer $B_a \circ B_a$ et vérifier que $B_a^{-1} = B_{-a}$
    (à rotation près).
  \item Soit $f \colon \D \to \D$ holomorphe avec $f(0) = 0$ et $f(a) = 0$
    pour un $a \neq 0$. Montrer que $\abs{f'(0)} \leq \abs{a}$.
\end{enumerate}
\end{exercice}

\begin{exercice}[Construction d'applications conformes]\label{exo:construction-conforme}
Construire explicitement des biholomorphismes entre les domaines suivants et $\D$~:
\begin{enumerate}[label=\alph*)]
  \item le quart de plan $\{z : \re(z) > 0,\, \im(z) > 0\}$~;
  \item la bande $\{z : 0 < \im(z) < 1\}$~;
  \item le domaine $\C \setminus [0, +\infty)$ (plan coupé)~;
  \item le demi-disque $\{z : \abs{z} < 1,\, \im(z) > 0\}$~;
  \item l'intersection des disques $\abs{z-1} < 2$ et $\abs{z+1} < 2$.
\end{enumerate}
\emph{Indication~: procéder par étapes, en composant des transformations du
tableau récapitulatif.}
\end{exercice}

\begin{exercice}[Formule de Schwarz--Christoffel]\label{exo:schwarz-christoffel}
La formule de Schwarz--Christoffel donne l'application conforme du demi-plan
supérieur sur un polygone à $n$ sommets. Pour un rectangle de sommets
$0, K, K+iK', iK'$, l'application conforme du demi-plan supérieur est~:
\[
  f(z) = C \int_0^z \frac{\dd\zeta}{\sqrt{(1-\zeta^2)(1-k^2\zeta^2)}}
\]
où $0 < k < 1$ et $C$ est une constante.
\begin{enumerate}[label=\alph*)]
  \item Vérifier que $f$ envoie l'axe réel sur le bord du rectangle
    (identifier les points $-1/k, -1, 0, 1, 1/k$ avec les sommets).
  \item En déduire que $K = \int_0^1 \frac{\dd t}{\sqrt{(1-t^2)(1-k^2t^2)}}$
    (intégrale elliptique complète de première espèce).
\end{enumerate}
\end{exercice}

\begin{exercice}[Application à la physique]\label{exo:phys-conforme}
Le potentiel complexe d'un écoulement de fluide parfait incompressible est une
fonction holomorphe $\Phi(z) = \phi(x,y) + i\psi(x,y)$ où $\phi$ est le potentiel
de vitesse et $\psi$ la fonction de courant.

\begin{enumerate}[label=\alph*)]
  \item On considère l'écoulement uniforme $\Phi(z) = Vz$ autour d'un cylindre
    circulaire. En utilisant la transformation de Joukowski $J(z) = z + 1/z$,
    montrer que le potentiel complexe de l'écoulement autour d'un profil de
    Joukowski est $\Phi(J^{-1}(w))$.
  \item Déterminer les lignes de courant de l'écoulement $\Phi(z) = z^2$
    (écoulement dans un coin de $90°$).
  \item En utilisant $f(z) = \cosh(z)$, décrire l'écoulement dans une bande
    infinie.
\end{enumerate}
\end{exercice}
%%%%%%%%%%%%%%%%%%%%%%%%%%%%%%%%%%%%%%%%%%%%%%%%%%%%%%%%%%%%%%%%%%%%
%  CHAPITRE 9 — Théorème de Rouché, Principe de l'Argument
%%%%%%%%%%%%%%%%%%%%%%%%%%%%%%%%%%%%%%%%%%%%%%%%%%%%%%%%%%%%%%%%%%%%
\chapter{Théorème de Rouché et Principe de l'Argument}\label{ch:rouche-argument}

\section{Zéros des fonctions holomorphes}\label{sec:zeros-holomorphes}

Nous commençons ce chapitre par une étude systématique des zéros des fonctions
holomorphes, en précisant la notion fondamentale d'\emph{ordre} d'un zéro.

\begin{definition}[Zéro d'ordre $m$]\label{def:zero-ordre}
Soit $f$ une fonction holomorphe sur un ouvert $\Omega \subseteq \C$ et soit
$a \in \Omega$ tel que $f(a) = 0$. On dit que $a$ est un \textbf{zéro d'ordre}
(ou de \textbf{multiplicité}) $m \geq 1$\index{zéro!ordre}\index{multiplicité d'un zéro}
de $f$ si
\[
  f(a) = f'(a) = f''(a) = \cdots = f^{(m-1)}(a) = 0
  \quad \text{et} \quad f^{(m)}(a) \neq 0.
\]
De manière équivalente, $a$ est un zéro d'ordre $m$ si et seulement s'il existe
une fonction holomorphe $g$ sur $\Omega$ avec $g(a) \neq 0$ telle que
\[
  f(z) = (z - a)^m g(z) \quad \text{pour tout } z \in \Omega.
\]
\end{definition}

\begin{proposition}[Isolement des zéros]\label{prop:zeros-isoles}
Soit $f$ holomorphe et non identiquement nulle sur un domaine (ouvert connexe)
$\Omega \subseteq \C$. Alors les zéros de $f$ sont \textbf{isolés} : pour tout
zéro $a$ de $f$, il existe $r > 0$ tel que $f(z) \neq 0$ pour
$0 < \abs{z - a} < r$.
\end{proposition}

\begin{proof}
Soit $a$ un zéro d'ordre $m$ de $f$. Par la définition~\ref{def:zero-ordre},
on écrit $f(z) = (z-a)^m g(z)$ avec $g$ holomorphe et $g(a) \neq 0$. Puisque
$g$ est continue et $g(a) \neq 0$, il existe $r > 0$ tel que $g(z) \neq 0$ pour
$\abs{z-a} < r$. Ainsi, pour $0 < \abs{z-a} < r$, on a
$f(z) = (z-a)^m g(z) \neq 0$.
\end{proof}

\begin{corollaire}[Principe des zéros isolés]\label{cor:zeros-isoles}
Soient $f$ et $g$ holomorphes sur un domaine $\Omega$. Si l'ensemble
$\{z \in \Omega : f(z) = g(z)\}$ admet un point d'accumulation dans $\Omega$,
alors $f \equiv g$ sur $\Omega$.
\end{corollaire}

\begin{proof}
La fonction $h = f - g$ est holomorphe sur $\Omega$ et ses zéros admettent un
point d'accumulation. Si $h$ n'était pas identiquement nulle, ses zéros seraient
isolés par la proposition~\ref{prop:zeros-isoles}, contradiction. Donc
$h \equiv 0$, c'est-à-dire $f \equiv g$.
\end{proof}

\begin{proposition}[Nombre de zéros par formule intégrale]\label{prop:nb-zeros-integrale}
Soit $f$ holomorphe sur un ouvert contenant le disque fermé $\overline{D(a,r)}$
et sans zéros sur le cercle $\abs{z - a} = r$. Alors le nombre de zéros de $f$
dans $D(a,r)$, comptés avec multiplicité, est donné par
\[
  N = \frac{1}{2\pi i} \oint_{\abs{z-a}=r} \frac{f'(z)}{f(z)} \, \dd z.
\]
\end{proposition}

\begin{proof}
Soient $a_1, \ldots, a_k$ les zéros de $f$ dans $D(a,r)$, de multiplicités
respectives $m_1, \ldots, m_k$. Alors $f(z) = (z - a_j)^{m_j} g_j(z)$ où
$g_j$ est holomorphe et $g_j(a_j) \neq 0$. On calcule la dérivée logarithmique :
\[
  \frac{f'(z)}{f(z)} = \sum_{j=1}^{k} \frac{m_j}{z - a_j} + h(z),
\]
où $h(z) = \sum_{j=1}^k g_j'(z)/g_j(z)$ est holomorphe dans $D(a,r)$
(plus précisément, chaque terme $g_j'/g_j$ est holomorphe au voisinage de $a_j$,
et la décomposition complète montre que $h$ est holomorphe sur $D(a,r)$).
Par le théorème des résidus :
\[
  \frac{1}{2\pi i} \oint_{\abs{z-a}=r} \frac{f'(z)}{f(z)} \, \dd z
  = \sum_{j=1}^{k} m_j = N. \qedhere
\]
\end{proof}

\section{Le principe de l'argument}\label{sec:principe-argument}

Le principe de l'argument est l'une des applications les plus élégantes et
puissantes du théorème des résidus. Il relie le nombre de zéros et de pôles
d'une fonction méromorphe à une intégrale de contour de sa dérivée logarithmique.

\begin{theoreme}[Principe de l'argument]\label{thm:principe-argument}
\index{principe de l'argument}
Soit $f$ une fonction méromorphe sur un ouvert contenant le disque fermé
$\overline{D(a,R)}$, n'ayant ni zéros ni pôles sur le cercle
$\gamma : \abs{z - a} = R$. Notons $N$ le nombre de zéros et $P$ le nombre de
pôles de $f$ dans $D(a,R)$, comptés avec multiplicité. Alors
\begin{equation}\label{eq:principe-argument}
  \frac{1}{2\pi i} \oint_{\gamma} \frac{f'(z)}{f(z)} \, \dd z = N - P.
\end{equation}
Plus généralement, si $\gamma$ est une courbe de Jordan rectifiable, positivement
orientée, bordant un domaine $\Omega$, et si $f$ est méromorphe sur un ouvert
contenant $\overline{\Omega}$, sans zéros ni pôles sur $\gamma$, alors la même
formule vaut avec $N$ et $P$ comptés dans $\Omega$.
\end{theoreme}

\begin{proof}
Soient $a_1, \ldots, a_k$ les zéros de $f$ dans $D(a,R)$, de multiplicités
$m_1, \ldots, m_k$, et $b_1, \ldots, b_\ell$ les pôles de $f$ dans $D(a,R)$,
d'ordres $n_1, \ldots, n_\ell$.

Au voisinage d'un zéro $a_j$ d'ordre $m_j$, on écrit $f(z) = (z - a_j)^{m_j}g_j(z)$
avec $g_j$ holomorphe et $g_j(a_j) \neq 0$. Alors :
\[
  \frac{f'(z)}{f(z)} = \frac{m_j}{z - a_j} + \frac{g_j'(z)}{g_j(z)},
\]
et donc $\Res_{z=a_j} \frac{f'(z)}{f(z)} = m_j$.

Au voisinage d'un pôle $b_j$ d'ordre $n_j$, on écrit
$f(z) = (z - b_j)^{-n_j} h_j(z)$ avec $h_j$ holomorphe et $h_j(b_j) \neq 0$. Alors :
\[
  \frac{f'(z)}{f(z)} = \frac{-n_j}{z - b_j} + \frac{h_j'(z)}{h_j(z)},
\]
et donc $\Res_{z=b_j} \frac{f'(z)}{f(z)} = -n_j$.

Par le théorème des résidus :
\[
  \frac{1}{2\pi i} \oint_{\gamma} \frac{f'(z)}{f(z)} \, \dd z
  = \sum_{j=1}^{k} m_j + \sum_{j=1}^{\ell} (-n_j)
  = N - P. \qedhere
\]
\end{proof}

\begin{remarque}[Interprétation géométrique]\label{rem:argument-geometrique}
L'intégrale $\frac{1}{2\pi i} \oint_{\gamma} \frac{f'(z)}{f(z)} \, \dd z$
représente la \textbf{variation totale} de l'argument de $f(z)$ lorsque $z$
parcourt $\gamma$, divisée par $2\pi$. En effet, formellement :
\[
  \frac{1}{2\pi i} \oint_{\gamma} \frac{f'(z)}{f(z)} \, \dd z
  = \frac{1}{2\pi i} \oint_{\gamma} \dd(\Log f(z))
  = \frac{1}{2\pi i} \Delta_{\gamma} \log f(z)
  = \frac{\Delta_{\gamma} \Arg f}{2\pi},
\]
où $\Delta_\gamma \Arg f$ désigne la variation totale de l'argument de $f$ le
long de $\gamma$. Géométriquement, $N - P$ est le nombre de tours que la courbe
image $f \circ \gamma$ effectue autour de l'origine, c'est-à-dire l'indice de
$f \circ \gamma$ par rapport à~$0$ :
\[
  N - P = \ind(f \circ \gamma, 0).
\]
\end{remarque}

\begin{center}
\begin{tikzpicture}[scale=1.0]
  % Domaine de départ
  \begin{scope}[xshift=-5cm]
    \draw[thick, blue!70!black, ->] (0,0) circle (2);
    \node[blue!70!black] at (0,-2.5) {$\gamma$};

    % Zéros
    \fill[red] (0.5, 0.7) circle (2pt) node[above right, font=\small] {$a_1$};
    \fill[red] (-0.8, -0.3) circle (2pt) node[below left, font=\small] {$a_2$};

    % Pôle
    \node[blue, font=\Large] at (0.3, -0.8) {$\times$};
    \node[blue, below right, font=\small] at (0.3, -0.8) {$b_1$};

    \node at (0, -3.3) {Plan $z$};
  \end{scope}

  % Flèche
  \draw[->, thick, >=stealth] (-2.2, 0) -- (-0.8, 0) node[midway, above] {$f$};

  % Domaine d'arrivée
  \begin{scope}[xshift=3cm]
    \draw[thick, red!70!black, ->]
      plot[smooth, domain=0:720, samples=200]
      ({1.5*cos(\x) + 0.3*cos(2*\x)}, {1.5*sin(\x) + 0.3*sin(2*\x)});

    \fill (0,0) circle (2pt) node[below right, font=\small] {$0$};
    \node at (0, -3.3) {Plan $w$};
    \node[red!70!black] at (2.5, 1.5) {$f \circ \gamma$};
    \node at (0, -4) {$\ind(f\circ\gamma, 0) = N - P = 1$};
  \end{scope}
\end{tikzpicture}
\end{center}

\begin{exemple}[Application directe du principe de l'argument]\label{ex:principe-argument-direct}
Considérons $f(z) = \frac{z^3 - 1}{z - 2}$ et le cercle $\gamma : \abs{z} = 3$.
Les zéros de $f$ dans $D(0,3)$ sont les racines cubiques de l'unité :
$z = 1, e^{2\pi i/3}, e^{-2\pi i/3}$, chacune de multiplicité~$1$, donc $N = 3$.
Le seul pôle dans $D(0,3)$ est $z = 2$, simple, donc $P = 1$. Par le principe
de l'argument :
\[
  \frac{1}{2\pi i} \oint_{\abs{z}=3} \frac{f'(z)}{f(z)} \, \dd z = 3 - 1 = 2.
\]
\end{exemple}

\begin{exemple}[Variation de l'argument]\label{ex:variation-argument}
Soit $f(z) = z^2 + 1$ et $\gamma : \abs{z} = 2$. Les zéros de $f$ sont
$z = \pm i$, chacun de multiplicité~$1$. Il n'y a pas de pôle. Donc :
\[
  \frac{1}{2\pi i} \oint_{\abs{z}=2} \frac{2z}{z^2+1} \, \dd z = 2.
\]
La courbe $f \circ \gamma$ fait $2$ tours autour de l'origine, ce qui correspond
à la variation totale $\Delta_\gamma \Arg f = 4\pi$.
\end{exemple}

\section{Le théorème de Rouché}\label{sec:rouche}

Le théorème de Rouché est un outil remarquable pour déterminer le nombre de
zéros d'une fonction holomorphe dans un domaine. Il repose sur le principe de
l'argument et l'idée qu'une <<~petite perturbation~>> ne change pas le nombre
de zéros.

\begin{theoreme}[Théorème de Rouché]\label{thm:rouche}
\index{théorème de Rouché}
Soient $f$ et $g$ holomorphes sur un ouvert contenant $\overline{\Omega}$, où
$\Omega$ est un domaine borné dont le bord $\partial\Omega$ est une courbe de
Jordan rectifiable. Si
\begin{equation}\label{eq:rouche-condition}
  \abs{g(z)} < \abs{f(z)} \quad \text{pour tout } z \in \partial\Omega,
\end{equation}
alors $f$ et $f + g$ ont le même nombre de zéros dans $\Omega$, comptés avec
multiplicité.
\end{theoreme}

\begin{proof}
La condition~\eqref{eq:rouche-condition} implique que $f(z) \neq 0$ pour tout
$z \in \partial\Omega$ (car $\abs{g(z)} < \abs{f(z)}$ force
$\abs{f(z)} > 0$). De plus, pour $z \in \partial\Omega$ :
\[
  \abs{(f+g)(z)} \geq \abs{f(z)} - \abs{g(z)} > 0,
\]
donc $f + g$ n'a pas de zéro sur $\partial\Omega$ non plus.

Considérons la famille $F_t(z) = f(z) + tg(z)$ pour $t \in [0,1]$.
Pour tout $z \in \partial\Omega$ et tout $t \in [0,1]$ :
\[
  \abs{F_t(z)} = \abs{f(z) + tg(z)} \geq \abs{f(z)} - t\abs{g(z)}
  \geq \abs{f(z)} - \abs{g(z)} > 0.
\]
Donc $F_t$ n'a pas de zéro sur $\partial\Omega$ pour tout $t \in [0,1]$.

Par le principe de l'argument, le nombre de zéros de $F_t$ dans $\Omega$ est :
\[
  N(t) = \frac{1}{2\pi i} \oint_{\partial\Omega} \frac{F_t'(z)}{F_t(z)} \, \dd z
  = \frac{1}{2\pi i} \oint_{\partial\Omega} \frac{f'(z) + tg'(z)}{f(z) + tg(z)} \, \dd z.
\]
La fonction $t \mapsto N(t)$ est continue (l'intégrande dépend continûment de $t$
et le dénominateur ne s'annule pas sur $\partial\Omega$). Comme $N(t)$ est à
valeurs entières, il est constant. En particulier :
\[
  N(1) = N(0),
\]
c'est-à-dire que le nombre de zéros de $f + g$ dans $\Omega$ est égal au nombre
de zéros de $f$ dans $\Omega$.
\end{proof}

\begin{remarque}[Forme symétrique du théorème de Rouché]\label{rem:rouche-symetrique}
Une version équivalente, parfois plus commode, est la suivante : si $f$ et $g$
sont holomorphes sur un voisinage de $\overline{\Omega}$ et si
\[
  \abs{f(z) - g(z)} < \abs{f(z)} + \abs{g(z)}
  \quad \text{pour tout } z \in \partial\Omega,
\]
alors $f$ et $g$ ont le même nombre de zéros dans $\Omega$. Cette version,
due à Estermann, est strictement plus forte que la version classique.
\end{remarque}

\begin{center}
\begin{tikzpicture}[scale=1.1]
  % Illustration de la condition de Rouché
  \draw[->] (-3.5,0) -- (3.5,0) node[right] {$\re$};
  \draw[->] (0,-2.5) -- (0,2.5) node[above] {$\im$};

  % Le cercle |f(z)|
  \draw[thick, blue!70!black, dashed] (0,0) circle (2);
  \node[blue!70!black] at (2.5, 1.8) {$\abs{f(z)}$};

  % Le vecteur f(z)
  \draw[->, ultra thick, blue!70!black] (0,0) -- (1.4, 1.4)
    node[midway, above left] {$f(z)$};

  % Le vecteur g(z) — plus court
  \draw[->, ultra thick, red!70!black] (1.4,1.4) -- (1.9, 0.7)
    node[midway, right] {$g(z)$};

  % Le vecteur f(z)+g(z)
  \draw[->, ultra thick, green!50!black] (0,0) -- (1.9, 0.7)
    node[below right] {$f(z)+g(z)$};

  % Condition
  \node[align=center] at (0, -3.2) {Condition de Rouché : $\abs{g(z)} < \abs{f(z)}$\\
  $\Rightarrow$ $f(z) + g(z) \neq 0$ sur $\partial\Omega$};
\end{tikzpicture}
\end{center}

\section{Applications du théorème de Rouché}\label{sec:applications-rouche}

\subsection{Le théorème fondamental de l'algèbre}

\begin{theoreme}[Théorème fondamental de l'algèbre]\label{thm:fta-rouche}
\index{théorème fondamental de l'algèbre}
Tout polynôme non constant $P(z) = a_n z^n + a_{n-1}z^{n-1} + \cdots + a_0$
à coefficients complexes ($a_n \neq 0$, $n \geq 1$) possède exactement $n$
racines dans $\C$, comptées avec multiplicité.
\end{theoreme}

\begin{proof}
Posons $f(z) = a_n z^n$ et $g(z) = a_{n-1}z^{n-1} + \cdots + a_0$, de sorte
que $P(z) = f(z) + g(z)$.

Pour $\abs{z} = R$ :
\[
  \abs{g(z)} \leq \abs{a_{n-1}} R^{n-1} + \cdots + \abs{a_0}
\]
et $\abs{f(z)} = \abs{a_n} R^n$. Pour $R$ suffisamment grand :
\[
  \frac{\abs{g(z)}}{\abs{f(z)}} \leq \frac{\abs{a_{n-1}}}{\abs{a_n} R}
  + \frac{\abs{a_{n-2}}}{\abs{a_n} R^2} + \cdots + \frac{\abs{a_0}}{\abs{a_n} R^n}
  < 1.
\]
Plus précisément, il suffit de choisir
\[
  R > \max\!\left(1, \; \frac{n(\abs{a_{n-1}} + \cdots + \abs{a_0})}{\abs{a_n}}\right).
\]
Par le théorème de Rouché~\ref{thm:rouche} appliqué à $\Omega = D(0,R)$,
$P = f + g$ a le même nombre de zéros que $f(z) = a_n z^n$ dans $D(0,R)$,
c'est-à-dire exactement $n$ zéros (comptés avec multiplicité), puisque
$f$ a un unique zéro en $z = 0$ de multiplicité~$n$.
\end{proof}

\begin{exemple}[Localisation des zéros]\label{ex:rouche-localisation}
Montrons que le polynôme $P(z) = z^5 + 3z + 1$ a exactement un zéro dans le
disque $D(0,1)$.

Posons $f(z) = 3z$ et $g(z) = z^5 + 1$. Sur $\abs{z} = 1$ :
\[
  \abs{f(z)} = 3, \qquad \abs{g(z)} \leq \abs{z}^5 + 1 = 2 < 3 = \abs{f(z)}.
\]
Par Rouché, $P = f + g$ a le même nombre de zéros que $f(z) = 3z$ dans $D(0,1)$,
soit exactement~$1$.
\end{exemple}

\begin{exemple}[Zéros dans une couronne]\label{ex:rouche-couronne}
Montrons que $P(z) = z^4 - 6z + 3$ a exactement trois zéros dans la couronne
$1 < \abs{z} < 2$.

\textbf{Étape 1 :} Sur $\abs{z} = 2$, posons $f(z) = z^4$ et $g(z) = -6z + 3$.
Alors $\abs{f(z)} = 16$ et $\abs{g(z)} \leq 12 + 3 = 15 < 16$. Par Rouché,
$P$ a $4$ zéros dans $D(0,2)$.

\textbf{Étape 2 :} Sur $\abs{z} = 1$, posons $f(z) = -6z$ et $g(z) = z^4 + 3$.
Alors $\abs{f(z)} = 6$ et $\abs{g(z)} \leq 1 + 3 = 4 < 6$. Par Rouché,
$P$ a $1$ zéro dans $D(0,1)$.

Donc $P$ a $4 - 1 = 3$ zéros dans la couronne $1 < \abs{z} < 2$.
\end{exemple}

\subsection{Perturbation des racines}

\begin{proposition}[Stabilité du nombre de zéros]\label{prop:stabilite-zeros}
Soit $f$ holomorphe sur un voisinage de $\overline{D(a,r)}$ avec exactement
$n$ zéros dans $D(a,r)$ et $f(z) \neq 0$ sur $\abs{z-a} = r$. Alors il
existe $\varepsilon > 0$ tel que toute fonction $h$ holomorphe sur un voisinage
de $\overline{D(a,r)}$ avec $\sup_{\abs{z-a}=r} \abs{h(z) - f(z)} < \varepsilon$
possède également exactement $n$ zéros dans $D(a,r)$.
\end{proposition}

\begin{proof}
Posons $\varepsilon = \min_{\abs{z-a}=r} \abs{f(z)} > 0$ (le minimum est atteint
car $f$ est continue et non nulle sur le cercle compact). Si
$\abs{h(z) - f(z)} < \varepsilon \leq \abs{f(z)}$ sur $\abs{z-a} = r$,
le théorème de Rouché avec $g = h - f$ donne le résultat.
\end{proof}

\section{Le théorème de Hurwitz}\label{sec:hurwitz}

\begin{theoreme}[Théorème de Hurwitz]\label{thm:hurwitz}
\index{théorème de Hurwitz}
Soit $(f_n)_{n \geq 1}$ une suite de fonctions holomorphes sur un domaine
$\Omega$ convergeant uniformément sur tout compact vers une fonction $f$.
Si $f$ n'est pas identiquement nulle, alors pour tout $a \in \Omega$ tel que
$f(a) = 0$ d'ordre $m$, et pour tout $r > 0$ assez petit
(avec $\overline{D(a,r)} \subset \Omega$ et $f(z) \neq 0$ pour
$0 < \abs{z-a} \leq r$), il existe $N \in \N$ tel que pour tout $n \geq N$,
$f_n$ possède exactement $m$ zéros dans $D(a,r)$, comptés avec multiplicité.

En particulier, si chaque $f_n$ est sans zéro dans $\Omega$ et $f_n \to f$
uniformément sur tout compact, alors $f$ est soit identiquement nulle, soit
sans zéro dans $\Omega$.
\end{theoreme}

\begin{proof}
Choisissons $r > 0$ tel que $\overline{D(a,r)} \subset \Omega$ et $f$ n'a pas
de zéro sur le cercle $\abs{z-a} = r$ ni dans le disque épointé
$0 < \abs{z-a} \leq r$ (c'est possible car les zéros de $f$ sont isolés).

Le nombre de zéros de $f$ dans $D(a,r)$ est exactement $m$ (le zéro en $a$
d'ordre~$m$). Posons
\[
  \varepsilon = \min_{\abs{z-a}=r} \abs{f(z)} > 0.
\]
Par convergence uniforme sur le compact $\{\abs{z-a} = r\}$, il existe $N$
tel que pour tout $n \geq N$ :
\[
  \max_{\abs{z-a}=r} \abs{f_n(z) - f(z)} < \varepsilon \leq \abs{f(z)}.
\]
Par le théorème de Rouché, $f_n$ et $f$ ont le même nombre de zéros dans
$D(a,r)$, soit $m$.

Pour la seconde assertion : si chaque $f_n$ est sans zéro, alors le nombre
de zéros de $f_n$ dans tout $D(a,r)$ est~$0$. Par ce qui précède, si $f$
n'est pas identiquement nulle, tout zéro de $f$ aurait multiplicité~$0$,
ce qui est impossible. Donc $f \equiv 0$ ou $f$ est sans zéro.
\end{proof}

\section{Le théorème de l'application ouverte}\label{sec:application-ouverte}

\begin{theoreme}[Théorème de l'application ouverte]\label{thm:application-ouverte}
\index{théorème de l'application ouverte}
Soit $f$ une fonction holomorphe et non constante sur un domaine
$\Omega \subseteq \C$. Alors $f$ est une \textbf{application ouverte} :
l'image de tout ouvert de $\Omega$ par $f$ est un ouvert de~$\C$.
\end{theoreme}

\begin{proof}
Soit $U \subseteq \Omega$ un ouvert et soit $w_0 = f(z_0)$ avec $z_0 \in U$.
Nous devons montrer que $w_0$ est un point intérieur de $f(U)$.

Puisque $f$ n'est pas constante, $z_0$ est un zéro isolé de $f(z) - w_0$.
Soit $m \geq 1$ l'ordre de ce zéro. Choisissons $r > 0$ tel que
$\overline{D(z_0, r)} \subset U$ et $f(z) - w_0 \neq 0$ pour
$0 < \abs{z - z_0} \leq r$.

Posons $\delta = \min_{\abs{z - z_0} = r} \abs{f(z) - w_0} > 0$.
Nous affirmons que $D(w_0, \delta) \subset f(U)$.

En effet, soit $w \in D(w_0, \delta)$, c'est-à-dire $\abs{w - w_0} < \delta$.
Considérons la fonction $F(z) = f(z) - w_0$ et la perturbation
$G(z) = w_0 - w$ (constante). Pour $\abs{z - z_0} = r$ :
\[
  \abs{G(z)} = \abs{w - w_0} < \delta \leq \abs{f(z) - w_0} = \abs{F(z)}.
\]
Par le théorème de Rouché, $F + G = f(z) - w$ a le même nombre de zéros que
$F = f(z) - w_0$ dans $D(z_0, r)$, soit $m \geq 1$.

Donc il existe $z_1 \in D(z_0, r) \subset U$ avec $f(z_1) = w$.
Ceci montre que $w \in f(U)$, et donc $D(w_0, \delta) \subset f(U)$.
\end{proof}

\begin{corollaire}[Principe du maximum (variante)]\label{cor:max-ouvert}
Soit $f$ holomorphe et non constante sur un domaine $\Omega$. Alors $\abs{f}$
n'atteint pas son maximum dans $\Omega$.
\end{corollaire}

\begin{proof}
Supposons par l'absurde que $\abs{f}$ atteint son maximum en $z_0 \in \Omega$,
c'est-à-dire $\abs{f(z)} \leq \abs{f(z_0)}$ pour tout $z \in \Omega$. Puisque
$f$ est une application ouverte, $f(\Omega)$ est un ouvert contenant
$w_0 = f(z_0)$. Il existe donc $\varepsilon > 0$ avec
$D(w_0, \varepsilon) \subset f(\Omega)$, et en particulier il existe $z_1 \in \Omega$
avec $\abs{f(z_1)} > \abs{f(z_0)}$, contradiction.
\end{proof}

\begin{corollaire}[Inverse d'une bijection holomorphe]\label{cor:inverse-holomorphe}
Si $f : \Omega_1 \to \Omega_2$ est une bijection holomorphe entre domaines, alors
$f^{-1}$ est holomorphe (et $f'(z) \neq 0$ pour tout $z \in \Omega_1$).
\end{corollaire}

\begin{proof}
Puisque $f$ est une application ouverte, $f^{-1}$ est continue. Pour montrer
que $f^{-1}$ est holomorphe, il suffit de montrer que $f'(z) \neq 0$ pour tout
$z \in \Omega_1$.

Supposons que $f'(z_0) = 0$ pour un $z_0 \in \Omega_1$, et soit $w_0 = f(z_0)$.
Alors $f(z) - w_0$ a un zéro d'ordre $m \geq 2$ en $z_0$. Par la preuve du
théorème de l'application ouverte, pour $w$ proche de $w_0$ (mais $w \neq w_0$),
l'équation $f(z) = w$ a $m \geq 2$ solutions (comptées avec multiplicité) dans
un voisinage de $z_0$. Pour $w$ générique (hors d'un ensemble discret), ces
solutions sont distinctes, contredisant l'injectivité de~$f$.

Donc $f'(z) \neq 0$ partout, et $f^{-1}$ est holomorphe par le théorème
d'inversion locale (ou directement : $(f^{-1})'(w) = 1/f'(f^{-1}(w))$).
\end{proof}

\begin{center}
\begin{tikzpicture}[scale=0.9]
  % Illustration du théorème de l'application ouverte
  \begin{scope}[xshift=-4.5cm]
    \draw[thick, fill=blue!8, draw=blue!60!black, rounded corners=8pt]
      (-1.8,-1.5) rectangle (1.8,1.5);
    \node[blue!60!black] at (0,1.9) {$\Omega$};

    % Petit disque ouvert
    \draw[thick, fill=green!20, draw=green!60!black] (0,0) circle (0.8);
    \node[font=\small] at (0,0) {$U$};
    \fill (0,0) circle (1.5pt) node[below right, font=\tiny] {$z_0$};
  \end{scope}

  \draw[->, thick, >=stealth] (-1.8,0) -- (-0.5,0) node[midway, above] {$f$};

  \begin{scope}[xshift=3cm]
    % Image — forme ouverte non régulière
    \draw[thick, fill=red!8, draw=red!60!black]
      plot[smooth cycle, tension=0.7] coordinates
      {(-2,0.5) (-1.2,1.6) (0.5,1.8) (1.8,1) (2,-0.3) (1.2,-1.5) (-0.5,-1.8) (-1.8,-0.8)};
    \node[red!60!black] at (0,2.3) {$f(\Omega)$};

    % Image du petit ouvert — aussi ouvert
    \draw[thick, fill=green!20, draw=green!60!black]
      plot[smooth cycle, tension=0.8] coordinates
      {(-0.5,0.3) (0.2,0.7) (0.8,0.2) (0.5,-0.5) (-0.3,-0.4)};
    \node[font=\small] at (0.2,0.1) {$f(U)$};
    \fill (0.2,0.1) circle (1.5pt) node[below right, font=\tiny] {$w_0$};
  \end{scope}
\end{tikzpicture}
\end{center}

\section{Exercices}\label{sec:exercices-ch9}

\begin{exercice}[Dénombrement de zéros]\label{exo:denombrement-zeros}
Déterminer le nombre de zéros (comptés avec multiplicité) des fonctions
suivantes dans le disque $D(0,1)$ :
\begin{enumerate}[label=(\alph*)]
  \item $f(z) = z^7 - 5z^3 + z - 2$,
  \item $g(z) = z^6 + 4z^2 + 1$,
  \item $h(z) = e^z - 4z^3$.
\end{enumerate}
\textit{Indication :} appliquer le théorème de Rouché avec un choix judicieux
de décomposition $f = F + G$.
\end{exercice}

\begin{exercice}[Zéros dans le demi-plan]\label{exo:zeros-demi-plan}
Montrer que le polynôme $P(z) = z^4 + z^3 + 4z^2 + 2z + 3$ a exactement deux
zéros dans le demi-plan $\re(z) > 0$.

\textit{Indication :} montrer d'abord que $P$ n'a pas de zéro réel ni de zéro
purement imaginaire, puis utiliser le principe de l'argument sur un
grand demi-cercle dans le demi-plan droit.
\end{exercice}

\begin{exercice}[Théorème de Rouché généralisé]\label{exo:rouche-generalise}
Soient $f$ et $g$ holomorphes sur un voisinage de $\overline{D(0,1)}$ telles que
$\abs{f(z) + g(z)} < \abs{f(z)} + \abs{g(z)}$ pour tout $z$ avec $\abs{z} = 1$.
Montrer que $f$ et $g$ ont le même nombre de zéros dans $D(0,1)$.

\textit{Indication :} montrer que $f/g$ n'a pas de valeur réelle négative ni
nulle sur $\abs{z} = 1$, puis considérer $\ind((f/g) \circ \gamma, 0)$.
\end{exercice}

\begin{exercice}[Application du principe de l'argument]\label{exo:principe-arg-calcul}
Soit $f(z) = \frac{z^2 e^z}{(z-1)(z+3)}$. Calculer
\[
  \frac{1}{2\pi i} \oint_{\abs{z}=2} \frac{f'(z)}{f(z)} \, \dd z.
\]
\end{exercice}

\begin{exercice}[Théorème de Hurwitz et injectivité]\label{exo:hurwitz-injectivite}
Soit $(f_n)$ une suite de fonctions holomorphes injectives sur un domaine
$\Omega$, convergeant uniformément sur tout compact vers $f$. Montrer que $f$
est soit injective, soit constante.

\textit{Indication :} si $f(a) = f(b)$ pour $a \neq b$, appliquer Hurwitz
à $f_n(z) - f_n(a)$ au voisinage de $b$.
\end{exercice}

\begin{exercice}[Localisation fine par Rouché]\label{exo:rouche-fine}
Soit $P(z) = z^5 + z + 16$. Montrer que :
\begin{enumerate}[label=(\alph*)]
  \item $P$ a exactement un zéro dans chacun des quatre disques
    $D(2e^{ik\pi/2}, 1)$ pour $k = 0, 1, 2, 3$,
  \item le cinquième zéro est situé dans $D(0, 1/16)$.
\end{enumerate}
\textit{Indication :} pour (b), écrire $P(z) = 16(1 + z/16 + z^5/16)$ et
montrer que $z + z^5$ est petit devant $16$ sur $\abs{z} = 1/16$.
\end{exercice}

\begin{exercice}[Nombre de solutions d'une équation transcendante]\label{exo:eq-transcendante}
Montrer que l'équation $e^z = az^n$ (où $a \in \C \setminus \{0\}$ et $n \geq 1$)
a exactement $n$ solutions dans un disque $D(0,R)$ pour $R$ suffisamment grand.

\textit{Indication :} réécrire sous la forme $e^z - az^n = 0$ et comparer
avec $-az^n$ sur $\abs{z} = R$.
\end{exercice}

\begin{exercice}[Application ouverte et non-constance]\label{exo:app-ouverte}
Soit $f$ holomorphe sur $\C$ telle que $f(\C) \subseteq \R$. Montrer que $f$
est constante, en utilisant le théorème de l'application ouverte.
\end{exercice}

\begin{exercice}[Variation de l'argument — calcul explicite]\label{exo:variation-arg}
Soit $f(z) = z^3 - 2z + 2$.
\begin{enumerate}[label=(\alph*)]
  \item Calculer le nombre de zéros de $f$ dans le premier quadrant
    $\{z : \re(z) > 0, \im(z) > 0\}$ en calculant la variation de l'argument
    de $f$ le long du contour constitué du segment $[0, R]$, de l'arc de cercle
    de rayon $R$ dans le premier quadrant, et du segment $[iR, 0]$, pour
    $R \to +\infty$.
  \item Vérifier le résultat avec le théorème de Rouché.
\end{enumerate}
\end{exercice}


%%%%%%%%%%%%%%%%%%%%%%%%%%%%%%%%%%%%%%%%%%%%%%%%%%%%%%%%%%%%%%%%%%%%
%  CHAPITRE 10 — Introduction aux Fonctions Entières et Méromorphes
%%%%%%%%%%%%%%%%%%%%%%%%%%%%%%%%%%%%%%%%%%%%%%%%%%%%%%%%%%%%%%%%%%%%
\chapter{Introduction aux Fonctions Entières et Méromorphes}\label{ch:entieres-meromorphes}

\section{Fonctions entières}\label{sec:fonctions-entieres}

\begin{definition}[Fonction entière]\label{def:fonction-entiere}
\index{fonction entière}
Une fonction $f : \C \to \C$ est dite \textbf{entière} si elle est holomorphe
sur tout $\C$. De manière équivalente, $f$ est représentable par une série
entière de rayon de convergence infini :
\[
  f(z) = \sum_{n=0}^{\infty} a_n z^n, \qquad a_n \in \C,
  \quad \text{convergente pour tout } z \in \C.
\]
\end{definition}

\begin{exemple}[Fonctions entières classiques]\label{ex:entieres-classiques}
Les fonctions suivantes sont entières :
\begin{enumerate}[label=(\roman*)]
  \item \textbf{Polynômes :} $P(z) = a_n z^n + \cdots + a_0$.
  \item \textbf{L'exponentielle :} $e^z = \sum_{n=0}^{\infty} \frac{z^n}{n!}$.
  \item \textbf{Les fonctions trigonométriques :}
    $\sin z = \sum_{n=0}^{\infty} \frac{(-1)^n z^{2n+1}}{(2n+1)!}$, \;
    $\cos z = \sum_{n=0}^{\infty} \frac{(-1)^n z^{2n}}{(2n)!}$.
  \item \textbf{Les fonctions hyperboliques :} $\sinh z$, $\cosh z$.
  \item \textbf{La fonction d'Airy :}
    $\mathrm{Ai}(z) = \frac{1}{\pi} \int_0^{\infty}
    \cos\!\left(\frac{t^3}{3} + zt\right) \dd t$.
  \item \textbf{Compositions et produits :} $e^{e^z}$, $\sin(z^2)$,
    $z \, e^{z}$, etc.
\end{enumerate}
\end{exemple}

\begin{theoreme}[Théorème de Liouville]\label{thm:liouville-revisit}
\index{théorème de Liouville}
Toute fonction entière bornée est constante. Plus généralement, si $f$ est
entière et s'il existe $C > 0$ et $n \in \N$ tels que
$\abs{f(z)} \leq C(1 + \abs{z}^n)$ pour tout $z \in \C$, alors $f$ est un
polynôme de degré au plus~$n$.
\end{theoreme}

\begin{proof}
Le cas $n = 0$ est le théorème de Liouville classique. Pour le cas général,
écrivons $f(z) = \sum_{k=0}^{\infty} a_k z^k$. Par les inégalités de Cauchy
sur le cercle de rayon $R$ :
\[
  \abs{a_k} \leq \frac{\max_{\abs{z}=R} \abs{f(z)}}{R^k}
  \leq \frac{C(1 + R^n)}{R^k}.
\]
Pour $k > n$, en faisant $R \to +\infty$ :
\[
  \abs{a_k} \leq \lim_{R \to \infty} \frac{C(1 + R^n)}{R^k} = 0.
\]
Donc $a_k = 0$ pour tout $k > n$, et $f$ est un polynôme de degré~$\leq n$.
\end{proof}

\begin{corollaire}[Théorème fondamental de l'algèbre — par Liouville]\label{cor:fta-liouville}
Tout polynôme non constant possède au moins une racine complexe.
\end{corollaire}

\begin{proof}
Si $P$ est un polynôme sans racine, alors $1/P$ est entière. Comme
$\abs{P(z)} \to \infty$ quand $\abs{z} \to \infty$, $1/P$ est bornée,
donc constante par Liouville, ce qui contredit le fait que $P$ est non constant.
\end{proof}

\section{Ordre de croissance}\label{sec:ordre-croissance}

L'une des questions fondamentales sur les fonctions entières est la vitesse à
laquelle elles croissent. Pour les polynômes, la croissance est polynomiale ;
pour $e^z$, elle est exponentielle. L'<<~ordre~>> d'une fonction entière
quantifie précisément cette croissance.

\begin{definition}[Ordre d'une fonction entière]\label{def:ordre-entiere}
\index{ordre!d'une fonction entière}
Soit $f$ une fonction entière. On définit la \textbf{fonction maximum} de $f$~par
\[
  M(r) = M_f(r) = \max_{\abs{z}=r} \abs{f(z)}, \qquad r \geq 0.
\]
L'\textbf{ordre} de $f$ est
\[
  \rho = \rho(f) = \limsup_{r \to +\infty} \frac{\log \log M(r)}{\log r}
  \in [0, +\infty].
\]
On dit que $f$ est d'\textbf{ordre fini} si $\rho < +\infty$.
\end{definition}

\begin{remarque}[Interprétation]\label{rem:ordre-interpretation}
L'ordre $\rho$ est le plus petit nombre (dans $[0,+\infty]$) tel que
$M(r) \leq e^{r^{\rho + \varepsilon}}$ pour tout $\varepsilon > 0$ et $r$
assez grand. Autrement dit, $f$ est d'ordre $\rho$ si et seulement si
\[
  \forall \varepsilon > 0, \quad M(r) = O(e^{r^{\rho+\varepsilon}})
  \quad \text{et} \quad
  M(r) \neq O(e^{r^{\rho-\varepsilon}}).
\]
\end{remarque}

\begin{exemple}[Ordres de fonctions classiques]\label{ex:ordres-classiques}
\begin{enumerate}[label=(\roman*)]
  \item Tout polynôme non nul est d'ordre $\rho = 0$.
  \item $e^z$ : on a $M(r) = e^r$, donc
    $\frac{\log \log M(r)}{\log r} = \frac{\log r}{\log r} = 1$, et $\rho = 1$.
  \item $e^{z^k}$ pour $k \in \N^*$ : $M(r) = e^{r^k}$, donc $\rho = k$.
  \item $\sin z$ : $M(r) \sim \frac{1}{2} e^r$ pour $r \to \infty$, donc $\rho = 1$.
  \item $\cos\!\sqrt{z}$ (entière car la série ne contient que des puissances
    entières) : on a $M(r) \sim \frac{1}{2} e^{\sqrt{r}}$, donc $\rho = 1/2$.
  \item $e^{e^z}$ : $M(r) = e^{e^r}$, donc $\rho = +\infty$ (ordre infini).
\end{enumerate}
\end{exemple}

\begin{definition}[Type d'une fonction entière]\label{def:type-entiere}
\index{type d'une fonction entière}
Soit $f$ une fonction entière d'ordre $\rho \in (0, +\infty)$. Le \textbf{type}
de $f$ est
\[
  \sigma = \sigma(f) = \limsup_{r \to +\infty} \frac{\log M(r)}{r^\rho}
  \in [0, +\infty].
\]
On dit que $f$ est de \textbf{type normal} si $0 < \sigma < +\infty$, de
\textbf{type minimal} si $\sigma = 0$, et de \textbf{type maximal} si
$\sigma = +\infty$.
\end{definition}

\begin{exemple}[Types]\label{ex:types}
\begin{enumerate}[label=(\roman*)]
  \item $e^{az}$ ($a \neq 0$) est d'ordre $1$ et de type $\abs{a}$.
  \item $\sin z$ est d'ordre $1$ et de type $1$.
  \item $e^{z^2}$ est d'ordre $2$ et de type $1$.
\end{enumerate}
\end{exemple}

\begin{proposition}[Ordre et coefficients de Taylor]\label{prop:ordre-coefficients}
Soit $f(z) = \sum_{n=0}^{\infty} a_n z^n$ une fonction entière d'ordre $\rho$.
Alors
\[
  \rho = \limsup_{n \to \infty} \frac{n \log n}{-\log \abs{a_n}}.
\]
Réciproquement, cette formule permet de calculer l'ordre directement à partir
des coefficients de Taylor.
\end{proposition}

\begin{proof}
Ceci est une conséquence de la formule de Hadamard reliant le rayon de
convergence aux coefficients. La preuve complète procède en deux étapes.

\textbf{Borne supérieure.} Soit $\rho' > \rho$. Alors $M(r) \leq e^{r^{\rho'}}$
pour $r$ assez grand. Par les inégalités de Cauchy, $\abs{a_n} \leq M(r)/r^n$
pour tout $r > 0$. Choisissons $r = (n/(\rho' e))^{1/\rho'}$ (qui minimise
approximativement $e^{r^{\rho'}}/r^n$). Un calcul direct donne
$\abs{a_n} \leq C \cdot (\rho' e / n)^{n/\rho'}$ pour une constante $C$, d'où
\[
  -\log\abs{a_n} \geq \frac{n}{\rho'}\left(\log n - \log(\rho' e)\right) + O(1),
\]
ce qui implique $\limsup \frac{n\log n}{-\log\abs{a_n}} \leq \rho'$. En faisant
$\rho' \to \rho$, on obtient la borne supérieure.

\textbf{Borne inférieure.} On montre de façon analogue, en utilisant
$M(r) \leq \sum \abs{a_n} r^n$ et le terme dominant, que
$\limsup \frac{n\log n}{-\log\abs{a_n}} \geq \rho$.
\end{proof}

\section{Le petit théorème de Picard}\label{sec:picard}

L'un des résultats les plus profonds de l'analyse complexe est le théorème de
Picard, qui renforce spectaculairement le théorème de Liouville.

\begin{theoreme}[Petit théorème de Picard]\label{thm:petit-picard}
\index{théorème de Picard!petit}
Toute fonction entière non constante prend toute valeur complexe, à une
exception au plus. Autrement dit, si $f : \C \to \C$ est entière et s'il
existe deux valeurs distinctes $a, b \in \C$ telles que $f(z) \neq a$ et
$f(z) \neq b$ pour tout $z \in \C$, alors $f$ est constante.
\end{theoreme}

\begin{remarque}[Sur la preuve]\label{rem:picard-preuve}
La preuve du petit théorème de Picard dépasse le cadre de ce cours introductif.
Elle repose sur la théorie du revêtement universel et la fonction modulaire
$\lambda : \mathbb{H} \to \C \setminus \{0, 1\}$. Une autre approche utilise
le lemme de Schottky et le théorème de Bloch.
\end{remarque}

\begin{remarque}[Optimalité]\label{rem:picard-optimal}
L'exception d'une valeur est nécessaire : $e^z$ est entière, non constante,
et ne prend jamais la valeur $0$. Mais $e^z$ prend toute autre valeur
$w \neq 0$ (en effet, $e^z = w \Leftrightarrow z = \log\abs{w} + i\Arg(w) + 2k\pi i$).
\end{remarque}

\begin{center}
\begin{tikzpicture}[scale=0.85]
  % Illustration du théorème de Picard
  \draw[->] (-4,0) -- (4,0) node[right] {$\re$};
  \draw[->] (0,-3) -- (0,3) node[above] {$\im$};

  % Plan d'arrivée : tout C sauf peut-être un point
  \fill[blue!8] (-3.8,-2.8) rectangle (3.8,2.8);

  % La valeur exceptionnelle
  \fill[white] (0,0) circle (5pt);
  \draw[red, thick] (0,0) circle (5pt);
  \node[red, below right] at (0.2, -0.2) {$a$};

  % Annotation
  \node[align=center] at (0, -3.8) {Image d'une fonction entière non constante\\
  $f(\C) = \C \setminus \{a\}$ (au plus une valeur omise)};

  % Flèches indiquant que f atteint tout le reste
  \foreach \angle in {30, 75, 120, 165, 210, 255, 300, 345} {
    \draw[->, green!50!black, thick] (0,0) -- ({2.2*cos(\angle)}, {2.2*sin(\angle)});
  }
\end{tikzpicture}
\end{center}

\begin{theoreme}[Grand théorème de Picard — énoncé]\label{thm:grand-picard}
\index{théorème de Picard!grand}
Soit $f$ holomorphe sur un voisinage épointé $D(a, r) \setminus \{a\}$ d'une
singularité essentielle isolée $a$. Alors $f$ prend toute valeur complexe, à
une exception au plus, dans tout voisinage de~$a$.

En d'autres termes, pour tout $\varepsilon > 0$, l'ensemble
$f(D(a, \varepsilon) \setminus \{a\})$ contient $\C$ privé d'au plus un point.
\end{theoreme}

\section{Produits infinis et théorème de Weierstrass}\label{sec:weierstrass}

La question fondamentale que nous abordons maintenant est celle de la
\emph{factorisation} des fonctions entières : peut-on exprimer une fonction
entière comme un produit (infini) faisant intervenir ses zéros, à l'image
de la factorisation des polynômes ?

\begin{definition}[Produit infini]\label{def:produit-infini}
\index{produit infini}
Soit $(a_n)_{n \geq 1}$ une suite de nombres complexes. On dit que le
\textbf{produit infini} $\prod_{n=1}^{\infty} a_n$ \textbf{converge} si les
produits partiels $P_N = \prod_{n=1}^N a_n$ tendent vers une limite
$P \neq 0$ quand $N \to \infty$. On écrit alors
$\prod_{n=1}^{\infty} a_n = P$.

On dit que $\prod_{n=1}^{\infty}(1 + u_n)$ converge \textbf{absolument}
si $\sum_{n=1}^{\infty} \abs{u_n}$ converge.
\end{definition}

\begin{proposition}[Convergence absolue des produits]\label{prop:produit-convergence}
Le produit $\prod_{n=1}^{\infty} (1 + u_n)$ converge absolument si et seulement
si la série $\sum_{n=1}^{\infty} \abs{u_n}$ converge. Dans ce cas, le produit
ne dépend pas de l'ordre des facteurs.
\end{proposition}

\begin{proof}
Si $\sum \abs{u_n} < \infty$, alors $\abs{u_n} < 1$ pour $n$ assez grand.
Pour ces $n$, on utilise l'inégalité $\abs{\log(1+u_n)} \leq 2\abs{u_n}$
(valable pour $\abs{u_n} \leq 1/2$). Donc $\sum \log(1 + u_n)$ converge
absolument, et par conséquent $\prod (1 + u_n) = \exp(\sum \log(1+u_n))$
converge vers une valeur non nulle.

Réciproquement, si $\prod (1 + u_n)$ converge absolument (c'est-à-dire
$\prod (1 + \abs{u_n})$ converge), alors $\sum \log(1 + \abs{u_n})$ converge,
ce qui entraîne $\sum \abs{u_n} < \infty$ puisque
$\log(1 + x) \geq x/2$ pour $0 \leq x \leq 1$.
\end{proof}

\begin{definition}[Facteurs élémentaires de Weierstrass]\label{def:facteurs-elementaires}
\index{facteur élémentaire}
Pour $p \in \N$, on définit le \textbf{facteur élémentaire} $E_p$ par
\[
  E_0(z) = 1 - z, \qquad
  E_p(z) = (1 - z) \exp\!\left(z + \frac{z^2}{2} + \cdots + \frac{z^p}{p}\right)
  \quad \text{pour } p \geq 1.
\]
\end{definition}

\begin{lemme}[Estimation du facteur élémentaire]\label{lem:estimation-Ep}
Pour tout $p \geq 0$ et tout $\abs{z} \leq 1$ :
\[
  \abs{1 - E_p(z)} \leq \abs{z}^{p+1}.
\]
\end{lemme}

\begin{proof}
Écrivons $E_p(z) = \exp(h_p(z))$ où
$h_p(z) = \log(1-z) + z + z^2/2 + \cdots + z^p/p$ (pour $\abs{z} < 1$).
Le développement en série de $\log(1-z) = -\sum_{k=1}^{\infty} z^k/k$ donne :
\[
  h_p(z) = -\sum_{k=p+1}^{\infty} \frac{z^k}{k}.
\]
Pour $\abs{z} \leq 1$ : $\abs{h_p(z)} \leq \sum_{k=p+1}^{\infty} \abs{z}^k/k
\leq \abs{z}^{p+1} \sum_{k=0}^{\infty} \abs{z}^k/(p+1+k) \leq \abs{z}^{p+1}$.
Puis $\abs{1 - E_p(z)} = \abs{1 - e^{h_p(z)}} \leq \abs{h_p(z)} e^{\abs{h_p(z)}}
\leq \abs{z}^{p+1} \cdot e \leq C\abs{z}^{p+1}$.
On peut affiner pour obtenir $\abs{1 - E_p(z)} \leq \abs{z}^{p+1}$ par un
argument plus soigné utilisant le fait que les coefficients de Taylor de
$1 - E_p(z)$ sont positifs.
\end{proof}

\begin{theoreme}[Théorème de factorisation de Weierstrass]\label{thm:weierstrass}
\index{théorème de Weierstrass!factorisation}
Soit $(a_n)_{n \geq 1}$ une suite de nombres complexes non nuls avec
$\abs{a_n} \to +\infty$, et soit $m \geq 0$ un entier. Alors il existe une
fonction entière dont les zéros sont exactement $z = 0$ (de multiplicité~$m$)
et les $a_n$ (chacun comptés avec multiplicité). Plus précisément, si $(p_n)$
est une suite d'entiers positifs telle que
\[
  \sum_{n=1}^{\infty} \left(\frac{r}{\abs{a_n}}\right)^{p_n+1} < \infty
  \quad \text{pour tout } r > 0,
\]
alors le produit
\[
  f(z) = z^m \prod_{n=1}^{\infty} E_{p_n}\!\left(\frac{z}{a_n}\right)
\]
converge uniformément sur tout compact et définit une fonction entière ayant
exactement les zéros prescrits.

En particulier, on peut toujours choisir $p_n = n - 1$.
\end{theoreme}

\begin{proof}[Esquisse de preuve]
Il suffit de montrer la convergence uniforme sur tout disque $\abs{z} \leq R$.
Pour $n$ assez grand, $\abs{a_n} > R$, donc $\abs{z/a_n} \leq R/\abs{a_n} < 1$.
Par le lemme~\ref{lem:estimation-Ep} :
\[
  \abs{1 - E_{p_n}(z/a_n)} \leq \left(\frac{R}{\abs{a_n}}\right)^{p_n+1}.
\]
La condition $\sum (R/\abs{a_n})^{p_n+1} < \infty$ assure la convergence
absolue du produit $\prod E_{p_n}(z/a_n)$, uniformément sur $\abs{z} \leq R$.

Un produit uniformément convergent de fonctions holomorphes est holomorphe
(par passage à la limite sous le signe $\exp\circ\log$), et les zéros du
produit sont exactement les $a_n$ avec les multiplicités souhaitées.
\end{proof}

\begin{corollaire}[Forme générale d'une fonction entière]\label{cor:hadamard-forme}
Toute fonction entière $f$ admet une factorisation de la forme
\[
  f(z) = z^m e^{g(z)} \prod_{n=1}^{\infty} E_{p_n}\!\left(\frac{z}{a_n}\right),
\]
où $m \geq 0$ est l'ordre du zéro en $0$, $(a_n)$ sont les zéros non nuls
de $f$, et $g$ est une fonction entière.
\end{corollaire}

\subsection{Factorisation du sinus}

L'exemple le plus classique de produit de Weierstrass est la factorisation
de $\sin(\pi z)$.

\begin{theoreme}[Produit d'Euler pour le sinus]\label{thm:euler-sinus}
\index{produit d'Euler!sinus}
Pour tout $z \in \C$ :
\begin{equation}\label{eq:euler-sinus}
  \sin(\pi z) = \pi z \prod_{n=1}^{\infty}
  \left(1 - \frac{z^2}{n^2}\right).
\end{equation}
\end{theoreme}

\begin{proof}
Les zéros de $\sin(\pi z)$ sont les entiers $z = n \in \Z$, chacun simple.
Par le théorème de Weierstrass (avec $p_n = 1$ pour assurer la convergence,
car $\sum 1/n^2 < \infty$) :
\[
  \sin(\pi z) = z \, e^{g(z)} \prod_{n=1}^{\infty}
  E_1\!\left(\frac{z}{n}\right) E_1\!\left(\frac{z}{-n}\right)
  = z \, e^{g(z)} \prod_{n=1}^{\infty}
  \left(1 - \frac{z}{n}\right)e^{z/n} \cdot
  \left(1 + \frac{z}{n}\right)e^{-z/n}.
\]
Les facteurs exponentiels se simplifient deux à deux, et on obtient :
\[
  \sin(\pi z) = z \, e^{g(z)} \prod_{n=1}^{\infty}
  \left(1 - \frac{z^2}{n^2}\right).
\]
Il reste à montrer que $g(z) = \log \pi$ (constante). En divisant par $z$
et prenant la limite $z \to 0$ :
\[
  \lim_{z \to 0} \frac{\sin(\pi z)}{z} = \pi
  = e^{g(0)} \prod_{n=1}^{\infty} 1 = e^{g(0)},
\]
donc $g(0) = \log \pi$. On peut montrer que $g'(z) = 0$ en comparant les
dérivées logarithmiques des deux membres : la dérivée logarithmique de
$\sin(\pi z)$ est $\pi\cot(\pi z)$, et celle du produit (avec le facteur $z$)
est $1/z + \sum_{n=1}^{\infty} 2z/(z^2 - n^2)$, ce qui donne exactement
le développement en série partielle de $\pi\cot(\pi z)$ (identité classique).
Donc $g'(z) = 0$ et $g$ est constante, d'où $e^{g(z)} = \pi$.
\end{proof}

\begin{corollaire}[Formule de Wallis]\label{cor:wallis}
\index{formule de Wallis}
En évaluant~\eqref{eq:euler-sinus} en $z = 1/2$ :
\[
  1 = \sin\!\left(\frac{\pi}{2}\right)
  = \frac{\pi}{2} \prod_{n=1}^{\infty}
  \left(1 - \frac{1}{4n^2}\right)
  = \frac{\pi}{2} \prod_{n=1}^{\infty}
  \frac{(2n-1)(2n+1)}{(2n)^2},
\]
ce qui donne la \textbf{formule de Wallis} :
\[
  \frac{\pi}{2} = \prod_{n=1}^{\infty} \frac{(2n)^2}{(2n-1)(2n+1)}
  = \frac{2 \cdot 2}{1 \cdot 3} \cdot \frac{4 \cdot 4}{3 \cdot 5}
  \cdot \frac{6 \cdot 6}{5 \cdot 7} \cdots
\]
\end{corollaire}

\begin{center}
\begin{tikzpicture}[scale=0.65]
  % Zéros de sin(pi z) dans le plan complexe
  \draw[->] (-5.5,0) -- (5.5,0) node[right] {$\re$};
  \draw[->] (0,-3) -- (0,3) node[above] {$\im$};

  % Les zéros sont les entiers
  \foreach \k in {-5,-4,...,5} {
    \fill[red] (\k, 0) circle (3pt);
  }

  % Labels
  \node[below, font=\small] at (-3, -0.3) {$-3$};
  \node[below, font=\small] at (-2, -0.3) {$-2$};
  \node[below, font=\small] at (-1, -0.3) {$-1$};
  \node[below left, font=\small] at (0, -0.3) {$0$};
  \node[below, font=\small] at (1, -0.3) {$1$};
  \node[below, font=\small] at (2, -0.3) {$2$};
  \node[below, font=\small] at (3, -0.3) {$3$};

  % Points de suspension
  \node at (-5, 0.5) {$\cdots$};
  \node at (5, 0.5) {$\cdots$};

  \node at (0, -4) {Zéros de $\sin(\pi z)$ : les entiers $\Z$};
\end{tikzpicture}
\end{center}

\section{Fonctions méromorphes et décomposition de Mittag-Leffler}%
\label{sec:mittag-leffler}

\begin{definition}[Fonction méromorphe]\label{def:fonction-meromorphe}
\index{fonction méromorphe}
Une fonction $f$ est \textbf{méromorphe} sur un ouvert $\Omega \subseteq \C$
si elle est holomorphe sur $\Omega$ privé d'un ensemble discret $S \subset \Omega$,
et si chaque point de $S$ est un pôle de~$f$.

L'ensemble des fonctions méromorphes sur $\Omega$ forme un \textbf{corps}
pour les opérations d'addition et de multiplication.
\end{definition}

\begin{exemple}[Fonctions méromorphes classiques]\label{ex:meromorphes}
\begin{enumerate}[label=(\roman*)]
  \item Toute fonction rationnelle $P(z)/Q(z)$ est méromorphe sur $\C$.
  \item $\tan z = \sin z / \cos z$ est méromorphe sur $\C$, avec des pôles
    simples en $z = \pi/2 + n\pi$ ($n \in \Z$).
  \item La fonction $\Gamma$ d'Euler est méromorphe sur $\C$, avec des pôles
    simples en $z = 0, -1, -2, \ldots$
  \item $\pi \cot(\pi z)$ est méromorphe avec des pôles simples en $z \in \Z$.
\end{enumerate}
\end{exemple}

Le théorème de Mittag-Leffler est l'analogue <<~additif~>> du théorème de
Weierstrass : là où Weierstrass prescrit les zéros d'une fonction entière,
Mittag-Leffler prescrit les parties principales aux pôles d'une fonction
méromorphe.

\begin{theoreme}[Théorème de Mittag-Leffler]\label{thm:mittag-leffler}
\index{théorème de Mittag-Leffler}
Soit $(b_n)_{n \geq 1}$ une suite de points distincts de $\C$ avec
$\abs{b_n} \to +\infty$, et pour chaque $n$, soit $P_n$ une \textbf{partie
principale} (c'est-à-dire un polynôme en $1/(z - b_n)$ sans terme constant) :
\[
  P_n(z) = \sum_{k=1}^{m_n} \frac{c_{n,k}}{(z - b_n)^k}.
\]
Alors il existe une fonction méromorphe $f$ sur $\C$ dont les seules
singularités sont des pôles en les $b_n$, avec partie principale $P_n$ en
chaque $b_n$. Plus précisément, il existe des polynômes $p_n(z)$ tels que
\[
  f(z) = \sum_{n=1}^{\infty} \bigl(P_n(z) - p_n(z)\bigr)
\]
converge uniformément sur tout compact de $\C \setminus \{b_1, b_2, \ldots\}$.

De plus, la fonction méromorphe la plus générale ayant ces parties principales
est de la forme $f(z) + g(z)$, où $g$ est une fonction entière arbitraire.
\end{theoreme}

\begin{proof}[Esquisse de preuve]
L'idée est de soustraire à chaque $P_n(z)$ un polynôme $p_n(z)$ (le début
de son développement de Taylor en $0$) pour assurer la convergence.

Pour $n$ assez grand, $\abs{b_n} > 0$ et $P_n(z)$ est holomorphe dans le disque
$D(0, \abs{b_n}/2)$. On développe $P_n(z)$ en série de Taylor au voisinage
de~$0$ :
\[
  P_n(z) = \sum_{k=0}^{\infty} \alpha_{n,k} z^k \quad
  \text{pour } \abs{z} < \abs{b_n}.
\]
On choisit $p_n(z) = \sum_{k=0}^{d_n} \alpha_{n,k} z^k$ (troncature à
l'ordre $d_n$). Pour $\abs{z} \leq R < \abs{b_n}$ :
\[
  \abs{P_n(z) - p_n(z)} \leq \sum_{k=d_n+1}^{\infty} \abs{\alpha_{n,k}} R^k
  \leq C_n \left(\frac{R}{\abs{b_n}}\right)^{d_n+1} \cdot \frac{1}{1 - R/\abs{b_n}}.
\]
En choisissant $d_n$ assez grand (par exemple $d_n = n$), on assure
$\abs{P_n(z) - p_n(z)} \leq 2^{-n}$ sur $\abs{z} \leq R$ pour tout $R$ fixé
et $n$ assez grand. La convergence uniforme sur les compacts en résulte.
\end{proof}

\begin{exemple}[Décomposition de $\pi\cot(\pi z)$]\label{ex:cot-mittag-leffler}
La fonction $\pi\cot(\pi z)$ est méromorphe sur $\C$ avec des pôles simples
en $z = n \in \Z$, chacun de résidu~$1$. La partie principale en $z = n$ est
$\frac{1}{z - n}$. Par le théorème de Mittag-Leffler (ou par un calcul direct) :
\begin{equation}\label{eq:cot-decomposition}
  \pi \cot(\pi z) = \frac{1}{z} + \sum_{n=1}^{\infty}
  \left(\frac{1}{z - n} + \frac{1}{z + n}\right)
  = \frac{1}{z} + \sum_{n=1}^{\infty} \frac{2z}{z^2 - n^2}.
\end{equation}
Cette série converge uniformément sur tout compact de $\C \setminus \Z$.
\end{exemple}

\begin{proof}[Justification de~\eqref{eq:cot-decomposition}]
Posons $g(z) = \pi\cot(\pi z) - \frac{1}{z} - \sum_{n=1}^{N}
\left(\frac{1}{z-n} + \frac{1}{z+n}\right)$. On montre que $g$ est bornée :
elle est holomorphe (les singularités sont levées) et périodique de période~$1$.
Sur la bande $\{z : 0 \leq \re(z) \leq 1, \abs{\im(z)} \geq 1\}$,
$\pi\cot(\pi z) \to \mp \pi i$ quand $\im(z) \to \pm\infty$, et les termes
$\frac{1}{z-n} + \frac{1}{z+n} = \frac{2z}{z^2-n^2} \to 0$. Donc la
différence $h(z) = \pi\cot(\pi z) - \frac{1}{z} - \sum_{n=1}^{\infty}
(\frac{1}{z-n} + \frac{1}{z+n})$ est entière (les pôles s'annulent) et
bornée, donc constante par Liouville. Comme $h(z) \to 0$ quand $z \to 0$
(par parité), cette constante est~$0$.
\end{proof}

\section{Fonctions elliptiques : un aperçu}\label{sec:elliptiques}

\begin{definition}[Fonction elliptique]\label{def:fonction-elliptique}
\index{fonction elliptique}
Une \textbf{fonction elliptique} est une fonction méromorphe sur $\C$ qui est
doublement périodique : il existe deux nombres complexes $\omega_1, \omega_2$,
$\R$-linéairement indépendants, tels que
\[
  f(z + \omega_1) = f(z) \quad \text{et} \quad f(z + \omega_2) = f(z)
  \quad \text{pour tout } z \in \C.
\]
Le \textbf{réseau} associé est $\Lambda = \Z\omega_1 + \Z\omega_2$, et le
\textbf{parallélogramme fondamental} est
\[
  \Pi = \{t_1 \omega_1 + t_2 \omega_2 : 0 \leq t_1, t_2 < 1\}.
\]
\end{definition}

\begin{proposition}[Propriétés élémentaires]\label{prop:elliptiques-elem}
Soit $f$ une fonction elliptique non constante de réseau $\Lambda$. Alors :
\begin{enumerate}[label=(\roman*)]
  \item $f$ a un nombre fini de pôles dans $\Pi$ (et de même pour les zéros).
  \item La somme des résidus de $f$ dans $\Pi$ est nulle.
  \item $f$ a au moins deux pôles dans $\Pi$ (comptés avec multiplicité) ;
    en particulier, il n'existe pas de fonction elliptique holomorphe non
    constante.
  \item Le nombre de zéros de $f$ dans $\Pi$ est égal au nombre de pôles
    (comptés avec multiplicité). Ce nombre commun est appelé l'\textbf{ordre}
    de~$f$.
\end{enumerate}
\end{proposition}

\begin{proof}
\textbf{(i)} L'adhérence $\overline{\Pi}$ est compacte et les pôles sont
isolés, donc il n'y en a qu'un nombre fini.

\textbf{(ii)} Par le théorème des résidus, $\sum \Res = \frac{1}{2\pi i}
\oint_{\partial\Pi} f(z)\,\dd z$. Par double périodicité, les contributions
des côtés opposés s'annulent.

\textbf{(iii)} Si $f$ n'a pas de pôle, elle est entière et doublement périodique,
donc bornée sur $\overline\Pi$, donc bornée sur $\C$ par périodicité. Par
Liouville, $f$ est constante. Si $f$ a un seul pôle simple, son résidu est non
nul, contredisant~(ii). Donc $f$ a au moins deux pôles (ou un pôle d'ordre
$\geq 2$).

\textbf{(iv)} Appliquer le principe de l'argument à $\partial\Pi$ (avec le
même argument de cancellation par périodicité).
\end{proof}

\begin{definition}[Fonction $\wp$ de Weierstrass]\label{def:weierstrass-p}
\index{fonction de Weierstrass $\wp$}
La \textbf{fonction $\wp$ de Weierstrass} associée au réseau $\Lambda$ est
\[
  \wp(z) = \frac{1}{z^2} + \sum_{\omega \in \Lambda \setminus \{0\}}
  \left(\frac{1}{(z - \omega)^2} - \frac{1}{\omega^2}\right).
\]
\end{definition}

\begin{proposition}[Propriétés de $\wp$]\label{prop:weierstrass-p}
La fonction $\wp$ est une fonction elliptique d'ordre~$2$ par rapport au
réseau $\Lambda$. Elle est paire ($\wp(-z) = \wp(z)$), et son unique pôle dans
$\Pi$ est un pôle double en $z = 0$. De plus :
\[
  \wp'(z) = -2 \sum_{\omega \in \Lambda} \frac{1}{(z - \omega)^3},
\]
et $\wp$ satisfait l'équation différentielle
\begin{equation}\label{eq:weierstrass-ode}
  (\wp')^2 = 4\wp^3 - g_2 \wp - g_3,
\end{equation}
où $g_2 = 60 \sum_{\omega \neq 0} \omega^{-4}$ et
$g_3 = 140 \sum_{\omega \neq 0} \omega^{-6}$ sont les \textbf{invariants}
du réseau.
\end{proposition}

\begin{center}
\begin{tikzpicture}[scale=1.2]
  % Réseau et parallélogramme fondamental
  \def\wa{2}    % omega_1 = 2
  \def\wbx{0.7} % Re(omega_2)
  \def\wby{1.8} % Im(omega_2)

  % Dessiner quelques points du réseau
  \foreach \m in {-1,0,1,2} {
    \foreach \n in {-1,0,1} {
      \fill[gray!60] ({\m*\wa + \n*\wbx}, {\n*\wby}) circle (2pt);
    }
  }

  % Parallélogramme fondamental
  \draw[thick, fill=blue!10, fill opacity=0.5, draw=blue!70!black]
    (0,0) -- (\wa, 0) -- ({\wa+\wbx}, \wby) -- (\wbx, \wby) -- cycle;

  % Labels
  \node[below left, font=\small] at (0,0) {$0$};
  \node[below, font=\small] at (\wa, 0) {$\omega_1$};
  \node[above left, font=\small] at (\wbx, \wby) {$\omega_2$};
  \node[above right, font=\small] at ({\wa+\wbx}, \wby) {$\omega_1+\omega_2$};

  % Pôle double en 0
  \node[red, font=\small] at (0.15, 0.25) {$\times\times$};

  \node[blue!70!black] at ({(\wa+\wbx)/2}, {\wby/2}) {$\Pi$};

  \node at ({\wa/2}, -0.8) {Parallélogramme fondamental et réseau $\Lambda$};
\end{tikzpicture}
\end{center}

\section{Exercices}\label{sec:exercices-ch10}

\begin{exercice}[Ordre de croissance]\label{exo:ordre-croissance}
Calculer l'ordre de croissance des fonctions entières suivantes :
\begin{enumerate}[label=(\alph*)]
  \item $f(z) = e^{z^3 + z}$,
  \item $g(z) = \sum_{n=0}^{\infty} \frac{z^n}{(n!)^2}$,
  \item $h(z) = \cos(\sqrt{z})$ (vérifier d'abord que $h$ est entière),
  \item $k(z) = \prod_{n=1}^{\infty} \left(1 + \frac{z}{n^2}\right)$.
\end{enumerate}
\end{exercice}

\begin{exercice}[Application de Liouville généralisé]\label{exo:liouville-gen}
Soit $f$ entière telle que $\abs{f(z)} \leq C e^{a\abs{z}}$ pour tout $z \in \C$,
avec $C, a > 0$. Montrer que $f$ est d'ordre au plus~$1$ et de type au plus~$a$.

Soit de plus $f$ telle que $f(n) = 0$ pour tout $n \in \N$. Montrer que
$f \equiv 0$.

\textit{Indication :} considérer $g(z) = f(z)/\sin(\pi z)$ et montrer que
$g$ est entière bornée.
\end{exercice}

\begin{exercice}[Formules pour $\zeta(2k)$]\label{exo:zeta-pairs}
En développant $\pi\cot(\pi z)$ en série de Laurent au voisinage de $z = 0$
et en utilisant la décomposition~\eqref{eq:cot-decomposition}, montrer que
\[
  \sum_{n=1}^{\infty} \frac{1}{n^{2k}} = \frac{(-1)^{k+1}(2\pi)^{2k}}{2(2k)!} B_{2k},
\]
où $B_{2k}$ sont les nombres de Bernoulli. En déduire
$\sum_{n=1}^{\infty} \frac{1}{n^2} = \frac{\pi^2}{6}$ et
$\sum_{n=1}^{\infty} \frac{1}{n^4} = \frac{\pi^4}{90}$.
\end{exercice}

\begin{exercice}[Produit de Weierstrass pour $1/\Gamma$]\label{exo:weierstrass-gamma}
La fonction $1/\Gamma(z)$ est entière avec des zéros simples en
$z = 0, -1, -2, \ldots$ Montrer que
\[
  \frac{1}{\Gamma(z)} = z e^{\gamma z} \prod_{n=1}^{\infty}
  \left(1 + \frac{z}{n}\right) e^{-z/n},
\]
où $\gamma = \lim_{n\to\infty}\left(\sum_{k=1}^n \frac{1}{k} - \log n\right)$
est la constante d'Euler--Mascheroni.

\textit{Indication :} identifier les zéros, appliquer Weierstrass avec $p_n = 1$,
et déterminer le facteur exponentiel $e^{g(z)}$ par des arguments de croissance.
\end{exercice}

\begin{exercice}[Théorème de Mittag-Leffler — application]\label{exo:mittag-leffler-app}
Construire une fonction méromorphe sur $\C$ dont les seuls pôles sont les
entiers positifs $n \geq 1$, avec partie principale $\frac{1}{(z-n)^2}$ en
chaque~$n$. Donner une expression en série.

\textit{Indication :} considérer $\sum_{n=1}^{\infty}
\left(\frac{1}{(z-n)^2} - \frac{1}{n^2}\right)$ et montrer la convergence.
\end{exercice}

\begin{exercice}[Petit théorème de Picard — applications]\label{exo:picard}
\begin{enumerate}[label=(\alph*)]
  \item Montrer que si $f$ est entière et $\re(f(z)) > 0$ pour tout $z$,
    alors $f$ est constante.
  \item Montrer que si $f$ est entière et $f(z) \notin \R^-$ pour tout $z$,
    alors $f$ est constante.
  \item Montrer que si $f$ est entière et $\abs{f(z)} \neq 1$ pour tout $z$,
    alors $f$ est constante. (\textit{Indication :} $f(\C)$ est connexe.)
\end{enumerate}
\end{exercice}

\begin{exercice}[Fonctions elliptiques élémentaires]\label{exo:elliptiques}
Soit $\Lambda = \Z + \Z i$ le réseau des entiers de Gauss.
\begin{enumerate}[label=(\alph*)]
  \item Montrer que la fonction $\wp$ associée satisfait $\wp(iz) = -\wp(z)$.
    En déduire que $g_3 = 0$ dans l'équation~\eqref{eq:weierstrass-ode}.
  \item Quelles sont les valeurs de $\wp$ aux demi-périodes
    $\omega_1/2 = 1/2$, $\omega_2/2 = i/2$ et $(\omega_1+\omega_2)/2 = (1+i)/2$ ?
\end{enumerate}
\end{exercice}

\begin{exercice}[Théorème de Casorati-Weierstrass]\label{exo:casorati}
Soit $a$ une singularité essentielle isolée de $f$. Montrer que pour tout
$\varepsilon > 0$, l'image $f(D(a,\varepsilon) \setminus \{a\})$ est dense
dans~$\C$.

\textit{Indication :} supposer le contraire, c'est-à-dire qu'il existe $w_0$
et $\delta > 0$ tels que $\abs{f(z) - w_0} \geq \delta$ au voisinage de $a$.
Considérer $g(z) = 1/(f(z) - w_0)$ et montrer que $a$ est un pôle ou une
singularité effaçable de $f$, contradiction.
\end{exercice}

\begin{exercice}[Produit de Blaschke]\label{exo:blaschke}
Soit $(a_n)_{n\geq 1}$ une suite dans le disque unité $\D$ avec
$\sum_{n=1}^{\infty}(1 - \abs{a_n}) < \infty$ et $a_n \neq 0$ pour tout $n$.
Montrer que le \textbf{produit de Blaschke}
\[
  B(z) = \prod_{n=1}^{\infty} \frac{\abs{a_n}}{a_n} \cdot \frac{a_n - z}{1 - \conj{a_n} z}
\]
converge uniformément sur tout compact de $\D$ et définit une fonction holomorphe
bornée sur $\D$ dont les zéros sont exactement les $a_n$.

\textit{Indication :} montrer que $\abs{1 - \frac{\abs{a_n}}{a_n} \cdot
\frac{a_n - z}{1 - \conj{a_n}z}} \leq C(1 - \abs{a_n})$ sur tout compact
de~$\D$.
\end{exercice}


%%%%%%%%%%%%%%%%%%%%%%%%%%%%%%%%%%%%%%%%%%%%%%%%%%%%%%%%%%%%%%%%%%%%
%  BIBLIOGRAPHIE
%%%%%%%%%%%%%%%%%%%%%%%%%%%%%%%%%%%%%%%%%%%%%%%%%%%%%%%%%%%%%%%%%%%%
\begin{thebibliography}{99}

\bibitem{ahlfors}
L.~V.~Ahlfors,
\textit{Complex Analysis: An Introduction to the Theory of Analytic Functions
of One Complex Variable},
3\ieme{} édition, McGraw-Hill, New York, 1979.
\index{Ahlfors}

\bibitem{conway}
J.~B.~Conway,
\textit{Functions of One Complex Variable},
2\ieme{} édition, Graduate Texts in Mathematics~11,
Springer-Verlag, New York, 1978.
\index{Conway}

\bibitem{stein-shakarchi}
E.~M.~Stein et R.~Shakarchi,
\textit{Complex Analysis},
Princeton Lectures in Analysis~II,
Princeton University Press, Princeton, 2003.
\index{Stein}\index{Shakarchi}

\bibitem{cartan}
H.~Cartan,
\textit{Théorie élémentaire des fonctions analytiques d'une ou plusieurs
variables complexes},
Hermann, Paris, 1961.
\index{Cartan}

\bibitem{rudin-rca}
W.~Rudin,
\textit{Real and Complex Analysis},
3\ieme{} édition, McGraw-Hill, New York, 1987.
\index{Rudin}

\bibitem{berenstein-gay}
C.~A.~Berenstein et R.~Gay,
\textit{Complex Variables: An Introduction},
Graduate Texts in Mathematics~125,
Springer-Verlag, New York, 1991.
\index{Berenstein}\index{Gay}

\bibitem{remmert}
R.~Remmert,
\textit{Classical Topics in Complex Function Theory},
traduit par L.~Kay, Graduate Texts in Mathematics~172,
Springer-Verlag, New York, 1998.
\index{Remmert}

\end{thebibliography}

%%%%%%%%%%%%%%%%%%%%%%%%%%%%%%%%%%%%%%%%%%%%%%%%%%%%%%%%%%%%%%%%%%%%
%  INDEX
%%%%%%%%%%%%%%%%%%%%%%%%%%%%%%%%%%%%%%%%%%%%%%%%%%%%%%%%%%%%%%%%%%%%
\printindex

\end{document}
